%%%%%%%%%%%%%%%%%%%%%%%%%%%%%%%%%%%%%%%%%
% Report of comparations of the program
% alignandcomparemolecules
% LaTeX Template
% Version 0.1 (10/12/17)
%
% Part of this template was downloaded from:
% http://www.LaTeXTemplates.com
%
% Modifications by:
% Isaias Morales Salazar (Shiver) isaiasms117@gmail.com
% Luis Alfredo Nuñez (Fdx) luis.alfredo.nu@gmail.com
%
%%%%%%%%%%%%%%%%%%%%%%%%%%%%%%%%%%%%%%%%%
%-----------------------------------------------------------------------------% 
% Packages
%-----------------------------------------------------------------------------%
\documentclass[12pt]{article} 	
\usepackage[utf8]{inputenc}
\usepackage[spanish]{babel}
\usepackage[usenames]{color}
\usepackage{float}
\usepackage{graphicx}
\usepackage{amsmath,mathtools}
\usepackage[font={small},format=plain,labelfont=bf]{caption} % Size and format of foot of figures and tables
\usepackage[dvipsnames]{xcolor}
\usepackage[most]{tcolorbox}
\usepackage{listings} % Add parts of code
\usepackage{multicol}
\usepackage{array,multirow}
%-----------------------------------------------------------------------------%
% Margin settings 
%-----------------------------------------------------------------------------%
\usepackage{geometry}
\geometry{
	paper=letterpaper, % Paper type
	inner=2.5cm, % Inner margin
	outer=2.5cm, % Outer margin
	top=3.5cm, % Top margin
	bottom=1.5cm, % Bottom margin
}
%-----------------------------------------------------------------------------% 
% New Commands
%-----------------------------------------------------------------------------%
\renewcommand{\tablename}{Tabla} % Name of table foot 
\renewcommand{\baselinestretch}{1.5} % space between lines 1.5
\newcommand\tab[1][0.5cm]{\hspace*{#1}}
\newcommand\vtab[1][0.5cm]{\vspace*{#1}}
\newcommand{\classC}[1]{\textcolor{MidnightBlue}{\bf {#1}}}
\newcommand{\alert}[1]{\begin{center}
\fcolorbox{RoyalBlue}{Mahogany!90}{\tab\begin{minipage}[c]{0.8\textwidth}
\vtab[-7mm]\textcolor{Goldenrod}{
\begin{flushleft}{\bf{#1}}\end{flushleft}
}\vtab[-3mm]
\end{minipage}\tab}
\end{center}}
%-----------------------------------------------------------------------------% 
% Path
%-----------------------------------------------------------------------------% 
\newcommand*{\srcpath}{src/}
%-----------------------------------------------------------------------------%
%-----------------------------------------------------------------------------%
% Document
%-----------------------------------------------------------------------------%
%-----------------------------------------------------------------------------% 

\begin{document}
%-----------------------------------------------------------------------------%
% Cover Page
%-----------------------------------------------------------------------------% 
\pagestyle{empty} 

\begin{center}
\includegraphics[width=4cm]{escudo-buap.pdf}

BENEMÉRITA UNIVERSIDAD AUTÓNOMA DE PUEBLA\\
\rule{150mm}{0.1mm}
FACULTAD DE CIENCIAS QUÍMICAS\\
LABORATORIO DE FISICOQUÍMICA TEÓRICA
\rule{150mm}{0.1mm}

\vtab
\Large{Reporte de resultados para el programa \\ alignandcomparemolecules}\\
\vtab[.3cm]

\large{Título del proyecto: \\ \textbf{Desarrollo de software para comparar y resolver estructuras isoméricas mediante el tensor de inercia}}\\
\vtab[0.5cm]

PRESENTA\\
Luis Alfredo Nuñez Meneses \\
\vtab[0.5cm]

\tab[-0.5cm]
\begin{minipage}[t]{0.5\textwidth}
\begin{center}
DIRECTOR DE TESIS  \\
Dr. Juan Manuel Solano Altamirano \\
FQ-BUAP \\
\end{center}
\end{minipage}

\vtab[0.5cm]
\today\\
\end{center}

\begin{flushright}
  \large{Jefe del Departamento de Fisicoquímica\\
  Dr. Juan Carlos Ramírez García}\\
\end{flushright}

\newpage

%-----------------------------------------------------------------------------% 
% Resultados
%-----------------------------------------------------------------------------% 
\section{Resultados}
\subsection{Pruebas realizadas}
Las Pruebas que ya se realizaron son las siguientes y los resultados son óptimos y satisfactorios en todas las pruebas.\\
Realicé todas las pruebas que muestro a continuación, para hacer todas las comparaciones programé unos scripts que me ayudaran a 
automatizar todo el trabajo.

		%\tab[-2cm]
		%\begin{tabular}{c|m{8cm}|c|c}
		%\# & Moléculas & Restultado esperado & Resultado programa \\ \hline\hline
		%\multirow{4}{*}{\tab 1 \tab} & r-2-chlorobutane\_out\_G09\_inversion & \multirow{4}{*}{Estereisómeros} & \multirow{4}{*}{Estereisómeros} \\
		%						 & E = -618.05708 \tab Freq = 115.2921   &											& 											 \\ \cline{2-2}
		%						 & r-2-chlorobutane\_out\_G09 			  &											& 											 \\
		%						 & E = -618.05708 \tab Freq = 115.2921   &											& 											 \\ \hline
		%\multirow{4}{*}{\tab 1 \tab} & r-2-chlorobutane\_out\_G09\_inversion & \multirow{4}{*}{Estereisómeros} & \multirow{4}{*}{Estereisómeros} \\
		%						 & E = -618.05708 \tab Freq = 115.2921   &											& 											 \\ \cline{2-2}
		%						 & r-2-chlorobutane\_out\_G09 			  &											& 											 \\
		%						 & E = -618.05708 \tab Freq = 115.2921   &											& 											 \\ \hline
		%\end{tabular}
		%
		%\vtab[2cm]
		%\begin{tabular}{c|m{5cm}|m{5cm}|m{3cm}|m{3cm}}
		%\# & Molécula A & Molécula B & Resultado esperado & Resultado programa \\ \hline
		%1  & r-2-chlorobutane\_out\_G09\_inversion & r-2-chlorobutane\_out\_G09 & Estereisómeros & Estereisómeros \\
		%	& E = -618.057081944 					 & E = -618.057081944 & & \\
		%	& Freq = 115.2921 						 & Freq = 115.2921 & & \\ \hline
		%
		%\end{tabular}



	%\begin{multicols}{2}
	%\begin{center}
	%Molecule A: \\ SymmetricEnantiomers\/Opt\_Geo\/s-tmp\_rotated\_out\_G09\_output.xyz \\
	% Inertia Tensor - Molecule A \\
	%\vtab
	%\begin{tabular}{|c c c|}
	%            103.7 &  -4.21263e-05 & 8.53829e-06  \\ 
	%     -4.21263e-05 &       630.074 & 2.17164e-05  \\ 
	%      8.53829e-06 &   2.17164e-05 &     722.449   
	%\end{tabular}
	%
	%\vtab
	% EingenVectors - Molecule A     \\ 
	%\vtab
	%\begin{tabular}{|c c c|}
	%               -1  & -8.00311e-08  &1.37993e-08  \\ 
	%      8.00312e-08  &           -1  &2.35089e-07  \\ 
	%      1.37993e-08  &  2.35089e-07  &          1   
	%\end{tabular}
	%
	%\vtab
	% EingenValues - Molecule A       \\
	%\vtab
	%\begin{tabular}{|c c c|}
	%            103.7  &      630.074  &    722.449   
	%\end{tabular}
	%
	%\columnbreak
	%Molecule B: \\ SymmetricEnantiomers\/Opt\_Geo\/s-tmp\_rotated\_out\_G09.xyz\\
	%Inertia Tensor - Molecule B\\
	%
	%\vtab
	%\begin{tabular}{|c c c|}
	%         103.707   &     1.51142    &  1.05403  \\
	%         1.51142   &     630.103    &  1.67204  \\
	%         1.05403   &     1.67204    &   722.42  
	%\end{tabular}
	%                                                
	%\vtab
	%EingenVectors - Molecule B\\
	%\vtab
	%\begin{tabular}{|c c c|}
	%       -0.999994   &  0.00286582  & 0.00169581  \\
	%      -0.0028346   &   -0.999832  &   0.018136  \\
	%       0.0017475   &   0.0181311  &   0.999834  
	%\end{tabular}
	%                                                
	%\vtab
	%EingenValues - Molecule B\\
	%\vtab
	%\begin{tabular}{|c c c|}
	%         103.701   &     630.077   &   722.452   
	%\end{tabular}
	%
	%\end{center}
	%\end{multicols}
	%\textcolor{NavyBlue}{\large Resultado: Are Isomers}


\vtab[-2cm]
\begin{center}
{\large AsymmetricEnantiomers}
\end{center}
\begin{multicols}{2}
\begin{center}
Molecule A \\ 
r-2-chlorobutane\_out\_G09
\includegraphics[width=8cm]{Images/ImagesFromVMD/r-2-chlorobutane_out_G09.jpg}
\\
Inertia Tensor - Molecule A \\
\vtab
\begin{tabular}{|c c c|}
113.621	 & 	-8.52391	 & 	-0.0867925	 \\
-8.52391	 & 	164.234	 & 	-0.885702	 \\
-0.0867925	 & 	-0.885702	 & 	257.008
\end{tabular}

\vtab
 EingenVectors - Molecule A     \\
\vtab
\begin{tabular}{|c c c|}
-0.986829	 & 	-0.161759	 & 	-0.00158111	 \\
0.161767	 & 	-0.986784	 & 	-0.00941026	 \\
-3.80259e-05	 & 	-0.00954208	 & 	0.999954
\end{tabular}

\vtab
 EingenValues - Molecule A     \\
\vtab
\begin{tabular}{|c c c|}
112.224	 & 	165.623	 & 	257.016
\end{tabular}
\columnbreak

Molecule B \\ 
r-2-chlorobutane\_rotated\_out\_G09
\includegraphics[width=8cm]{Images/ImagesFromVMD/r-2-chlorobutane_rotated_out_G09.jpg}
\\
Inertia Tensor - Molecule B \\
\vtab
\begin{tabular}{|c c c|}
113.621	 & 	-8.52391	 & 	-0.0867928	 \\
-8.52391	 & 	164.234	 & 	-0.885702	 \\
-0.0867928	 & 	-0.885702	 & 	257.008
\end{tabular}

\vtab
 EingenVectors - Molecule B     \\
\vtab
\begin{tabular}{|c c c|}
-0.986829	 & 	-0.161759	 & 	-0.00158112	 \\
0.161767	 & 	-0.986784	 & 	-0.00941025	 \\
-3.80279e-05	 & 	-0.00954208	 & 	0.999954
\end{tabular}

\vtab
 EingenValues - Molecule B     \\
\vtab
\begin{tabular}{|c c c|}
112.224	 & 	165.623	 & 	257.016
\end{tabular}

\end{center}
\end{multicols}
\begin{center}
\vtab
\vtab
\textcolor{NavyBlue}{\Large Are the same molecule}
\end{center}
%%%%%%%%%%%%%%%%%%%%%%%%%%%%%%%%%%%%%%%%%%%%%%%%%%%%%%%%
%%%%%%%%%%%%%%%%%%%%%%%%%%%%%%%%%%%%%%%%%%%%%%%%%%%%%%%%
\newpage
%%%%%%%%%%%%%%%%%%%%%%%%%%%%%%%%%%%%%%%%%%%%%%%%%%%%%%%%
%%%%%%%%%%%%%%%%%%%%%%%%%%%%%%%%%%%%%%%%%%%%%%%%%%%%%%%%

\vtab[-2cm]
\begin{center}
{\large AsymmetricEnantiomers}
\end{center}
\begin{multicols}{2}
\begin{center}
Molecule A \\ 
r-2-chlorobutane\_out\_G09
\includegraphics[width=8cm]{Images/ImagesFromVMD/r-2-chlorobutane_out_G09.jpg}
\\
Inertia Tensor - Molecule A \\
\vtab
\begin{tabular}{|c c c|}
113.621	 & 	-8.52391	 & 	-0.0867925	 \\
-8.52391	 & 	164.234	 & 	-0.885702	 \\
-0.0867925	 & 	-0.885702	 & 	257.008
\end{tabular}

\vtab
 EingenVectors - Molecule A     \\
\vtab
\begin{tabular}{|c c c|}
-0.986829	 & 	-0.161759	 & 	-0.00158111	 \\
0.161767	 & 	-0.986784	 & 	-0.00941026	 \\
-3.80259e-05	 & 	-0.00954208	 & 	0.999954
\end{tabular}

\vtab
 EingenValues - Molecule A     \\
\vtab
\begin{tabular}{|c c c|}
112.224	 & 	165.623	 & 	257.016
\end{tabular}
\columnbreak

Molecule B \\ 
r-2-chlorobutane\_tmp\_out\_G09
\includegraphics[width=8cm]{Images/ImagesFromVMD/r-2-chlorobutane_tmp_out_G09.jpg}
\\
Inertia Tensor - Molecule B \\
\vtab
\begin{tabular}{|c c c|}
113.621	 & 	-8.52391	 & 	-0.0867928	 \\
-8.52391	 & 	164.234	 & 	-0.885704	 \\
-0.0867928	 & 	-0.885704	 & 	257.008
\end{tabular}

\vtab
 EingenVectors - Molecule B     \\
\vtab
\begin{tabular}{|c c c|}
-0.986829	 & 	-0.161759	 & 	-0.00158112	 \\
0.161767	 & 	-0.986784	 & 	-0.00941027	 \\
-3.8027e-05	 & 	-0.0095421	 & 	0.999954
\end{tabular}

\vtab
 EingenValues - Molecule B     \\
\vtab
\begin{tabular}{|c c c|}
112.224	 & 	165.623	 & 	257.016
\end{tabular}

\end{center}
\end{multicols}
\begin{center}
\vtab
\vtab
\textcolor{NavyBlue}{\Large Are the same molecule}
\end{center}
%%%%%%%%%%%%%%%%%%%%%%%%%%%%%%%%%%%%%%%%%%%%%%%%%%%%%%%%
%%%%%%%%%%%%%%%%%%%%%%%%%%%%%%%%%%%%%%%%%%%%%%%%%%%%%%%%
\newpage
%%%%%%%%%%%%%%%%%%%%%%%%%%%%%%%%%%%%%%%%%%%%%%%%%%%%%%%%
%%%%%%%%%%%%%%%%%%%%%%%%%%%%%%%%%%%%%%%%%%%%%%%%%%%%%%%%

\vtab[-2cm]
\begin{center}
{\large AsymmetricEnantiomers}
\end{center}
\begin{multicols}{2}
\begin{center}
Molecule A \\ 
r-2-chlorobutane\_out\_G09
\includegraphics[width=8cm]{Images/ImagesFromVMD/r-2-chlorobutane_out_G09.jpg}
\\
Inertia Tensor - Molecule A \\
\vtab
\begin{tabular}{|c c c|}
113.621	 & 	-8.52391	 & 	-0.0867925	 \\
-8.52391	 & 	164.234	 & 	-0.885702	 \\
-0.0867925	 & 	-0.885702	 & 	257.008
\end{tabular}

\vtab
 EingenVectors - Molecule A     \\
\vtab
\begin{tabular}{|c c c|}
-0.986829	 & 	-0.161759	 & 	-0.00158111	 \\
0.161767	 & 	-0.986784	 & 	-0.00941026	 \\
-3.80259e-05	 & 	-0.00954208	 & 	0.999954
\end{tabular}

\vtab
 EingenValues - Molecule A     \\
\vtab
\begin{tabular}{|c c c|}
112.224	 & 	165.623	 & 	257.016
\end{tabular}
\columnbreak

Molecule B \\ 
s-2-chlorobutane\_out\_G09
\includegraphics[width=8cm]{Images/ImagesFromVMD/s-2-chlorobutane_out_G09.jpg}
\\
Inertia Tensor - Molecule B \\
\vtab
\begin{tabular}{|c c c|}
113.621	 & 	8.52396	 & 	0.0868155	 \\
8.52396	 & 	164.235	 & 	-0.885715	 \\
0.0868155	 & 	-0.885715	 & 	257.008
\end{tabular}

\vtab
 EingenVectors - Molecule B     \\
\vtab
\begin{tabular}{|c c c|}
-0.98683	 & 	0.161755	 & 	0.00158125	 \\
-0.161763	 & 	-0.986785	 & 	-0.00941041	 \\
3.81715e-05	 & 	-0.00954226	 & 	0.999954
\end{tabular}

\vtab
 EingenValues - Molecule B     \\
\vtab
\begin{tabular}{|c c c|}
112.224	 & 	165.624	 & 	257.017
\end{tabular}

\end{center}
\end{multicols}
\begin{center}
\vtab
\vtab
\textcolor{NavyBlue}{\Large Are Isomer and Stereoisomers}
\end{center}
%%%%%%%%%%%%%%%%%%%%%%%%%%%%%%%%%%%%%%%%%%%%%%%%%%%%%%%%
%%%%%%%%%%%%%%%%%%%%%%%%%%%%%%%%%%%%%%%%%%%%%%%%%%%%%%%%
\newpage
%%%%%%%%%%%%%%%%%%%%%%%%%%%%%%%%%%%%%%%%%%%%%%%%%%%%%%%%
%%%%%%%%%%%%%%%%%%%%%%%%%%%%%%%%%%%%%%%%%%%%%%%%%%%%%%%%

\vtab[-2cm]
\begin{center}
{\large AsymmetricEnantiomers}
\end{center}
\begin{multicols}{2}
\begin{center}
Molecule A \\ 
r-2-chlorobutane\_out\_G09
\includegraphics[width=8cm]{Images/ImagesFromVMD/r-2-chlorobutane_out_G09.jpg}
\\
Inertia Tensor - Molecule A \\
\vtab
\begin{tabular}{|c c c|}
113.621	 & 	-8.52391	 & 	-0.0867925	 \\
-8.52391	 & 	164.234	 & 	-0.885702	 \\
-0.0867925	 & 	-0.885702	 & 	257.008
\end{tabular}

\vtab
 EingenVectors - Molecule A     \\
\vtab
\begin{tabular}{|c c c|}
-0.986829	 & 	-0.161759	 & 	-0.00158111	 \\
0.161767	 & 	-0.986784	 & 	-0.00941026	 \\
-3.80259e-05	 & 	-0.00954208	 & 	0.999954
\end{tabular}

\vtab
 EingenValues - Molecule A     \\
\vtab
\begin{tabular}{|c c c|}
112.224	 & 	165.623	 & 	257.016
\end{tabular}
\columnbreak

Molecule B \\ 
s-2-chlorobutane\_rotated\_out\_G09
\includegraphics[width=8cm]{Images/ImagesFromVMD/s-2-chlorobutane_rotated_out_G09.jpg}
\\
Inertia Tensor - Molecule B \\
\vtab
\begin{tabular}{|c c c|}
113.621	 & 	8.52397	 & 	0.0867992	 \\
8.52397	 & 	164.235	 & 	-0.885706	 \\
0.0867992	 & 	-0.885706	 & 	257.008
\end{tabular}

\vtab
 EingenVectors - Molecule B     \\
\vtab
\begin{tabular}{|c c c|}
-0.98683	 & 	0.161756	 & 	0.00158113	 \\
-0.161763	 & 	-0.986785	 & 	-0.00941034	 \\
3.80628e-05	 & 	-0.00954217	 & 	0.999954
\end{tabular}

\vtab
 EingenValues - Molecule B     \\
\vtab
\begin{tabular}{|c c c|}
112.224	 & 	165.624	 & 	257.017
\end{tabular}

\end{center}
\end{multicols}
\begin{center}
\vtab
\vtab
\textcolor{NavyBlue}{\Large Are Isomer and Stereoisomers}
\end{center}
%%%%%%%%%%%%%%%%%%%%%%%%%%%%%%%%%%%%%%%%%%%%%%%%%%%%%%%%
%%%%%%%%%%%%%%%%%%%%%%%%%%%%%%%%%%%%%%%%%%%%%%%%%%%%%%%%
\newpage
%%%%%%%%%%%%%%%%%%%%%%%%%%%%%%%%%%%%%%%%%%%%%%%%%%%%%%%%
%%%%%%%%%%%%%%%%%%%%%%%%%%%%%%%%%%%%%%%%%%%%%%%%%%%%%%%%

\vtab[-2cm]
\begin{center}
{\large AsymmetricEnantiomers}
\end{center}
\begin{multicols}{2}
\begin{center}
Molecule A \\ 
r-2-chlorobutane\_out\_G09
\includegraphics[width=8cm]{Images/ImagesFromVMD/r-2-chlorobutane_out_G09.jpg}
\\
Inertia Tensor - Molecule A \\
\vtab
\begin{tabular}{|c c c|}
113.621	 & 	-8.52391	 & 	-0.0867925	 \\
-8.52391	 & 	164.234	 & 	-0.885702	 \\
-0.0867925	 & 	-0.885702	 & 	257.008
\end{tabular}

\vtab
 EingenVectors - Molecule A     \\
\vtab
\begin{tabular}{|c c c|}
-0.986829	 & 	-0.161759	 & 	-0.00158111	 \\
0.161767	 & 	-0.986784	 & 	-0.00941026	 \\
-3.80259e-05	 & 	-0.00954208	 & 	0.999954
\end{tabular}

\vtab
 EingenValues - Molecule A     \\
\vtab
\begin{tabular}{|c c c|}
112.224	 & 	165.623	 & 	257.016
\end{tabular}
\columnbreak

Molecule B \\ 
s-2-chlorobutane\_tmp\_out\_G09
\includegraphics[width=8cm]{Images/ImagesFromVMD/s-2-chlorobutane_tmp_out_G09.jpg}
\\
Inertia Tensor - Molecule B \\
\vtab
\begin{tabular}{|c c c|}
113.621	 & 	8.52397	 & 	0.086812	 \\
8.52397	 & 	164.235	 & 	-0.885714	 \\
0.086812	 & 	-0.885714	 & 	257.008
\end{tabular}

\vtab
 EingenVectors - Molecule B     \\
\vtab
\begin{tabular}{|c c c|}
-0.98683	 & 	0.161756	 & 	0.00158123	 \\
-0.161763	 & 	-0.986785	 & 	-0.00941041	 \\
3.81472e-05	 & 	-0.00954225	 & 	0.999954
\end{tabular}

\vtab
 EingenValues - Molecule B     \\
\vtab
\begin{tabular}{|c c c|}
112.224	 & 	165.624	 & 	257.017
\end{tabular}

\end{center}
\end{multicols}
\begin{center}
\vtab
\vtab
\textcolor{NavyBlue}{\Large Are Isomer and Stereoisomers}
\end{center}
%%%%%%%%%%%%%%%%%%%%%%%%%%%%%%%%%%%%%%%%%%%%%%%%%%%%%%%%
%%%%%%%%%%%%%%%%%%%%%%%%%%%%%%%%%%%%%%%%%%%%%%%%%%%%%%%%
\newpage
%%%%%%%%%%%%%%%%%%%%%%%%%%%%%%%%%%%%%%%%%%%%%%%%%%%%%%%%
%%%%%%%%%%%%%%%%%%%%%%%%%%%%%%%%%%%%%%%%%%%%%%%%%%%%%%%%

\vtab[-2cm]
\begin{center}
{\large AsymmetricEnantiomers}
\end{center}
\begin{multicols}{2}
\begin{center}
Molecule A \\ 
r-2-chlorobutane\_rotated\_out\_G09
\includegraphics[width=8cm]{Images/ImagesFromVMD/r-2-chlorobutane_rotated_out_G09.jpg}
\\
Inertia Tensor - Molecule A \\
\vtab
\begin{tabular}{|c c c|}
113.621	 & 	-8.52391	 & 	-0.0867928	 \\
-8.52391	 & 	164.234	 & 	-0.885702	 \\
-0.0867928	 & 	-0.885702	 & 	257.008
\end{tabular}

\vtab
 EingenVectors - Molecule A     \\
\vtab
\begin{tabular}{|c c c|}
-0.986829	 & 	-0.161759	 & 	-0.00158112	 \\
0.161767	 & 	-0.986784	 & 	-0.00941025	 \\
-3.80279e-05	 & 	-0.00954208	 & 	0.999954
\end{tabular}

\vtab
 EingenValues - Molecule A     \\
\vtab
\begin{tabular}{|c c c|}
112.224	 & 	165.623	 & 	257.016
\end{tabular}
\columnbreak

Molecule B \\ 
r-2-chlorobutane\_tmp\_out\_G09
\includegraphics[width=8cm]{Images/ImagesFromVMD/r-2-chlorobutane_tmp_out_G09.jpg}
\\
Inertia Tensor - Molecule B \\
\vtab
\begin{tabular}{|c c c|}
113.621	 & 	-8.52391	 & 	-0.0867928	 \\
-8.52391	 & 	164.234	 & 	-0.885704	 \\
-0.0867928	 & 	-0.885704	 & 	257.008
\end{tabular}

\vtab
 EingenVectors - Molecule B     \\
\vtab
\begin{tabular}{|c c c|}
-0.986829	 & 	-0.161759	 & 	-0.00158112	 \\
0.161767	 & 	-0.986784	 & 	-0.00941027	 \\
-3.8027e-05	 & 	-0.0095421	 & 	0.999954
\end{tabular}

\vtab
 EingenValues - Molecule B     \\
\vtab
\begin{tabular}{|c c c|}
112.224	 & 	165.623	 & 	257.016
\end{tabular}

\end{center}
\end{multicols}
\begin{center}
\vtab
\vtab
\textcolor{NavyBlue}{\Large Are the same molecule}
\end{center}
%%%%%%%%%%%%%%%%%%%%%%%%%%%%%%%%%%%%%%%%%%%%%%%%%%%%%%%%
%%%%%%%%%%%%%%%%%%%%%%%%%%%%%%%%%%%%%%%%%%%%%%%%%%%%%%%%
\newpage
%%%%%%%%%%%%%%%%%%%%%%%%%%%%%%%%%%%%%%%%%%%%%%%%%%%%%%%%
%%%%%%%%%%%%%%%%%%%%%%%%%%%%%%%%%%%%%%%%%%%%%%%%%%%%%%%%

\vtab[-2cm]
\begin{center}
{\large AsymmetricEnantiomers}
\end{center}
\begin{multicols}{2}
\begin{center}
Molecule A \\ 
r-2-chlorobutane\_rotated\_out\_G09
\includegraphics[width=8cm]{Images/ImagesFromVMD/r-2-chlorobutane_rotated_out_G09.jpg}
\\
Inertia Tensor - Molecule A \\
\vtab
\begin{tabular}{|c c c|}
113.621	 & 	-8.52391	 & 	-0.0867928	 \\
-8.52391	 & 	164.234	 & 	-0.885702	 \\
-0.0867928	 & 	-0.885702	 & 	257.008
\end{tabular}

\vtab
 EingenVectors - Molecule A     \\
\vtab
\begin{tabular}{|c c c|}
-0.986829	 & 	-0.161759	 & 	-0.00158112	 \\
0.161767	 & 	-0.986784	 & 	-0.00941025	 \\
-3.80279e-05	 & 	-0.00954208	 & 	0.999954
\end{tabular}

\vtab
 EingenValues - Molecule A     \\
\vtab
\begin{tabular}{|c c c|}
112.224	 & 	165.623	 & 	257.016
\end{tabular}
\columnbreak

Molecule B \\ 
s-2-chlorobutane\_out\_G09
\includegraphics[width=8cm]{Images/ImagesFromVMD/s-2-chlorobutane_out_G09.jpg}
\\
Inertia Tensor - Molecule B \\
\vtab
\begin{tabular}{|c c c|}
113.621	 & 	8.52396	 & 	0.0868155	 \\
8.52396	 & 	164.235	 & 	-0.885715	 \\
0.0868155	 & 	-0.885715	 & 	257.008
\end{tabular}

\vtab
 EingenVectors - Molecule B     \\
\vtab
\begin{tabular}{|c c c|}
-0.98683	 & 	0.161755	 & 	0.00158125	 \\
-0.161763	 & 	-0.986785	 & 	-0.00941041	 \\
3.81715e-05	 & 	-0.00954226	 & 	0.999954
\end{tabular}

\vtab
 EingenValues - Molecule B     \\
\vtab
\begin{tabular}{|c c c|}
112.224	 & 	165.624	 & 	257.017
\end{tabular}

\end{center}
\end{multicols}
\begin{center}
\vtab
\vtab
\textcolor{NavyBlue}{\Large Are Isomer and Stereoisomers}
\end{center}
%%%%%%%%%%%%%%%%%%%%%%%%%%%%%%%%%%%%%%%%%%%%%%%%%%%%%%%%
%%%%%%%%%%%%%%%%%%%%%%%%%%%%%%%%%%%%%%%%%%%%%%%%%%%%%%%%
\newpage
%%%%%%%%%%%%%%%%%%%%%%%%%%%%%%%%%%%%%%%%%%%%%%%%%%%%%%%%
%%%%%%%%%%%%%%%%%%%%%%%%%%%%%%%%%%%%%%%%%%%%%%%%%%%%%%%%

\vtab[-2cm]
\begin{center}
{\large AsymmetricEnantiomers}
\end{center}
\begin{multicols}{2}
\begin{center}
Molecule A \\ 
r-2-chlorobutane\_rotated\_out\_G09
\includegraphics[width=8cm]{Images/ImagesFromVMD/r-2-chlorobutane_rotated_out_G09.jpg}
\\
Inertia Tensor - Molecule A \\
\vtab
\begin{tabular}{|c c c|}
113.621	 & 	-8.52391	 & 	-0.0867928	 \\
-8.52391	 & 	164.234	 & 	-0.885702	 \\
-0.0867928	 & 	-0.885702	 & 	257.008
\end{tabular}

\vtab
 EingenVectors - Molecule A     \\
\vtab
\begin{tabular}{|c c c|}
-0.986829	 & 	-0.161759	 & 	-0.00158112	 \\
0.161767	 & 	-0.986784	 & 	-0.00941025	 \\
-3.80279e-05	 & 	-0.00954208	 & 	0.999954
\end{tabular}

\vtab
 EingenValues - Molecule A     \\
\vtab
\begin{tabular}{|c c c|}
112.224	 & 	165.623	 & 	257.016
\end{tabular}
\columnbreak

Molecule B \\ 
s-2-chlorobutane\_rotated\_out\_G09
\includegraphics[width=8cm]{Images/ImagesFromVMD/s-2-chlorobutane_rotated_out_G09.jpg}
\\
Inertia Tensor - Molecule B \\
\vtab
\begin{tabular}{|c c c|}
113.621	 & 	8.52397	 & 	0.0867992	 \\
8.52397	 & 	164.235	 & 	-0.885706	 \\
0.0867992	 & 	-0.885706	 & 	257.008
\end{tabular}

\vtab
 EingenVectors - Molecule B     \\
\vtab
\begin{tabular}{|c c c|}
-0.98683	 & 	0.161756	 & 	0.00158113	 \\
-0.161763	 & 	-0.986785	 & 	-0.00941034	 \\
3.80628e-05	 & 	-0.00954217	 & 	0.999954
\end{tabular}

\vtab
 EingenValues - Molecule B     \\
\vtab
\begin{tabular}{|c c c|}
112.224	 & 	165.624	 & 	257.017
\end{tabular}

\end{center}
\end{multicols}
\begin{center}
\vtab
\vtab
\textcolor{NavyBlue}{\Large Are Isomer and Stereoisomers}
\end{center}
%%%%%%%%%%%%%%%%%%%%%%%%%%%%%%%%%%%%%%%%%%%%%%%%%%%%%%%%
%%%%%%%%%%%%%%%%%%%%%%%%%%%%%%%%%%%%%%%%%%%%%%%%%%%%%%%%
\newpage
%%%%%%%%%%%%%%%%%%%%%%%%%%%%%%%%%%%%%%%%%%%%%%%%%%%%%%%%
%%%%%%%%%%%%%%%%%%%%%%%%%%%%%%%%%%%%%%%%%%%%%%%%%%%%%%%%

\vtab[-2cm]
\begin{center}
{\large AsymmetricEnantiomers}
\end{center}
\begin{multicols}{2}
\begin{center}
Molecule A \\ 
r-2-chlorobutane\_rotated\_out\_G09
\includegraphics[width=8cm]{Images/ImagesFromVMD/r-2-chlorobutane_rotated_out_G09.jpg}
\\
Inertia Tensor - Molecule A \\
\vtab
\begin{tabular}{|c c c|}
113.621	 & 	-8.52391	 & 	-0.0867928	 \\
-8.52391	 & 	164.234	 & 	-0.885702	 \\
-0.0867928	 & 	-0.885702	 & 	257.008
\end{tabular}

\vtab
 EingenVectors - Molecule A     \\
\vtab
\begin{tabular}{|c c c|}
-0.986829	 & 	-0.161759	 & 	-0.00158112	 \\
0.161767	 & 	-0.986784	 & 	-0.00941025	 \\
-3.80279e-05	 & 	-0.00954208	 & 	0.999954
\end{tabular}

\vtab
 EingenValues - Molecule A     \\
\vtab
\begin{tabular}{|c c c|}
112.224	 & 	165.623	 & 	257.016
\end{tabular}
\columnbreak

Molecule B \\ 
s-2-chlorobutane\_tmp\_out\_G09
\includegraphics[width=8cm]{Images/ImagesFromVMD/s-2-chlorobutane_tmp_out_G09.jpg}
\\
Inertia Tensor - Molecule B \\
\vtab
\begin{tabular}{|c c c|}
113.621	 & 	8.52397	 & 	0.086812	 \\
8.52397	 & 	164.235	 & 	-0.885714	 \\
0.086812	 & 	-0.885714	 & 	257.008
\end{tabular}

\vtab
 EingenVectors - Molecule B     \\
\vtab
\begin{tabular}{|c c c|}
-0.98683	 & 	0.161756	 & 	0.00158123	 \\
-0.161763	 & 	-0.986785	 & 	-0.00941041	 \\
3.81472e-05	 & 	-0.00954225	 & 	0.999954
\end{tabular}

\vtab
 EingenValues - Molecule B     \\
\vtab
\begin{tabular}{|c c c|}
112.224	 & 	165.624	 & 	257.017
\end{tabular}

\end{center}
\end{multicols}
\begin{center}
\vtab
\vtab
\textcolor{NavyBlue}{\Large Are Isomer and Stereoisomers}
\end{center}
%%%%%%%%%%%%%%%%%%%%%%%%%%%%%%%%%%%%%%%%%%%%%%%%%%%%%%%%
%%%%%%%%%%%%%%%%%%%%%%%%%%%%%%%%%%%%%%%%%%%%%%%%%%%%%%%%
\newpage
%%%%%%%%%%%%%%%%%%%%%%%%%%%%%%%%%%%%%%%%%%%%%%%%%%%%%%%%
%%%%%%%%%%%%%%%%%%%%%%%%%%%%%%%%%%%%%%%%%%%%%%%%%%%%%%%%

\vtab[-2cm]
\begin{center}
{\large AsymmetricEnantiomers}
\end{center}
\begin{multicols}{2}
\begin{center}
Molecule A \\ 
r-2-chlorobutane\_tmp\_out\_G09
\includegraphics[width=8cm]{Images/ImagesFromVMD/r-2-chlorobutane_tmp_out_G09.jpg}
\\
Inertia Tensor - Molecule A \\
\vtab
\begin{tabular}{|c c c|}
113.621	 & 	-8.52391	 & 	-0.0867928	 \\
-8.52391	 & 	164.234	 & 	-0.885704	 \\
-0.0867928	 & 	-0.885704	 & 	257.008
\end{tabular}

\vtab
 EingenVectors - Molecule A     \\
\vtab
\begin{tabular}{|c c c|}
-0.986829	 & 	-0.161759	 & 	-0.00158112	 \\
0.161767	 & 	-0.986784	 & 	-0.00941027	 \\
-3.8027e-05	 & 	-0.0095421	 & 	0.999954
\end{tabular}

\vtab
 EingenValues - Molecule A     \\
\vtab
\begin{tabular}{|c c c|}
112.224	 & 	165.623	 & 	257.016
\end{tabular}
\columnbreak

Molecule B \\ 
s-2-chlorobutane\_out\_G09
\includegraphics[width=8cm]{Images/ImagesFromVMD/s-2-chlorobutane_out_G09.jpg}
\\
Inertia Tensor - Molecule B \\
\vtab
\begin{tabular}{|c c c|}
113.621	 & 	8.52396	 & 	0.0868155	 \\
8.52396	 & 	164.235	 & 	-0.885715	 \\
0.0868155	 & 	-0.885715	 & 	257.008
\end{tabular}

\vtab
 EingenVectors - Molecule B     \\
\vtab
\begin{tabular}{|c c c|}
-0.98683	 & 	0.161755	 & 	0.00158125	 \\
-0.161763	 & 	-0.986785	 & 	-0.00941041	 \\
3.81715e-05	 & 	-0.00954226	 & 	0.999954
\end{tabular}

\vtab
 EingenValues - Molecule B     \\
\vtab
\begin{tabular}{|c c c|}
112.224	 & 	165.624	 & 	257.017
\end{tabular}

\end{center}
\end{multicols}
\begin{center}
\vtab
\vtab
\textcolor{NavyBlue}{\Large Are Isomer and Stereoisomers}
\end{center}
%%%%%%%%%%%%%%%%%%%%%%%%%%%%%%%%%%%%%%%%%%%%%%%%%%%%%%%%
%%%%%%%%%%%%%%%%%%%%%%%%%%%%%%%%%%%%%%%%%%%%%%%%%%%%%%%%
\newpage
%%%%%%%%%%%%%%%%%%%%%%%%%%%%%%%%%%%%%%%%%%%%%%%%%%%%%%%%
%%%%%%%%%%%%%%%%%%%%%%%%%%%%%%%%%%%%%%%%%%%%%%%%%%%%%%%%

\vtab[-2cm]
\begin{center}
{\large AsymmetricEnantiomers}
\end{center}
\begin{multicols}{2}
\begin{center}
Molecule A \\ 
r-2-chlorobutane\_tmp\_out\_G09
\includegraphics[width=8cm]{Images/ImagesFromVMD/r-2-chlorobutane_tmp_out_G09.jpg}
\\
Inertia Tensor - Molecule A \\
\vtab
\begin{tabular}{|c c c|}
113.621	 & 	-8.52391	 & 	-0.0867928	 \\
-8.52391	 & 	164.234	 & 	-0.885704	 \\
-0.0867928	 & 	-0.885704	 & 	257.008
\end{tabular}

\vtab
 EingenVectors - Molecule A     \\
\vtab
\begin{tabular}{|c c c|}
-0.986829	 & 	-0.161759	 & 	-0.00158112	 \\
0.161767	 & 	-0.986784	 & 	-0.00941027	 \\
-3.8027e-05	 & 	-0.0095421	 & 	0.999954
\end{tabular}

\vtab
 EingenValues - Molecule A     \\
\vtab
\begin{tabular}{|c c c|}
112.224	 & 	165.623	 & 	257.016
\end{tabular}
\columnbreak

Molecule B \\ 
s-2-chlorobutane\_rotated\_out\_G09
\includegraphics[width=8cm]{Images/ImagesFromVMD/s-2-chlorobutane_rotated_out_G09.jpg}
\\
Inertia Tensor - Molecule B \\
\vtab
\begin{tabular}{|c c c|}
113.621	 & 	8.52397	 & 	0.0867992	 \\
8.52397	 & 	164.235	 & 	-0.885706	 \\
0.0867992	 & 	-0.885706	 & 	257.008
\end{tabular}

\vtab
 EingenVectors - Molecule B     \\
\vtab
\begin{tabular}{|c c c|}
-0.98683	 & 	0.161756	 & 	0.00158113	 \\
-0.161763	 & 	-0.986785	 & 	-0.00941034	 \\
3.80628e-05	 & 	-0.00954217	 & 	0.999954
\end{tabular}

\vtab
 EingenValues - Molecule B     \\
\vtab
\begin{tabular}{|c c c|}
112.224	 & 	165.624	 & 	257.017
\end{tabular}

\end{center}
\end{multicols}
\begin{center}
\vtab
\vtab
\textcolor{NavyBlue}{\Large Are Isomer and Stereoisomers}
\end{center}
%%%%%%%%%%%%%%%%%%%%%%%%%%%%%%%%%%%%%%%%%%%%%%%%%%%%%%%%
%%%%%%%%%%%%%%%%%%%%%%%%%%%%%%%%%%%%%%%%%%%%%%%%%%%%%%%%
\newpage
%%%%%%%%%%%%%%%%%%%%%%%%%%%%%%%%%%%%%%%%%%%%%%%%%%%%%%%%
%%%%%%%%%%%%%%%%%%%%%%%%%%%%%%%%%%%%%%%%%%%%%%%%%%%%%%%%

\vtab[-2cm]
\begin{center}
{\large AsymmetricEnantiomers}
\end{center}
\begin{multicols}{2}
\begin{center}
Molecule A \\ 
r-2-chlorobutane\_tmp\_out\_G09
\includegraphics[width=8cm]{Images/ImagesFromVMD/r-2-chlorobutane_tmp_out_G09.jpg}
\\
Inertia Tensor - Molecule A \\
\vtab
\begin{tabular}{|c c c|}
113.621	 & 	-8.52391	 & 	-0.0867928	 \\
-8.52391	 & 	164.234	 & 	-0.885704	 \\
-0.0867928	 & 	-0.885704	 & 	257.008
\end{tabular}

\vtab
 EingenVectors - Molecule A     \\
\vtab
\begin{tabular}{|c c c|}
-0.986829	 & 	-0.161759	 & 	-0.00158112	 \\
0.161767	 & 	-0.986784	 & 	-0.00941027	 \\
-3.8027e-05	 & 	-0.0095421	 & 	0.999954
\end{tabular}

\vtab
 EingenValues - Molecule A     \\
\vtab
\begin{tabular}{|c c c|}
112.224	 & 	165.623	 & 	257.016
\end{tabular}
\columnbreak

Molecule B \\ 
s-2-chlorobutane\_tmp\_out\_G09
\includegraphics[width=8cm]{Images/ImagesFromVMD/s-2-chlorobutane_tmp_out_G09.jpg}
\\
Inertia Tensor - Molecule B \\
\vtab
\begin{tabular}{|c c c|}
113.621	 & 	8.52397	 & 	0.086812	 \\
8.52397	 & 	164.235	 & 	-0.885714	 \\
0.086812	 & 	-0.885714	 & 	257.008
\end{tabular}

\vtab
 EingenVectors - Molecule B     \\
\vtab
\begin{tabular}{|c c c|}
-0.98683	 & 	0.161756	 & 	0.00158123	 \\
-0.161763	 & 	-0.986785	 & 	-0.00941041	 \\
3.81472e-05	 & 	-0.00954225	 & 	0.999954
\end{tabular}

\vtab
 EingenValues - Molecule B     \\
\vtab
\begin{tabular}{|c c c|}
112.224	 & 	165.624	 & 	257.017
\end{tabular}

\end{center}
\end{multicols}
\begin{center}
\vtab
\vtab
\textcolor{NavyBlue}{\Large Are Isomer and Stereoisomers}
\end{center}
%%%%%%%%%%%%%%%%%%%%%%%%%%%%%%%%%%%%%%%%%%%%%%%%%%%%%%%%
%%%%%%%%%%%%%%%%%%%%%%%%%%%%%%%%%%%%%%%%%%%%%%%%%%%%%%%%
\newpage
%%%%%%%%%%%%%%%%%%%%%%%%%%%%%%%%%%%%%%%%%%%%%%%%%%%%%%%%
%%%%%%%%%%%%%%%%%%%%%%%%%%%%%%%%%%%%%%%%%%%%%%%%%%%%%%%%

\vtab[-2cm]
\begin{center}
{\large AsymmetricEnantiomers}
\end{center}
\begin{multicols}{2}
\begin{center}
Molecule A \\ 
s-2-chlorobutane\_out\_G09
\includegraphics[width=8cm]{Images/ImagesFromVMD/s-2-chlorobutane_out_G09.jpg}
\\
Inertia Tensor - Molecule A \\
\vtab
\begin{tabular}{|c c c|}
113.621	 & 	8.52396	 & 	0.0868155	 \\
8.52396	 & 	164.235	 & 	-0.885715	 \\
0.0868155	 & 	-0.885715	 & 	257.008
\end{tabular}

\vtab
 EingenVectors - Molecule A     \\
\vtab
\begin{tabular}{|c c c|}
-0.98683	 & 	0.161755	 & 	0.00158125	 \\
-0.161763	 & 	-0.986785	 & 	-0.00941041	 \\
3.81715e-05	 & 	-0.00954226	 & 	0.999954
\end{tabular}

\vtab
 EingenValues - Molecule A     \\
\vtab
\begin{tabular}{|c c c|}
112.224	 & 	165.624	 & 	257.017
\end{tabular}
\columnbreak

Molecule B \\ 
s-2-chlorobutane\_rotated\_out\_G09
\includegraphics[width=8cm]{Images/ImagesFromVMD/s-2-chlorobutane_rotated_out_G09.jpg}
\\
Inertia Tensor - Molecule B \\
\vtab
\begin{tabular}{|c c c|}
113.621	 & 	8.52397	 & 	0.0867992	 \\
8.52397	 & 	164.235	 & 	-0.885706	 \\
0.0867992	 & 	-0.885706	 & 	257.008
\end{tabular}

\vtab
 EingenVectors - Molecule B     \\
\vtab
\begin{tabular}{|c c c|}
-0.98683	 & 	0.161756	 & 	0.00158113	 \\
-0.161763	 & 	-0.986785	 & 	-0.00941034	 \\
3.80628e-05	 & 	-0.00954217	 & 	0.999954
\end{tabular}

\vtab
 EingenValues - Molecule B     \\
\vtab
\begin{tabular}{|c c c|}
112.224	 & 	165.624	 & 	257.017
\end{tabular}

\end{center}
\end{multicols}
\begin{center}
\vtab
\vtab
\textcolor{NavyBlue}{\Large Are the same molecule}
\end{center}
%%%%%%%%%%%%%%%%%%%%%%%%%%%%%%%%%%%%%%%%%%%%%%%%%%%%%%%%
%%%%%%%%%%%%%%%%%%%%%%%%%%%%%%%%%%%%%%%%%%%%%%%%%%%%%%%%
\newpage
%%%%%%%%%%%%%%%%%%%%%%%%%%%%%%%%%%%%%%%%%%%%%%%%%%%%%%%%
%%%%%%%%%%%%%%%%%%%%%%%%%%%%%%%%%%%%%%%%%%%%%%%%%%%%%%%%

\vtab[-2cm]
\begin{center}
{\large AsymmetricEnantiomers}
\end{center}
\begin{multicols}{2}
\begin{center}
Molecule A \\ 
s-2-chlorobutane\_out\_G09
\includegraphics[width=8cm]{Images/ImagesFromVMD/s-2-chlorobutane_out_G09.jpg}
\\
Inertia Tensor - Molecule A \\
\vtab
\begin{tabular}{|c c c|}
113.621	 & 	8.52396	 & 	0.0868155	 \\
8.52396	 & 	164.235	 & 	-0.885715	 \\
0.0868155	 & 	-0.885715	 & 	257.008
\end{tabular}

\vtab
 EingenVectors - Molecule A     \\
\vtab
\begin{tabular}{|c c c|}
-0.98683	 & 	0.161755	 & 	0.00158125	 \\
-0.161763	 & 	-0.986785	 & 	-0.00941041	 \\
3.81715e-05	 & 	-0.00954226	 & 	0.999954
\end{tabular}

\vtab
 EingenValues - Molecule A     \\
\vtab
\begin{tabular}{|c c c|}
112.224	 & 	165.624	 & 	257.017
\end{tabular}
\columnbreak

Molecule B \\ 
s-2-chlorobutane\_tmp\_out\_G09
\includegraphics[width=8cm]{Images/ImagesFromVMD/s-2-chlorobutane_tmp_out_G09.jpg}
\\
Inertia Tensor - Molecule B \\
\vtab
\begin{tabular}{|c c c|}
113.621	 & 	8.52397	 & 	0.086812	 \\
8.52397	 & 	164.235	 & 	-0.885714	 \\
0.086812	 & 	-0.885714	 & 	257.008
\end{tabular}

\vtab
 EingenVectors - Molecule B     \\
\vtab
\begin{tabular}{|c c c|}
-0.98683	 & 	0.161756	 & 	0.00158123	 \\
-0.161763	 & 	-0.986785	 & 	-0.00941041	 \\
3.81472e-05	 & 	-0.00954225	 & 	0.999954
\end{tabular}

\vtab
 EingenValues - Molecule B     \\
\vtab
\begin{tabular}{|c c c|}
112.224	 & 	165.624	 & 	257.017
\end{tabular}

\end{center}
\end{multicols}
\begin{center}
\vtab
\vtab
\textcolor{NavyBlue}{\Large Are the same molecule}
\end{center}
%%%%%%%%%%%%%%%%%%%%%%%%%%%%%%%%%%%%%%%%%%%%%%%%%%%%%%%%
%%%%%%%%%%%%%%%%%%%%%%%%%%%%%%%%%%%%%%%%%%%%%%%%%%%%%%%%
\newpage
%%%%%%%%%%%%%%%%%%%%%%%%%%%%%%%%%%%%%%%%%%%%%%%%%%%%%%%%
%%%%%%%%%%%%%%%%%%%%%%%%%%%%%%%%%%%%%%%%%%%%%%%%%%%%%%%%

\vtab[-2cm]
\begin{center}
{\large AsymmetricEnantiomers}
\end{center}
\begin{multicols}{2}
\begin{center}
Molecule A \\ 
s-2-chlorobutane\_rotated\_out\_G09
\includegraphics[width=8cm]{Images/ImagesFromVMD/s-2-chlorobutane_rotated_out_G09.jpg}
\\
Inertia Tensor - Molecule A \\
\vtab
\begin{tabular}{|c c c|}
113.621	 & 	8.52397	 & 	0.0867992	 \\
8.52397	 & 	164.235	 & 	-0.885706	 \\
0.0867992	 & 	-0.885706	 & 	257.008
\end{tabular}

\vtab
 EingenVectors - Molecule A     \\
\vtab
\begin{tabular}{|c c c|}
-0.98683	 & 	0.161756	 & 	0.00158113	 \\
-0.161763	 & 	-0.986785	 & 	-0.00941034	 \\
3.80628e-05	 & 	-0.00954217	 & 	0.999954
\end{tabular}

\vtab
 EingenValues - Molecule A     \\
\vtab
\begin{tabular}{|c c c|}
112.224	 & 	165.624	 & 	257.017
\end{tabular}
\columnbreak

Molecule B \\ 
s-2-chlorobutane\_tmp\_out\_G09
\includegraphics[width=8cm]{Images/ImagesFromVMD/s-2-chlorobutane_tmp_out_G09.jpg}
\\
Inertia Tensor - Molecule B \\
\vtab
\begin{tabular}{|c c c|}
113.621	 & 	8.52397	 & 	0.086812	 \\
8.52397	 & 	164.235	 & 	-0.885714	 \\
0.086812	 & 	-0.885714	 & 	257.008
\end{tabular}

\vtab
 EingenVectors - Molecule B     \\
\vtab
\begin{tabular}{|c c c|}
-0.98683	 & 	0.161756	 & 	0.00158123	 \\
-0.161763	 & 	-0.986785	 & 	-0.00941041	 \\
3.81472e-05	 & 	-0.00954225	 & 	0.999954
\end{tabular}

\vtab
 EingenValues - Molecule B     \\
\vtab
\begin{tabular}{|c c c|}
112.224	 & 	165.624	 & 	257.017
\end{tabular}

\end{center}
\end{multicols}
\begin{center}
\vtab
\vtab
\textcolor{NavyBlue}{\Large Are the same molecule}
\end{center}
%%%%%%%%%%%%%%%%%%%%%%%%%%%%%%%%%%%%%%%%%%%%%%%%%%%%%%%%
%%%%%%%%%%%%%%%%%%%%%%%%%%%%%%%%%%%%%%%%%%%%%%%%%%%%%%%%
\newpage
%%%%%%%%%%%%%%%%%%%%%%%%%%%%%%%%%%%%%%%%%%%%%%%%%%%%%%%%
%%%%%%%%%%%%%%%%%%%%%%%%%%%%%%%%%%%%%%%%%%%%%%%%%%%%%%%%

\vtab[-2cm]
\begin{center}
{\large IsomerC6H14}
\end{center}
\begin{multicols}{2}
\begin{center}
Molecule A \\ 
2,2-dimethyl-butane\_out\_G09
\includegraphics[width=8cm]{Images/ImagesFromVMD/2,2-dimethyl-butane_out_G09.jpg}
\\
Inertia Tensor - Molecule A \\
\vtab
\begin{tabular}{|c c c|}
139.709	 & 	36.9725	 & 	-0	 \\
36.9725	 & 	183.699	 & 	-0	 \\
-0	 & 	-0	 & 	205.419
\end{tabular}

\vtab
 EingenVectors - Molecule A     \\
\vtab
\begin{tabular}{|c c c|}
-0.869274	 & 	0.494331	 & 	0	 \\
0.494331	 & 	0.869274	 & 	0	 \\
0	 & 	0	 & 	1
\end{tabular}

\vtab
 EingenValues - Molecule A     \\
\vtab
\begin{tabular}{|c c c|}
118.683	 & 	204.724	 & 	205.419
\end{tabular}
\columnbreak

Molecule B \\ 
2,3-dimethyl-butane\_out\_G09
\includegraphics[width=8cm]{Images/ImagesFromVMD/2,3-dimethyl-butane_out_G09.jpg}
\\
Inertia Tensor - Molecule B \\
\vtab
\begin{tabular}{|c c c|}
270.146	 & 	-16.9873	 & 	-0	 \\
-16.9873	 & 	122.541	 & 	-0	 \\
-0	 & 	-0	 & 	177.08
\end{tabular}

\vtab
 EingenVectors - Molecule B     \\
\vtab
\begin{tabular}{|c c c|}
-0.112876	 & 	-0.993609	 & 	0	 \\
0	 & 	0	 & 	1	 \\
0.993609	 & 	-0.112876	 & 	-0
\end{tabular}

\vtab
 EingenValues - Molecule B     \\
\vtab
\begin{tabular}{|c c c|}
120.611	 & 	177.08	 & 	272.075
\end{tabular}

\end{center}
\end{multicols}
\begin{center}
\vtab
\vtab
\textcolor{NavyBlue}{\Large Are Isomers}
\end{center}
%%%%%%%%%%%%%%%%%%%%%%%%%%%%%%%%%%%%%%%%%%%%%%%%%%%%%%%%
%%%%%%%%%%%%%%%%%%%%%%%%%%%%%%%%%%%%%%%%%%%%%%%%%%%%%%%%
\newpage
%%%%%%%%%%%%%%%%%%%%%%%%%%%%%%%%%%%%%%%%%%%%%%%%%%%%%%%%
%%%%%%%%%%%%%%%%%%%%%%%%%%%%%%%%%%%%%%%%%%%%%%%%%%%%%%%%

\vtab[-2cm]
\begin{center}
{\large IsomerC6H14}
\end{center}
\begin{multicols}{2}
\begin{center}
Molecule A \\ 
2,2-dimethyl-butane\_out\_G09
\includegraphics[width=8cm]{Images/ImagesFromVMD/2,2-dimethyl-butane_out_G09.jpg}
\\
Inertia Tensor - Molecule A \\
\vtab
\begin{tabular}{|c c c|}
139.709	 & 	36.9725	 & 	-0	 \\
36.9725	 & 	183.699	 & 	-0	 \\
-0	 & 	-0	 & 	205.419
\end{tabular}

\vtab
 EingenVectors - Molecule A     \\
\vtab
\begin{tabular}{|c c c|}
-0.869274	 & 	0.494331	 & 	0	 \\
0.494331	 & 	0.869274	 & 	0	 \\
0	 & 	0	 & 	1
\end{tabular}

\vtab
 EingenValues - Molecule A     \\
\vtab
\begin{tabular}{|c c c|}
118.683	 & 	204.724	 & 	205.419
\end{tabular}
\columnbreak

Molecule B \\ 
2-methyl-pentane\_out\_G09
\includegraphics[width=8cm]{Images/ImagesFromVMD/2-methyl-pentane_out_G09.jpg}
\\
Inertia Tensor - Molecule B \\
\vtab
\begin{tabular}{|c c c|}
76.7701	 & 	-0.234564	 & 	-0.159008	 \\
-0.234564	 & 	299.061	 & 	-0.102167	 \\
-0.159008	 & 	-0.102167	 & 	347.557
\end{tabular}

\vtab
 EingenVectors - Molecule B     \\
\vtab
\begin{tabular}{|c c c|}
-0.999999	 & 	-0.00105548	 & 	-0.000587604	 \\
-0.00105671	 & 	0.999997	 & 	0.00210321	 \\
-0.000585382	 & 	-0.00210383	 & 	0.999998
\end{tabular}

\vtab
 EingenValues - Molecule B     \\
\vtab
\begin{tabular}{|c c c|}
76.7698	 & 	299.061	 & 	347.558
\end{tabular}

\end{center}
\end{multicols}
\begin{center}
\vtab
\vtab
\textcolor{NavyBlue}{\Large Are Isomers}
\end{center}
%%%%%%%%%%%%%%%%%%%%%%%%%%%%%%%%%%%%%%%%%%%%%%%%%%%%%%%%
%%%%%%%%%%%%%%%%%%%%%%%%%%%%%%%%%%%%%%%%%%%%%%%%%%%%%%%%
\newpage
%%%%%%%%%%%%%%%%%%%%%%%%%%%%%%%%%%%%%%%%%%%%%%%%%%%%%%%%
%%%%%%%%%%%%%%%%%%%%%%%%%%%%%%%%%%%%%%%%%%%%%%%%%%%%%%%%

\vtab[-2cm]
\begin{center}
{\large IsomerC6H14}
\end{center}
\begin{multicols}{2}
\begin{center}
Molecule A \\ 
2,2-dimethyl-butane\_out\_G09
\includegraphics[width=8cm]{Images/ImagesFromVMD/2,2-dimethyl-butane_out_G09.jpg}
\\
Inertia Tensor - Molecule A \\
\vtab
\begin{tabular}{|c c c|}
139.709	 & 	36.9725	 & 	-0	 \\
36.9725	 & 	183.699	 & 	-0	 \\
-0	 & 	-0	 & 	205.419
\end{tabular}

\vtab
 EingenVectors - Molecule A     \\
\vtab
\begin{tabular}{|c c c|}
-0.869274	 & 	0.494331	 & 	0	 \\
0.494331	 & 	0.869274	 & 	0	 \\
0	 & 	0	 & 	1
\end{tabular}

\vtab
 EingenValues - Molecule A     \\
\vtab
\begin{tabular}{|c c c|}
118.683	 & 	204.724	 & 	205.419
\end{tabular}
\columnbreak

Molecule B \\ 
3-methyl-pentane\_out\_G09
\includegraphics[width=8cm]{Images/ImagesFromVMD/3-methyl-pentane_out_G09.jpg}
\\
Inertia Tensor - Molecule B \\
\vtab
\begin{tabular}{|c c c|}
294.318	 & 	21.9352	 & 	-0	 \\
21.9352	 & 	302.288	 & 	-0	 \\
-0	 & 	-0	 & 	77.8079
\end{tabular}

\vtab
 EingenVectors - Molecule B     \\
\vtab
\begin{tabular}{|c c c|}
0	 & 	0	 & 	1	 \\
-0.767709	 & 	0.640799	 & 	0	 \\
0.640799	 & 	0.767709	 & 	0
\end{tabular}

\vtab
 EingenValues - Molecule B     \\
\vtab
\begin{tabular}{|c c c|}
77.8079	 & 	276.009	 & 	320.597
\end{tabular}

\end{center}
\end{multicols}
\begin{center}
\vtab
\vtab
\textcolor{NavyBlue}{\Large Are Isomers}
\end{center}
%%%%%%%%%%%%%%%%%%%%%%%%%%%%%%%%%%%%%%%%%%%%%%%%%%%%%%%%
%%%%%%%%%%%%%%%%%%%%%%%%%%%%%%%%%%%%%%%%%%%%%%%%%%%%%%%%
\newpage
%%%%%%%%%%%%%%%%%%%%%%%%%%%%%%%%%%%%%%%%%%%%%%%%%%%%%%%%
%%%%%%%%%%%%%%%%%%%%%%%%%%%%%%%%%%%%%%%%%%%%%%%%%%%%%%%%

\vtab[-2cm]
\begin{center}
{\large IsomerC6H14}
\end{center}
\begin{multicols}{2}
\begin{center}
Molecule A \\ 
2,2-dimethyl-butane\_out\_G09
\includegraphics[width=8cm]{Images/ImagesFromVMD/2,2-dimethyl-butane_out_G09.jpg}
\\
Inertia Tensor - Molecule A \\
\vtab
\begin{tabular}{|c c c|}
139.709	 & 	36.9725	 & 	-0	 \\
36.9725	 & 	183.699	 & 	-0	 \\
-0	 & 	-0	 & 	205.419
\end{tabular}

\vtab
 EingenVectors - Molecule A     \\
\vtab
\begin{tabular}{|c c c|}
-0.869274	 & 	0.494331	 & 	0	 \\
0.494331	 & 	0.869274	 & 	0	 \\
0	 & 	0	 & 	1
\end{tabular}

\vtab
 EingenValues - Molecule A     \\
\vtab
\begin{tabular}{|c c c|}
118.683	 & 	204.724	 & 	205.419
\end{tabular}
\columnbreak

Molecule B \\ 
hexane\_out\_G09
\includegraphics[width=8cm]{Images/ImagesFromVMD/hexane_out_G09.jpg}
\\
Inertia Tensor - Molecule B \\
\vtab
\begin{tabular}{|c c c|}
346.244	 & 	-179.129	 & 	-0	 \\
-179.129	 & 	137.676	 & 	-0	 \\
-0	 & 	-0	 & 	465.177
\end{tabular}

\vtab
 EingenVectors - Molecule B     \\
\vtab
\begin{tabular}{|c c c|}
-0.498436	 & 	-0.866927	 & 	0	 \\
0.866927	 & 	-0.498436	 & 	-0	 \\
0	 & 	0	 & 	1
\end{tabular}

\vtab
 EingenValues - Molecule B     \\
\vtab
\begin{tabular}{|c c c|}
34.6869	 & 	449.234	 & 	465.177
\end{tabular}

\end{center}
\end{multicols}
\begin{center}
\vtab
\vtab
\textcolor{NavyBlue}{\Large Are Isomers}
\end{center}
%%%%%%%%%%%%%%%%%%%%%%%%%%%%%%%%%%%%%%%%%%%%%%%%%%%%%%%%
%%%%%%%%%%%%%%%%%%%%%%%%%%%%%%%%%%%%%%%%%%%%%%%%%%%%%%%%
\newpage
%%%%%%%%%%%%%%%%%%%%%%%%%%%%%%%%%%%%%%%%%%%%%%%%%%%%%%%%
%%%%%%%%%%%%%%%%%%%%%%%%%%%%%%%%%%%%%%%%%%%%%%%%%%%%%%%%

\vtab[-2cm]
\begin{center}
{\large IsomerC6H14}
\end{center}
\begin{multicols}{2}
\begin{center}
Molecule A \\ 
2,3-dimethyl-butane\_out\_G09
\includegraphics[width=8cm]{Images/ImagesFromVMD/2,3-dimethyl-butane_out_G09.jpg}
\\
Inertia Tensor - Molecule A \\
\vtab
\begin{tabular}{|c c c|}
270.146	 & 	-16.9873	 & 	-0	 \\
-16.9873	 & 	122.541	 & 	-0	 \\
-0	 & 	-0	 & 	177.08
\end{tabular}

\vtab
 EingenVectors - Molecule A     \\
\vtab
\begin{tabular}{|c c c|}
-0.112876	 & 	-0.993609	 & 	0	 \\
0	 & 	0	 & 	1	 \\
0.993609	 & 	-0.112876	 & 	-0
\end{tabular}

\vtab
 EingenValues - Molecule A     \\
\vtab
\begin{tabular}{|c c c|}
120.611	 & 	177.08	 & 	272.075
\end{tabular}
\columnbreak

Molecule B \\ 
2-methyl-pentane\_out\_G09
\includegraphics[width=8cm]{Images/ImagesFromVMD/2-methyl-pentane_out_G09.jpg}
\\
Inertia Tensor - Molecule B \\
\vtab
\begin{tabular}{|c c c|}
76.7701	 & 	-0.234564	 & 	-0.159008	 \\
-0.234564	 & 	299.061	 & 	-0.102167	 \\
-0.159008	 & 	-0.102167	 & 	347.557
\end{tabular}

\vtab
 EingenVectors - Molecule B     \\
\vtab
\begin{tabular}{|c c c|}
-0.999999	 & 	-0.00105548	 & 	-0.000587604	 \\
-0.00105671	 & 	0.999997	 & 	0.00210321	 \\
-0.000585382	 & 	-0.00210383	 & 	0.999998
\end{tabular}

\vtab
 EingenValues - Molecule B     \\
\vtab
\begin{tabular}{|c c c|}
76.7698	 & 	299.061	 & 	347.558
\end{tabular}

\end{center}
\end{multicols}
\begin{center}
\vtab
\vtab
\textcolor{NavyBlue}{\Large Are Isomers}
\end{center}
%%%%%%%%%%%%%%%%%%%%%%%%%%%%%%%%%%%%%%%%%%%%%%%%%%%%%%%%
%%%%%%%%%%%%%%%%%%%%%%%%%%%%%%%%%%%%%%%%%%%%%%%%%%%%%%%%
\newpage
%%%%%%%%%%%%%%%%%%%%%%%%%%%%%%%%%%%%%%%%%%%%%%%%%%%%%%%%
%%%%%%%%%%%%%%%%%%%%%%%%%%%%%%%%%%%%%%%%%%%%%%%%%%%%%%%%

\vtab[-2cm]
\begin{center}
{\large IsomerC6H14}
\end{center}
\begin{multicols}{2}
\begin{center}
Molecule A \\ 
2,3-dimethyl-butane\_out\_G09
\includegraphics[width=8cm]{Images/ImagesFromVMD/2,3-dimethyl-butane_out_G09.jpg}
\\
Inertia Tensor - Molecule A \\
\vtab
\begin{tabular}{|c c c|}
270.146	 & 	-16.9873	 & 	-0	 \\
-16.9873	 & 	122.541	 & 	-0	 \\
-0	 & 	-0	 & 	177.08
\end{tabular}

\vtab
 EingenVectors - Molecule A     \\
\vtab
\begin{tabular}{|c c c|}
-0.112876	 & 	-0.993609	 & 	0	 \\
0	 & 	0	 & 	1	 \\
0.993609	 & 	-0.112876	 & 	-0
\end{tabular}

\vtab
 EingenValues - Molecule A     \\
\vtab
\begin{tabular}{|c c c|}
120.611	 & 	177.08	 & 	272.075
\end{tabular}
\columnbreak

Molecule B \\ 
3-methyl-pentane\_out\_G09
\includegraphics[width=8cm]{Images/ImagesFromVMD/3-methyl-pentane_out_G09.jpg}
\\
Inertia Tensor - Molecule B \\
\vtab
\begin{tabular}{|c c c|}
294.318	 & 	21.9352	 & 	-0	 \\
21.9352	 & 	302.288	 & 	-0	 \\
-0	 & 	-0	 & 	77.8079
\end{tabular}

\vtab
 EingenVectors - Molecule B     \\
\vtab
\begin{tabular}{|c c c|}
0	 & 	0	 & 	1	 \\
-0.767709	 & 	0.640799	 & 	0	 \\
0.640799	 & 	0.767709	 & 	0
\end{tabular}

\vtab
 EingenValues - Molecule B     \\
\vtab
\begin{tabular}{|c c c|}
77.8079	 & 	276.009	 & 	320.597
\end{tabular}

\end{center}
\end{multicols}
\begin{center}
\vtab
\vtab
\textcolor{NavyBlue}{\Large Are Isomers}
\end{center}
%%%%%%%%%%%%%%%%%%%%%%%%%%%%%%%%%%%%%%%%%%%%%%%%%%%%%%%%
%%%%%%%%%%%%%%%%%%%%%%%%%%%%%%%%%%%%%%%%%%%%%%%%%%%%%%%%
\newpage
%%%%%%%%%%%%%%%%%%%%%%%%%%%%%%%%%%%%%%%%%%%%%%%%%%%%%%%%
%%%%%%%%%%%%%%%%%%%%%%%%%%%%%%%%%%%%%%%%%%%%%%%%%%%%%%%%

\vtab[-2cm]
\begin{center}
{\large IsomerC6H14}
\end{center}
\begin{multicols}{2}
\begin{center}
Molecule A \\ 
2,3-dimethyl-butane\_out\_G09
\includegraphics[width=8cm]{Images/ImagesFromVMD/2,3-dimethyl-butane_out_G09.jpg}
\\
Inertia Tensor - Molecule A \\
\vtab
\begin{tabular}{|c c c|}
270.146	 & 	-16.9873	 & 	-0	 \\
-16.9873	 & 	122.541	 & 	-0	 \\
-0	 & 	-0	 & 	177.08
\end{tabular}

\vtab
 EingenVectors - Molecule A     \\
\vtab
\begin{tabular}{|c c c|}
-0.112876	 & 	-0.993609	 & 	0	 \\
0	 & 	0	 & 	1	 \\
0.993609	 & 	-0.112876	 & 	-0
\end{tabular}

\vtab
 EingenValues - Molecule A     \\
\vtab
\begin{tabular}{|c c c|}
120.611	 & 	177.08	 & 	272.075
\end{tabular}
\columnbreak

Molecule B \\ 
hexane\_out\_G09
\includegraphics[width=8cm]{Images/ImagesFromVMD/hexane_out_G09.jpg}
\\
Inertia Tensor - Molecule B \\
\vtab
\begin{tabular}{|c c c|}
346.244	 & 	-179.129	 & 	-0	 \\
-179.129	 & 	137.676	 & 	-0	 \\
-0	 & 	-0	 & 	465.177
\end{tabular}

\vtab
 EingenVectors - Molecule B     \\
\vtab
\begin{tabular}{|c c c|}
-0.498436	 & 	-0.866927	 & 	0	 \\
0.866927	 & 	-0.498436	 & 	-0	 \\
0	 & 	0	 & 	1
\end{tabular}

\vtab
 EingenValues - Molecule B     \\
\vtab
\begin{tabular}{|c c c|}
34.6869	 & 	449.234	 & 	465.177
\end{tabular}

\end{center}
\end{multicols}
\begin{center}
\vtab
\vtab
\textcolor{NavyBlue}{\Large Are Isomers}
\end{center}
%%%%%%%%%%%%%%%%%%%%%%%%%%%%%%%%%%%%%%%%%%%%%%%%%%%%%%%%
%%%%%%%%%%%%%%%%%%%%%%%%%%%%%%%%%%%%%%%%%%%%%%%%%%%%%%%%
\newpage
%%%%%%%%%%%%%%%%%%%%%%%%%%%%%%%%%%%%%%%%%%%%%%%%%%%%%%%%
%%%%%%%%%%%%%%%%%%%%%%%%%%%%%%%%%%%%%%%%%%%%%%%%%%%%%%%%

\vtab[-2cm]
\begin{center}
{\large IsomerC6H14}
\end{center}
\begin{multicols}{2}
\begin{center}
Molecule A \\ 
2-methyl-pentane\_out\_G09
\includegraphics[width=8cm]{Images/ImagesFromVMD/2-methyl-pentane_out_G09.jpg}
\\
Inertia Tensor - Molecule A \\
\vtab
\begin{tabular}{|c c c|}
76.7701	 & 	-0.234564	 & 	-0.159008	 \\
-0.234564	 & 	299.061	 & 	-0.102167	 \\
-0.159008	 & 	-0.102167	 & 	347.557
\end{tabular}

\vtab
 EingenVectors - Molecule A     \\
\vtab
\begin{tabular}{|c c c|}
-0.999999	 & 	-0.00105548	 & 	-0.000587604	 \\
-0.00105671	 & 	0.999997	 & 	0.00210321	 \\
-0.000585382	 & 	-0.00210383	 & 	0.999998
\end{tabular}

\vtab
 EingenValues - Molecule A     \\
\vtab
\begin{tabular}{|c c c|}
76.7698	 & 	299.061	 & 	347.558
\end{tabular}
\columnbreak

Molecule B \\ 
3-methyl-pentane\_out\_G09
\includegraphics[width=8cm]{Images/ImagesFromVMD/3-methyl-pentane_out_G09.jpg}
\\
Inertia Tensor - Molecule B \\
\vtab
\begin{tabular}{|c c c|}
294.318	 & 	21.9352	 & 	-0	 \\
21.9352	 & 	302.288	 & 	-0	 \\
-0	 & 	-0	 & 	77.8079
\end{tabular}

\vtab
 EingenVectors - Molecule B     \\
\vtab
\begin{tabular}{|c c c|}
0	 & 	0	 & 	1	 \\
-0.767709	 & 	0.640799	 & 	0	 \\
0.640799	 & 	0.767709	 & 	0
\end{tabular}

\vtab
 EingenValues - Molecule B     \\
\vtab
\begin{tabular}{|c c c|}
77.8079	 & 	276.009	 & 	320.597
\end{tabular}

\end{center}
\end{multicols}
\begin{center}
\vtab
\vtab
\textcolor{NavyBlue}{\Large Are Isomers}
\end{center}
%%%%%%%%%%%%%%%%%%%%%%%%%%%%%%%%%%%%%%%%%%%%%%%%%%%%%%%%
%%%%%%%%%%%%%%%%%%%%%%%%%%%%%%%%%%%%%%%%%%%%%%%%%%%%%%%%
\newpage
%%%%%%%%%%%%%%%%%%%%%%%%%%%%%%%%%%%%%%%%%%%%%%%%%%%%%%%%
%%%%%%%%%%%%%%%%%%%%%%%%%%%%%%%%%%%%%%%%%%%%%%%%%%%%%%%%

\vtab[-2cm]
\begin{center}
{\large IsomerC6H14}
\end{center}
\begin{multicols}{2}
\begin{center}
Molecule A \\ 
2-methyl-pentane\_out\_G09
\includegraphics[width=8cm]{Images/ImagesFromVMD/2-methyl-pentane_out_G09.jpg}
\\
Inertia Tensor - Molecule A \\
\vtab
\begin{tabular}{|c c c|}
76.7701	 & 	-0.234564	 & 	-0.159008	 \\
-0.234564	 & 	299.061	 & 	-0.102167	 \\
-0.159008	 & 	-0.102167	 & 	347.557
\end{tabular}

\vtab
 EingenVectors - Molecule A     \\
\vtab
\begin{tabular}{|c c c|}
-0.999999	 & 	-0.00105548	 & 	-0.000587604	 \\
-0.00105671	 & 	0.999997	 & 	0.00210321	 \\
-0.000585382	 & 	-0.00210383	 & 	0.999998
\end{tabular}

\vtab
 EingenValues - Molecule A     \\
\vtab
\begin{tabular}{|c c c|}
76.7698	 & 	299.061	 & 	347.558
\end{tabular}
\columnbreak

Molecule B \\ 
hexane\_out\_G09
\includegraphics[width=8cm]{Images/ImagesFromVMD/hexane_out_G09.jpg}
\\
Inertia Tensor - Molecule B \\
\vtab
\begin{tabular}{|c c c|}
346.244	 & 	-179.129	 & 	-0	 \\
-179.129	 & 	137.676	 & 	-0	 \\
-0	 & 	-0	 & 	465.177
\end{tabular}

\vtab
 EingenVectors - Molecule B     \\
\vtab
\begin{tabular}{|c c c|}
-0.498436	 & 	-0.866927	 & 	0	 \\
0.866927	 & 	-0.498436	 & 	-0	 \\
0	 & 	0	 & 	1
\end{tabular}

\vtab
 EingenValues - Molecule B     \\
\vtab
\begin{tabular}{|c c c|}
34.6869	 & 	449.234	 & 	465.177
\end{tabular}

\end{center}
\end{multicols}
\begin{center}
\vtab
\vtab
\textcolor{NavyBlue}{\Large Are Isomers}
\end{center}
%%%%%%%%%%%%%%%%%%%%%%%%%%%%%%%%%%%%%%%%%%%%%%%%%%%%%%%%
%%%%%%%%%%%%%%%%%%%%%%%%%%%%%%%%%%%%%%%%%%%%%%%%%%%%%%%%
\newpage
%%%%%%%%%%%%%%%%%%%%%%%%%%%%%%%%%%%%%%%%%%%%%%%%%%%%%%%%
%%%%%%%%%%%%%%%%%%%%%%%%%%%%%%%%%%%%%%%%%%%%%%%%%%%%%%%%

\vtab[-2cm]
\begin{center}
{\large IsomerC6H14}
\end{center}
\begin{multicols}{2}
\begin{center}
Molecule A \\ 
3-methyl-pentane\_out\_G09
\includegraphics[width=8cm]{Images/ImagesFromVMD/3-methyl-pentane_out_G09.jpg}
\\
Inertia Tensor - Molecule A \\
\vtab
\begin{tabular}{|c c c|}
294.318	 & 	21.9352	 & 	-0	 \\
21.9352	 & 	302.288	 & 	-0	 \\
-0	 & 	-0	 & 	77.8079
\end{tabular}

\vtab
 EingenVectors - Molecule A     \\
\vtab
\begin{tabular}{|c c c|}
0	 & 	0	 & 	1	 \\
-0.767709	 & 	0.640799	 & 	0	 \\
0.640799	 & 	0.767709	 & 	0
\end{tabular}

\vtab
 EingenValues - Molecule A     \\
\vtab
\begin{tabular}{|c c c|}
77.8079	 & 	276.009	 & 	320.597
\end{tabular}
\columnbreak

Molecule B \\ 
hexane\_out\_G09
\includegraphics[width=8cm]{Images/ImagesFromVMD/hexane_out_G09.jpg}
\\
Inertia Tensor - Molecule B \\
\vtab
\begin{tabular}{|c c c|}
346.244	 & 	-179.129	 & 	-0	 \\
-179.129	 & 	137.676	 & 	-0	 \\
-0	 & 	-0	 & 	465.177
\end{tabular}

\vtab
 EingenVectors - Molecule B     \\
\vtab
\begin{tabular}{|c c c|}
-0.498436	 & 	-0.866927	 & 	0	 \\
0.866927	 & 	-0.498436	 & 	-0	 \\
0	 & 	0	 & 	1
\end{tabular}

\vtab
 EingenValues - Molecule B     \\
\vtab
\begin{tabular}{|c c c|}
34.6869	 & 	449.234	 & 	465.177
\end{tabular}

\end{center}
\end{multicols}
\begin{center}
\vtab
\vtab
\textcolor{NavyBlue}{\Large Are Isomers}
\end{center}
%%%%%%%%%%%%%%%%%%%%%%%%%%%%%%%%%%%%%%%%%%%%%%%%%%%%%%%%
%%%%%%%%%%%%%%%%%%%%%%%%%%%%%%%%%%%%%%%%%%%%%%%%%%%%%%%%
\newpage
%%%%%%%%%%%%%%%%%%%%%%%%%%%%%%%%%%%%%%%%%%%%%%%%%%%%%%%%
%%%%%%%%%%%%%%%%%%%%%%%%%%%%%%%%%%%%%%%%%%%%%%%%%%%%%%%%
{\Large NeoConformation}
\begin{multicols}{2}
\begin{center}
Molecule A \\ 
neopentane\_out\_G09
\includegraphics[width=8cm]{Images/ImagesFromVMD/neopentane_out_G09.jpg}
\vtab

\columnbreak

Molecule B \\ 
tert-butylamine\_out\_G09
\includegraphics[width=8cm]{Images/ImagesFromVMD/tert-butylamine_out_G09.jpg}
\vtab
\end{center}
\end{multicols}
\begin{center}
\vtab
\vtab
\textcolor{NavyBlue}{\large Are different molecules}
\end{center}
%%%%%%%%%%%%%%%%%%%%%%%%%%%%%%%%%%%%%%%%%%%%%%%%%%%%%%%%
%%%%%%%%%%%%%%%%%%%%%%%%%%%%%%%%%%%%%%%%%%%%%%%%%%%%%%%%
\newpage
%%%%%%%%%%%%%%%%%%%%%%%%%%%%%%%%%%%%%%%%%%%%%%%%%%%%%%%%
%%%%%%%%%%%%%%%%%%%%%%%%%%%%%%%%%%%%%%%%%%%%%%%%%%%%%%%%
{\Large NeoConformation}
\begin{multicols}{2}
\begin{center}
Molecule A \\ 
neopentane\_out\_G09
\includegraphics[width=8cm]{Images/ImagesFromVMD/neopentane_out_G09.jpg}
\vtab

\columnbreak

Molecule B \\ 
tetramethylsilane\_out\_G09
\includegraphics[width=8cm]{Images/ImagesFromVMD/tetramethylsilane_out_G09.jpg}
\vtab
\end{center}
\end{multicols}
\begin{center}
\vtab
\vtab
\textcolor{NavyBlue}{\large Are different molecules}
\end{center}
%%%%%%%%%%%%%%%%%%%%%%%%%%%%%%%%%%%%%%%%%%%%%%%%%%%%%%%%
%%%%%%%%%%%%%%%%%%%%%%%%%%%%%%%%%%%%%%%%%%%%%%%%%%%%%%%%
\newpage
%%%%%%%%%%%%%%%%%%%%%%%%%%%%%%%%%%%%%%%%%%%%%%%%%%%%%%%%
%%%%%%%%%%%%%%%%%%%%%%%%%%%%%%%%%%%%%%%%%%%%%%%%%%%%%%%%
{\Large NeoConformation}
\begin{multicols}{2}
\begin{center}
Molecule A \\ 
tert-butylamine\_out\_G09
\includegraphics[width=8cm]{Images/ImagesFromVMD/tert-butylamine_out_G09.jpg}
\vtab

\columnbreak

Molecule B \\ 
tetramethylsilane\_out\_G09
\includegraphics[width=8cm]{Images/ImagesFromVMD/tetramethylsilane_out_G09.jpg}
\vtab
\end{center}
\end{multicols}
\begin{center}
\vtab
\vtab
\textcolor{NavyBlue}{\large Are different molecules}
\end{center}
%%%%%%%%%%%%%%%%%%%%%%%%%%%%%%%%%%%%%%%%%%%%%%%%%%%%%%%%
%%%%%%%%%%%%%%%%%%%%%%%%%%%%%%%%%%%%%%%%%%%%%%%%%%%%%%%%
\newpage
%%%%%%%%%%%%%%%%%%%%%%%%%%%%%%%%%%%%%%%%%%%%%%%%%%%%%%%%
%%%%%%%%%%%%%%%%%%%%%%%%%%%%%%%%%%%%%%%%%%%%%%%%%%%%%%%%
{\Large RotationSameMolecule}
\begin{multicols}{2}
\begin{center}
Molecule A \\ 
Porphin\_out\_G09
\includegraphics[width=8cm]{Images/ImagesFromVMD/Porphin_out_G09.jpg}
\vtab

\columnbreak

Molecule B \\ 
Tetrabenzoporphyrin\_out\_G09
\includegraphics[width=8cm]{Images/ImagesFromVMD/Tetrabenzoporphyrin_out_G09.jpg}
\vtab
\end{center}
\end{multicols}
\begin{center}
\vtab
\vtab
\textcolor{NavyBlue}{\large Are different molecules}
\end{center}
%%%%%%%%%%%%%%%%%%%%%%%%%%%%%%%%%%%%%%%%%%%%%%%%%%%%%%%%
%%%%%%%%%%%%%%%%%%%%%%%%%%%%%%%%%%%%%%%%%%%%%%%%%%%%%%%%
\newpage
%%%%%%%%%%%%%%%%%%%%%%%%%%%%%%%%%%%%%%%%%%%%%%%%%%%%%%%%
%%%%%%%%%%%%%%%%%%%%%%%%%%%%%%%%%%%%%%%%%%%%%%%%%%%%%%%%
{\Large RotationSameMolecule}
\begin{multicols}{2}
\begin{center}
Molecule A \\ 
Porphin\_out\_G09
\includegraphics[width=8cm]{Images/ImagesFromVMD/Porphin_out_G09.jpg}
\vtab

\columnbreak

Molecule B \\ 
Tetrabenzoporphyrin\_rotated02\_out\_G09
\includegraphics[width=8cm]{Images/ImagesFromVMD/Tetrabenzoporphyrin_rotated02_out_G09.jpg}
\vtab
\end{center}
\end{multicols}
\begin{center}
\vtab
\vtab
\textcolor{NavyBlue}{\large Are different molecules}
\end{center}
%%%%%%%%%%%%%%%%%%%%%%%%%%%%%%%%%%%%%%%%%%%%%%%%%%%%%%%%
%%%%%%%%%%%%%%%%%%%%%%%%%%%%%%%%%%%%%%%%%%%%%%%%%%%%%%%%
\newpage
%%%%%%%%%%%%%%%%%%%%%%%%%%%%%%%%%%%%%%%%%%%%%%%%%%%%%%%%
%%%%%%%%%%%%%%%%%%%%%%%%%%%%%%%%%%%%%%%%%%%%%%%%%%%%%%%%
{\Large RotationSameMolecule}
\begin{multicols}{2}
\begin{center}
Molecule A \\ 
Porphin\_out\_G09
\includegraphics[width=8cm]{Images/ImagesFromVMD/Porphin_out_G09.jpg}
\vtab

\columnbreak

Molecule B \\ 
Tetrabenzoporphyrin\_rotated\_out\_G09
\includegraphics[width=8cm]{Images/ImagesFromVMD/Tetrabenzoporphyrin_rotated_out_G09.jpg}
\vtab
\end{center}
\end{multicols}
\begin{center}
\vtab
\vtab
\textcolor{NavyBlue}{\large Are different molecules}
\end{center}
%%%%%%%%%%%%%%%%%%%%%%%%%%%%%%%%%%%%%%%%%%%%%%%%%%%%%%%%
%%%%%%%%%%%%%%%%%%%%%%%%%%%%%%%%%%%%%%%%%%%%%%%%%%%%%%%%
\newpage
%%%%%%%%%%%%%%%%%%%%%%%%%%%%%%%%%%%%%%%%%%%%%%%%%%%%%%%%
%%%%%%%%%%%%%%%%%%%%%%%%%%%%%%%%%%%%%%%%%%%%%%%%%%%%%%%%

\vtab[-2cm]
\begin{center}
{\large RotationSameMolecule}
\end{center}
\begin{multicols}{2}
\begin{center}
Molecule A \\ 
Tetrabenzoporphyrin\_out\_G09
\includegraphics[width=8cm]{Images/ImagesFromVMD/Tetrabenzoporphyrin_out_G09.jpg}
\\
Inertia Tensor - Molecule A \\
\vtab
\begin{tabular}{|c c c|}
11864.8	 & 	-0	 & 	-0	 \\
-0	 & 	5941.95	 & 	-0	 \\
-0	 & 	-0	 & 	5922.88
\end{tabular}

\vtab
 EingenVectors - Molecule A     \\
\vtab
\begin{tabular}{|c c c|}
0	 & 	0	 & 	1	 \\
0	 & 	1	 & 	0	 \\
1	 & 	0	 & 	0
\end{tabular}

\vtab
 EingenValues - Molecule A     \\
\vtab
\begin{tabular}{|c c c|}
5922.88	 & 	5941.95	 & 	11864.8
\end{tabular}
\columnbreak

Molecule B \\ 
Tetrabenzoporphyrin\_rotated02\_out\_G09
\includegraphics[width=8cm]{Images/ImagesFromVMD/Tetrabenzoporphyrin_rotated02_out_G09.jpg}
\\
Inertia Tensor - Molecule B \\
\vtab
\begin{tabular}{|c c c|}
7791.57	 & 	2641.52	 & 	-788.168	 \\
2641.52	 & 	9673.67	 & 	-1103.82	 \\
-788.168	 & 	-1103.82	 & 	6265.27
\end{tabular}

\vtab
 EingenVectors - Molecule B     \\
\vtab
\begin{tabular}{|c c c|}
-0.771476	 & 	0.397136	 & 	-0.497099	 \\
0.301236	 & 	-0.460189	 & 	-0.835154	 \\
-0.56043	 & 	-0.794046	 & 	0.235394
\end{tabular}

\vtab
 EingenValues - Molecule B     \\
\vtab
\begin{tabular}{|c c c|}
5923.93	 & 	5941.33	 & 	11865.3
\end{tabular}

\end{center}
\end{multicols}
\begin{center}
\vtab
\vtab
\textcolor{NavyBlue}{\Large Are Isomers}
\end{center}
%%%%%%%%%%%%%%%%%%%%%%%%%%%%%%%%%%%%%%%%%%%%%%%%%%%%%%%%
%%%%%%%%%%%%%%%%%%%%%%%%%%%%%%%%%%%%%%%%%%%%%%%%%%%%%%%%
\newpage
%%%%%%%%%%%%%%%%%%%%%%%%%%%%%%%%%%%%%%%%%%%%%%%%%%%%%%%%
%%%%%%%%%%%%%%%%%%%%%%%%%%%%%%%%%%%%%%%%%%%%%%%%%%%%%%%%

\vtab[-2cm]
\begin{center}
{\large RotationSameMolecule}
\end{center}
\begin{multicols}{2}
\begin{center}
Molecule A \\ 
Tetrabenzoporphyrin\_out\_G09
\includegraphics[width=8cm]{Images/ImagesFromVMD/Tetrabenzoporphyrin_out_G09.jpg}
\\
Inertia Tensor - Molecule A \\
\vtab
\begin{tabular}{|c c c|}
11864.8	 & 	-0	 & 	-0	 \\
-0	 & 	5941.95	 & 	-0	 \\
-0	 & 	-0	 & 	5922.88
\end{tabular}

\vtab
 EingenVectors - Molecule A     \\
\vtab
\begin{tabular}{|c c c|}
0	 & 	0	 & 	1	 \\
0	 & 	1	 & 	0	 \\
1	 & 	0	 & 	0
\end{tabular}

\vtab
 EingenValues - Molecule A     \\
\vtab
\begin{tabular}{|c c c|}
5922.88	 & 	5941.95	 & 	11864.8
\end{tabular}
\columnbreak

Molecule B \\ 
Tetrabenzoporphyrin\_rotated\_out\_G09
\includegraphics[width=8cm]{Images/ImagesFromVMD/Tetrabenzoporphyrin_rotated_out_G09.jpg}
\\
Inertia Tensor - Molecule B \\
\vtab
\begin{tabular}{|c c c|}
11864.8	 & 	-0	 & 	-0	 \\
-0	 & 	5941.95	 & 	-0	 \\
-0	 & 	-0	 & 	5922.88
\end{tabular}

\vtab
 EingenVectors - Molecule B     \\
\vtab
\begin{tabular}{|c c c|}
0	 & 	0	 & 	1	 \\
0	 & 	1	 & 	0	 \\
1	 & 	0	 & 	0
\end{tabular}

\vtab
 EingenValues - Molecule B     \\
\vtab
\begin{tabular}{|c c c|}
5922.88	 & 	5941.95	 & 	11864.8
\end{tabular}

\end{center}
\end{multicols}
\begin{center}
\vtab
\vtab
\textcolor{NavyBlue}{\Large Are the same molecule}
\end{center}
%%%%%%%%%%%%%%%%%%%%%%%%%%%%%%%%%%%%%%%%%%%%%%%%%%%%%%%%
%%%%%%%%%%%%%%%%%%%%%%%%%%%%%%%%%%%%%%%%%%%%%%%%%%%%%%%%
\newpage
%%%%%%%%%%%%%%%%%%%%%%%%%%%%%%%%%%%%%%%%%%%%%%%%%%%%%%%%
%%%%%%%%%%%%%%%%%%%%%%%%%%%%%%%%%%%%%%%%%%%%%%%%%%%%%%%%

\vtab[-2cm]
\begin{center}
{\large RotationSameMolecule}
\end{center}
\begin{multicols}{2}
\begin{center}
Molecule A \\ 
Tetrabenzoporphyrin\_rotated02\_out\_G09
\includegraphics[width=8cm]{Images/ImagesFromVMD/Tetrabenzoporphyrin_rotated02_out_G09.jpg}
\\
Inertia Tensor - Molecule A \\
\vtab
\begin{tabular}{|c c c|}
7791.57	 & 	2641.52	 & 	-788.168	 \\
2641.52	 & 	9673.67	 & 	-1103.82	 \\
-788.168	 & 	-1103.82	 & 	6265.27
\end{tabular}

\vtab
 EingenVectors - Molecule A     \\
\vtab
\begin{tabular}{|c c c|}
-0.771476	 & 	0.397136	 & 	-0.497099	 \\
0.301236	 & 	-0.460189	 & 	-0.835154	 \\
-0.56043	 & 	-0.794046	 & 	0.235394
\end{tabular}

\vtab
 EingenValues - Molecule A     \\
\vtab
\begin{tabular}{|c c c|}
5923.93	 & 	5941.33	 & 	11865.3
\end{tabular}
\columnbreak

Molecule B \\ 
Tetrabenzoporphyrin\_rotated\_out\_G09
\includegraphics[width=8cm]{Images/ImagesFromVMD/Tetrabenzoporphyrin_rotated_out_G09.jpg}
\\
Inertia Tensor - Molecule B \\
\vtab
\begin{tabular}{|c c c|}
11864.8	 & 	-0	 & 	-0	 \\
-0	 & 	5941.95	 & 	-0	 \\
-0	 & 	-0	 & 	5922.88
\end{tabular}

\vtab
 EingenVectors - Molecule B     \\
\vtab
\begin{tabular}{|c c c|}
0	 & 	0	 & 	1	 \\
0	 & 	1	 & 	0	 \\
1	 & 	0	 & 	0
\end{tabular}

\vtab
 EingenValues - Molecule B     \\
\vtab
\begin{tabular}{|c c c|}
5922.88	 & 	5941.95	 & 	11864.8
\end{tabular}

\end{center}
\end{multicols}
\begin{center}
\vtab
\vtab
\textcolor{NavyBlue}{\Large Are Isomers}
\end{center}
%%%%%%%%%%%%%%%%%%%%%%%%%%%%%%%%%%%%%%%%%%%%%%%%%%%%%%%%
%%%%%%%%%%%%%%%%%%%%%%%%%%%%%%%%%%%%%%%%%%%%%%%%%%%%%%%%
\newpage
%%%%%%%%%%%%%%%%%%%%%%%%%%%%%%%%%%%%%%%%%%%%%%%%%%%%%%%%
%%%%%%%%%%%%%%%%%%%%%%%%%%%%%%%%%%%%%%%%%%%%%%%%%%%%%%%%

\vtab[-2cm]
\begin{center}
{\large Sugars}
\end{center}
\begin{multicols}{2}
\begin{center}
Molecule A \\ 
Base-Fructuose\_out\_G09
\includegraphics[width=8cm]{Images/ImagesFromVMD/Base-Fructuose_out_G09.jpg}
\\
Inertia Tensor - Molecule A \\
\vtab
\begin{tabular}{|c c c|}
228.962	 & 	-17.4299	 & 	2.01118	 \\
-17.4299	 & 	1222.64	 & 	-0.103419	 \\
2.01118	 & 	-0.103419	 & 	1316.38
\end{tabular}

\vtab
 EingenVectors - Molecule A     \\
\vtab
\begin{tabular}{|c c c|}
0.999845	 & 	0.0175325	 & 	-0.00184703	 \\
0.0175298	 & 	-0.999845	 & 	-0.00148399	 \\
0.00187276	 & 	-0.00145138	 & 	0.999997
\end{tabular}

\vtab
 EingenValues - Molecule A     \\
\vtab
\begin{tabular}{|c c c|}
228.653	 & 	1222.94	 & 	1316.38
\end{tabular}
\columnbreak

Molecule B \\ 
D\_-\_Fructuose\_out\_G09
\includegraphics[width=8cm]{Images/ImagesFromVMD/D_-_Fructuose_out_G09.jpg}
\\
Inertia Tensor - Molecule B \\
\vtab
\begin{tabular}{|c c c|}
263.78	 & 	3.63756	 & 	5.48351	 \\
3.63756	 & 	1174.62	 & 	1.15892	 \\
5.48351	 & 	1.15892	 & 	1226.18
\end{tabular}

\vtab
 EingenVectors - Molecule B     \\
\vtab
\begin{tabular}{|c c c|}
0.999976	 & 	-0.00398608	 & 	-0.00569256	 \\
0.00385485	 & 	0.999731	 & 	-0.0228795	 \\
0.00578223	 & 	0.022857	 & 	0.999722
\end{tabular}

\vtab
 EingenValues - Molecule B     \\
\vtab
\begin{tabular}{|c c c|}
263.734	 & 	1174.61	 & 	1226.23
\end{tabular}

\end{center}
\end{multicols}
\begin{center}
\vtab
\vtab
\textcolor{NavyBlue}{\Large Are Isomers}
\end{center}
%%%%%%%%%%%%%%%%%%%%%%%%%%%%%%%%%%%%%%%%%%%%%%%%%%%%%%%%
%%%%%%%%%%%%%%%%%%%%%%%%%%%%%%%%%%%%%%%%%%%%%%%%%%%%%%%%
\newpage
%%%%%%%%%%%%%%%%%%%%%%%%%%%%%%%%%%%%%%%%%%%%%%%%%%%%%%%%
%%%%%%%%%%%%%%%%%%%%%%%%%%%%%%%%%%%%%%%%%%%%%%%%%%%%%%%%

\vtab[-2cm]
\begin{center}
{\large Sugars}
\end{center}
\begin{multicols}{2}
\begin{center}
Molecule A \\ 
Base-Fructuose\_out\_G09
\includegraphics[width=8cm]{Images/ImagesFromVMD/Base-Fructuose_out_G09.jpg}
\\
Inertia Tensor - Molecule A \\
\vtab
\begin{tabular}{|c c c|}
228.962	 & 	-17.4299	 & 	2.01118	 \\
-17.4299	 & 	1222.64	 & 	-0.103419	 \\
2.01118	 & 	-0.103419	 & 	1316.38
\end{tabular}

\vtab
 EingenVectors - Molecule A     \\
\vtab
\begin{tabular}{|c c c|}
0.999845	 & 	0.0175325	 & 	-0.00184703	 \\
0.0175298	 & 	-0.999845	 & 	-0.00148399	 \\
0.00187276	 & 	-0.00145138	 & 	0.999997
\end{tabular}

\vtab
 EingenValues - Molecule A     \\
\vtab
\begin{tabular}{|c c c|}
228.653	 & 	1222.94	 & 	1316.38
\end{tabular}
\columnbreak

Molecule B \\ 
R-R-R-Fructuose\_out\_G09
\includegraphics[width=8cm]{Images/ImagesFromVMD/R-R-R-Fructuose_out_G09.jpg}
\\
Inertia Tensor - Molecule B \\
\vtab
\begin{tabular}{|c c c|}
270.815	 & 	9.92006	 & 	1.58147	 \\
9.92006	 & 	1105.83	 & 	-0.519818	 \\
1.58147	 & 	-0.519818	 & 	1246.52
\end{tabular}

\vtab
 EingenVectors - Molecule B     \\
\vtab
\begin{tabular}{|c c c|}
0.999928	 & 	-0.0118786	 & 	-0.00162686	 \\
-0.0118843	 & 	-0.999923	 & 	-0.00356382	 \\
0.0015844	 & 	-0.0035829	 & 	0.999992
\end{tabular}

\vtab
 EingenValues - Molecule B     \\
\vtab
\begin{tabular}{|c c c|}
270.695	 & 	1105.94	 & 	1246.52
\end{tabular}

\end{center}
\end{multicols}
\begin{center}
\vtab
\vtab
\textcolor{NavyBlue}{\Large Are Isomers}
\end{center}
%%%%%%%%%%%%%%%%%%%%%%%%%%%%%%%%%%%%%%%%%%%%%%%%%%%%%%%%
%%%%%%%%%%%%%%%%%%%%%%%%%%%%%%%%%%%%%%%%%%%%%%%%%%%%%%%%
\newpage
%%%%%%%%%%%%%%%%%%%%%%%%%%%%%%%%%%%%%%%%%%%%%%%%%%%%%%%%
%%%%%%%%%%%%%%%%%%%%%%%%%%%%%%%%%%%%%%%%%%%%%%%%%%%%%%%%

\vtab[-2cm]
\begin{center}
{\large Sugars}
\end{center}
\begin{multicols}{2}
\begin{center}
Molecule A \\ 
Base-Fructuose\_out\_G09
\includegraphics[width=8cm]{Images/ImagesFromVMD/Base-Fructuose_out_G09.jpg}
\\
Inertia Tensor - Molecule A \\
\vtab
\begin{tabular}{|c c c|}
228.962	 & 	-17.4299	 & 	2.01118	 \\
-17.4299	 & 	1222.64	 & 	-0.103419	 \\
2.01118	 & 	-0.103419	 & 	1316.38
\end{tabular}

\vtab
 EingenVectors - Molecule A     \\
\vtab
\begin{tabular}{|c c c|}
0.999845	 & 	0.0175325	 & 	-0.00184703	 \\
0.0175298	 & 	-0.999845	 & 	-0.00148399	 \\
0.00187276	 & 	-0.00145138	 & 	0.999997
\end{tabular}

\vtab
 EingenValues - Molecule A     \\
\vtab
\begin{tabular}{|c c c|}
228.653	 & 	1222.94	 & 	1316.38
\end{tabular}
\columnbreak

Molecule B \\ 
R-R-S-Fructuose\_out\_G09
\includegraphics[width=8cm]{Images/ImagesFromVMD/R-R-S-Fructuose_out_G09.jpg}
\\
Inertia Tensor - Molecule B \\
\vtab
\begin{tabular}{|c c c|}
258.077	 & 	6.73578	 & 	-10.2403	 \\
6.73578	 & 	1157.51	 & 	-1.31161	 \\
-10.2403	 & 	-1.31161	 & 	1279.98
\end{tabular}

\vtab
 EingenVectors - Molecule B     \\
\vtab
\begin{tabular}{|c c c|}
-0.999922	 & 	0.00747247	 & 	-0.0100089	 \\
0.00735899	 & 	0.999909	 & 	0.0113273	 \\
-0.0100927	 & 	-0.0112527	 & 	0.999886
\end{tabular}

\vtab
 EingenValues - Molecule B     \\
\vtab
\begin{tabular}{|c c c|}
257.924	 & 	1157.54	 & 	1280.1
\end{tabular}

\end{center}
\end{multicols}
\begin{center}
\vtab
\vtab
\textcolor{NavyBlue}{\Large Are Isomers}
\end{center}
%%%%%%%%%%%%%%%%%%%%%%%%%%%%%%%%%%%%%%%%%%%%%%%%%%%%%%%%
%%%%%%%%%%%%%%%%%%%%%%%%%%%%%%%%%%%%%%%%%%%%%%%%%%%%%%%%
\newpage
%%%%%%%%%%%%%%%%%%%%%%%%%%%%%%%%%%%%%%%%%%%%%%%%%%%%%%%%
%%%%%%%%%%%%%%%%%%%%%%%%%%%%%%%%%%%%%%%%%%%%%%%%%%%%%%%%

\vtab[-2cm]
\begin{center}
{\large Sugars}
\end{center}
\begin{multicols}{2}
\begin{center}
Molecule A \\ 
Base-Fructuose\_out\_G09
\includegraphics[width=8cm]{Images/ImagesFromVMD/Base-Fructuose_out_G09.jpg}
\\
Inertia Tensor - Molecule A \\
\vtab
\begin{tabular}{|c c c|}
228.962	 & 	-17.4299	 & 	2.01118	 \\
-17.4299	 & 	1222.64	 & 	-0.103419	 \\
2.01118	 & 	-0.103419	 & 	1316.38
\end{tabular}

\vtab
 EingenVectors - Molecule A     \\
\vtab
\begin{tabular}{|c c c|}
0.999845	 & 	0.0175325	 & 	-0.00184703	 \\
0.0175298	 & 	-0.999845	 & 	-0.00148399	 \\
0.00187276	 & 	-0.00145138	 & 	0.999997
\end{tabular}

\vtab
 EingenValues - Molecule A     \\
\vtab
\begin{tabular}{|c c c|}
228.653	 & 	1222.94	 & 	1316.38
\end{tabular}
\columnbreak

Molecule B \\ 
R-S-R-Fructuose\_out\_G09
\includegraphics[width=8cm]{Images/ImagesFromVMD/R-S-R-Fructuose_out_G09.jpg}
\\
Inertia Tensor - Molecule B \\
\vtab
\begin{tabular}{|c c c|}
202.772	 & 	14.8494	 & 	0.184396	 \\
14.8494	 & 	1218.4	 & 	-0.831783	 \\
0.184396	 & 	-0.831783	 & 	1335.51
\end{tabular}

\vtab
 EingenVectors - Molecule B     \\
\vtab
\begin{tabular}{|c c c|}
-0.999893	 & 	0.0146163	 & 	0.00017347	 \\
-0.0146172	 & 	-0.999868	 & 	-0.00709174	 \\
6.97925e-05	 & 	-0.00709352	 & 	0.999975
\end{tabular}

\vtab
 EingenValues - Molecule B     \\
\vtab
\begin{tabular}{|c c c|}
202.555	 & 	1218.61	 & 	1335.51
\end{tabular}

\end{center}
\end{multicols}
\begin{center}
\vtab
\vtab
\textcolor{NavyBlue}{\Large Are Isomers}
\end{center}
%%%%%%%%%%%%%%%%%%%%%%%%%%%%%%%%%%%%%%%%%%%%%%%%%%%%%%%%
%%%%%%%%%%%%%%%%%%%%%%%%%%%%%%%%%%%%%%%%%%%%%%%%%%%%%%%%
\newpage
%%%%%%%%%%%%%%%%%%%%%%%%%%%%%%%%%%%%%%%%%%%%%%%%%%%%%%%%
%%%%%%%%%%%%%%%%%%%%%%%%%%%%%%%%%%%%%%%%%%%%%%%%%%%%%%%%

\vtab[-2cm]
\begin{center}
{\large Sugars}
\end{center}
\begin{multicols}{2}
\begin{center}
Molecule A \\ 
Base-Fructuose\_out\_G09
\includegraphics[width=8cm]{Images/ImagesFromVMD/Base-Fructuose_out_G09.jpg}
\\
Inertia Tensor - Molecule A \\
\vtab
\begin{tabular}{|c c c|}
228.962	 & 	-17.4299	 & 	2.01118	 \\
-17.4299	 & 	1222.64	 & 	-0.103419	 \\
2.01118	 & 	-0.103419	 & 	1316.38
\end{tabular}

\vtab
 EingenVectors - Molecule A     \\
\vtab
\begin{tabular}{|c c c|}
0.999845	 & 	0.0175325	 & 	-0.00184703	 \\
0.0175298	 & 	-0.999845	 & 	-0.00148399	 \\
0.00187276	 & 	-0.00145138	 & 	0.999997
\end{tabular}

\vtab
 EingenValues - Molecule A     \\
\vtab
\begin{tabular}{|c c c|}
228.653	 & 	1222.94	 & 	1316.38
\end{tabular}
\columnbreak

Molecule B \\ 
R-S-S-Fructuose\_out\_G09
\includegraphics[width=8cm]{Images/ImagesFromVMD/R-S-S-Fructuose_out_G09.jpg}
\\
Inertia Tensor - Molecule B \\
\vtab
\begin{tabular}{|c c c|}
215.024	 & 	-14.0942	 & 	-5.97678	 \\
-14.0942	 & 	1187.1	 & 	0.653719	 \\
-5.97678	 & 	0.653719	 & 	1326.43
\end{tabular}

\vtab
 EingenVectors - Molecule B     \\
\vtab
\begin{tabular}{|c c c|}
0.999881	 & 	0.0144902	 & 	0.00536739	 \\
0.0144617	 & 	-0.999881	 & 	0.0053195	 \\
-0.00544384	 & 	0.00524124	 & 	0.999971
\end{tabular}

\vtab
 EingenValues - Molecule B     \\
\vtab
\begin{tabular}{|c c c|}
214.787	 & 	1187.3	 & 	1326.46
\end{tabular}

\end{center}
\end{multicols}
\begin{center}
\vtab
\vtab
\textcolor{NavyBlue}{\Large Are Isomers}
\end{center}
%%%%%%%%%%%%%%%%%%%%%%%%%%%%%%%%%%%%%%%%%%%%%%%%%%%%%%%%
%%%%%%%%%%%%%%%%%%%%%%%%%%%%%%%%%%%%%%%%%%%%%%%%%%%%%%%%
\newpage
%%%%%%%%%%%%%%%%%%%%%%%%%%%%%%%%%%%%%%%%%%%%%%%%%%%%%%%%
%%%%%%%%%%%%%%%%%%%%%%%%%%%%%%%%%%%%%%%%%%%%%%%%%%%%%%%%

\vtab[-2cm]
\begin{center}
{\large Sugars}
\end{center}
\begin{multicols}{2}
\begin{center}
Molecule A \\ 
Base-Fructuose\_out\_G09
\includegraphics[width=8cm]{Images/ImagesFromVMD/Base-Fructuose_out_G09.jpg}
\\
Inertia Tensor - Molecule A \\
\vtab
\begin{tabular}{|c c c|}
228.962	 & 	-17.4299	 & 	2.01118	 \\
-17.4299	 & 	1222.64	 & 	-0.103419	 \\
2.01118	 & 	-0.103419	 & 	1316.38
\end{tabular}

\vtab
 EingenVectors - Molecule A     \\
\vtab
\begin{tabular}{|c c c|}
0.999845	 & 	0.0175325	 & 	-0.00184703	 \\
0.0175298	 & 	-0.999845	 & 	-0.00148399	 \\
0.00187276	 & 	-0.00145138	 & 	0.999997
\end{tabular}

\vtab
 EingenValues - Molecule A     \\
\vtab
\begin{tabular}{|c c c|}
228.653	 & 	1222.94	 & 	1316.38
\end{tabular}
\columnbreak

Molecule B \\ 
S-R-R-Fructuose\_out\_G09
\includegraphics[width=8cm]{Images/ImagesFromVMD/S-R-R-Fructuose_out_G09.jpg}
\\
Inertia Tensor - Molecule B \\
\vtab
\begin{tabular}{|c c c|}
267.669	 & 	-2.83987	 & 	8.0125	 \\
-2.83987	 & 	1173.13	 & 	0.735007	 \\
8.0125	 & 	0.735007	 & 	1229.07
\end{tabular}

\vtab
 EingenVectors - Molecule B     \\
\vtab
\begin{tabular}{|c c c|}
0.99996	 & 	0.00314275	 & 	-0.00833561	 \\
-0.00324825	 & 	0.999914	 & 	-0.0126739	 \\
0.00829507	 & 	0.0127005	 & 	0.999885
\end{tabular}

\vtab
 EingenValues - Molecule B     \\
\vtab
\begin{tabular}{|c c c|}
267.593	 & 	1173.13	 & 	1229.15
\end{tabular}

\end{center}
\end{multicols}
\begin{center}
\vtab
\vtab
\textcolor{NavyBlue}{\Large Are Isomers}
\end{center}
%%%%%%%%%%%%%%%%%%%%%%%%%%%%%%%%%%%%%%%%%%%%%%%%%%%%%%%%
%%%%%%%%%%%%%%%%%%%%%%%%%%%%%%%%%%%%%%%%%%%%%%%%%%%%%%%%
\newpage
%%%%%%%%%%%%%%%%%%%%%%%%%%%%%%%%%%%%%%%%%%%%%%%%%%%%%%%%
%%%%%%%%%%%%%%%%%%%%%%%%%%%%%%%%%%%%%%%%%%%%%%%%%%%%%%%%

\vtab[-2cm]
\begin{center}
{\large Sugars}
\end{center}
\begin{multicols}{2}
\begin{center}
Molecule A \\ 
Base-Fructuose\_out\_G09
\includegraphics[width=8cm]{Images/ImagesFromVMD/Base-Fructuose_out_G09.jpg}
\\
Inertia Tensor - Molecule A \\
\vtab
\begin{tabular}{|c c c|}
228.962	 & 	-17.4299	 & 	2.01118	 \\
-17.4299	 & 	1222.64	 & 	-0.103419	 \\
2.01118	 & 	-0.103419	 & 	1316.38
\end{tabular}

\vtab
 EingenVectors - Molecule A     \\
\vtab
\begin{tabular}{|c c c|}
0.999845	 & 	0.0175325	 & 	-0.00184703	 \\
0.0175298	 & 	-0.999845	 & 	-0.00148399	 \\
0.00187276	 & 	-0.00145138	 & 	0.999997
\end{tabular}

\vtab
 EingenValues - Molecule A     \\
\vtab
\begin{tabular}{|c c c|}
228.653	 & 	1222.94	 & 	1316.38
\end{tabular}
\columnbreak

Molecule B \\ 
S-R-S-Fructuose\_out\_G09
\includegraphics[width=8cm]{Images/ImagesFromVMD/S-R-S-Fructuose_out_G09.jpg}
\\
Inertia Tensor - Molecule B \\
\vtab
\begin{tabular}{|c c c|}
202.782	 & 	-14.8472	 & 	-0.161586	 \\
-14.8472	 & 	1218.41	 & 	-0.836204	 \\
-0.161586	 & 	-0.836204	 & 	1335.52
\end{tabular}

\vtab
 EingenVectors - Molecule B     \\
\vtab
\begin{tabular}{|c c c|}
-0.999893	 & 	-0.0146142	 & 	-0.000153395	 \\
0.0146149	 & 	-0.999868	 & 	-0.00713203	 \\
-4.91458e-05	 & 	-0.00713351	 & 	0.999975
\end{tabular}

\vtab
 EingenValues - Molecule B     \\
\vtab
\begin{tabular}{|c c c|}
202.565	 & 	1218.62	 & 	1335.52
\end{tabular}

\end{center}
\end{multicols}
\begin{center}
\vtab
\vtab
\textcolor{NavyBlue}{\Large Are Isomers}
\end{center}
%%%%%%%%%%%%%%%%%%%%%%%%%%%%%%%%%%%%%%%%%%%%%%%%%%%%%%%%
%%%%%%%%%%%%%%%%%%%%%%%%%%%%%%%%%%%%%%%%%%%%%%%%%%%%%%%%
\newpage
%%%%%%%%%%%%%%%%%%%%%%%%%%%%%%%%%%%%%%%%%%%%%%%%%%%%%%%%
%%%%%%%%%%%%%%%%%%%%%%%%%%%%%%%%%%%%%%%%%%%%%%%%%%%%%%%%

\vtab[-2cm]
\begin{center}
{\large Sugars}
\end{center}
\begin{multicols}{2}
\begin{center}
Molecule A \\ 
Base-Fructuose\_out\_G09
\includegraphics[width=8cm]{Images/ImagesFromVMD/Base-Fructuose_out_G09.jpg}
\\
Inertia Tensor - Molecule A \\
\vtab
\begin{tabular}{|c c c|}
228.962	 & 	-17.4299	 & 	2.01118	 \\
-17.4299	 & 	1222.64	 & 	-0.103419	 \\
2.01118	 & 	-0.103419	 & 	1316.38
\end{tabular}

\vtab
 EingenVectors - Molecule A     \\
\vtab
\begin{tabular}{|c c c|}
0.999845	 & 	0.0175325	 & 	-0.00184703	 \\
0.0175298	 & 	-0.999845	 & 	-0.00148399	 \\
0.00187276	 & 	-0.00145138	 & 	0.999997
\end{tabular}

\vtab
 EingenValues - Molecule A     \\
\vtab
\begin{tabular}{|c c c|}
228.653	 & 	1222.94	 & 	1316.38
\end{tabular}
\columnbreak

Molecule B \\ 
S-S-R-Fructuose\_out\_G09
\includegraphics[width=8cm]{Images/ImagesFromVMD/S-S-R-Fructuose_out_G09.jpg}
\\
Inertia Tensor - Molecule B \\
\vtab
\begin{tabular}{|c c c|}
258.093	 & 	-6.7331	 & 	10.2577	 \\
-6.7331	 & 	1157.52	 & 	-1.31924	 \\
10.2577	 & 	-1.31924	 & 	1279.94
\end{tabular}

\vtab
 EingenVectors - Molecule B     \\
\vtab
\begin{tabular}{|c c c|}
0.999922	 & 	0.00746944	 & 	-0.0100264	 \\
0.00735508	 & 	-0.999908	 & 	-0.0113945	 \\
0.0101106	 & 	-0.0113198	 & 	0.999885
\end{tabular}

\vtab
 EingenValues - Molecule B     \\
\vtab
\begin{tabular}{|c c c|}
257.939	 & 	1157.55	 & 	1280.06
\end{tabular}

\end{center}
\end{multicols}
\begin{center}
\vtab
\vtab
\textcolor{NavyBlue}{\Large Are Isomers}
\end{center}
%%%%%%%%%%%%%%%%%%%%%%%%%%%%%%%%%%%%%%%%%%%%%%%%%%%%%%%%
%%%%%%%%%%%%%%%%%%%%%%%%%%%%%%%%%%%%%%%%%%%%%%%%%%%%%%%%
\newpage
%%%%%%%%%%%%%%%%%%%%%%%%%%%%%%%%%%%%%%%%%%%%%%%%%%%%%%%%
%%%%%%%%%%%%%%%%%%%%%%%%%%%%%%%%%%%%%%%%%%%%%%%%%%%%%%%%

\vtab[-2cm]
\begin{center}
{\large Sugars}
\end{center}
\begin{multicols}{2}
\begin{center}
Molecule A \\ 
Base-Fructuose\_out\_G09
\includegraphics[width=8cm]{Images/ImagesFromVMD/Base-Fructuose_out_G09.jpg}
\\
Inertia Tensor - Molecule A \\
\vtab
\begin{tabular}{|c c c|}
228.962	 & 	-17.4299	 & 	2.01118	 \\
-17.4299	 & 	1222.64	 & 	-0.103419	 \\
2.01118	 & 	-0.103419	 & 	1316.38
\end{tabular}

\vtab
 EingenVectors - Molecule A     \\
\vtab
\begin{tabular}{|c c c|}
0.999845	 & 	0.0175325	 & 	-0.00184703	 \\
0.0175298	 & 	-0.999845	 & 	-0.00148399	 \\
0.00187276	 & 	-0.00145138	 & 	0.999997
\end{tabular}

\vtab
 EingenValues - Molecule A     \\
\vtab
\begin{tabular}{|c c c|}
228.653	 & 	1222.94	 & 	1316.38
\end{tabular}
\columnbreak

Molecule B \\ 
S-S-S-Fructuose\_out\_G09
\includegraphics[width=8cm]{Images/ImagesFromVMD/S-S-S-Fructuose_out_G09.jpg}
\\
Inertia Tensor - Molecule B \\
\vtab
\begin{tabular}{|c c c|}
226.452	 & 	13.0925	 & 	-2.23012	 \\
13.0925	 & 	1177.96	 & 	-0.925239	 \\
-2.23012	 & 	-0.925239	 & 	1359.88
\end{tabular}

\vtab
 EingenVectors - Molecule B     \\
\vtab
\begin{tabular}{|c c c|}
-0.999903	 & 	0.0137538	 & 	-0.00195585	 \\
0.0137434	 & 	0.999892	 & 	0.00525895	 \\
-0.00202797	 & 	-0.00523156	 & 	0.999984
\end{tabular}

\vtab
 EingenValues - Molecule B     \\
\vtab
\begin{tabular}{|c c c|}
226.267	 & 	1178.14	 & 	1359.89
\end{tabular}

\end{center}
\end{multicols}
\begin{center}
\vtab
\vtab
\textcolor{NavyBlue}{\Large Are Isomers}
\end{center}
%%%%%%%%%%%%%%%%%%%%%%%%%%%%%%%%%%%%%%%%%%%%%%%%%%%%%%%%
%%%%%%%%%%%%%%%%%%%%%%%%%%%%%%%%%%%%%%%%%%%%%%%%%%%%%%%%
\newpage
%%%%%%%%%%%%%%%%%%%%%%%%%%%%%%%%%%%%%%%%%%%%%%%%%%%%%%%%
%%%%%%%%%%%%%%%%%%%%%%%%%%%%%%%%%%%%%%%%%%%%%%%%%%%%%%%%

\vtab[-2cm]
\begin{center}
{\large Sugars}
\end{center}
\begin{multicols}{2}
\begin{center}
Molecule A \\ 
D\_-\_Fructuose\_out\_G09
\includegraphics[width=8cm]{Images/ImagesFromVMD/D_-_Fructuose_out_G09.jpg}
\\
Inertia Tensor - Molecule A \\
\vtab
\begin{tabular}{|c c c|}
263.78	 & 	3.63756	 & 	5.48351	 \\
3.63756	 & 	1174.62	 & 	1.15892	 \\
5.48351	 & 	1.15892	 & 	1226.18
\end{tabular}

\vtab
 EingenVectors - Molecule A     \\
\vtab
\begin{tabular}{|c c c|}
0.999976	 & 	-0.00398608	 & 	-0.00569256	 \\
0.00385485	 & 	0.999731	 & 	-0.0228795	 \\
0.00578223	 & 	0.022857	 & 	0.999722
\end{tabular}

\vtab
 EingenValues - Molecule A     \\
\vtab
\begin{tabular}{|c c c|}
263.734	 & 	1174.61	 & 	1226.23
\end{tabular}
\columnbreak

Molecule B \\ 
R-R-R-Fructuose\_out\_G09
\includegraphics[width=8cm]{Images/ImagesFromVMD/R-R-R-Fructuose_out_G09.jpg}
\\
Inertia Tensor - Molecule B \\
\vtab
\begin{tabular}{|c c c|}
270.815	 & 	9.92006	 & 	1.58147	 \\
9.92006	 & 	1105.83	 & 	-0.519818	 \\
1.58147	 & 	-0.519818	 & 	1246.52
\end{tabular}

\vtab
 EingenVectors - Molecule B     \\
\vtab
\begin{tabular}{|c c c|}
0.999928	 & 	-0.0118786	 & 	-0.00162686	 \\
-0.0118843	 & 	-0.999923	 & 	-0.00356382	 \\
0.0015844	 & 	-0.0035829	 & 	0.999992
\end{tabular}

\vtab
 EingenValues - Molecule B     \\
\vtab
\begin{tabular}{|c c c|}
270.695	 & 	1105.94	 & 	1246.52
\end{tabular}

\end{center}
\end{multicols}
\begin{center}
\vtab
\vtab
\textcolor{NavyBlue}{\Large Are Isomers}
\end{center}
%%%%%%%%%%%%%%%%%%%%%%%%%%%%%%%%%%%%%%%%%%%%%%%%%%%%%%%%
%%%%%%%%%%%%%%%%%%%%%%%%%%%%%%%%%%%%%%%%%%%%%%%%%%%%%%%%
\newpage
%%%%%%%%%%%%%%%%%%%%%%%%%%%%%%%%%%%%%%%%%%%%%%%%%%%%%%%%
%%%%%%%%%%%%%%%%%%%%%%%%%%%%%%%%%%%%%%%%%%%%%%%%%%%%%%%%

\vtab[-2cm]
\begin{center}
{\large Sugars}
\end{center}
\begin{multicols}{2}
\begin{center}
Molecule A \\ 
D\_-\_Fructuose\_out\_G09
\includegraphics[width=8cm]{Images/ImagesFromVMD/D_-_Fructuose_out_G09.jpg}
\\
Inertia Tensor - Molecule A \\
\vtab
\begin{tabular}{|c c c|}
263.78	 & 	3.63756	 & 	5.48351	 \\
3.63756	 & 	1174.62	 & 	1.15892	 \\
5.48351	 & 	1.15892	 & 	1226.18
\end{tabular}

\vtab
 EingenVectors - Molecule A     \\
\vtab
\begin{tabular}{|c c c|}
0.999976	 & 	-0.00398608	 & 	-0.00569256	 \\
0.00385485	 & 	0.999731	 & 	-0.0228795	 \\
0.00578223	 & 	0.022857	 & 	0.999722
\end{tabular}

\vtab
 EingenValues - Molecule A     \\
\vtab
\begin{tabular}{|c c c|}
263.734	 & 	1174.61	 & 	1226.23
\end{tabular}
\columnbreak

Molecule B \\ 
R-R-S-Fructuose\_out\_G09
\includegraphics[width=8cm]{Images/ImagesFromVMD/R-R-S-Fructuose_out_G09.jpg}
\\
Inertia Tensor - Molecule B \\
\vtab
\begin{tabular}{|c c c|}
258.077	 & 	6.73578	 & 	-10.2403	 \\
6.73578	 & 	1157.51	 & 	-1.31161	 \\
-10.2403	 & 	-1.31161	 & 	1279.98
\end{tabular}

\vtab
 EingenVectors - Molecule B     \\
\vtab
\begin{tabular}{|c c c|}
-0.999922	 & 	0.00747247	 & 	-0.0100089	 \\
0.00735899	 & 	0.999909	 & 	0.0113273	 \\
-0.0100927	 & 	-0.0112527	 & 	0.999886
\end{tabular}

\vtab
 EingenValues - Molecule B     \\
\vtab
\begin{tabular}{|c c c|}
257.924	 & 	1157.54	 & 	1280.1
\end{tabular}

\end{center}
\end{multicols}
\begin{center}
\vtab
\vtab
\textcolor{NavyBlue}{\Large Are Isomers}
\end{center}
%%%%%%%%%%%%%%%%%%%%%%%%%%%%%%%%%%%%%%%%%%%%%%%%%%%%%%%%
%%%%%%%%%%%%%%%%%%%%%%%%%%%%%%%%%%%%%%%%%%%%%%%%%%%%%%%%
\newpage
%%%%%%%%%%%%%%%%%%%%%%%%%%%%%%%%%%%%%%%%%%%%%%%%%%%%%%%%
%%%%%%%%%%%%%%%%%%%%%%%%%%%%%%%%%%%%%%%%%%%%%%%%%%%%%%%%

\vtab[-2cm]
\begin{center}
{\large Sugars}
\end{center}
\begin{multicols}{2}
\begin{center}
Molecule A \\ 
D\_-\_Fructuose\_out\_G09
\includegraphics[width=8cm]{Images/ImagesFromVMD/D_-_Fructuose_out_G09.jpg}
\\
Inertia Tensor - Molecule A \\
\vtab
\begin{tabular}{|c c c|}
263.78	 & 	3.63756	 & 	5.48351	 \\
3.63756	 & 	1174.62	 & 	1.15892	 \\
5.48351	 & 	1.15892	 & 	1226.18
\end{tabular}

\vtab
 EingenVectors - Molecule A     \\
\vtab
\begin{tabular}{|c c c|}
0.999976	 & 	-0.00398608	 & 	-0.00569256	 \\
0.00385485	 & 	0.999731	 & 	-0.0228795	 \\
0.00578223	 & 	0.022857	 & 	0.999722
\end{tabular}

\vtab
 EingenValues - Molecule A     \\
\vtab
\begin{tabular}{|c c c|}
263.734	 & 	1174.61	 & 	1226.23
\end{tabular}
\columnbreak

Molecule B \\ 
R-S-R-Fructuose\_out\_G09
\includegraphics[width=8cm]{Images/ImagesFromVMD/R-S-R-Fructuose_out_G09.jpg}
\\
Inertia Tensor - Molecule B \\
\vtab
\begin{tabular}{|c c c|}
202.772	 & 	14.8494	 & 	0.184396	 \\
14.8494	 & 	1218.4	 & 	-0.831783	 \\
0.184396	 & 	-0.831783	 & 	1335.51
\end{tabular}

\vtab
 EingenVectors - Molecule B     \\
\vtab
\begin{tabular}{|c c c|}
-0.999893	 & 	0.0146163	 & 	0.00017347	 \\
-0.0146172	 & 	-0.999868	 & 	-0.00709174	 \\
6.97925e-05	 & 	-0.00709352	 & 	0.999975
\end{tabular}

\vtab
 EingenValues - Molecule B     \\
\vtab
\begin{tabular}{|c c c|}
202.555	 & 	1218.61	 & 	1335.51
\end{tabular}

\end{center}
\end{multicols}
\begin{center}
\vtab
\vtab
\textcolor{NavyBlue}{\Large Are Isomers}
\end{center}
%%%%%%%%%%%%%%%%%%%%%%%%%%%%%%%%%%%%%%%%%%%%%%%%%%%%%%%%
%%%%%%%%%%%%%%%%%%%%%%%%%%%%%%%%%%%%%%%%%%%%%%%%%%%%%%%%
\newpage
%%%%%%%%%%%%%%%%%%%%%%%%%%%%%%%%%%%%%%%%%%%%%%%%%%%%%%%%
%%%%%%%%%%%%%%%%%%%%%%%%%%%%%%%%%%%%%%%%%%%%%%%%%%%%%%%%

\vtab[-2cm]
\begin{center}
{\large Sugars}
\end{center}
\begin{multicols}{2}
\begin{center}
Molecule A \\ 
D\_-\_Fructuose\_out\_G09
\includegraphics[width=8cm]{Images/ImagesFromVMD/D_-_Fructuose_out_G09.jpg}
\\
Inertia Tensor - Molecule A \\
\vtab
\begin{tabular}{|c c c|}
263.78	 & 	3.63756	 & 	5.48351	 \\
3.63756	 & 	1174.62	 & 	1.15892	 \\
5.48351	 & 	1.15892	 & 	1226.18
\end{tabular}

\vtab
 EingenVectors - Molecule A     \\
\vtab
\begin{tabular}{|c c c|}
0.999976	 & 	-0.00398608	 & 	-0.00569256	 \\
0.00385485	 & 	0.999731	 & 	-0.0228795	 \\
0.00578223	 & 	0.022857	 & 	0.999722
\end{tabular}

\vtab
 EingenValues - Molecule A     \\
\vtab
\begin{tabular}{|c c c|}
263.734	 & 	1174.61	 & 	1226.23
\end{tabular}
\columnbreak

Molecule B \\ 
R-S-S-Fructuose\_out\_G09
\includegraphics[width=8cm]{Images/ImagesFromVMD/R-S-S-Fructuose_out_G09.jpg}
\\
Inertia Tensor - Molecule B \\
\vtab
\begin{tabular}{|c c c|}
215.024	 & 	-14.0942	 & 	-5.97678	 \\
-14.0942	 & 	1187.1	 & 	0.653719	 \\
-5.97678	 & 	0.653719	 & 	1326.43
\end{tabular}

\vtab
 EingenVectors - Molecule B     \\
\vtab
\begin{tabular}{|c c c|}
0.999881	 & 	0.0144902	 & 	0.00536739	 \\
0.0144617	 & 	-0.999881	 & 	0.0053195	 \\
-0.00544384	 & 	0.00524124	 & 	0.999971
\end{tabular}

\vtab
 EingenValues - Molecule B     \\
\vtab
\begin{tabular}{|c c c|}
214.787	 & 	1187.3	 & 	1326.46
\end{tabular}

\end{center}
\end{multicols}
\begin{center}
\vtab
\vtab
\textcolor{NavyBlue}{\Large Are Isomers}
\end{center}
%%%%%%%%%%%%%%%%%%%%%%%%%%%%%%%%%%%%%%%%%%%%%%%%%%%%%%%%
%%%%%%%%%%%%%%%%%%%%%%%%%%%%%%%%%%%%%%%%%%%%%%%%%%%%%%%%
\newpage
%%%%%%%%%%%%%%%%%%%%%%%%%%%%%%%%%%%%%%%%%%%%%%%%%%%%%%%%
%%%%%%%%%%%%%%%%%%%%%%%%%%%%%%%%%%%%%%%%%%%%%%%%%%%%%%%%

\vtab[-2cm]
\begin{center}
{\large Sugars}
\end{center}
\begin{multicols}{2}
\begin{center}
Molecule A \\ 
D\_-\_Fructuose\_out\_G09
\includegraphics[width=8cm]{Images/ImagesFromVMD/D_-_Fructuose_out_G09.jpg}
\\
Inertia Tensor - Molecule A \\
\vtab
\begin{tabular}{|c c c|}
263.78	 & 	3.63756	 & 	5.48351	 \\
3.63756	 & 	1174.62	 & 	1.15892	 \\
5.48351	 & 	1.15892	 & 	1226.18
\end{tabular}

\vtab
 EingenVectors - Molecule A     \\
\vtab
\begin{tabular}{|c c c|}
0.999976	 & 	-0.00398608	 & 	-0.00569256	 \\
0.00385485	 & 	0.999731	 & 	-0.0228795	 \\
0.00578223	 & 	0.022857	 & 	0.999722
\end{tabular}

\vtab
 EingenValues - Molecule A     \\
\vtab
\begin{tabular}{|c c c|}
263.734	 & 	1174.61	 & 	1226.23
\end{tabular}
\columnbreak

Molecule B \\ 
S-R-R-Fructuose\_out\_G09
\includegraphics[width=8cm]{Images/ImagesFromVMD/S-R-R-Fructuose_out_G09.jpg}
\\
Inertia Tensor - Molecule B \\
\vtab
\begin{tabular}{|c c c|}
267.669	 & 	-2.83987	 & 	8.0125	 \\
-2.83987	 & 	1173.13	 & 	0.735007	 \\
8.0125	 & 	0.735007	 & 	1229.07
\end{tabular}

\vtab
 EingenVectors - Molecule B     \\
\vtab
\begin{tabular}{|c c c|}
0.99996	 & 	0.00314275	 & 	-0.00833561	 \\
-0.00324825	 & 	0.999914	 & 	-0.0126739	 \\
0.00829507	 & 	0.0127005	 & 	0.999885
\end{tabular}

\vtab
 EingenValues - Molecule B     \\
\vtab
\begin{tabular}{|c c c|}
267.593	 & 	1173.13	 & 	1229.15
\end{tabular}

\end{center}
\end{multicols}
\begin{center}
\vtab
\vtab
\textcolor{NavyBlue}{\Large Are Isomers}
\end{center}
%%%%%%%%%%%%%%%%%%%%%%%%%%%%%%%%%%%%%%%%%%%%%%%%%%%%%%%%
%%%%%%%%%%%%%%%%%%%%%%%%%%%%%%%%%%%%%%%%%%%%%%%%%%%%%%%%
\newpage
%%%%%%%%%%%%%%%%%%%%%%%%%%%%%%%%%%%%%%%%%%%%%%%%%%%%%%%%
%%%%%%%%%%%%%%%%%%%%%%%%%%%%%%%%%%%%%%%%%%%%%%%%%%%%%%%%

\vtab[-2cm]
\begin{center}
{\large Sugars}
\end{center}
\begin{multicols}{2}
\begin{center}
Molecule A \\ 
D\_-\_Fructuose\_out\_G09
\includegraphics[width=8cm]{Images/ImagesFromVMD/D_-_Fructuose_out_G09.jpg}
\\
Inertia Tensor - Molecule A \\
\vtab
\begin{tabular}{|c c c|}
263.78	 & 	3.63756	 & 	5.48351	 \\
3.63756	 & 	1174.62	 & 	1.15892	 \\
5.48351	 & 	1.15892	 & 	1226.18
\end{tabular}

\vtab
 EingenVectors - Molecule A     \\
\vtab
\begin{tabular}{|c c c|}
0.999976	 & 	-0.00398608	 & 	-0.00569256	 \\
0.00385485	 & 	0.999731	 & 	-0.0228795	 \\
0.00578223	 & 	0.022857	 & 	0.999722
\end{tabular}

\vtab
 EingenValues - Molecule A     \\
\vtab
\begin{tabular}{|c c c|}
263.734	 & 	1174.61	 & 	1226.23
\end{tabular}
\columnbreak

Molecule B \\ 
S-R-S-Fructuose\_out\_G09
\includegraphics[width=8cm]{Images/ImagesFromVMD/S-R-S-Fructuose_out_G09.jpg}
\\
Inertia Tensor - Molecule B \\
\vtab
\begin{tabular}{|c c c|}
202.782	 & 	-14.8472	 & 	-0.161586	 \\
-14.8472	 & 	1218.41	 & 	-0.836204	 \\
-0.161586	 & 	-0.836204	 & 	1335.52
\end{tabular}

\vtab
 EingenVectors - Molecule B     \\
\vtab
\begin{tabular}{|c c c|}
-0.999893	 & 	-0.0146142	 & 	-0.000153395	 \\
0.0146149	 & 	-0.999868	 & 	-0.00713203	 \\
-4.91458e-05	 & 	-0.00713351	 & 	0.999975
\end{tabular}

\vtab
 EingenValues - Molecule B     \\
\vtab
\begin{tabular}{|c c c|}
202.565	 & 	1218.62	 & 	1335.52
\end{tabular}

\end{center}
\end{multicols}
\begin{center}
\vtab
\vtab
\textcolor{NavyBlue}{\Large Are Isomers}
\end{center}
%%%%%%%%%%%%%%%%%%%%%%%%%%%%%%%%%%%%%%%%%%%%%%%%%%%%%%%%
%%%%%%%%%%%%%%%%%%%%%%%%%%%%%%%%%%%%%%%%%%%%%%%%%%%%%%%%
\newpage
%%%%%%%%%%%%%%%%%%%%%%%%%%%%%%%%%%%%%%%%%%%%%%%%%%%%%%%%
%%%%%%%%%%%%%%%%%%%%%%%%%%%%%%%%%%%%%%%%%%%%%%%%%%%%%%%%

\vtab[-2cm]
\begin{center}
{\large Sugars}
\end{center}
\begin{multicols}{2}
\begin{center}
Molecule A \\ 
D\_-\_Fructuose\_out\_G09
\includegraphics[width=8cm]{Images/ImagesFromVMD/D_-_Fructuose_out_G09.jpg}
\\
Inertia Tensor - Molecule A \\
\vtab
\begin{tabular}{|c c c|}
263.78	 & 	3.63756	 & 	5.48351	 \\
3.63756	 & 	1174.62	 & 	1.15892	 \\
5.48351	 & 	1.15892	 & 	1226.18
\end{tabular}

\vtab
 EingenVectors - Molecule A     \\
\vtab
\begin{tabular}{|c c c|}
0.999976	 & 	-0.00398608	 & 	-0.00569256	 \\
0.00385485	 & 	0.999731	 & 	-0.0228795	 \\
0.00578223	 & 	0.022857	 & 	0.999722
\end{tabular}

\vtab
 EingenValues - Molecule A     \\
\vtab
\begin{tabular}{|c c c|}
263.734	 & 	1174.61	 & 	1226.23
\end{tabular}
\columnbreak

Molecule B \\ 
S-S-R-Fructuose\_out\_G09
\includegraphics[width=8cm]{Images/ImagesFromVMD/S-S-R-Fructuose_out_G09.jpg}
\\
Inertia Tensor - Molecule B \\
\vtab
\begin{tabular}{|c c c|}
258.093	 & 	-6.7331	 & 	10.2577	 \\
-6.7331	 & 	1157.52	 & 	-1.31924	 \\
10.2577	 & 	-1.31924	 & 	1279.94
\end{tabular}

\vtab
 EingenVectors - Molecule B     \\
\vtab
\begin{tabular}{|c c c|}
0.999922	 & 	0.00746944	 & 	-0.0100264	 \\
0.00735508	 & 	-0.999908	 & 	-0.0113945	 \\
0.0101106	 & 	-0.0113198	 & 	0.999885
\end{tabular}

\vtab
 EingenValues - Molecule B     \\
\vtab
\begin{tabular}{|c c c|}
257.939	 & 	1157.55	 & 	1280.06
\end{tabular}

\end{center}
\end{multicols}
\begin{center}
\vtab
\vtab
\textcolor{NavyBlue}{\Large Are Isomers}
\end{center}
%%%%%%%%%%%%%%%%%%%%%%%%%%%%%%%%%%%%%%%%%%%%%%%%%%%%%%%%
%%%%%%%%%%%%%%%%%%%%%%%%%%%%%%%%%%%%%%%%%%%%%%%%%%%%%%%%
\newpage
%%%%%%%%%%%%%%%%%%%%%%%%%%%%%%%%%%%%%%%%%%%%%%%%%%%%%%%%
%%%%%%%%%%%%%%%%%%%%%%%%%%%%%%%%%%%%%%%%%%%%%%%%%%%%%%%%

\vtab[-2cm]
\begin{center}
{\large Sugars}
\end{center}
\begin{multicols}{2}
\begin{center}
Molecule A \\ 
D\_-\_Fructuose\_out\_G09
\includegraphics[width=8cm]{Images/ImagesFromVMD/D_-_Fructuose_out_G09.jpg}
\\
Inertia Tensor - Molecule A \\
\vtab
\begin{tabular}{|c c c|}
263.78	 & 	3.63756	 & 	5.48351	 \\
3.63756	 & 	1174.62	 & 	1.15892	 \\
5.48351	 & 	1.15892	 & 	1226.18
\end{tabular}

\vtab
 EingenVectors - Molecule A     \\
\vtab
\begin{tabular}{|c c c|}
0.999976	 & 	-0.00398608	 & 	-0.00569256	 \\
0.00385485	 & 	0.999731	 & 	-0.0228795	 \\
0.00578223	 & 	0.022857	 & 	0.999722
\end{tabular}

\vtab
 EingenValues - Molecule A     \\
\vtab
\begin{tabular}{|c c c|}
263.734	 & 	1174.61	 & 	1226.23
\end{tabular}
\columnbreak

Molecule B \\ 
S-S-S-Fructuose\_out\_G09
\includegraphics[width=8cm]{Images/ImagesFromVMD/S-S-S-Fructuose_out_G09.jpg}
\\
Inertia Tensor - Molecule B \\
\vtab
\begin{tabular}{|c c c|}
226.452	 & 	13.0925	 & 	-2.23012	 \\
13.0925	 & 	1177.96	 & 	-0.925239	 \\
-2.23012	 & 	-0.925239	 & 	1359.88
\end{tabular}

\vtab
 EingenVectors - Molecule B     \\
\vtab
\begin{tabular}{|c c c|}
-0.999903	 & 	0.0137538	 & 	-0.00195585	 \\
0.0137434	 & 	0.999892	 & 	0.00525895	 \\
-0.00202797	 & 	-0.00523156	 & 	0.999984
\end{tabular}

\vtab
 EingenValues - Molecule B     \\
\vtab
\begin{tabular}{|c c c|}
226.267	 & 	1178.14	 & 	1359.89
\end{tabular}

\end{center}
\end{multicols}
\begin{center}
\vtab
\vtab
\textcolor{NavyBlue}{\Large Are Isomers}
\end{center}
%%%%%%%%%%%%%%%%%%%%%%%%%%%%%%%%%%%%%%%%%%%%%%%%%%%%%%%%
%%%%%%%%%%%%%%%%%%%%%%%%%%%%%%%%%%%%%%%%%%%%%%%%%%%%%%%%
\newpage
%%%%%%%%%%%%%%%%%%%%%%%%%%%%%%%%%%%%%%%%%%%%%%%%%%%%%%%%
%%%%%%%%%%%%%%%%%%%%%%%%%%%%%%%%%%%%%%%%%%%%%%%%%%%%%%%%

\vtab[-2cm]
\begin{center}
{\large Sugars}
\end{center}
\begin{multicols}{2}
\begin{center}
Molecule A \\ 
R-R-R-Fructuose\_out\_G09
\includegraphics[width=8cm]{Images/ImagesFromVMD/R-R-R-Fructuose_out_G09.jpg}
\\
Inertia Tensor - Molecule A \\
\vtab
\begin{tabular}{|c c c|}
270.815	 & 	9.92006	 & 	1.58147	 \\
9.92006	 & 	1105.83	 & 	-0.519818	 \\
1.58147	 & 	-0.519818	 & 	1246.52
\end{tabular}

\vtab
 EingenVectors - Molecule A     \\
\vtab
\begin{tabular}{|c c c|}
0.999928	 & 	-0.0118786	 & 	-0.00162686	 \\
-0.0118843	 & 	-0.999923	 & 	-0.00356382	 \\
0.0015844	 & 	-0.0035829	 & 	0.999992
\end{tabular}

\vtab
 EingenValues - Molecule A     \\
\vtab
\begin{tabular}{|c c c|}
270.695	 & 	1105.94	 & 	1246.52
\end{tabular}
\columnbreak

Molecule B \\ 
R-R-S-Fructuose\_out\_G09
\includegraphics[width=8cm]{Images/ImagesFromVMD/R-R-S-Fructuose_out_G09.jpg}
\\
Inertia Tensor - Molecule B \\
\vtab
\begin{tabular}{|c c c|}
258.077	 & 	6.73578	 & 	-10.2403	 \\
6.73578	 & 	1157.51	 & 	-1.31161	 \\
-10.2403	 & 	-1.31161	 & 	1279.98
\end{tabular}

\vtab
 EingenVectors - Molecule B     \\
\vtab
\begin{tabular}{|c c c|}
-0.999922	 & 	0.00747247	 & 	-0.0100089	 \\
0.00735899	 & 	0.999909	 & 	0.0113273	 \\
-0.0100927	 & 	-0.0112527	 & 	0.999886
\end{tabular}

\vtab
 EingenValues - Molecule B     \\
\vtab
\begin{tabular}{|c c c|}
257.924	 & 	1157.54	 & 	1280.1
\end{tabular}

\end{center}
\end{multicols}
\begin{center}
\vtab
\vtab
\textcolor{NavyBlue}{\Large Are Isomers}
\end{center}
%%%%%%%%%%%%%%%%%%%%%%%%%%%%%%%%%%%%%%%%%%%%%%%%%%%%%%%%
%%%%%%%%%%%%%%%%%%%%%%%%%%%%%%%%%%%%%%%%%%%%%%%%%%%%%%%%
\newpage
%%%%%%%%%%%%%%%%%%%%%%%%%%%%%%%%%%%%%%%%%%%%%%%%%%%%%%%%
%%%%%%%%%%%%%%%%%%%%%%%%%%%%%%%%%%%%%%%%%%%%%%%%%%%%%%%%

\vtab[-2cm]
\begin{center}
{\large Sugars}
\end{center}
\begin{multicols}{2}
\begin{center}
Molecule A \\ 
R-R-R-Fructuose\_out\_G09
\includegraphics[width=8cm]{Images/ImagesFromVMD/R-R-R-Fructuose_out_G09.jpg}
\\
Inertia Tensor - Molecule A \\
\vtab
\begin{tabular}{|c c c|}
270.815	 & 	9.92006	 & 	1.58147	 \\
9.92006	 & 	1105.83	 & 	-0.519818	 \\
1.58147	 & 	-0.519818	 & 	1246.52
\end{tabular}

\vtab
 EingenVectors - Molecule A     \\
\vtab
\begin{tabular}{|c c c|}
0.999928	 & 	-0.0118786	 & 	-0.00162686	 \\
-0.0118843	 & 	-0.999923	 & 	-0.00356382	 \\
0.0015844	 & 	-0.0035829	 & 	0.999992
\end{tabular}

\vtab
 EingenValues - Molecule A     \\
\vtab
\begin{tabular}{|c c c|}
270.695	 & 	1105.94	 & 	1246.52
\end{tabular}
\columnbreak

Molecule B \\ 
R-S-R-Fructuose\_out\_G09
\includegraphics[width=8cm]{Images/ImagesFromVMD/R-S-R-Fructuose_out_G09.jpg}
\\
Inertia Tensor - Molecule B \\
\vtab
\begin{tabular}{|c c c|}
202.772	 & 	14.8494	 & 	0.184396	 \\
14.8494	 & 	1218.4	 & 	-0.831783	 \\
0.184396	 & 	-0.831783	 & 	1335.51
\end{tabular}

\vtab
 EingenVectors - Molecule B     \\
\vtab
\begin{tabular}{|c c c|}
-0.999893	 & 	0.0146163	 & 	0.00017347	 \\
-0.0146172	 & 	-0.999868	 & 	-0.00709174	 \\
6.97925e-05	 & 	-0.00709352	 & 	0.999975
\end{tabular}

\vtab
 EingenValues - Molecule B     \\
\vtab
\begin{tabular}{|c c c|}
202.555	 & 	1218.61	 & 	1335.51
\end{tabular}

\end{center}
\end{multicols}
\begin{center}
\vtab
\vtab
\textcolor{NavyBlue}{\Large Are Isomers}
\end{center}
%%%%%%%%%%%%%%%%%%%%%%%%%%%%%%%%%%%%%%%%%%%%%%%%%%%%%%%%
%%%%%%%%%%%%%%%%%%%%%%%%%%%%%%%%%%%%%%%%%%%%%%%%%%%%%%%%
\newpage
%%%%%%%%%%%%%%%%%%%%%%%%%%%%%%%%%%%%%%%%%%%%%%%%%%%%%%%%
%%%%%%%%%%%%%%%%%%%%%%%%%%%%%%%%%%%%%%%%%%%%%%%%%%%%%%%%

\vtab[-2cm]
\begin{center}
{\large Sugars}
\end{center}
\begin{multicols}{2}
\begin{center}
Molecule A \\ 
R-R-R-Fructuose\_out\_G09
\includegraphics[width=8cm]{Images/ImagesFromVMD/R-R-R-Fructuose_out_G09.jpg}
\\
Inertia Tensor - Molecule A \\
\vtab
\begin{tabular}{|c c c|}
270.815	 & 	9.92006	 & 	1.58147	 \\
9.92006	 & 	1105.83	 & 	-0.519818	 \\
1.58147	 & 	-0.519818	 & 	1246.52
\end{tabular}

\vtab
 EingenVectors - Molecule A     \\
\vtab
\begin{tabular}{|c c c|}
0.999928	 & 	-0.0118786	 & 	-0.00162686	 \\
-0.0118843	 & 	-0.999923	 & 	-0.00356382	 \\
0.0015844	 & 	-0.0035829	 & 	0.999992
\end{tabular}

\vtab
 EingenValues - Molecule A     \\
\vtab
\begin{tabular}{|c c c|}
270.695	 & 	1105.94	 & 	1246.52
\end{tabular}
\columnbreak

Molecule B \\ 
R-S-S-Fructuose\_out\_G09
\includegraphics[width=8cm]{Images/ImagesFromVMD/R-S-S-Fructuose_out_G09.jpg}
\\
Inertia Tensor - Molecule B \\
\vtab
\begin{tabular}{|c c c|}
215.024	 & 	-14.0942	 & 	-5.97678	 \\
-14.0942	 & 	1187.1	 & 	0.653719	 \\
-5.97678	 & 	0.653719	 & 	1326.43
\end{tabular}

\vtab
 EingenVectors - Molecule B     \\
\vtab
\begin{tabular}{|c c c|}
0.999881	 & 	0.0144902	 & 	0.00536739	 \\
0.0144617	 & 	-0.999881	 & 	0.0053195	 \\
-0.00544384	 & 	0.00524124	 & 	0.999971
\end{tabular}

\vtab
 EingenValues - Molecule B     \\
\vtab
\begin{tabular}{|c c c|}
214.787	 & 	1187.3	 & 	1326.46
\end{tabular}

\end{center}
\end{multicols}
\begin{center}
\vtab
\vtab
\textcolor{NavyBlue}{\Large Are Isomers}
\end{center}
%%%%%%%%%%%%%%%%%%%%%%%%%%%%%%%%%%%%%%%%%%%%%%%%%%%%%%%%
%%%%%%%%%%%%%%%%%%%%%%%%%%%%%%%%%%%%%%%%%%%%%%%%%%%%%%%%
\newpage
%%%%%%%%%%%%%%%%%%%%%%%%%%%%%%%%%%%%%%%%%%%%%%%%%%%%%%%%
%%%%%%%%%%%%%%%%%%%%%%%%%%%%%%%%%%%%%%%%%%%%%%%%%%%%%%%%

\vtab[-2cm]
\begin{center}
{\large Sugars}
\end{center}
\begin{multicols}{2}
\begin{center}
Molecule A \\ 
R-R-R-Fructuose\_out\_G09
\includegraphics[width=8cm]{Images/ImagesFromVMD/R-R-R-Fructuose_out_G09.jpg}
\\
Inertia Tensor - Molecule A \\
\vtab
\begin{tabular}{|c c c|}
270.815	 & 	9.92006	 & 	1.58147	 \\
9.92006	 & 	1105.83	 & 	-0.519818	 \\
1.58147	 & 	-0.519818	 & 	1246.52
\end{tabular}

\vtab
 EingenVectors - Molecule A     \\
\vtab
\begin{tabular}{|c c c|}
0.999928	 & 	-0.0118786	 & 	-0.00162686	 \\
-0.0118843	 & 	-0.999923	 & 	-0.00356382	 \\
0.0015844	 & 	-0.0035829	 & 	0.999992
\end{tabular}

\vtab
 EingenValues - Molecule A     \\
\vtab
\begin{tabular}{|c c c|}
270.695	 & 	1105.94	 & 	1246.52
\end{tabular}
\columnbreak

Molecule B \\ 
S-R-R-Fructuose\_out\_G09
\includegraphics[width=8cm]{Images/ImagesFromVMD/S-R-R-Fructuose_out_G09.jpg}
\\
Inertia Tensor - Molecule B \\
\vtab
\begin{tabular}{|c c c|}
267.669	 & 	-2.83987	 & 	8.0125	 \\
-2.83987	 & 	1173.13	 & 	0.735007	 \\
8.0125	 & 	0.735007	 & 	1229.07
\end{tabular}

\vtab
 EingenVectors - Molecule B     \\
\vtab
\begin{tabular}{|c c c|}
0.99996	 & 	0.00314275	 & 	-0.00833561	 \\
-0.00324825	 & 	0.999914	 & 	-0.0126739	 \\
0.00829507	 & 	0.0127005	 & 	0.999885
\end{tabular}

\vtab
 EingenValues - Molecule B     \\
\vtab
\begin{tabular}{|c c c|}
267.593	 & 	1173.13	 & 	1229.15
\end{tabular}

\end{center}
\end{multicols}
\begin{center}
\vtab
\vtab
\textcolor{NavyBlue}{\Large Are Isomers}
\end{center}
%%%%%%%%%%%%%%%%%%%%%%%%%%%%%%%%%%%%%%%%%%%%%%%%%%%%%%%%
%%%%%%%%%%%%%%%%%%%%%%%%%%%%%%%%%%%%%%%%%%%%%%%%%%%%%%%%
\newpage
%%%%%%%%%%%%%%%%%%%%%%%%%%%%%%%%%%%%%%%%%%%%%%%%%%%%%%%%
%%%%%%%%%%%%%%%%%%%%%%%%%%%%%%%%%%%%%%%%%%%%%%%%%%%%%%%%

\vtab[-2cm]
\begin{center}
{\large Sugars}
\end{center}
\begin{multicols}{2}
\begin{center}
Molecule A \\ 
R-R-R-Fructuose\_out\_G09
\includegraphics[width=8cm]{Images/ImagesFromVMD/R-R-R-Fructuose_out_G09.jpg}
\\
Inertia Tensor - Molecule A \\
\vtab
\begin{tabular}{|c c c|}
270.815	 & 	9.92006	 & 	1.58147	 \\
9.92006	 & 	1105.83	 & 	-0.519818	 \\
1.58147	 & 	-0.519818	 & 	1246.52
\end{tabular}

\vtab
 EingenVectors - Molecule A     \\
\vtab
\begin{tabular}{|c c c|}
0.999928	 & 	-0.0118786	 & 	-0.00162686	 \\
-0.0118843	 & 	-0.999923	 & 	-0.00356382	 \\
0.0015844	 & 	-0.0035829	 & 	0.999992
\end{tabular}

\vtab
 EingenValues - Molecule A     \\
\vtab
\begin{tabular}{|c c c|}
270.695	 & 	1105.94	 & 	1246.52
\end{tabular}
\columnbreak

Molecule B \\ 
S-R-S-Fructuose\_out\_G09
\includegraphics[width=8cm]{Images/ImagesFromVMD/S-R-S-Fructuose_out_G09.jpg}
\\
Inertia Tensor - Molecule B \\
\vtab
\begin{tabular}{|c c c|}
202.782	 & 	-14.8472	 & 	-0.161586	 \\
-14.8472	 & 	1218.41	 & 	-0.836204	 \\
-0.161586	 & 	-0.836204	 & 	1335.52
\end{tabular}

\vtab
 EingenVectors - Molecule B     \\
\vtab
\begin{tabular}{|c c c|}
-0.999893	 & 	-0.0146142	 & 	-0.000153395	 \\
0.0146149	 & 	-0.999868	 & 	-0.00713203	 \\
-4.91458e-05	 & 	-0.00713351	 & 	0.999975
\end{tabular}

\vtab
 EingenValues - Molecule B     \\
\vtab
\begin{tabular}{|c c c|}
202.565	 & 	1218.62	 & 	1335.52
\end{tabular}

\end{center}
\end{multicols}
\begin{center}
\vtab
\vtab
\textcolor{NavyBlue}{\Large Are Isomers}
\end{center}
%%%%%%%%%%%%%%%%%%%%%%%%%%%%%%%%%%%%%%%%%%%%%%%%%%%%%%%%
%%%%%%%%%%%%%%%%%%%%%%%%%%%%%%%%%%%%%%%%%%%%%%%%%%%%%%%%
\newpage
%%%%%%%%%%%%%%%%%%%%%%%%%%%%%%%%%%%%%%%%%%%%%%%%%%%%%%%%
%%%%%%%%%%%%%%%%%%%%%%%%%%%%%%%%%%%%%%%%%%%%%%%%%%%%%%%%

\vtab[-2cm]
\begin{center}
{\large Sugars}
\end{center}
\begin{multicols}{2}
\begin{center}
Molecule A \\ 
R-R-R-Fructuose\_out\_G09
\includegraphics[width=8cm]{Images/ImagesFromVMD/R-R-R-Fructuose_out_G09.jpg}
\\
Inertia Tensor - Molecule A \\
\vtab
\begin{tabular}{|c c c|}
270.815	 & 	9.92006	 & 	1.58147	 \\
9.92006	 & 	1105.83	 & 	-0.519818	 \\
1.58147	 & 	-0.519818	 & 	1246.52
\end{tabular}

\vtab
 EingenVectors - Molecule A     \\
\vtab
\begin{tabular}{|c c c|}
0.999928	 & 	-0.0118786	 & 	-0.00162686	 \\
-0.0118843	 & 	-0.999923	 & 	-0.00356382	 \\
0.0015844	 & 	-0.0035829	 & 	0.999992
\end{tabular}

\vtab
 EingenValues - Molecule A     \\
\vtab
\begin{tabular}{|c c c|}
270.695	 & 	1105.94	 & 	1246.52
\end{tabular}
\columnbreak

Molecule B \\ 
S-S-R-Fructuose\_out\_G09
\includegraphics[width=8cm]{Images/ImagesFromVMD/S-S-R-Fructuose_out_G09.jpg}
\\
Inertia Tensor - Molecule B \\
\vtab
\begin{tabular}{|c c c|}
258.093	 & 	-6.7331	 & 	10.2577	 \\
-6.7331	 & 	1157.52	 & 	-1.31924	 \\
10.2577	 & 	-1.31924	 & 	1279.94
\end{tabular}

\vtab
 EingenVectors - Molecule B     \\
\vtab
\begin{tabular}{|c c c|}
0.999922	 & 	0.00746944	 & 	-0.0100264	 \\
0.00735508	 & 	-0.999908	 & 	-0.0113945	 \\
0.0101106	 & 	-0.0113198	 & 	0.999885
\end{tabular}

\vtab
 EingenValues - Molecule B     \\
\vtab
\begin{tabular}{|c c c|}
257.939	 & 	1157.55	 & 	1280.06
\end{tabular}

\end{center}
\end{multicols}
\begin{center}
\vtab
\vtab
\textcolor{NavyBlue}{\Large Are Isomers}
\end{center}
%%%%%%%%%%%%%%%%%%%%%%%%%%%%%%%%%%%%%%%%%%%%%%%%%%%%%%%%
%%%%%%%%%%%%%%%%%%%%%%%%%%%%%%%%%%%%%%%%%%%%%%%%%%%%%%%%
\newpage
%%%%%%%%%%%%%%%%%%%%%%%%%%%%%%%%%%%%%%%%%%%%%%%%%%%%%%%%
%%%%%%%%%%%%%%%%%%%%%%%%%%%%%%%%%%%%%%%%%%%%%%%%%%%%%%%%

\vtab[-2cm]
\begin{center}
{\large Sugars}
\end{center}
\begin{multicols}{2}
\begin{center}
Molecule A \\ 
R-R-R-Fructuose\_out\_G09
\includegraphics[width=8cm]{Images/ImagesFromVMD/R-R-R-Fructuose_out_G09.jpg}
\\
Inertia Tensor - Molecule A \\
\vtab
\begin{tabular}{|c c c|}
270.815	 & 	9.92006	 & 	1.58147	 \\
9.92006	 & 	1105.83	 & 	-0.519818	 \\
1.58147	 & 	-0.519818	 & 	1246.52
\end{tabular}

\vtab
 EingenVectors - Molecule A     \\
\vtab
\begin{tabular}{|c c c|}
0.999928	 & 	-0.0118786	 & 	-0.00162686	 \\
-0.0118843	 & 	-0.999923	 & 	-0.00356382	 \\
0.0015844	 & 	-0.0035829	 & 	0.999992
\end{tabular}

\vtab
 EingenValues - Molecule A     \\
\vtab
\begin{tabular}{|c c c|}
270.695	 & 	1105.94	 & 	1246.52
\end{tabular}
\columnbreak

Molecule B \\ 
S-S-S-Fructuose\_out\_G09
\includegraphics[width=8cm]{Images/ImagesFromVMD/S-S-S-Fructuose_out_G09.jpg}
\\
Inertia Tensor - Molecule B \\
\vtab
\begin{tabular}{|c c c|}
226.452	 & 	13.0925	 & 	-2.23012	 \\
13.0925	 & 	1177.96	 & 	-0.925239	 \\
-2.23012	 & 	-0.925239	 & 	1359.88
\end{tabular}

\vtab
 EingenVectors - Molecule B     \\
\vtab
\begin{tabular}{|c c c|}
-0.999903	 & 	0.0137538	 & 	-0.00195585	 \\
0.0137434	 & 	0.999892	 & 	0.00525895	 \\
-0.00202797	 & 	-0.00523156	 & 	0.999984
\end{tabular}

\vtab
 EingenValues - Molecule B     \\
\vtab
\begin{tabular}{|c c c|}
226.267	 & 	1178.14	 & 	1359.89
\end{tabular}

\end{center}
\end{multicols}
\begin{center}
\vtab
\vtab
\textcolor{NavyBlue}{\Large Are Isomers}
\end{center}
%%%%%%%%%%%%%%%%%%%%%%%%%%%%%%%%%%%%%%%%%%%%%%%%%%%%%%%%
%%%%%%%%%%%%%%%%%%%%%%%%%%%%%%%%%%%%%%%%%%%%%%%%%%%%%%%%
\newpage
%%%%%%%%%%%%%%%%%%%%%%%%%%%%%%%%%%%%%%%%%%%%%%%%%%%%%%%%
%%%%%%%%%%%%%%%%%%%%%%%%%%%%%%%%%%%%%%%%%%%%%%%%%%%%%%%%

\vtab[-2cm]
\begin{center}
{\large Sugars}
\end{center}
\begin{multicols}{2}
\begin{center}
Molecule A \\ 
R-R-S-Fructuose\_out\_G09
\includegraphics[width=8cm]{Images/ImagesFromVMD/R-R-S-Fructuose_out_G09.jpg}
\\
Inertia Tensor - Molecule A \\
\vtab
\begin{tabular}{|c c c|}
258.077	 & 	6.73578	 & 	-10.2403	 \\
6.73578	 & 	1157.51	 & 	-1.31161	 \\
-10.2403	 & 	-1.31161	 & 	1279.98
\end{tabular}

\vtab
 EingenVectors - Molecule A     \\
\vtab
\begin{tabular}{|c c c|}
-0.999922	 & 	0.00747247	 & 	-0.0100089	 \\
0.00735899	 & 	0.999909	 & 	0.0113273	 \\
-0.0100927	 & 	-0.0112527	 & 	0.999886
\end{tabular}

\vtab
 EingenValues - Molecule A     \\
\vtab
\begin{tabular}{|c c c|}
257.924	 & 	1157.54	 & 	1280.1
\end{tabular}
\columnbreak

Molecule B \\ 
R-S-R-Fructuose\_out\_G09
\includegraphics[width=8cm]{Images/ImagesFromVMD/R-S-R-Fructuose_out_G09.jpg}
\\
Inertia Tensor - Molecule B \\
\vtab
\begin{tabular}{|c c c|}
202.772	 & 	14.8494	 & 	0.184396	 \\
14.8494	 & 	1218.4	 & 	-0.831783	 \\
0.184396	 & 	-0.831783	 & 	1335.51
\end{tabular}

\vtab
 EingenVectors - Molecule B     \\
\vtab
\begin{tabular}{|c c c|}
-0.999893	 & 	0.0146163	 & 	0.00017347	 \\
-0.0146172	 & 	-0.999868	 & 	-0.00709174	 \\
6.97925e-05	 & 	-0.00709352	 & 	0.999975
\end{tabular}

\vtab
 EingenValues - Molecule B     \\
\vtab
\begin{tabular}{|c c c|}
202.555	 & 	1218.61	 & 	1335.51
\end{tabular}

\end{center}
\end{multicols}
\begin{center}
\vtab
\vtab
\textcolor{NavyBlue}{\Large Are Isomers}
\end{center}
%%%%%%%%%%%%%%%%%%%%%%%%%%%%%%%%%%%%%%%%%%%%%%%%%%%%%%%%
%%%%%%%%%%%%%%%%%%%%%%%%%%%%%%%%%%%%%%%%%%%%%%%%%%%%%%%%
\newpage
%%%%%%%%%%%%%%%%%%%%%%%%%%%%%%%%%%%%%%%%%%%%%%%%%%%%%%%%
%%%%%%%%%%%%%%%%%%%%%%%%%%%%%%%%%%%%%%%%%%%%%%%%%%%%%%%%

\vtab[-2cm]
\begin{center}
{\large Sugars}
\end{center}
\begin{multicols}{2}
\begin{center}
Molecule A \\ 
R-R-S-Fructuose\_out\_G09
\includegraphics[width=8cm]{Images/ImagesFromVMD/R-R-S-Fructuose_out_G09.jpg}
\\
Inertia Tensor - Molecule A \\
\vtab
\begin{tabular}{|c c c|}
258.077	 & 	6.73578	 & 	-10.2403	 \\
6.73578	 & 	1157.51	 & 	-1.31161	 \\
-10.2403	 & 	-1.31161	 & 	1279.98
\end{tabular}

\vtab
 EingenVectors - Molecule A     \\
\vtab
\begin{tabular}{|c c c|}
-0.999922	 & 	0.00747247	 & 	-0.0100089	 \\
0.00735899	 & 	0.999909	 & 	0.0113273	 \\
-0.0100927	 & 	-0.0112527	 & 	0.999886
\end{tabular}

\vtab
 EingenValues - Molecule A     \\
\vtab
\begin{tabular}{|c c c|}
257.924	 & 	1157.54	 & 	1280.1
\end{tabular}
\columnbreak

Molecule B \\ 
R-S-S-Fructuose\_out\_G09
\includegraphics[width=8cm]{Images/ImagesFromVMD/R-S-S-Fructuose_out_G09.jpg}
\\
Inertia Tensor - Molecule B \\
\vtab
\begin{tabular}{|c c c|}
215.024	 & 	-14.0942	 & 	-5.97678	 \\
-14.0942	 & 	1187.1	 & 	0.653719	 \\
-5.97678	 & 	0.653719	 & 	1326.43
\end{tabular}

\vtab
 EingenVectors - Molecule B     \\
\vtab
\begin{tabular}{|c c c|}
0.999881	 & 	0.0144902	 & 	0.00536739	 \\
0.0144617	 & 	-0.999881	 & 	0.0053195	 \\
-0.00544384	 & 	0.00524124	 & 	0.999971
\end{tabular}

\vtab
 EingenValues - Molecule B     \\
\vtab
\begin{tabular}{|c c c|}
214.787	 & 	1187.3	 & 	1326.46
\end{tabular}

\end{center}
\end{multicols}
\begin{center}
\vtab
\vtab
\textcolor{NavyBlue}{\Large Are Isomers}
\end{center}
%%%%%%%%%%%%%%%%%%%%%%%%%%%%%%%%%%%%%%%%%%%%%%%%%%%%%%%%
%%%%%%%%%%%%%%%%%%%%%%%%%%%%%%%%%%%%%%%%%%%%%%%%%%%%%%%%
\newpage
%%%%%%%%%%%%%%%%%%%%%%%%%%%%%%%%%%%%%%%%%%%%%%%%%%%%%%%%
%%%%%%%%%%%%%%%%%%%%%%%%%%%%%%%%%%%%%%%%%%%%%%%%%%%%%%%%

\vtab[-2cm]
\begin{center}
{\large Sugars}
\end{center}
\begin{multicols}{2}
\begin{center}
Molecule A \\ 
R-R-S-Fructuose\_out\_G09
\includegraphics[width=8cm]{Images/ImagesFromVMD/R-R-S-Fructuose_out_G09.jpg}
\\
Inertia Tensor - Molecule A \\
\vtab
\begin{tabular}{|c c c|}
258.077	 & 	6.73578	 & 	-10.2403	 \\
6.73578	 & 	1157.51	 & 	-1.31161	 \\
-10.2403	 & 	-1.31161	 & 	1279.98
\end{tabular}

\vtab
 EingenVectors - Molecule A     \\
\vtab
\begin{tabular}{|c c c|}
-0.999922	 & 	0.00747247	 & 	-0.0100089	 \\
0.00735899	 & 	0.999909	 & 	0.0113273	 \\
-0.0100927	 & 	-0.0112527	 & 	0.999886
\end{tabular}

\vtab
 EingenValues - Molecule A     \\
\vtab
\begin{tabular}{|c c c|}
257.924	 & 	1157.54	 & 	1280.1
\end{tabular}
\columnbreak

Molecule B \\ 
S-R-R-Fructuose\_out\_G09
\includegraphics[width=8cm]{Images/ImagesFromVMD/S-R-R-Fructuose_out_G09.jpg}
\\
Inertia Tensor - Molecule B \\
\vtab
\begin{tabular}{|c c c|}
267.669	 & 	-2.83987	 & 	8.0125	 \\
-2.83987	 & 	1173.13	 & 	0.735007	 \\
8.0125	 & 	0.735007	 & 	1229.07
\end{tabular}

\vtab
 EingenVectors - Molecule B     \\
\vtab
\begin{tabular}{|c c c|}
0.99996	 & 	0.00314275	 & 	-0.00833561	 \\
-0.00324825	 & 	0.999914	 & 	-0.0126739	 \\
0.00829507	 & 	0.0127005	 & 	0.999885
\end{tabular}

\vtab
 EingenValues - Molecule B     \\
\vtab
\begin{tabular}{|c c c|}
267.593	 & 	1173.13	 & 	1229.15
\end{tabular}

\end{center}
\end{multicols}
\begin{center}
\vtab
\vtab
\textcolor{NavyBlue}{\Large Are Isomers}
\end{center}
%%%%%%%%%%%%%%%%%%%%%%%%%%%%%%%%%%%%%%%%%%%%%%%%%%%%%%%%
%%%%%%%%%%%%%%%%%%%%%%%%%%%%%%%%%%%%%%%%%%%%%%%%%%%%%%%%
\newpage
%%%%%%%%%%%%%%%%%%%%%%%%%%%%%%%%%%%%%%%%%%%%%%%%%%%%%%%%
%%%%%%%%%%%%%%%%%%%%%%%%%%%%%%%%%%%%%%%%%%%%%%%%%%%%%%%%

\vtab[-2cm]
\begin{center}
{\large Sugars}
\end{center}
\begin{multicols}{2}
\begin{center}
Molecule A \\ 
R-R-S-Fructuose\_out\_G09
\includegraphics[width=8cm]{Images/ImagesFromVMD/R-R-S-Fructuose_out_G09.jpg}
\\
Inertia Tensor - Molecule A \\
\vtab
\begin{tabular}{|c c c|}
258.077	 & 	6.73578	 & 	-10.2403	 \\
6.73578	 & 	1157.51	 & 	-1.31161	 \\
-10.2403	 & 	-1.31161	 & 	1279.98
\end{tabular}

\vtab
 EingenVectors - Molecule A     \\
\vtab
\begin{tabular}{|c c c|}
-0.999922	 & 	0.00747247	 & 	-0.0100089	 \\
0.00735899	 & 	0.999909	 & 	0.0113273	 \\
-0.0100927	 & 	-0.0112527	 & 	0.999886
\end{tabular}

\vtab
 EingenValues - Molecule A     \\
\vtab
\begin{tabular}{|c c c|}
257.924	 & 	1157.54	 & 	1280.1
\end{tabular}
\columnbreak

Molecule B \\ 
S-R-S-Fructuose\_out\_G09
\includegraphics[width=8cm]{Images/ImagesFromVMD/S-R-S-Fructuose_out_G09.jpg}
\\
Inertia Tensor - Molecule B \\
\vtab
\begin{tabular}{|c c c|}
202.782	 & 	-14.8472	 & 	-0.161586	 \\
-14.8472	 & 	1218.41	 & 	-0.836204	 \\
-0.161586	 & 	-0.836204	 & 	1335.52
\end{tabular}

\vtab
 EingenVectors - Molecule B     \\
\vtab
\begin{tabular}{|c c c|}
-0.999893	 & 	-0.0146142	 & 	-0.000153395	 \\
0.0146149	 & 	-0.999868	 & 	-0.00713203	 \\
-4.91458e-05	 & 	-0.00713351	 & 	0.999975
\end{tabular}

\vtab
 EingenValues - Molecule B     \\
\vtab
\begin{tabular}{|c c c|}
202.565	 & 	1218.62	 & 	1335.52
\end{tabular}

\end{center}
\end{multicols}
\begin{center}
\vtab
\vtab
\textcolor{NavyBlue}{\Large Are Isomers}
\end{center}
%%%%%%%%%%%%%%%%%%%%%%%%%%%%%%%%%%%%%%%%%%%%%%%%%%%%%%%%
%%%%%%%%%%%%%%%%%%%%%%%%%%%%%%%%%%%%%%%%%%%%%%%%%%%%%%%%
\newpage
%%%%%%%%%%%%%%%%%%%%%%%%%%%%%%%%%%%%%%%%%%%%%%%%%%%%%%%%
%%%%%%%%%%%%%%%%%%%%%%%%%%%%%%%%%%%%%%%%%%%%%%%%%%%%%%%%

\vtab[-2cm]
\begin{center}
{\large Sugars}
\end{center}
\begin{multicols}{2}
\begin{center}
Molecule A \\ 
R-R-S-Fructuose\_out\_G09
\includegraphics[width=8cm]{Images/ImagesFromVMD/R-R-S-Fructuose_out_G09.jpg}
\\
Inertia Tensor - Molecule A \\
\vtab
\begin{tabular}{|c c c|}
258.077	 & 	6.73578	 & 	-10.2403	 \\
6.73578	 & 	1157.51	 & 	-1.31161	 \\
-10.2403	 & 	-1.31161	 & 	1279.98
\end{tabular}

\vtab
 EingenVectors - Molecule A     \\
\vtab
\begin{tabular}{|c c c|}
-0.999922	 & 	0.00747247	 & 	-0.0100089	 \\
0.00735899	 & 	0.999909	 & 	0.0113273	 \\
-0.0100927	 & 	-0.0112527	 & 	0.999886
\end{tabular}

\vtab
 EingenValues - Molecule A     \\
\vtab
\begin{tabular}{|c c c|}
257.924	 & 	1157.54	 & 	1280.1
\end{tabular}
\columnbreak

Molecule B \\ 
S-S-R-Fructuose\_out\_G09
\includegraphics[width=8cm]{Images/ImagesFromVMD/S-S-R-Fructuose_out_G09.jpg}
\\
Inertia Tensor - Molecule B \\
\vtab
\begin{tabular}{|c c c|}
258.093	 & 	-6.7331	 & 	10.2577	 \\
-6.7331	 & 	1157.52	 & 	-1.31924	 \\
10.2577	 & 	-1.31924	 & 	1279.94
\end{tabular}

\vtab
 EingenVectors - Molecule B     \\
\vtab
\begin{tabular}{|c c c|}
0.999922	 & 	0.00746944	 & 	-0.0100264	 \\
0.00735508	 & 	-0.999908	 & 	-0.0113945	 \\
0.0101106	 & 	-0.0113198	 & 	0.999885
\end{tabular}

\vtab
 EingenValues - Molecule B     \\
\vtab
\begin{tabular}{|c c c|}
257.939	 & 	1157.55	 & 	1280.06
\end{tabular}

\end{center}
\end{multicols}
\begin{center}
\vtab
\vtab
\textcolor{NavyBlue}{\Large Are Isomer and Stereoisomers}
\end{center}
%%%%%%%%%%%%%%%%%%%%%%%%%%%%%%%%%%%%%%%%%%%%%%%%%%%%%%%%
%%%%%%%%%%%%%%%%%%%%%%%%%%%%%%%%%%%%%%%%%%%%%%%%%%%%%%%%
\newpage
%%%%%%%%%%%%%%%%%%%%%%%%%%%%%%%%%%%%%%%%%%%%%%%%%%%%%%%%
%%%%%%%%%%%%%%%%%%%%%%%%%%%%%%%%%%%%%%%%%%%%%%%%%%%%%%%%

\vtab[-2cm]
\begin{center}
{\large Sugars}
\end{center}
\begin{multicols}{2}
\begin{center}
Molecule A \\ 
R-R-S-Fructuose\_out\_G09
\includegraphics[width=8cm]{Images/ImagesFromVMD/R-R-S-Fructuose_out_G09.jpg}
\\
Inertia Tensor - Molecule A \\
\vtab
\begin{tabular}{|c c c|}
258.077	 & 	6.73578	 & 	-10.2403	 \\
6.73578	 & 	1157.51	 & 	-1.31161	 \\
-10.2403	 & 	-1.31161	 & 	1279.98
\end{tabular}

\vtab
 EingenVectors - Molecule A     \\
\vtab
\begin{tabular}{|c c c|}
-0.999922	 & 	0.00747247	 & 	-0.0100089	 \\
0.00735899	 & 	0.999909	 & 	0.0113273	 \\
-0.0100927	 & 	-0.0112527	 & 	0.999886
\end{tabular}

\vtab
 EingenValues - Molecule A     \\
\vtab
\begin{tabular}{|c c c|}
257.924	 & 	1157.54	 & 	1280.1
\end{tabular}
\columnbreak

Molecule B \\ 
S-S-S-Fructuose\_out\_G09
\includegraphics[width=8cm]{Images/ImagesFromVMD/S-S-S-Fructuose_out_G09.jpg}
\\
Inertia Tensor - Molecule B \\
\vtab
\begin{tabular}{|c c c|}
226.452	 & 	13.0925	 & 	-2.23012	 \\
13.0925	 & 	1177.96	 & 	-0.925239	 \\
-2.23012	 & 	-0.925239	 & 	1359.88
\end{tabular}

\vtab
 EingenVectors - Molecule B     \\
\vtab
\begin{tabular}{|c c c|}
-0.999903	 & 	0.0137538	 & 	-0.00195585	 \\
0.0137434	 & 	0.999892	 & 	0.00525895	 \\
-0.00202797	 & 	-0.00523156	 & 	0.999984
\end{tabular}

\vtab
 EingenValues - Molecule B     \\
\vtab
\begin{tabular}{|c c c|}
226.267	 & 	1178.14	 & 	1359.89
\end{tabular}

\end{center}
\end{multicols}
\begin{center}
\vtab
\vtab
\textcolor{NavyBlue}{\Large Are Isomers}
\end{center}
%%%%%%%%%%%%%%%%%%%%%%%%%%%%%%%%%%%%%%%%%%%%%%%%%%%%%%%%
%%%%%%%%%%%%%%%%%%%%%%%%%%%%%%%%%%%%%%%%%%%%%%%%%%%%%%%%
\newpage
%%%%%%%%%%%%%%%%%%%%%%%%%%%%%%%%%%%%%%%%%%%%%%%%%%%%%%%%
%%%%%%%%%%%%%%%%%%%%%%%%%%%%%%%%%%%%%%%%%%%%%%%%%%%%%%%%

\vtab[-2cm]
\begin{center}
{\large Sugars}
\end{center}
\begin{multicols}{2}
\begin{center}
Molecule A \\ 
R-S-R-Fructuose\_out\_G09
\includegraphics[width=8cm]{Images/ImagesFromVMD/R-S-R-Fructuose_out_G09.jpg}
\\
Inertia Tensor - Molecule A \\
\vtab
\begin{tabular}{|c c c|}
202.772	 & 	14.8494	 & 	0.184396	 \\
14.8494	 & 	1218.4	 & 	-0.831783	 \\
0.184396	 & 	-0.831783	 & 	1335.51
\end{tabular}

\vtab
 EingenVectors - Molecule A     \\
\vtab
\begin{tabular}{|c c c|}
-0.999893	 & 	0.0146163	 & 	0.00017347	 \\
-0.0146172	 & 	-0.999868	 & 	-0.00709174	 \\
6.97925e-05	 & 	-0.00709352	 & 	0.999975
\end{tabular}

\vtab
 EingenValues - Molecule A     \\
\vtab
\begin{tabular}{|c c c|}
202.555	 & 	1218.61	 & 	1335.51
\end{tabular}
\columnbreak

Molecule B \\ 
R-S-S-Fructuose\_out\_G09
\includegraphics[width=8cm]{Images/ImagesFromVMD/R-S-S-Fructuose_out_G09.jpg}
\\
Inertia Tensor - Molecule B \\
\vtab
\begin{tabular}{|c c c|}
215.024	 & 	-14.0942	 & 	-5.97678	 \\
-14.0942	 & 	1187.1	 & 	0.653719	 \\
-5.97678	 & 	0.653719	 & 	1326.43
\end{tabular}

\vtab
 EingenVectors - Molecule B     \\
\vtab
\begin{tabular}{|c c c|}
0.999881	 & 	0.0144902	 & 	0.00536739	 \\
0.0144617	 & 	-0.999881	 & 	0.0053195	 \\
-0.00544384	 & 	0.00524124	 & 	0.999971
\end{tabular}

\vtab
 EingenValues - Molecule B     \\
\vtab
\begin{tabular}{|c c c|}
214.787	 & 	1187.3	 & 	1326.46
\end{tabular}

\end{center}
\end{multicols}
\begin{center}
\vtab
\vtab
\textcolor{NavyBlue}{\Large Are Isomers}
\end{center}
%%%%%%%%%%%%%%%%%%%%%%%%%%%%%%%%%%%%%%%%%%%%%%%%%%%%%%%%
%%%%%%%%%%%%%%%%%%%%%%%%%%%%%%%%%%%%%%%%%%%%%%%%%%%%%%%%
\newpage
%%%%%%%%%%%%%%%%%%%%%%%%%%%%%%%%%%%%%%%%%%%%%%%%%%%%%%%%
%%%%%%%%%%%%%%%%%%%%%%%%%%%%%%%%%%%%%%%%%%%%%%%%%%%%%%%%

\vtab[-2cm]
\begin{center}
{\large Sugars}
\end{center}
\begin{multicols}{2}
\begin{center}
Molecule A \\ 
R-S-R-Fructuose\_out\_G09
\includegraphics[width=8cm]{Images/ImagesFromVMD/R-S-R-Fructuose_out_G09.jpg}
\\
Inertia Tensor - Molecule A \\
\vtab
\begin{tabular}{|c c c|}
202.772	 & 	14.8494	 & 	0.184396	 \\
14.8494	 & 	1218.4	 & 	-0.831783	 \\
0.184396	 & 	-0.831783	 & 	1335.51
\end{tabular}

\vtab
 EingenVectors - Molecule A     \\
\vtab
\begin{tabular}{|c c c|}
-0.999893	 & 	0.0146163	 & 	0.00017347	 \\
-0.0146172	 & 	-0.999868	 & 	-0.00709174	 \\
6.97925e-05	 & 	-0.00709352	 & 	0.999975
\end{tabular}

\vtab
 EingenValues - Molecule A     \\
\vtab
\begin{tabular}{|c c c|}
202.555	 & 	1218.61	 & 	1335.51
\end{tabular}
\columnbreak

Molecule B \\ 
S-R-R-Fructuose\_out\_G09
\includegraphics[width=8cm]{Images/ImagesFromVMD/S-R-R-Fructuose_out_G09.jpg}
\\
Inertia Tensor - Molecule B \\
\vtab
\begin{tabular}{|c c c|}
267.669	 & 	-2.83987	 & 	8.0125	 \\
-2.83987	 & 	1173.13	 & 	0.735007	 \\
8.0125	 & 	0.735007	 & 	1229.07
\end{tabular}

\vtab
 EingenVectors - Molecule B     \\
\vtab
\begin{tabular}{|c c c|}
0.99996	 & 	0.00314275	 & 	-0.00833561	 \\
-0.00324825	 & 	0.999914	 & 	-0.0126739	 \\
0.00829507	 & 	0.0127005	 & 	0.999885
\end{tabular}

\vtab
 EingenValues - Molecule B     \\
\vtab
\begin{tabular}{|c c c|}
267.593	 & 	1173.13	 & 	1229.15
\end{tabular}

\end{center}
\end{multicols}
\begin{center}
\vtab
\vtab
\textcolor{NavyBlue}{\Large Are Isomers}
\end{center}
%%%%%%%%%%%%%%%%%%%%%%%%%%%%%%%%%%%%%%%%%%%%%%%%%%%%%%%%
%%%%%%%%%%%%%%%%%%%%%%%%%%%%%%%%%%%%%%%%%%%%%%%%%%%%%%%%
\newpage
%%%%%%%%%%%%%%%%%%%%%%%%%%%%%%%%%%%%%%%%%%%%%%%%%%%%%%%%
%%%%%%%%%%%%%%%%%%%%%%%%%%%%%%%%%%%%%%%%%%%%%%%%%%%%%%%%

\vtab[-2cm]
\begin{center}
{\large Sugars}
\end{center}
\begin{multicols}{2}
\begin{center}
Molecule A \\ 
R-S-R-Fructuose\_out\_G09
\includegraphics[width=8cm]{Images/ImagesFromVMD/R-S-R-Fructuose_out_G09.jpg}
\\
Inertia Tensor - Molecule A \\
\vtab
\begin{tabular}{|c c c|}
202.772	 & 	14.8494	 & 	0.184396	 \\
14.8494	 & 	1218.4	 & 	-0.831783	 \\
0.184396	 & 	-0.831783	 & 	1335.51
\end{tabular}

\vtab
 EingenVectors - Molecule A     \\
\vtab
\begin{tabular}{|c c c|}
-0.999893	 & 	0.0146163	 & 	0.00017347	 \\
-0.0146172	 & 	-0.999868	 & 	-0.00709174	 \\
6.97925e-05	 & 	-0.00709352	 & 	0.999975
\end{tabular}

\vtab
 EingenValues - Molecule A     \\
\vtab
\begin{tabular}{|c c c|}
202.555	 & 	1218.61	 & 	1335.51
\end{tabular}
\columnbreak

Molecule B \\ 
S-R-S-Fructuose\_out\_G09
\includegraphics[width=8cm]{Images/ImagesFromVMD/S-R-S-Fructuose_out_G09.jpg}
\\
Inertia Tensor - Molecule B \\
\vtab
\begin{tabular}{|c c c|}
202.782	 & 	-14.8472	 & 	-0.161586	 \\
-14.8472	 & 	1218.41	 & 	-0.836204	 \\
-0.161586	 & 	-0.836204	 & 	1335.52
\end{tabular}

\vtab
 EingenVectors - Molecule B     \\
\vtab
\begin{tabular}{|c c c|}
-0.999893	 & 	-0.0146142	 & 	-0.000153395	 \\
0.0146149	 & 	-0.999868	 & 	-0.00713203	 \\
-4.91458e-05	 & 	-0.00713351	 & 	0.999975
\end{tabular}

\vtab
 EingenValues - Molecule B     \\
\vtab
\begin{tabular}{|c c c|}
202.565	 & 	1218.62	 & 	1335.52
\end{tabular}

\end{center}
\end{multicols}
\begin{center}
\vtab
\vtab
\textcolor{NavyBlue}{\Large Are Isomer and Stereoisomers}
\end{center}
%%%%%%%%%%%%%%%%%%%%%%%%%%%%%%%%%%%%%%%%%%%%%%%%%%%%%%%%
%%%%%%%%%%%%%%%%%%%%%%%%%%%%%%%%%%%%%%%%%%%%%%%%%%%%%%%%
\newpage
%%%%%%%%%%%%%%%%%%%%%%%%%%%%%%%%%%%%%%%%%%%%%%%%%%%%%%%%
%%%%%%%%%%%%%%%%%%%%%%%%%%%%%%%%%%%%%%%%%%%%%%%%%%%%%%%%

\vtab[-2cm]
\begin{center}
{\large Sugars}
\end{center}
\begin{multicols}{2}
\begin{center}
Molecule A \\ 
R-S-R-Fructuose\_out\_G09
\includegraphics[width=8cm]{Images/ImagesFromVMD/R-S-R-Fructuose_out_G09.jpg}
\\
Inertia Tensor - Molecule A \\
\vtab
\begin{tabular}{|c c c|}
202.772	 & 	14.8494	 & 	0.184396	 \\
14.8494	 & 	1218.4	 & 	-0.831783	 \\
0.184396	 & 	-0.831783	 & 	1335.51
\end{tabular}

\vtab
 EingenVectors - Molecule A     \\
\vtab
\begin{tabular}{|c c c|}
-0.999893	 & 	0.0146163	 & 	0.00017347	 \\
-0.0146172	 & 	-0.999868	 & 	-0.00709174	 \\
6.97925e-05	 & 	-0.00709352	 & 	0.999975
\end{tabular}

\vtab
 EingenValues - Molecule A     \\
\vtab
\begin{tabular}{|c c c|}
202.555	 & 	1218.61	 & 	1335.51
\end{tabular}
\columnbreak

Molecule B \\ 
S-S-R-Fructuose\_out\_G09
\includegraphics[width=8cm]{Images/ImagesFromVMD/S-S-R-Fructuose_out_G09.jpg}
\\
Inertia Tensor - Molecule B \\
\vtab
\begin{tabular}{|c c c|}
258.093	 & 	-6.7331	 & 	10.2577	 \\
-6.7331	 & 	1157.52	 & 	-1.31924	 \\
10.2577	 & 	-1.31924	 & 	1279.94
\end{tabular}

\vtab
 EingenVectors - Molecule B     \\
\vtab
\begin{tabular}{|c c c|}
0.999922	 & 	0.00746944	 & 	-0.0100264	 \\
0.00735508	 & 	-0.999908	 & 	-0.0113945	 \\
0.0101106	 & 	-0.0113198	 & 	0.999885
\end{tabular}

\vtab
 EingenValues - Molecule B     \\
\vtab
\begin{tabular}{|c c c|}
257.939	 & 	1157.55	 & 	1280.06
\end{tabular}

\end{center}
\end{multicols}
\begin{center}
\vtab
\vtab
\textcolor{NavyBlue}{\Large Are Isomers}
\end{center}
%%%%%%%%%%%%%%%%%%%%%%%%%%%%%%%%%%%%%%%%%%%%%%%%%%%%%%%%
%%%%%%%%%%%%%%%%%%%%%%%%%%%%%%%%%%%%%%%%%%%%%%%%%%%%%%%%
\newpage
%%%%%%%%%%%%%%%%%%%%%%%%%%%%%%%%%%%%%%%%%%%%%%%%%%%%%%%%
%%%%%%%%%%%%%%%%%%%%%%%%%%%%%%%%%%%%%%%%%%%%%%%%%%%%%%%%

\vtab[-2cm]
\begin{center}
{\large Sugars}
\end{center}
\begin{multicols}{2}
\begin{center}
Molecule A \\ 
R-S-R-Fructuose\_out\_G09
\includegraphics[width=8cm]{Images/ImagesFromVMD/R-S-R-Fructuose_out_G09.jpg}
\\
Inertia Tensor - Molecule A \\
\vtab
\begin{tabular}{|c c c|}
202.772	 & 	14.8494	 & 	0.184396	 \\
14.8494	 & 	1218.4	 & 	-0.831783	 \\
0.184396	 & 	-0.831783	 & 	1335.51
\end{tabular}

\vtab
 EingenVectors - Molecule A     \\
\vtab
\begin{tabular}{|c c c|}
-0.999893	 & 	0.0146163	 & 	0.00017347	 \\
-0.0146172	 & 	-0.999868	 & 	-0.00709174	 \\
6.97925e-05	 & 	-0.00709352	 & 	0.999975
\end{tabular}

\vtab
 EingenValues - Molecule A     \\
\vtab
\begin{tabular}{|c c c|}
202.555	 & 	1218.61	 & 	1335.51
\end{tabular}
\columnbreak

Molecule B \\ 
S-S-S-Fructuose\_out\_G09
\includegraphics[width=8cm]{Images/ImagesFromVMD/S-S-S-Fructuose_out_G09.jpg}
\\
Inertia Tensor - Molecule B \\
\vtab
\begin{tabular}{|c c c|}
226.452	 & 	13.0925	 & 	-2.23012	 \\
13.0925	 & 	1177.96	 & 	-0.925239	 \\
-2.23012	 & 	-0.925239	 & 	1359.88
\end{tabular}

\vtab
 EingenVectors - Molecule B     \\
\vtab
\begin{tabular}{|c c c|}
-0.999903	 & 	0.0137538	 & 	-0.00195585	 \\
0.0137434	 & 	0.999892	 & 	0.00525895	 \\
-0.00202797	 & 	-0.00523156	 & 	0.999984
\end{tabular}

\vtab
 EingenValues - Molecule B     \\
\vtab
\begin{tabular}{|c c c|}
226.267	 & 	1178.14	 & 	1359.89
\end{tabular}

\end{center}
\end{multicols}
\begin{center}
\vtab
\vtab
\textcolor{NavyBlue}{\Large Are Isomers}
\end{center}
%%%%%%%%%%%%%%%%%%%%%%%%%%%%%%%%%%%%%%%%%%%%%%%%%%%%%%%%
%%%%%%%%%%%%%%%%%%%%%%%%%%%%%%%%%%%%%%%%%%%%%%%%%%%%%%%%
\newpage
%%%%%%%%%%%%%%%%%%%%%%%%%%%%%%%%%%%%%%%%%%%%%%%%%%%%%%%%
%%%%%%%%%%%%%%%%%%%%%%%%%%%%%%%%%%%%%%%%%%%%%%%%%%%%%%%%

\vtab[-2cm]
\begin{center}
{\large Sugars}
\end{center}
\begin{multicols}{2}
\begin{center}
Molecule A \\ 
R-S-S-Fructuose\_out\_G09
\includegraphics[width=8cm]{Images/ImagesFromVMD/R-S-S-Fructuose_out_G09.jpg}
\\
Inertia Tensor - Molecule A \\
\vtab
\begin{tabular}{|c c c|}
215.024	 & 	-14.0942	 & 	-5.97678	 \\
-14.0942	 & 	1187.1	 & 	0.653719	 \\
-5.97678	 & 	0.653719	 & 	1326.43
\end{tabular}

\vtab
 EingenVectors - Molecule A     \\
\vtab
\begin{tabular}{|c c c|}
0.999881	 & 	0.0144902	 & 	0.00536739	 \\
0.0144617	 & 	-0.999881	 & 	0.0053195	 \\
-0.00544384	 & 	0.00524124	 & 	0.999971
\end{tabular}

\vtab
 EingenValues - Molecule A     \\
\vtab
\begin{tabular}{|c c c|}
214.787	 & 	1187.3	 & 	1326.46
\end{tabular}
\columnbreak

Molecule B \\ 
S-R-R-Fructuose\_out\_G09
\includegraphics[width=8cm]{Images/ImagesFromVMD/S-R-R-Fructuose_out_G09.jpg}
\\
Inertia Tensor - Molecule B \\
\vtab
\begin{tabular}{|c c c|}
267.669	 & 	-2.83987	 & 	8.0125	 \\
-2.83987	 & 	1173.13	 & 	0.735007	 \\
8.0125	 & 	0.735007	 & 	1229.07
\end{tabular}

\vtab
 EingenVectors - Molecule B     \\
\vtab
\begin{tabular}{|c c c|}
0.99996	 & 	0.00314275	 & 	-0.00833561	 \\
-0.00324825	 & 	0.999914	 & 	-0.0126739	 \\
0.00829507	 & 	0.0127005	 & 	0.999885
\end{tabular}

\vtab
 EingenValues - Molecule B     \\
\vtab
\begin{tabular}{|c c c|}
267.593	 & 	1173.13	 & 	1229.15
\end{tabular}

\end{center}
\end{multicols}
\begin{center}
\vtab
\vtab
\textcolor{NavyBlue}{\Large Are Isomers}
\end{center}
%%%%%%%%%%%%%%%%%%%%%%%%%%%%%%%%%%%%%%%%%%%%%%%%%%%%%%%%
%%%%%%%%%%%%%%%%%%%%%%%%%%%%%%%%%%%%%%%%%%%%%%%%%%%%%%%%
\newpage
%%%%%%%%%%%%%%%%%%%%%%%%%%%%%%%%%%%%%%%%%%%%%%%%%%%%%%%%
%%%%%%%%%%%%%%%%%%%%%%%%%%%%%%%%%%%%%%%%%%%%%%%%%%%%%%%%

\vtab[-2cm]
\begin{center}
{\large Sugars}
\end{center}
\begin{multicols}{2}
\begin{center}
Molecule A \\ 
R-S-S-Fructuose\_out\_G09
\includegraphics[width=8cm]{Images/ImagesFromVMD/R-S-S-Fructuose_out_G09.jpg}
\\
Inertia Tensor - Molecule A \\
\vtab
\begin{tabular}{|c c c|}
215.024	 & 	-14.0942	 & 	-5.97678	 \\
-14.0942	 & 	1187.1	 & 	0.653719	 \\
-5.97678	 & 	0.653719	 & 	1326.43
\end{tabular}

\vtab
 EingenVectors - Molecule A     \\
\vtab
\begin{tabular}{|c c c|}
0.999881	 & 	0.0144902	 & 	0.00536739	 \\
0.0144617	 & 	-0.999881	 & 	0.0053195	 \\
-0.00544384	 & 	0.00524124	 & 	0.999971
\end{tabular}

\vtab
 EingenValues - Molecule A     \\
\vtab
\begin{tabular}{|c c c|}
214.787	 & 	1187.3	 & 	1326.46
\end{tabular}
\columnbreak

Molecule B \\ 
S-R-S-Fructuose\_out\_G09
\includegraphics[width=8cm]{Images/ImagesFromVMD/S-R-S-Fructuose_out_G09.jpg}
\\
Inertia Tensor - Molecule B \\
\vtab
\begin{tabular}{|c c c|}
202.782	 & 	-14.8472	 & 	-0.161586	 \\
-14.8472	 & 	1218.41	 & 	-0.836204	 \\
-0.161586	 & 	-0.836204	 & 	1335.52
\end{tabular}

\vtab
 EingenVectors - Molecule B     \\
\vtab
\begin{tabular}{|c c c|}
-0.999893	 & 	-0.0146142	 & 	-0.000153395	 \\
0.0146149	 & 	-0.999868	 & 	-0.00713203	 \\
-4.91458e-05	 & 	-0.00713351	 & 	0.999975
\end{tabular}

\vtab
 EingenValues - Molecule B     \\
\vtab
\begin{tabular}{|c c c|}
202.565	 & 	1218.62	 & 	1335.52
\end{tabular}

\end{center}
\end{multicols}
\begin{center}
\vtab
\vtab
\textcolor{NavyBlue}{\Large Are Isomers}
\end{center}
%%%%%%%%%%%%%%%%%%%%%%%%%%%%%%%%%%%%%%%%%%%%%%%%%%%%%%%%
%%%%%%%%%%%%%%%%%%%%%%%%%%%%%%%%%%%%%%%%%%%%%%%%%%%%%%%%
\newpage
%%%%%%%%%%%%%%%%%%%%%%%%%%%%%%%%%%%%%%%%%%%%%%%%%%%%%%%%
%%%%%%%%%%%%%%%%%%%%%%%%%%%%%%%%%%%%%%%%%%%%%%%%%%%%%%%%

\vtab[-2cm]
\begin{center}
{\large Sugars}
\end{center}
\begin{multicols}{2}
\begin{center}
Molecule A \\ 
R-S-S-Fructuose\_out\_G09
\includegraphics[width=8cm]{Images/ImagesFromVMD/R-S-S-Fructuose_out_G09.jpg}
\\
Inertia Tensor - Molecule A \\
\vtab
\begin{tabular}{|c c c|}
215.024	 & 	-14.0942	 & 	-5.97678	 \\
-14.0942	 & 	1187.1	 & 	0.653719	 \\
-5.97678	 & 	0.653719	 & 	1326.43
\end{tabular}

\vtab
 EingenVectors - Molecule A     \\
\vtab
\begin{tabular}{|c c c|}
0.999881	 & 	0.0144902	 & 	0.00536739	 \\
0.0144617	 & 	-0.999881	 & 	0.0053195	 \\
-0.00544384	 & 	0.00524124	 & 	0.999971
\end{tabular}

\vtab
 EingenValues - Molecule A     \\
\vtab
\begin{tabular}{|c c c|}
214.787	 & 	1187.3	 & 	1326.46
\end{tabular}
\columnbreak

Molecule B \\ 
S-S-R-Fructuose\_out\_G09
\includegraphics[width=8cm]{Images/ImagesFromVMD/S-S-R-Fructuose_out_G09.jpg}
\\
Inertia Tensor - Molecule B \\
\vtab
\begin{tabular}{|c c c|}
258.093	 & 	-6.7331	 & 	10.2577	 \\
-6.7331	 & 	1157.52	 & 	-1.31924	 \\
10.2577	 & 	-1.31924	 & 	1279.94
\end{tabular}

\vtab
 EingenVectors - Molecule B     \\
\vtab
\begin{tabular}{|c c c|}
0.999922	 & 	0.00746944	 & 	-0.0100264	 \\
0.00735508	 & 	-0.999908	 & 	-0.0113945	 \\
0.0101106	 & 	-0.0113198	 & 	0.999885
\end{tabular}

\vtab
 EingenValues - Molecule B     \\
\vtab
\begin{tabular}{|c c c|}
257.939	 & 	1157.55	 & 	1280.06
\end{tabular}

\end{center}
\end{multicols}
\begin{center}
\vtab
\vtab
\textcolor{NavyBlue}{\Large Are Isomers}
\end{center}
%%%%%%%%%%%%%%%%%%%%%%%%%%%%%%%%%%%%%%%%%%%%%%%%%%%%%%%%
%%%%%%%%%%%%%%%%%%%%%%%%%%%%%%%%%%%%%%%%%%%%%%%%%%%%%%%%
\newpage
%%%%%%%%%%%%%%%%%%%%%%%%%%%%%%%%%%%%%%%%%%%%%%%%%%%%%%%%
%%%%%%%%%%%%%%%%%%%%%%%%%%%%%%%%%%%%%%%%%%%%%%%%%%%%%%%%

\vtab[-2cm]
\begin{center}
{\large Sugars}
\end{center}
\begin{multicols}{2}
\begin{center}
Molecule A \\ 
R-S-S-Fructuose\_out\_G09
\includegraphics[width=8cm]{Images/ImagesFromVMD/R-S-S-Fructuose_out_G09.jpg}
\\
Inertia Tensor - Molecule A \\
\vtab
\begin{tabular}{|c c c|}
215.024	 & 	-14.0942	 & 	-5.97678	 \\
-14.0942	 & 	1187.1	 & 	0.653719	 \\
-5.97678	 & 	0.653719	 & 	1326.43
\end{tabular}

\vtab
 EingenVectors - Molecule A     \\
\vtab
\begin{tabular}{|c c c|}
0.999881	 & 	0.0144902	 & 	0.00536739	 \\
0.0144617	 & 	-0.999881	 & 	0.0053195	 \\
-0.00544384	 & 	0.00524124	 & 	0.999971
\end{tabular}

\vtab
 EingenValues - Molecule A     \\
\vtab
\begin{tabular}{|c c c|}
214.787	 & 	1187.3	 & 	1326.46
\end{tabular}
\columnbreak

Molecule B \\ 
S-S-S-Fructuose\_out\_G09
\includegraphics[width=8cm]{Images/ImagesFromVMD/S-S-S-Fructuose_out_G09.jpg}
\\
Inertia Tensor - Molecule B \\
\vtab
\begin{tabular}{|c c c|}
226.452	 & 	13.0925	 & 	-2.23012	 \\
13.0925	 & 	1177.96	 & 	-0.925239	 \\
-2.23012	 & 	-0.925239	 & 	1359.88
\end{tabular}

\vtab
 EingenVectors - Molecule B     \\
\vtab
\begin{tabular}{|c c c|}
-0.999903	 & 	0.0137538	 & 	-0.00195585	 \\
0.0137434	 & 	0.999892	 & 	0.00525895	 \\
-0.00202797	 & 	-0.00523156	 & 	0.999984
\end{tabular}

\vtab
 EingenValues - Molecule B     \\
\vtab
\begin{tabular}{|c c c|}
226.267	 & 	1178.14	 & 	1359.89
\end{tabular}

\end{center}
\end{multicols}
\begin{center}
\vtab
\vtab
\textcolor{NavyBlue}{\Large Are Isomers}
\end{center}
%%%%%%%%%%%%%%%%%%%%%%%%%%%%%%%%%%%%%%%%%%%%%%%%%%%%%%%%
%%%%%%%%%%%%%%%%%%%%%%%%%%%%%%%%%%%%%%%%%%%%%%%%%%%%%%%%
\newpage
%%%%%%%%%%%%%%%%%%%%%%%%%%%%%%%%%%%%%%%%%%%%%%%%%%%%%%%%
%%%%%%%%%%%%%%%%%%%%%%%%%%%%%%%%%%%%%%%%%%%%%%%%%%%%%%%%

\vtab[-2cm]
\begin{center}
{\large Sugars}
\end{center}
\begin{multicols}{2}
\begin{center}
Molecule A \\ 
S-R-R-Fructuose\_out\_G09
\includegraphics[width=8cm]{Images/ImagesFromVMD/S-R-R-Fructuose_out_G09.jpg}
\\
Inertia Tensor - Molecule A \\
\vtab
\begin{tabular}{|c c c|}
267.669	 & 	-2.83987	 & 	8.0125	 \\
-2.83987	 & 	1173.13	 & 	0.735007	 \\
8.0125	 & 	0.735007	 & 	1229.07
\end{tabular}

\vtab
 EingenVectors - Molecule A     \\
\vtab
\begin{tabular}{|c c c|}
0.99996	 & 	0.00314275	 & 	-0.00833561	 \\
-0.00324825	 & 	0.999914	 & 	-0.0126739	 \\
0.00829507	 & 	0.0127005	 & 	0.999885
\end{tabular}

\vtab
 EingenValues - Molecule A     \\
\vtab
\begin{tabular}{|c c c|}
267.593	 & 	1173.13	 & 	1229.15
\end{tabular}
\columnbreak

Molecule B \\ 
S-R-S-Fructuose\_out\_G09
\includegraphics[width=8cm]{Images/ImagesFromVMD/S-R-S-Fructuose_out_G09.jpg}
\\
Inertia Tensor - Molecule B \\
\vtab
\begin{tabular}{|c c c|}
202.782	 & 	-14.8472	 & 	-0.161586	 \\
-14.8472	 & 	1218.41	 & 	-0.836204	 \\
-0.161586	 & 	-0.836204	 & 	1335.52
\end{tabular}

\vtab
 EingenVectors - Molecule B     \\
\vtab
\begin{tabular}{|c c c|}
-0.999893	 & 	-0.0146142	 & 	-0.000153395	 \\
0.0146149	 & 	-0.999868	 & 	-0.00713203	 \\
-4.91458e-05	 & 	-0.00713351	 & 	0.999975
\end{tabular}

\vtab
 EingenValues - Molecule B     \\
\vtab
\begin{tabular}{|c c c|}
202.565	 & 	1218.62	 & 	1335.52
\end{tabular}

\end{center}
\end{multicols}
\begin{center}
\vtab
\vtab
\textcolor{NavyBlue}{\Large Are Isomers}
\end{center}
%%%%%%%%%%%%%%%%%%%%%%%%%%%%%%%%%%%%%%%%%%%%%%%%%%%%%%%%
%%%%%%%%%%%%%%%%%%%%%%%%%%%%%%%%%%%%%%%%%%%%%%%%%%%%%%%%
\newpage
%%%%%%%%%%%%%%%%%%%%%%%%%%%%%%%%%%%%%%%%%%%%%%%%%%%%%%%%
%%%%%%%%%%%%%%%%%%%%%%%%%%%%%%%%%%%%%%%%%%%%%%%%%%%%%%%%

\vtab[-2cm]
\begin{center}
{\large Sugars}
\end{center}
\begin{multicols}{2}
\begin{center}
Molecule A \\ 
S-R-R-Fructuose\_out\_G09
\includegraphics[width=8cm]{Images/ImagesFromVMD/S-R-R-Fructuose_out_G09.jpg}
\\
Inertia Tensor - Molecule A \\
\vtab
\begin{tabular}{|c c c|}
267.669	 & 	-2.83987	 & 	8.0125	 \\
-2.83987	 & 	1173.13	 & 	0.735007	 \\
8.0125	 & 	0.735007	 & 	1229.07
\end{tabular}

\vtab
 EingenVectors - Molecule A     \\
\vtab
\begin{tabular}{|c c c|}
0.99996	 & 	0.00314275	 & 	-0.00833561	 \\
-0.00324825	 & 	0.999914	 & 	-0.0126739	 \\
0.00829507	 & 	0.0127005	 & 	0.999885
\end{tabular}

\vtab
 EingenValues - Molecule A     \\
\vtab
\begin{tabular}{|c c c|}
267.593	 & 	1173.13	 & 	1229.15
\end{tabular}
\columnbreak

Molecule B \\ 
S-S-R-Fructuose\_out\_G09
\includegraphics[width=8cm]{Images/ImagesFromVMD/S-S-R-Fructuose_out_G09.jpg}
\\
Inertia Tensor - Molecule B \\
\vtab
\begin{tabular}{|c c c|}
258.093	 & 	-6.7331	 & 	10.2577	 \\
-6.7331	 & 	1157.52	 & 	-1.31924	 \\
10.2577	 & 	-1.31924	 & 	1279.94
\end{tabular}

\vtab
 EingenVectors - Molecule B     \\
\vtab
\begin{tabular}{|c c c|}
0.999922	 & 	0.00746944	 & 	-0.0100264	 \\
0.00735508	 & 	-0.999908	 & 	-0.0113945	 \\
0.0101106	 & 	-0.0113198	 & 	0.999885
\end{tabular}

\vtab
 EingenValues - Molecule B     \\
\vtab
\begin{tabular}{|c c c|}
257.939	 & 	1157.55	 & 	1280.06
\end{tabular}

\end{center}
\end{multicols}
\begin{center}
\vtab
\vtab
\textcolor{NavyBlue}{\Large Are Isomers}
\end{center}
%%%%%%%%%%%%%%%%%%%%%%%%%%%%%%%%%%%%%%%%%%%%%%%%%%%%%%%%
%%%%%%%%%%%%%%%%%%%%%%%%%%%%%%%%%%%%%%%%%%%%%%%%%%%%%%%%
\newpage
%%%%%%%%%%%%%%%%%%%%%%%%%%%%%%%%%%%%%%%%%%%%%%%%%%%%%%%%
%%%%%%%%%%%%%%%%%%%%%%%%%%%%%%%%%%%%%%%%%%%%%%%%%%%%%%%%

\vtab[-2cm]
\begin{center}
{\large Sugars}
\end{center}
\begin{multicols}{2}
\begin{center}
Molecule A \\ 
S-R-R-Fructuose\_out\_G09
\includegraphics[width=8cm]{Images/ImagesFromVMD/S-R-R-Fructuose_out_G09.jpg}
\\
Inertia Tensor - Molecule A \\
\vtab
\begin{tabular}{|c c c|}
267.669	 & 	-2.83987	 & 	8.0125	 \\
-2.83987	 & 	1173.13	 & 	0.735007	 \\
8.0125	 & 	0.735007	 & 	1229.07
\end{tabular}

\vtab
 EingenVectors - Molecule A     \\
\vtab
\begin{tabular}{|c c c|}
0.99996	 & 	0.00314275	 & 	-0.00833561	 \\
-0.00324825	 & 	0.999914	 & 	-0.0126739	 \\
0.00829507	 & 	0.0127005	 & 	0.999885
\end{tabular}

\vtab
 EingenValues - Molecule A     \\
\vtab
\begin{tabular}{|c c c|}
267.593	 & 	1173.13	 & 	1229.15
\end{tabular}
\columnbreak

Molecule B \\ 
S-S-S-Fructuose\_out\_G09
\includegraphics[width=8cm]{Images/ImagesFromVMD/S-S-S-Fructuose_out_G09.jpg}
\\
Inertia Tensor - Molecule B \\
\vtab
\begin{tabular}{|c c c|}
226.452	 & 	13.0925	 & 	-2.23012	 \\
13.0925	 & 	1177.96	 & 	-0.925239	 \\
-2.23012	 & 	-0.925239	 & 	1359.88
\end{tabular}

\vtab
 EingenVectors - Molecule B     \\
\vtab
\begin{tabular}{|c c c|}
-0.999903	 & 	0.0137538	 & 	-0.00195585	 \\
0.0137434	 & 	0.999892	 & 	0.00525895	 \\
-0.00202797	 & 	-0.00523156	 & 	0.999984
\end{tabular}

\vtab
 EingenValues - Molecule B     \\
\vtab
\begin{tabular}{|c c c|}
226.267	 & 	1178.14	 & 	1359.89
\end{tabular}

\end{center}
\end{multicols}
\begin{center}
\vtab
\vtab
\textcolor{NavyBlue}{\Large Are Isomers}
\end{center}
%%%%%%%%%%%%%%%%%%%%%%%%%%%%%%%%%%%%%%%%%%%%%%%%%%%%%%%%
%%%%%%%%%%%%%%%%%%%%%%%%%%%%%%%%%%%%%%%%%%%%%%%%%%%%%%%%
\newpage
%%%%%%%%%%%%%%%%%%%%%%%%%%%%%%%%%%%%%%%%%%%%%%%%%%%%%%%%
%%%%%%%%%%%%%%%%%%%%%%%%%%%%%%%%%%%%%%%%%%%%%%%%%%%%%%%%

\vtab[-2cm]
\begin{center}
{\large Sugars}
\end{center}
\begin{multicols}{2}
\begin{center}
Molecule A \\ 
S-R-S-Fructuose\_out\_G09
\includegraphics[width=8cm]{Images/ImagesFromVMD/S-R-S-Fructuose_out_G09.jpg}
\\
Inertia Tensor - Molecule A \\
\vtab
\begin{tabular}{|c c c|}
202.782	 & 	-14.8472	 & 	-0.161586	 \\
-14.8472	 & 	1218.41	 & 	-0.836204	 \\
-0.161586	 & 	-0.836204	 & 	1335.52
\end{tabular}

\vtab
 EingenVectors - Molecule A     \\
\vtab
\begin{tabular}{|c c c|}
-0.999893	 & 	-0.0146142	 & 	-0.000153395	 \\
0.0146149	 & 	-0.999868	 & 	-0.00713203	 \\
-4.91458e-05	 & 	-0.00713351	 & 	0.999975
\end{tabular}

\vtab
 EingenValues - Molecule A     \\
\vtab
\begin{tabular}{|c c c|}
202.565	 & 	1218.62	 & 	1335.52
\end{tabular}
\columnbreak

Molecule B \\ 
S-S-R-Fructuose\_out\_G09
\includegraphics[width=8cm]{Images/ImagesFromVMD/S-S-R-Fructuose_out_G09.jpg}
\\
Inertia Tensor - Molecule B \\
\vtab
\begin{tabular}{|c c c|}
258.093	 & 	-6.7331	 & 	10.2577	 \\
-6.7331	 & 	1157.52	 & 	-1.31924	 \\
10.2577	 & 	-1.31924	 & 	1279.94
\end{tabular}

\vtab
 EingenVectors - Molecule B     \\
\vtab
\begin{tabular}{|c c c|}
0.999922	 & 	0.00746944	 & 	-0.0100264	 \\
0.00735508	 & 	-0.999908	 & 	-0.0113945	 \\
0.0101106	 & 	-0.0113198	 & 	0.999885
\end{tabular}

\vtab
 EingenValues - Molecule B     \\
\vtab
\begin{tabular}{|c c c|}
257.939	 & 	1157.55	 & 	1280.06
\end{tabular}

\end{center}
\end{multicols}
\begin{center}
\vtab
\vtab
\textcolor{NavyBlue}{\Large Are Isomers}
\end{center}
%%%%%%%%%%%%%%%%%%%%%%%%%%%%%%%%%%%%%%%%%%%%%%%%%%%%%%%%
%%%%%%%%%%%%%%%%%%%%%%%%%%%%%%%%%%%%%%%%%%%%%%%%%%%%%%%%
\newpage
%%%%%%%%%%%%%%%%%%%%%%%%%%%%%%%%%%%%%%%%%%%%%%%%%%%%%%%%
%%%%%%%%%%%%%%%%%%%%%%%%%%%%%%%%%%%%%%%%%%%%%%%%%%%%%%%%

\vtab[-2cm]
\begin{center}
{\large Sugars}
\end{center}
\begin{multicols}{2}
\begin{center}
Molecule A \\ 
S-R-S-Fructuose\_out\_G09
\includegraphics[width=8cm]{Images/ImagesFromVMD/S-R-S-Fructuose_out_G09.jpg}
\\
Inertia Tensor - Molecule A \\
\vtab
\begin{tabular}{|c c c|}
202.782	 & 	-14.8472	 & 	-0.161586	 \\
-14.8472	 & 	1218.41	 & 	-0.836204	 \\
-0.161586	 & 	-0.836204	 & 	1335.52
\end{tabular}

\vtab
 EingenVectors - Molecule A     \\
\vtab
\begin{tabular}{|c c c|}
-0.999893	 & 	-0.0146142	 & 	-0.000153395	 \\
0.0146149	 & 	-0.999868	 & 	-0.00713203	 \\
-4.91458e-05	 & 	-0.00713351	 & 	0.999975
\end{tabular}

\vtab
 EingenValues - Molecule A     \\
\vtab
\begin{tabular}{|c c c|}
202.565	 & 	1218.62	 & 	1335.52
\end{tabular}
\columnbreak

Molecule B \\ 
S-S-S-Fructuose\_out\_G09
\includegraphics[width=8cm]{Images/ImagesFromVMD/S-S-S-Fructuose_out_G09.jpg}
\\
Inertia Tensor - Molecule B \\
\vtab
\begin{tabular}{|c c c|}
226.452	 & 	13.0925	 & 	-2.23012	 \\
13.0925	 & 	1177.96	 & 	-0.925239	 \\
-2.23012	 & 	-0.925239	 & 	1359.88
\end{tabular}

\vtab
 EingenVectors - Molecule B     \\
\vtab
\begin{tabular}{|c c c|}
-0.999903	 & 	0.0137538	 & 	-0.00195585	 \\
0.0137434	 & 	0.999892	 & 	0.00525895	 \\
-0.00202797	 & 	-0.00523156	 & 	0.999984
\end{tabular}

\vtab
 EingenValues - Molecule B     \\
\vtab
\begin{tabular}{|c c c|}
226.267	 & 	1178.14	 & 	1359.89
\end{tabular}

\end{center}
\end{multicols}
\begin{center}
\vtab
\vtab
\textcolor{NavyBlue}{\Large Are Isomers}
\end{center}
%%%%%%%%%%%%%%%%%%%%%%%%%%%%%%%%%%%%%%%%%%%%%%%%%%%%%%%%
%%%%%%%%%%%%%%%%%%%%%%%%%%%%%%%%%%%%%%%%%%%%%%%%%%%%%%%%
\newpage
%%%%%%%%%%%%%%%%%%%%%%%%%%%%%%%%%%%%%%%%%%%%%%%%%%%%%%%%
%%%%%%%%%%%%%%%%%%%%%%%%%%%%%%%%%%%%%%%%%%%%%%%%%%%%%%%%

\vtab[-2cm]
\begin{center}
{\large Sugars}
\end{center}
\begin{multicols}{2}
\begin{center}
Molecule A \\ 
S-S-R-Fructuose\_out\_G09
\includegraphics[width=8cm]{Images/ImagesFromVMD/S-S-R-Fructuose_out_G09.jpg}
\\
Inertia Tensor - Molecule A \\
\vtab
\begin{tabular}{|c c c|}
258.093	 & 	-6.7331	 & 	10.2577	 \\
-6.7331	 & 	1157.52	 & 	-1.31924	 \\
10.2577	 & 	-1.31924	 & 	1279.94
\end{tabular}

\vtab
 EingenVectors - Molecule A     \\
\vtab
\begin{tabular}{|c c c|}
0.999922	 & 	0.00746944	 & 	-0.0100264	 \\
0.00735508	 & 	-0.999908	 & 	-0.0113945	 \\
0.0101106	 & 	-0.0113198	 & 	0.999885
\end{tabular}

\vtab
 EingenValues - Molecule A     \\
\vtab
\begin{tabular}{|c c c|}
257.939	 & 	1157.55	 & 	1280.06
\end{tabular}
\columnbreak

Molecule B \\ 
S-S-S-Fructuose\_out\_G09
\includegraphics[width=8cm]{Images/ImagesFromVMD/S-S-S-Fructuose_out_G09.jpg}
\\
Inertia Tensor - Molecule B \\
\vtab
\begin{tabular}{|c c c|}
226.452	 & 	13.0925	 & 	-2.23012	 \\
13.0925	 & 	1177.96	 & 	-0.925239	 \\
-2.23012	 & 	-0.925239	 & 	1359.88
\end{tabular}

\vtab
 EingenVectors - Molecule B     \\
\vtab
\begin{tabular}{|c c c|}
-0.999903	 & 	0.0137538	 & 	-0.00195585	 \\
0.0137434	 & 	0.999892	 & 	0.00525895	 \\
-0.00202797	 & 	-0.00523156	 & 	0.999984
\end{tabular}

\vtab
 EingenValues - Molecule B     \\
\vtab
\begin{tabular}{|c c c|}
226.267	 & 	1178.14	 & 	1359.89
\end{tabular}

\end{center}
\end{multicols}
\begin{center}
\vtab
\vtab
\textcolor{NavyBlue}{\Large Are Isomers}
\end{center}
%%%%%%%%%%%%%%%%%%%%%%%%%%%%%%%%%%%%%%%%%%%%%%%%%%%%%%%%
%%%%%%%%%%%%%%%%%%%%%%%%%%%%%%%%%%%%%%%%%%%%%%%%%%%%%%%%
\newpage
%%%%%%%%%%%%%%%%%%%%%%%%%%%%%%%%%%%%%%%%%%%%%%%%%%%%%%%%
%%%%%%%%%%%%%%%%%%%%%%%%%%%%%%%%%%%%%%%%%%%%%%%%%%%%%%%%

\vtab[-2cm]
\begin{center}
{\large SymmetricEnantiomers}
\end{center}
\begin{multicols}{2}
\begin{center}
Molecule A \\ 
r-enan\_out\_G09
\includegraphics[width=8cm]{Images/ImagesFromVMD/r-enan_out_G09.jpg}
\\
Inertia Tensor - Molecule A \\
\vtab
\begin{tabular}{|c c c|}
103.745	 & 	1.51033	 & 	1.05456	 \\
1.51033	 & 	629.842	 & 	1.67376	 \\
1.05456	 & 	1.67376	 & 	722.189
\end{tabular}

\vtab
 EingenVectors - Molecule A     \\
\vtab
\begin{tabular}{|c c c|}
-0.999994	 & 	0.00286537	 & 	0.0016974	 \\
-0.00283409	 & 	-0.999831	 & 	0.0181489	 \\
0.00174912	 & 	0.018144	 & 	0.999834
\end{tabular}

\vtab
 EingenValues - Molecule A     \\
\vtab
\begin{tabular}{|c c c|}
103.739	 & 	629.816	 & 	722.221
\end{tabular}
\columnbreak

Molecule B \\ 
r-enan\_rotated\_out\_G09
\includegraphics[width=8cm]{Images/ImagesFromVMD/r-enan_rotated_out_G09.jpg}
\\
Inertia Tensor - Molecule B \\
\vtab
\begin{tabular}{|c c c|}
103.745	 & 	1.51033	 & 	1.05456	 \\
1.51033	 & 	629.842	 & 	1.67376	 \\
1.05456	 & 	1.67376	 & 	722.189
\end{tabular}

\vtab
 EingenVectors - Molecule B     \\
\vtab
\begin{tabular}{|c c c|}
-0.999994	 & 	0.00286537	 & 	0.0016974	 \\
-0.00283409	 & 	-0.999831	 & 	0.0181489	 \\
0.00174912	 & 	0.018144	 & 	0.999834
\end{tabular}

\vtab
 EingenValues - Molecule B     \\
\vtab
\begin{tabular}{|c c c|}
103.739	 & 	629.816	 & 	722.221
\end{tabular}

\end{center}
\end{multicols}
\begin{center}
\vtab
\vtab
\textcolor{NavyBlue}{\Large Are the same molecule}
\end{center}
%%%%%%%%%%%%%%%%%%%%%%%%%%%%%%%%%%%%%%%%%%%%%%%%%%%%%%%%
%%%%%%%%%%%%%%%%%%%%%%%%%%%%%%%%%%%%%%%%%%%%%%%%%%%%%%%%
\newpage
%%%%%%%%%%%%%%%%%%%%%%%%%%%%%%%%%%%%%%%%%%%%%%%%%%%%%%%%
%%%%%%%%%%%%%%%%%%%%%%%%%%%%%%%%%%%%%%%%%%%%%%%%%%%%%%%%

\vtab[-2cm]
\begin{center}
{\large SymmetricEnantiomers}
\end{center}
\begin{multicols}{2}
\begin{center}
Molecule A \\ 
r-enan\_out\_G09
\includegraphics[width=8cm]{Images/ImagesFromVMD/r-enan_out_G09.jpg}
\\
Inertia Tensor - Molecule A \\
\vtab
\begin{tabular}{|c c c|}
103.745	 & 	1.51033	 & 	1.05456	 \\
1.51033	 & 	629.842	 & 	1.67376	 \\
1.05456	 & 	1.67376	 & 	722.189
\end{tabular}

\vtab
 EingenVectors - Molecule A     \\
\vtab
\begin{tabular}{|c c c|}
-0.999994	 & 	0.00286537	 & 	0.0016974	 \\
-0.00283409	 & 	-0.999831	 & 	0.0181489	 \\
0.00174912	 & 	0.018144	 & 	0.999834
\end{tabular}

\vtab
 EingenValues - Molecule A     \\
\vtab
\begin{tabular}{|c c c|}
103.739	 & 	629.816	 & 	722.221
\end{tabular}
\columnbreak

Molecule B \\ 
r-tmp\_out\_G09
\includegraphics[width=8cm]{Images/ImagesFromVMD/r-tmp_out_G09.jpg}
\\
Inertia Tensor - Molecule B \\
\vtab
\begin{tabular}{|c c c|}
103.745	 & 	-1.51033	 & 	-1.05456	 \\
-1.51033	 & 	629.842	 & 	1.67376	 \\
-1.05456	 & 	1.67376	 & 	722.189
\end{tabular}

\vtab
 EingenVectors - Molecule B     \\
\vtab
\begin{tabular}{|c c c|}
-0.999994	 & 	-0.00286537	 & 	-0.0016974	 \\
0.00283409	 & 	-0.999831	 & 	0.0181489	 \\
-0.00174912	 & 	0.018144	 & 	0.999834
\end{tabular}

\vtab
 EingenValues - Molecule B     \\
\vtab
\begin{tabular}{|c c c|}
103.739	 & 	629.816	 & 	722.221
\end{tabular}

\end{center}
\end{multicols}
\begin{center}
\vtab
\vtab
\textcolor{NavyBlue}{\Large Are Isomer and Stereoisomers}
\end{center}
%%%%%%%%%%%%%%%%%%%%%%%%%%%%%%%%%%%%%%%%%%%%%%%%%%%%%%%%
%%%%%%%%%%%%%%%%%%%%%%%%%%%%%%%%%%%%%%%%%%%%%%%%%%%%%%%%
\newpage
%%%%%%%%%%%%%%%%%%%%%%%%%%%%%%%%%%%%%%%%%%%%%%%%%%%%%%%%
%%%%%%%%%%%%%%%%%%%%%%%%%%%%%%%%%%%%%%%%%%%%%%%%%%%%%%%%

\vtab[-2cm]
\begin{center}
{\large SymmetricEnantiomers}
\end{center}
\begin{multicols}{2}
\begin{center}
Molecule A \\ 
r-enan\_out\_G09
\includegraphics[width=8cm]{Images/ImagesFromVMD/r-enan_out_G09.jpg}
\\
Inertia Tensor - Molecule A \\
\vtab
\begin{tabular}{|c c c|}
103.745	 & 	1.51033	 & 	1.05456	 \\
1.51033	 & 	629.842	 & 	1.67376	 \\
1.05456	 & 	1.67376	 & 	722.189
\end{tabular}

\vtab
 EingenVectors - Molecule A     \\
\vtab
\begin{tabular}{|c c c|}
-0.999994	 & 	0.00286537	 & 	0.0016974	 \\
-0.00283409	 & 	-0.999831	 & 	0.0181489	 \\
0.00174912	 & 	0.018144	 & 	0.999834
\end{tabular}

\vtab
 EingenValues - Molecule A     \\
\vtab
\begin{tabular}{|c c c|}
103.739	 & 	629.816	 & 	722.221
\end{tabular}
\columnbreak

Molecule B \\ 
r-tmp\_rotated\_out\_G09
\includegraphics[width=8cm]{Images/ImagesFromVMD/r-tmp_rotated_out_G09.jpg}
\\
Inertia Tensor - Molecule B \\
\vtab
\begin{tabular}{|c c c|}
103.745	 & 	-1.51033	 & 	-1.05456	 \\
-1.51033	 & 	629.842	 & 	1.67376	 \\
-1.05456	 & 	1.67376	 & 	722.189
\end{tabular}

\vtab
 EingenVectors - Molecule B     \\
\vtab
\begin{tabular}{|c c c|}
-0.999994	 & 	-0.00286537	 & 	-0.0016974	 \\
0.00283409	 & 	-0.999831	 & 	0.0181489	 \\
-0.00174912	 & 	0.018144	 & 	0.999834
\end{tabular}

\vtab
 EingenValues - Molecule B     \\
\vtab
\begin{tabular}{|c c c|}
103.739	 & 	629.816	 & 	722.221
\end{tabular}

\end{center}
\end{multicols}
\begin{center}
\vtab
\vtab
\textcolor{NavyBlue}{\Large Are Isomer and Stereoisomers}
\end{center}
%%%%%%%%%%%%%%%%%%%%%%%%%%%%%%%%%%%%%%%%%%%%%%%%%%%%%%%%
%%%%%%%%%%%%%%%%%%%%%%%%%%%%%%%%%%%%%%%%%%%%%%%%%%%%%%%%
\newpage
%%%%%%%%%%%%%%%%%%%%%%%%%%%%%%%%%%%%%%%%%%%%%%%%%%%%%%%%
%%%%%%%%%%%%%%%%%%%%%%%%%%%%%%%%%%%%%%%%%%%%%%%%%%%%%%%%

\vtab[-2cm]
\begin{center}
{\large SymmetricEnantiomers}
\end{center}
\begin{multicols}{2}
\begin{center}
Molecule A \\ 
r-enan\_out\_G09
\includegraphics[width=8cm]{Images/ImagesFromVMD/r-enan_out_G09.jpg}
\\
Inertia Tensor - Molecule A \\
\vtab
\begin{tabular}{|c c c|}
103.745	 & 	1.51033	 & 	1.05456	 \\
1.51033	 & 	629.842	 & 	1.67376	 \\
1.05456	 & 	1.67376	 & 	722.189
\end{tabular}

\vtab
 EingenVectors - Molecule A     \\
\vtab
\begin{tabular}{|c c c|}
-0.999994	 & 	0.00286537	 & 	0.0016974	 \\
-0.00283409	 & 	-0.999831	 & 	0.0181489	 \\
0.00174912	 & 	0.018144	 & 	0.999834
\end{tabular}

\vtab
 EingenValues - Molecule A     \\
\vtab
\begin{tabular}{|c c c|}
103.739	 & 	629.816	 & 	722.221
\end{tabular}
\columnbreak

Molecule B \\ 
s-enan\_out\_G09
\includegraphics[width=8cm]{Images/ImagesFromVMD/s-enan_out_G09.jpg}
\\
Inertia Tensor - Molecule B \\
\vtab
\begin{tabular}{|c c c|}
103.707	 & 	-1.51142	 & 	-1.05403	 \\
-1.51142	 & 	630.103	 & 	1.67204	 \\
-1.05403	 & 	1.67204	 & 	722.42
\end{tabular}

\vtab
 EingenVectors - Molecule B     \\
\vtab
\begin{tabular}{|c c c|}
-0.999994	 & 	-0.00286582	 & 	-0.00169581	 \\
0.0028346	 & 	-0.999832	 & 	0.018136	 \\
-0.0017475	 & 	0.0181311	 & 	0.999834
\end{tabular}

\vtab
 EingenValues - Molecule B     \\
\vtab
\begin{tabular}{|c c c|}
103.701	 & 	630.077	 & 	722.452
\end{tabular}

\end{center}
\end{multicols}
\begin{center}
\vtab
\vtab
\textcolor{NavyBlue}{\Large Are Isomer and Stereoisomers}
\end{center}
%%%%%%%%%%%%%%%%%%%%%%%%%%%%%%%%%%%%%%%%%%%%%%%%%%%%%%%%
%%%%%%%%%%%%%%%%%%%%%%%%%%%%%%%%%%%%%%%%%%%%%%%%%%%%%%%%
\newpage
%%%%%%%%%%%%%%%%%%%%%%%%%%%%%%%%%%%%%%%%%%%%%%%%%%%%%%%%
%%%%%%%%%%%%%%%%%%%%%%%%%%%%%%%%%%%%%%%%%%%%%%%%%%%%%%%%

\vtab[-2cm]
\begin{center}
{\large SymmetricEnantiomers}
\end{center}
\begin{multicols}{2}
\begin{center}
Molecule A \\ 
r-enan\_out\_G09
\includegraphics[width=8cm]{Images/ImagesFromVMD/r-enan_out_G09.jpg}
\\
Inertia Tensor - Molecule A \\
\vtab
\begin{tabular}{|c c c|}
103.745	 & 	1.51033	 & 	1.05456	 \\
1.51033	 & 	629.842	 & 	1.67376	 \\
1.05456	 & 	1.67376	 & 	722.189
\end{tabular}

\vtab
 EingenVectors - Molecule A     \\
\vtab
\begin{tabular}{|c c c|}
-0.999994	 & 	0.00286537	 & 	0.0016974	 \\
-0.00283409	 & 	-0.999831	 & 	0.0181489	 \\
0.00174912	 & 	0.018144	 & 	0.999834
\end{tabular}

\vtab
 EingenValues - Molecule A     \\
\vtab
\begin{tabular}{|c c c|}
103.739	 & 	629.816	 & 	722.221
\end{tabular}
\columnbreak

Molecule B \\ 
s-enan\_rotated\_out\_G09
\includegraphics[width=8cm]{Images/ImagesFromVMD/s-enan_rotated_out_G09.jpg}
\\
Inertia Tensor - Molecule B \\
\vtab
\begin{tabular}{|c c c|}
103.707	 & 	-1.51138	 & 	-1.05403	 \\
-1.51138	 & 	630.103	 & 	1.67204	 \\
-1.05403	 & 	1.67204	 & 	722.42
\end{tabular}

\vtab
 EingenVectors - Molecule B     \\
\vtab
\begin{tabular}{|c c c|}
-0.999994	 & 	-0.00286575	 & 	-0.00169582	 \\
0.00283453	 & 	-0.999832	 & 	0.018136	 \\
-0.0017475	 & 	0.0181311	 & 	0.999834
\end{tabular}

\vtab
 EingenValues - Molecule B     \\
\vtab
\begin{tabular}{|c c c|}
103.701	 & 	630.077	 & 	722.452
\end{tabular}

\end{center}
\end{multicols}
\begin{center}
\vtab
\vtab
\textcolor{NavyBlue}{\Large Are Isomer and Stereoisomers}
\end{center}
%%%%%%%%%%%%%%%%%%%%%%%%%%%%%%%%%%%%%%%%%%%%%%%%%%%%%%%%
%%%%%%%%%%%%%%%%%%%%%%%%%%%%%%%%%%%%%%%%%%%%%%%%%%%%%%%%
\newpage
%%%%%%%%%%%%%%%%%%%%%%%%%%%%%%%%%%%%%%%%%%%%%%%%%%%%%%%%
%%%%%%%%%%%%%%%%%%%%%%%%%%%%%%%%%%%%%%%%%%%%%%%%%%%%%%%%

\vtab[-2cm]
\begin{center}
{\large SymmetricEnantiomers}
\end{center}
\begin{multicols}{2}
\begin{center}
Molecule A \\ 
r-enan\_out\_G09
\includegraphics[width=8cm]{Images/ImagesFromVMD/r-enan_out_G09.jpg}
\\
Inertia Tensor - Molecule A \\
\vtab
\begin{tabular}{|c c c|}
103.745	 & 	1.51033	 & 	1.05456	 \\
1.51033	 & 	629.842	 & 	1.67376	 \\
1.05456	 & 	1.67376	 & 	722.189
\end{tabular}

\vtab
 EingenVectors - Molecule A     \\
\vtab
\begin{tabular}{|c c c|}
-0.999994	 & 	0.00286537	 & 	0.0016974	 \\
-0.00283409	 & 	-0.999831	 & 	0.0181489	 \\
0.00174912	 & 	0.018144	 & 	0.999834
\end{tabular}

\vtab
 EingenValues - Molecule A     \\
\vtab
\begin{tabular}{|c c c|}
103.739	 & 	629.816	 & 	722.221
\end{tabular}
\columnbreak

Molecule B \\ 
s-tmp\_out\_G09
\includegraphics[width=8cm]{Images/ImagesFromVMD/s-tmp_out_G09.jpg}
\\
Inertia Tensor - Molecule B \\
\vtab
\begin{tabular}{|c c c|}
103.707	 & 	1.51142	 & 	1.05403	 \\
1.51142	 & 	630.103	 & 	1.67204	 \\
1.05403	 & 	1.67204	 & 	722.42
\end{tabular}

\vtab
 EingenVectors - Molecule B     \\
\vtab
\begin{tabular}{|c c c|}
-0.999994	 & 	0.00286582	 & 	0.00169581	 \\
-0.0028346	 & 	-0.999832	 & 	0.018136	 \\
0.0017475	 & 	0.0181311	 & 	0.999834
\end{tabular}

\vtab
 EingenValues - Molecule B     \\
\vtab
\begin{tabular}{|c c c|}
103.701	 & 	630.077	 & 	722.452
\end{tabular}

\end{center}
\end{multicols}
\begin{center}
\vtab
\vtab
\textcolor{NavyBlue}{\Large Are the same molecule}
\end{center}
%%%%%%%%%%%%%%%%%%%%%%%%%%%%%%%%%%%%%%%%%%%%%%%%%%%%%%%%
%%%%%%%%%%%%%%%%%%%%%%%%%%%%%%%%%%%%%%%%%%%%%%%%%%%%%%%%
\newpage
%%%%%%%%%%%%%%%%%%%%%%%%%%%%%%%%%%%%%%%%%%%%%%%%%%%%%%%%
%%%%%%%%%%%%%%%%%%%%%%%%%%%%%%%%%%%%%%%%%%%%%%%%%%%%%%%%

\vtab[-2cm]
\begin{center}
{\large SymmetricEnantiomers}
\end{center}
\begin{multicols}{2}
\begin{center}
Molecule A \\ 
r-enan\_out\_G09
\includegraphics[width=8cm]{Images/ImagesFromVMD/r-enan_out_G09.jpg}
\\
Inertia Tensor - Molecule A \\
\vtab
\begin{tabular}{|c c c|}
103.745	 & 	1.51033	 & 	1.05456	 \\
1.51033	 & 	629.842	 & 	1.67376	 \\
1.05456	 & 	1.67376	 & 	722.189
\end{tabular}

\vtab
 EingenVectors - Molecule A     \\
\vtab
\begin{tabular}{|c c c|}
-0.999994	 & 	0.00286537	 & 	0.0016974	 \\
-0.00283409	 & 	-0.999831	 & 	0.0181489	 \\
0.00174912	 & 	0.018144	 & 	0.999834
\end{tabular}

\vtab
 EingenValues - Molecule A     \\
\vtab
\begin{tabular}{|c c c|}
103.739	 & 	629.816	 & 	722.221
\end{tabular}
\columnbreak

Molecule B \\ 
s-tmp\_rotated\_out\_G09
\includegraphics[width=8cm]{Images/ImagesFromVMD/s-tmp_rotated_out_G09.jpg}
\\
Inertia Tensor - Molecule B \\
\vtab
\begin{tabular}{|c c c|}
103.707	 & 	1.51142	 & 	1.05403	 \\
1.51142	 & 	630.103	 & 	1.67204	 \\
1.05403	 & 	1.67204	 & 	722.42
\end{tabular}

\vtab
 EingenVectors - Molecule B     \\
\vtab
\begin{tabular}{|c c c|}
-0.999994	 & 	0.00286582	 & 	0.00169581	 \\
-0.0028346	 & 	-0.999832	 & 	0.018136	 \\
0.0017475	 & 	0.0181311	 & 	0.999834
\end{tabular}

\vtab
 EingenValues - Molecule B     \\
\vtab
\begin{tabular}{|c c c|}
103.701	 & 	630.077	 & 	722.452
\end{tabular}

\end{center}
\end{multicols}
\begin{center}
\vtab
\vtab
\textcolor{NavyBlue}{\Large Are the same molecule}
\end{center}
%%%%%%%%%%%%%%%%%%%%%%%%%%%%%%%%%%%%%%%%%%%%%%%%%%%%%%%%
%%%%%%%%%%%%%%%%%%%%%%%%%%%%%%%%%%%%%%%%%%%%%%%%%%%%%%%%
\newpage
%%%%%%%%%%%%%%%%%%%%%%%%%%%%%%%%%%%%%%%%%%%%%%%%%%%%%%%%
%%%%%%%%%%%%%%%%%%%%%%%%%%%%%%%%%%%%%%%%%%%%%%%%%%%%%%%%

\vtab[-2cm]
\begin{center}
{\large SymmetricEnantiomers}
\end{center}
\begin{multicols}{2}
\begin{center}
Molecule A \\ 
r-enan\_rotated\_out\_G09
\includegraphics[width=8cm]{Images/ImagesFromVMD/r-enan_rotated_out_G09.jpg}
\\
Inertia Tensor - Molecule A \\
\vtab
\begin{tabular}{|c c c|}
103.745	 & 	1.51033	 & 	1.05456	 \\
1.51033	 & 	629.842	 & 	1.67376	 \\
1.05456	 & 	1.67376	 & 	722.189
\end{tabular}

\vtab
 EingenVectors - Molecule A     \\
\vtab
\begin{tabular}{|c c c|}
-0.999994	 & 	0.00286537	 & 	0.0016974	 \\
-0.00283409	 & 	-0.999831	 & 	0.0181489	 \\
0.00174912	 & 	0.018144	 & 	0.999834
\end{tabular}

\vtab
 EingenValues - Molecule A     \\
\vtab
\begin{tabular}{|c c c|}
103.739	 & 	629.816	 & 	722.221
\end{tabular}
\columnbreak

Molecule B \\ 
r-tmp\_out\_G09
\includegraphics[width=8cm]{Images/ImagesFromVMD/r-tmp_out_G09.jpg}
\\
Inertia Tensor - Molecule B \\
\vtab
\begin{tabular}{|c c c|}
103.745	 & 	-1.51033	 & 	-1.05456	 \\
-1.51033	 & 	629.842	 & 	1.67376	 \\
-1.05456	 & 	1.67376	 & 	722.189
\end{tabular}

\vtab
 EingenVectors - Molecule B     \\
\vtab
\begin{tabular}{|c c c|}
-0.999994	 & 	-0.00286537	 & 	-0.0016974	 \\
0.00283409	 & 	-0.999831	 & 	0.0181489	 \\
-0.00174912	 & 	0.018144	 & 	0.999834
\end{tabular}

\vtab
 EingenValues - Molecule B     \\
\vtab
\begin{tabular}{|c c c|}
103.739	 & 	629.816	 & 	722.221
\end{tabular}

\end{center}
\end{multicols}
\begin{center}
\vtab
\vtab
\textcolor{NavyBlue}{\Large Are Isomer and Stereoisomers}
\end{center}
%%%%%%%%%%%%%%%%%%%%%%%%%%%%%%%%%%%%%%%%%%%%%%%%%%%%%%%%
%%%%%%%%%%%%%%%%%%%%%%%%%%%%%%%%%%%%%%%%%%%%%%%%%%%%%%%%
\newpage
%%%%%%%%%%%%%%%%%%%%%%%%%%%%%%%%%%%%%%%%%%%%%%%%%%%%%%%%
%%%%%%%%%%%%%%%%%%%%%%%%%%%%%%%%%%%%%%%%%%%%%%%%%%%%%%%%

\vtab[-2cm]
\begin{center}
{\large SymmetricEnantiomers}
\end{center}
\begin{multicols}{2}
\begin{center}
Molecule A \\ 
r-enan\_rotated\_out\_G09
\includegraphics[width=8cm]{Images/ImagesFromVMD/r-enan_rotated_out_G09.jpg}
\\
Inertia Tensor - Molecule A \\
\vtab
\begin{tabular}{|c c c|}
103.745	 & 	1.51033	 & 	1.05456	 \\
1.51033	 & 	629.842	 & 	1.67376	 \\
1.05456	 & 	1.67376	 & 	722.189
\end{tabular}

\vtab
 EingenVectors - Molecule A     \\
\vtab
\begin{tabular}{|c c c|}
-0.999994	 & 	0.00286537	 & 	0.0016974	 \\
-0.00283409	 & 	-0.999831	 & 	0.0181489	 \\
0.00174912	 & 	0.018144	 & 	0.999834
\end{tabular}

\vtab
 EingenValues - Molecule A     \\
\vtab
\begin{tabular}{|c c c|}
103.739	 & 	629.816	 & 	722.221
\end{tabular}
\columnbreak

Molecule B \\ 
r-tmp\_rotated\_out\_G09
\includegraphics[width=8cm]{Images/ImagesFromVMD/r-tmp_rotated_out_G09.jpg}
\\
Inertia Tensor - Molecule B \\
\vtab
\begin{tabular}{|c c c|}
103.745	 & 	-1.51033	 & 	-1.05456	 \\
-1.51033	 & 	629.842	 & 	1.67376	 \\
-1.05456	 & 	1.67376	 & 	722.189
\end{tabular}

\vtab
 EingenVectors - Molecule B     \\
\vtab
\begin{tabular}{|c c c|}
-0.999994	 & 	-0.00286537	 & 	-0.0016974	 \\
0.00283409	 & 	-0.999831	 & 	0.0181489	 \\
-0.00174912	 & 	0.018144	 & 	0.999834
\end{tabular}

\vtab
 EingenValues - Molecule B     \\
\vtab
\begin{tabular}{|c c c|}
103.739	 & 	629.816	 & 	722.221
\end{tabular}

\end{center}
\end{multicols}
\begin{center}
\vtab
\vtab
\textcolor{NavyBlue}{\Large Are Isomer and Stereoisomers}
\end{center}
%%%%%%%%%%%%%%%%%%%%%%%%%%%%%%%%%%%%%%%%%%%%%%%%%%%%%%%%
%%%%%%%%%%%%%%%%%%%%%%%%%%%%%%%%%%%%%%%%%%%%%%%%%%%%%%%%
\newpage
%%%%%%%%%%%%%%%%%%%%%%%%%%%%%%%%%%%%%%%%%%%%%%%%%%%%%%%%
%%%%%%%%%%%%%%%%%%%%%%%%%%%%%%%%%%%%%%%%%%%%%%%%%%%%%%%%

\vtab[-2cm]
\begin{center}
{\large SymmetricEnantiomers}
\end{center}
\begin{multicols}{2}
\begin{center}
Molecule A \\ 
r-enan\_rotated\_out\_G09
\includegraphics[width=8cm]{Images/ImagesFromVMD/r-enan_rotated_out_G09.jpg}
\\
Inertia Tensor - Molecule A \\
\vtab
\begin{tabular}{|c c c|}
103.745	 & 	1.51033	 & 	1.05456	 \\
1.51033	 & 	629.842	 & 	1.67376	 \\
1.05456	 & 	1.67376	 & 	722.189
\end{tabular}

\vtab
 EingenVectors - Molecule A     \\
\vtab
\begin{tabular}{|c c c|}
-0.999994	 & 	0.00286537	 & 	0.0016974	 \\
-0.00283409	 & 	-0.999831	 & 	0.0181489	 \\
0.00174912	 & 	0.018144	 & 	0.999834
\end{tabular}

\vtab
 EingenValues - Molecule A     \\
\vtab
\begin{tabular}{|c c c|}
103.739	 & 	629.816	 & 	722.221
\end{tabular}
\columnbreak

Molecule B \\ 
s-enan\_out\_G09
\includegraphics[width=8cm]{Images/ImagesFromVMD/s-enan_out_G09.jpg}
\\
Inertia Tensor - Molecule B \\
\vtab
\begin{tabular}{|c c c|}
103.707	 & 	-1.51142	 & 	-1.05403	 \\
-1.51142	 & 	630.103	 & 	1.67204	 \\
-1.05403	 & 	1.67204	 & 	722.42
\end{tabular}

\vtab
 EingenVectors - Molecule B     \\
\vtab
\begin{tabular}{|c c c|}
-0.999994	 & 	-0.00286582	 & 	-0.00169581	 \\
0.0028346	 & 	-0.999832	 & 	0.018136	 \\
-0.0017475	 & 	0.0181311	 & 	0.999834
\end{tabular}

\vtab
 EingenValues - Molecule B     \\
\vtab
\begin{tabular}{|c c c|}
103.701	 & 	630.077	 & 	722.452
\end{tabular}

\end{center}
\end{multicols}
\begin{center}
\vtab
\vtab
\textcolor{NavyBlue}{\Large Are Isomer and Stereoisomers}
\end{center}
%%%%%%%%%%%%%%%%%%%%%%%%%%%%%%%%%%%%%%%%%%%%%%%%%%%%%%%%
%%%%%%%%%%%%%%%%%%%%%%%%%%%%%%%%%%%%%%%%%%%%%%%%%%%%%%%%
\newpage
%%%%%%%%%%%%%%%%%%%%%%%%%%%%%%%%%%%%%%%%%%%%%%%%%%%%%%%%
%%%%%%%%%%%%%%%%%%%%%%%%%%%%%%%%%%%%%%%%%%%%%%%%%%%%%%%%

\vtab[-2cm]
\begin{center}
{\large SymmetricEnantiomers}
\end{center}
\begin{multicols}{2}
\begin{center}
Molecule A \\ 
r-enan\_rotated\_out\_G09
\includegraphics[width=8cm]{Images/ImagesFromVMD/r-enan_rotated_out_G09.jpg}
\\
Inertia Tensor - Molecule A \\
\vtab
\begin{tabular}{|c c c|}
103.745	 & 	1.51033	 & 	1.05456	 \\
1.51033	 & 	629.842	 & 	1.67376	 \\
1.05456	 & 	1.67376	 & 	722.189
\end{tabular}

\vtab
 EingenVectors - Molecule A     \\
\vtab
\begin{tabular}{|c c c|}
-0.999994	 & 	0.00286537	 & 	0.0016974	 \\
-0.00283409	 & 	-0.999831	 & 	0.0181489	 \\
0.00174912	 & 	0.018144	 & 	0.999834
\end{tabular}

\vtab
 EingenValues - Molecule A     \\
\vtab
\begin{tabular}{|c c c|}
103.739	 & 	629.816	 & 	722.221
\end{tabular}
\columnbreak

Molecule B \\ 
s-enan\_rotated\_out\_G09
\includegraphics[width=8cm]{Images/ImagesFromVMD/s-enan_rotated_out_G09.jpg}
\\
Inertia Tensor - Molecule B \\
\vtab
\begin{tabular}{|c c c|}
103.707	 & 	-1.51138	 & 	-1.05403	 \\
-1.51138	 & 	630.103	 & 	1.67204	 \\
-1.05403	 & 	1.67204	 & 	722.42
\end{tabular}

\vtab
 EingenVectors - Molecule B     \\
\vtab
\begin{tabular}{|c c c|}
-0.999994	 & 	-0.00286575	 & 	-0.00169582	 \\
0.00283453	 & 	-0.999832	 & 	0.018136	 \\
-0.0017475	 & 	0.0181311	 & 	0.999834
\end{tabular}

\vtab
 EingenValues - Molecule B     \\
\vtab
\begin{tabular}{|c c c|}
103.701	 & 	630.077	 & 	722.452
\end{tabular}

\end{center}
\end{multicols}
\begin{center}
\vtab
\vtab
\textcolor{NavyBlue}{\Large Are Isomer and Stereoisomers}
\end{center}
%%%%%%%%%%%%%%%%%%%%%%%%%%%%%%%%%%%%%%%%%%%%%%%%%%%%%%%%
%%%%%%%%%%%%%%%%%%%%%%%%%%%%%%%%%%%%%%%%%%%%%%%%%%%%%%%%
\newpage
%%%%%%%%%%%%%%%%%%%%%%%%%%%%%%%%%%%%%%%%%%%%%%%%%%%%%%%%
%%%%%%%%%%%%%%%%%%%%%%%%%%%%%%%%%%%%%%%%%%%%%%%%%%%%%%%%

\vtab[-2cm]
\begin{center}
{\large SymmetricEnantiomers}
\end{center}
\begin{multicols}{2}
\begin{center}
Molecule A \\ 
r-enan\_rotated\_out\_G09
\includegraphics[width=8cm]{Images/ImagesFromVMD/r-enan_rotated_out_G09.jpg}
\\
Inertia Tensor - Molecule A \\
\vtab
\begin{tabular}{|c c c|}
103.745	 & 	1.51033	 & 	1.05456	 \\
1.51033	 & 	629.842	 & 	1.67376	 \\
1.05456	 & 	1.67376	 & 	722.189
\end{tabular}

\vtab
 EingenVectors - Molecule A     \\
\vtab
\begin{tabular}{|c c c|}
-0.999994	 & 	0.00286537	 & 	0.0016974	 \\
-0.00283409	 & 	-0.999831	 & 	0.0181489	 \\
0.00174912	 & 	0.018144	 & 	0.999834
\end{tabular}

\vtab
 EingenValues - Molecule A     \\
\vtab
\begin{tabular}{|c c c|}
103.739	 & 	629.816	 & 	722.221
\end{tabular}
\columnbreak

Molecule B \\ 
s-tmp\_out\_G09
\includegraphics[width=8cm]{Images/ImagesFromVMD/s-tmp_out_G09.jpg}
\\
Inertia Tensor - Molecule B \\
\vtab
\begin{tabular}{|c c c|}
103.707	 & 	1.51142	 & 	1.05403	 \\
1.51142	 & 	630.103	 & 	1.67204	 \\
1.05403	 & 	1.67204	 & 	722.42
\end{tabular}

\vtab
 EingenVectors - Molecule B     \\
\vtab
\begin{tabular}{|c c c|}
-0.999994	 & 	0.00286582	 & 	0.00169581	 \\
-0.0028346	 & 	-0.999832	 & 	0.018136	 \\
0.0017475	 & 	0.0181311	 & 	0.999834
\end{tabular}

\vtab
 EingenValues - Molecule B     \\
\vtab
\begin{tabular}{|c c c|}
103.701	 & 	630.077	 & 	722.452
\end{tabular}

\end{center}
\end{multicols}
\begin{center}
\vtab
\vtab
\textcolor{NavyBlue}{\Large Are the same molecule}
\end{center}
%%%%%%%%%%%%%%%%%%%%%%%%%%%%%%%%%%%%%%%%%%%%%%%%%%%%%%%%
%%%%%%%%%%%%%%%%%%%%%%%%%%%%%%%%%%%%%%%%%%%%%%%%%%%%%%%%
\newpage
%%%%%%%%%%%%%%%%%%%%%%%%%%%%%%%%%%%%%%%%%%%%%%%%%%%%%%%%
%%%%%%%%%%%%%%%%%%%%%%%%%%%%%%%%%%%%%%%%%%%%%%%%%%%%%%%%

\vtab[-2cm]
\begin{center}
{\large SymmetricEnantiomers}
\end{center}
\begin{multicols}{2}
\begin{center}
Molecule A \\ 
r-enan\_rotated\_out\_G09
\includegraphics[width=8cm]{Images/ImagesFromVMD/r-enan_rotated_out_G09.jpg}
\\
Inertia Tensor - Molecule A \\
\vtab
\begin{tabular}{|c c c|}
103.745	 & 	1.51033	 & 	1.05456	 \\
1.51033	 & 	629.842	 & 	1.67376	 \\
1.05456	 & 	1.67376	 & 	722.189
\end{tabular}

\vtab
 EingenVectors - Molecule A     \\
\vtab
\begin{tabular}{|c c c|}
-0.999994	 & 	0.00286537	 & 	0.0016974	 \\
-0.00283409	 & 	-0.999831	 & 	0.0181489	 \\
0.00174912	 & 	0.018144	 & 	0.999834
\end{tabular}

\vtab
 EingenValues - Molecule A     \\
\vtab
\begin{tabular}{|c c c|}
103.739	 & 	629.816	 & 	722.221
\end{tabular}
\columnbreak

Molecule B \\ 
s-tmp\_rotated\_out\_G09
\includegraphics[width=8cm]{Images/ImagesFromVMD/s-tmp_rotated_out_G09.jpg}
\\
Inertia Tensor - Molecule B \\
\vtab
\begin{tabular}{|c c c|}
103.707	 & 	1.51142	 & 	1.05403	 \\
1.51142	 & 	630.103	 & 	1.67204	 \\
1.05403	 & 	1.67204	 & 	722.42
\end{tabular}

\vtab
 EingenVectors - Molecule B     \\
\vtab
\begin{tabular}{|c c c|}
-0.999994	 & 	0.00286582	 & 	0.00169581	 \\
-0.0028346	 & 	-0.999832	 & 	0.018136	 \\
0.0017475	 & 	0.0181311	 & 	0.999834
\end{tabular}

\vtab
 EingenValues - Molecule B     \\
\vtab
\begin{tabular}{|c c c|}
103.701	 & 	630.077	 & 	722.452
\end{tabular}

\end{center}
\end{multicols}
\begin{center}
\vtab
\vtab
\textcolor{NavyBlue}{\Large Are the same molecule}
\end{center}
%%%%%%%%%%%%%%%%%%%%%%%%%%%%%%%%%%%%%%%%%%%%%%%%%%%%%%%%
%%%%%%%%%%%%%%%%%%%%%%%%%%%%%%%%%%%%%%%%%%%%%%%%%%%%%%%%
\newpage
%%%%%%%%%%%%%%%%%%%%%%%%%%%%%%%%%%%%%%%%%%%%%%%%%%%%%%%%
%%%%%%%%%%%%%%%%%%%%%%%%%%%%%%%%%%%%%%%%%%%%%%%%%%%%%%%%

\vtab[-2cm]
\begin{center}
{\large SymmetricEnantiomers}
\end{center}
\begin{multicols}{2}
\begin{center}
Molecule A \\ 
r-tmp\_out\_G09
\includegraphics[width=8cm]{Images/ImagesFromVMD/r-tmp_out_G09.jpg}
\\
Inertia Tensor - Molecule A \\
\vtab
\begin{tabular}{|c c c|}
103.745	 & 	-1.51033	 & 	-1.05456	 \\
-1.51033	 & 	629.842	 & 	1.67376	 \\
-1.05456	 & 	1.67376	 & 	722.189
\end{tabular}

\vtab
 EingenVectors - Molecule A     \\
\vtab
\begin{tabular}{|c c c|}
-0.999994	 & 	-0.00286537	 & 	-0.0016974	 \\
0.00283409	 & 	-0.999831	 & 	0.0181489	 \\
-0.00174912	 & 	0.018144	 & 	0.999834
\end{tabular}

\vtab
 EingenValues - Molecule A     \\
\vtab
\begin{tabular}{|c c c|}
103.739	 & 	629.816	 & 	722.221
\end{tabular}
\columnbreak

Molecule B \\ 
r-tmp\_rotated\_out\_G09
\includegraphics[width=8cm]{Images/ImagesFromVMD/r-tmp_rotated_out_G09.jpg}
\\
Inertia Tensor - Molecule B \\
\vtab
\begin{tabular}{|c c c|}
103.745	 & 	-1.51033	 & 	-1.05456	 \\
-1.51033	 & 	629.842	 & 	1.67376	 \\
-1.05456	 & 	1.67376	 & 	722.189
\end{tabular}

\vtab
 EingenVectors - Molecule B     \\
\vtab
\begin{tabular}{|c c c|}
-0.999994	 & 	-0.00286537	 & 	-0.0016974	 \\
0.00283409	 & 	-0.999831	 & 	0.0181489	 \\
-0.00174912	 & 	0.018144	 & 	0.999834
\end{tabular}

\vtab
 EingenValues - Molecule B     \\
\vtab
\begin{tabular}{|c c c|}
103.739	 & 	629.816	 & 	722.221
\end{tabular}

\end{center}
\end{multicols}
\begin{center}
\vtab
\vtab
\textcolor{NavyBlue}{\Large Are the same molecule}
\end{center}
%%%%%%%%%%%%%%%%%%%%%%%%%%%%%%%%%%%%%%%%%%%%%%%%%%%%%%%%
%%%%%%%%%%%%%%%%%%%%%%%%%%%%%%%%%%%%%%%%%%%%%%%%%%%%%%%%
\newpage
%%%%%%%%%%%%%%%%%%%%%%%%%%%%%%%%%%%%%%%%%%%%%%%%%%%%%%%%
%%%%%%%%%%%%%%%%%%%%%%%%%%%%%%%%%%%%%%%%%%%%%%%%%%%%%%%%

\vtab[-2cm]
\begin{center}
{\large SymmetricEnantiomers}
\end{center}
\begin{multicols}{2}
\begin{center}
Molecule A \\ 
r-tmp\_out\_G09
\includegraphics[width=8cm]{Images/ImagesFromVMD/r-tmp_out_G09.jpg}
\\
Inertia Tensor - Molecule A \\
\vtab
\begin{tabular}{|c c c|}
103.745	 & 	-1.51033	 & 	-1.05456	 \\
-1.51033	 & 	629.842	 & 	1.67376	 \\
-1.05456	 & 	1.67376	 & 	722.189
\end{tabular}

\vtab
 EingenVectors - Molecule A     \\
\vtab
\begin{tabular}{|c c c|}
-0.999994	 & 	-0.00286537	 & 	-0.0016974	 \\
0.00283409	 & 	-0.999831	 & 	0.0181489	 \\
-0.00174912	 & 	0.018144	 & 	0.999834
\end{tabular}

\vtab
 EingenValues - Molecule A     \\
\vtab
\begin{tabular}{|c c c|}
103.739	 & 	629.816	 & 	722.221
\end{tabular}
\columnbreak

Molecule B \\ 
s-enan\_out\_G09
\includegraphics[width=8cm]{Images/ImagesFromVMD/s-enan_out_G09.jpg}
\\
Inertia Tensor - Molecule B \\
\vtab
\begin{tabular}{|c c c|}
103.707	 & 	-1.51142	 & 	-1.05403	 \\
-1.51142	 & 	630.103	 & 	1.67204	 \\
-1.05403	 & 	1.67204	 & 	722.42
\end{tabular}

\vtab
 EingenVectors - Molecule B     \\
\vtab
\begin{tabular}{|c c c|}
-0.999994	 & 	-0.00286582	 & 	-0.00169581	 \\
0.0028346	 & 	-0.999832	 & 	0.018136	 \\
-0.0017475	 & 	0.0181311	 & 	0.999834
\end{tabular}

\vtab
 EingenValues - Molecule B     \\
\vtab
\begin{tabular}{|c c c|}
103.701	 & 	630.077	 & 	722.452
\end{tabular}

\end{center}
\end{multicols}
\begin{center}
\vtab
\vtab
\textcolor{NavyBlue}{\Large Are the same molecule}
\end{center}
%%%%%%%%%%%%%%%%%%%%%%%%%%%%%%%%%%%%%%%%%%%%%%%%%%%%%%%%
%%%%%%%%%%%%%%%%%%%%%%%%%%%%%%%%%%%%%%%%%%%%%%%%%%%%%%%%
\newpage
%%%%%%%%%%%%%%%%%%%%%%%%%%%%%%%%%%%%%%%%%%%%%%%%%%%%%%%%
%%%%%%%%%%%%%%%%%%%%%%%%%%%%%%%%%%%%%%%%%%%%%%%%%%%%%%%%

\vtab[-2cm]
\begin{center}
{\large SymmetricEnantiomers}
\end{center}
\begin{multicols}{2}
\begin{center}
Molecule A \\ 
r-tmp\_out\_G09
\includegraphics[width=8cm]{Images/ImagesFromVMD/r-tmp_out_G09.jpg}
\\
Inertia Tensor - Molecule A \\
\vtab
\begin{tabular}{|c c c|}
103.745	 & 	-1.51033	 & 	-1.05456	 \\
-1.51033	 & 	629.842	 & 	1.67376	 \\
-1.05456	 & 	1.67376	 & 	722.189
\end{tabular}

\vtab
 EingenVectors - Molecule A     \\
\vtab
\begin{tabular}{|c c c|}
-0.999994	 & 	-0.00286537	 & 	-0.0016974	 \\
0.00283409	 & 	-0.999831	 & 	0.0181489	 \\
-0.00174912	 & 	0.018144	 & 	0.999834
\end{tabular}

\vtab
 EingenValues - Molecule A     \\
\vtab
\begin{tabular}{|c c c|}
103.739	 & 	629.816	 & 	722.221
\end{tabular}
\columnbreak

Molecule B \\ 
s-enan\_rotated\_out\_G09
\includegraphics[width=8cm]{Images/ImagesFromVMD/s-enan_rotated_out_G09.jpg}
\\
Inertia Tensor - Molecule B \\
\vtab
\begin{tabular}{|c c c|}
103.707	 & 	-1.51138	 & 	-1.05403	 \\
-1.51138	 & 	630.103	 & 	1.67204	 \\
-1.05403	 & 	1.67204	 & 	722.42
\end{tabular}

\vtab
 EingenVectors - Molecule B     \\
\vtab
\begin{tabular}{|c c c|}
-0.999994	 & 	-0.00286575	 & 	-0.00169582	 \\
0.00283453	 & 	-0.999832	 & 	0.018136	 \\
-0.0017475	 & 	0.0181311	 & 	0.999834
\end{tabular}

\vtab
 EingenValues - Molecule B     \\
\vtab
\begin{tabular}{|c c c|}
103.701	 & 	630.077	 & 	722.452
\end{tabular}

\end{center}
\end{multicols}
\begin{center}
\vtab
\vtab
\textcolor{NavyBlue}{\Large Are the same molecule}
\end{center}
%%%%%%%%%%%%%%%%%%%%%%%%%%%%%%%%%%%%%%%%%%%%%%%%%%%%%%%%
%%%%%%%%%%%%%%%%%%%%%%%%%%%%%%%%%%%%%%%%%%%%%%%%%%%%%%%%
\newpage
%%%%%%%%%%%%%%%%%%%%%%%%%%%%%%%%%%%%%%%%%%%%%%%%%%%%%%%%
%%%%%%%%%%%%%%%%%%%%%%%%%%%%%%%%%%%%%%%%%%%%%%%%%%%%%%%%

\vtab[-2cm]
\begin{center}
{\large SymmetricEnantiomers}
\end{center}
\begin{multicols}{2}
\begin{center}
Molecule A \\ 
r-tmp\_out\_G09
\includegraphics[width=8cm]{Images/ImagesFromVMD/r-tmp_out_G09.jpg}
\\
Inertia Tensor - Molecule A \\
\vtab
\begin{tabular}{|c c c|}
103.745	 & 	-1.51033	 & 	-1.05456	 \\
-1.51033	 & 	629.842	 & 	1.67376	 \\
-1.05456	 & 	1.67376	 & 	722.189
\end{tabular}

\vtab
 EingenVectors - Molecule A     \\
\vtab
\begin{tabular}{|c c c|}
-0.999994	 & 	-0.00286537	 & 	-0.0016974	 \\
0.00283409	 & 	-0.999831	 & 	0.0181489	 \\
-0.00174912	 & 	0.018144	 & 	0.999834
\end{tabular}

\vtab
 EingenValues - Molecule A     \\
\vtab
\begin{tabular}{|c c c|}
103.739	 & 	629.816	 & 	722.221
\end{tabular}
\columnbreak

Molecule B \\ 
s-tmp\_out\_G09
\includegraphics[width=8cm]{Images/ImagesFromVMD/s-tmp_out_G09.jpg}
\\
Inertia Tensor - Molecule B \\
\vtab
\begin{tabular}{|c c c|}
103.707	 & 	1.51142	 & 	1.05403	 \\
1.51142	 & 	630.103	 & 	1.67204	 \\
1.05403	 & 	1.67204	 & 	722.42
\end{tabular}

\vtab
 EingenVectors - Molecule B     \\
\vtab
\begin{tabular}{|c c c|}
-0.999994	 & 	0.00286582	 & 	0.00169581	 \\
-0.0028346	 & 	-0.999832	 & 	0.018136	 \\
0.0017475	 & 	0.0181311	 & 	0.999834
\end{tabular}

\vtab
 EingenValues - Molecule B     \\
\vtab
\begin{tabular}{|c c c|}
103.701	 & 	630.077	 & 	722.452
\end{tabular}

\end{center}
\end{multicols}
\begin{center}
\vtab
\vtab
\textcolor{NavyBlue}{\Large Are Isomer and Stereoisomers}
\end{center}
%%%%%%%%%%%%%%%%%%%%%%%%%%%%%%%%%%%%%%%%%%%%%%%%%%%%%%%%
%%%%%%%%%%%%%%%%%%%%%%%%%%%%%%%%%%%%%%%%%%%%%%%%%%%%%%%%
\newpage
%%%%%%%%%%%%%%%%%%%%%%%%%%%%%%%%%%%%%%%%%%%%%%%%%%%%%%%%
%%%%%%%%%%%%%%%%%%%%%%%%%%%%%%%%%%%%%%%%%%%%%%%%%%%%%%%%

\vtab[-2cm]
\begin{center}
{\large SymmetricEnantiomers}
\end{center}
\begin{multicols}{2}
\begin{center}
Molecule A \\ 
r-tmp\_out\_G09
\includegraphics[width=8cm]{Images/ImagesFromVMD/r-tmp_out_G09.jpg}
\\
Inertia Tensor - Molecule A \\
\vtab
\begin{tabular}{|c c c|}
103.745	 & 	-1.51033	 & 	-1.05456	 \\
-1.51033	 & 	629.842	 & 	1.67376	 \\
-1.05456	 & 	1.67376	 & 	722.189
\end{tabular}

\vtab
 EingenVectors - Molecule A     \\
\vtab
\begin{tabular}{|c c c|}
-0.999994	 & 	-0.00286537	 & 	-0.0016974	 \\
0.00283409	 & 	-0.999831	 & 	0.0181489	 \\
-0.00174912	 & 	0.018144	 & 	0.999834
\end{tabular}

\vtab
 EingenValues - Molecule A     \\
\vtab
\begin{tabular}{|c c c|}
103.739	 & 	629.816	 & 	722.221
\end{tabular}
\columnbreak

Molecule B \\ 
s-tmp\_rotated\_out\_G09
\includegraphics[width=8cm]{Images/ImagesFromVMD/s-tmp_rotated_out_G09.jpg}
\\
Inertia Tensor - Molecule B \\
\vtab
\begin{tabular}{|c c c|}
103.707	 & 	1.51142	 & 	1.05403	 \\
1.51142	 & 	630.103	 & 	1.67204	 \\
1.05403	 & 	1.67204	 & 	722.42
\end{tabular}

\vtab
 EingenVectors - Molecule B     \\
\vtab
\begin{tabular}{|c c c|}
-0.999994	 & 	0.00286582	 & 	0.00169581	 \\
-0.0028346	 & 	-0.999832	 & 	0.018136	 \\
0.0017475	 & 	0.0181311	 & 	0.999834
\end{tabular}

\vtab
 EingenValues - Molecule B     \\
\vtab
\begin{tabular}{|c c c|}
103.701	 & 	630.077	 & 	722.452
\end{tabular}

\end{center}
\end{multicols}
\begin{center}
\vtab
\vtab
\textcolor{NavyBlue}{\Large Are Isomer and Stereoisomers}
\end{center}
%%%%%%%%%%%%%%%%%%%%%%%%%%%%%%%%%%%%%%%%%%%%%%%%%%%%%%%%
%%%%%%%%%%%%%%%%%%%%%%%%%%%%%%%%%%%%%%%%%%%%%%%%%%%%%%%%
\newpage
%%%%%%%%%%%%%%%%%%%%%%%%%%%%%%%%%%%%%%%%%%%%%%%%%%%%%%%%
%%%%%%%%%%%%%%%%%%%%%%%%%%%%%%%%%%%%%%%%%%%%%%%%%%%%%%%%

\vtab[-2cm]
\begin{center}
{\large SymmetricEnantiomers}
\end{center}
\begin{multicols}{2}
\begin{center}
Molecule A \\ 
r-tmp\_rotated\_out\_G09
\includegraphics[width=8cm]{Images/ImagesFromVMD/r-tmp_rotated_out_G09.jpg}
\\
Inertia Tensor - Molecule A \\
\vtab
\begin{tabular}{|c c c|}
103.745	 & 	-1.51033	 & 	-1.05456	 \\
-1.51033	 & 	629.842	 & 	1.67376	 \\
-1.05456	 & 	1.67376	 & 	722.189
\end{tabular}

\vtab
 EingenVectors - Molecule A     \\
\vtab
\begin{tabular}{|c c c|}
-0.999994	 & 	-0.00286537	 & 	-0.0016974	 \\
0.00283409	 & 	-0.999831	 & 	0.0181489	 \\
-0.00174912	 & 	0.018144	 & 	0.999834
\end{tabular}

\vtab
 EingenValues - Molecule A     \\
\vtab
\begin{tabular}{|c c c|}
103.739	 & 	629.816	 & 	722.221
\end{tabular}
\columnbreak

Molecule B \\ 
s-enan\_out\_G09
\includegraphics[width=8cm]{Images/ImagesFromVMD/s-enan_out_G09.jpg}
\\
Inertia Tensor - Molecule B \\
\vtab
\begin{tabular}{|c c c|}
103.707	 & 	-1.51142	 & 	-1.05403	 \\
-1.51142	 & 	630.103	 & 	1.67204	 \\
-1.05403	 & 	1.67204	 & 	722.42
\end{tabular}

\vtab
 EingenVectors - Molecule B     \\
\vtab
\begin{tabular}{|c c c|}
-0.999994	 & 	-0.00286582	 & 	-0.00169581	 \\
0.0028346	 & 	-0.999832	 & 	0.018136	 \\
-0.0017475	 & 	0.0181311	 & 	0.999834
\end{tabular}

\vtab
 EingenValues - Molecule B     \\
\vtab
\begin{tabular}{|c c c|}
103.701	 & 	630.077	 & 	722.452
\end{tabular}

\end{center}
\end{multicols}
\begin{center}
\vtab
\vtab
\textcolor{NavyBlue}{\Large Are the same molecule}
\end{center}
%%%%%%%%%%%%%%%%%%%%%%%%%%%%%%%%%%%%%%%%%%%%%%%%%%%%%%%%
%%%%%%%%%%%%%%%%%%%%%%%%%%%%%%%%%%%%%%%%%%%%%%%%%%%%%%%%
\newpage
%%%%%%%%%%%%%%%%%%%%%%%%%%%%%%%%%%%%%%%%%%%%%%%%%%%%%%%%
%%%%%%%%%%%%%%%%%%%%%%%%%%%%%%%%%%%%%%%%%%%%%%%%%%%%%%%%

\vtab[-2cm]
\begin{center}
{\large SymmetricEnantiomers}
\end{center}
\begin{multicols}{2}
\begin{center}
Molecule A \\ 
r-tmp\_rotated\_out\_G09
\includegraphics[width=8cm]{Images/ImagesFromVMD/r-tmp_rotated_out_G09.jpg}
\\
Inertia Tensor - Molecule A \\
\vtab
\begin{tabular}{|c c c|}
103.745	 & 	-1.51033	 & 	-1.05456	 \\
-1.51033	 & 	629.842	 & 	1.67376	 \\
-1.05456	 & 	1.67376	 & 	722.189
\end{tabular}

\vtab
 EingenVectors - Molecule A     \\
\vtab
\begin{tabular}{|c c c|}
-0.999994	 & 	-0.00286537	 & 	-0.0016974	 \\
0.00283409	 & 	-0.999831	 & 	0.0181489	 \\
-0.00174912	 & 	0.018144	 & 	0.999834
\end{tabular}

\vtab
 EingenValues - Molecule A     \\
\vtab
\begin{tabular}{|c c c|}
103.739	 & 	629.816	 & 	722.221
\end{tabular}
\columnbreak

Molecule B \\ 
s-enan\_rotated\_out\_G09
\includegraphics[width=8cm]{Images/ImagesFromVMD/s-enan_rotated_out_G09.jpg}
\\
Inertia Tensor - Molecule B \\
\vtab
\begin{tabular}{|c c c|}
103.707	 & 	-1.51138	 & 	-1.05403	 \\
-1.51138	 & 	630.103	 & 	1.67204	 \\
-1.05403	 & 	1.67204	 & 	722.42
\end{tabular}

\vtab
 EingenVectors - Molecule B     \\
\vtab
\begin{tabular}{|c c c|}
-0.999994	 & 	-0.00286575	 & 	-0.00169582	 \\
0.00283453	 & 	-0.999832	 & 	0.018136	 \\
-0.0017475	 & 	0.0181311	 & 	0.999834
\end{tabular}

\vtab
 EingenValues - Molecule B     \\
\vtab
\begin{tabular}{|c c c|}
103.701	 & 	630.077	 & 	722.452
\end{tabular}

\end{center}
\end{multicols}
\begin{center}
\vtab
\vtab
\textcolor{NavyBlue}{\Large Are the same molecule}
\end{center}
%%%%%%%%%%%%%%%%%%%%%%%%%%%%%%%%%%%%%%%%%%%%%%%%%%%%%%%%
%%%%%%%%%%%%%%%%%%%%%%%%%%%%%%%%%%%%%%%%%%%%%%%%%%%%%%%%
\newpage
%%%%%%%%%%%%%%%%%%%%%%%%%%%%%%%%%%%%%%%%%%%%%%%%%%%%%%%%
%%%%%%%%%%%%%%%%%%%%%%%%%%%%%%%%%%%%%%%%%%%%%%%%%%%%%%%%

\vtab[-2cm]
\begin{center}
{\large SymmetricEnantiomers}
\end{center}
\begin{multicols}{2}
\begin{center}
Molecule A \\ 
r-tmp\_rotated\_out\_G09
\includegraphics[width=8cm]{Images/ImagesFromVMD/r-tmp_rotated_out_G09.jpg}
\\
Inertia Tensor - Molecule A \\
\vtab
\begin{tabular}{|c c c|}
103.745	 & 	-1.51033	 & 	-1.05456	 \\
-1.51033	 & 	629.842	 & 	1.67376	 \\
-1.05456	 & 	1.67376	 & 	722.189
\end{tabular}

\vtab
 EingenVectors - Molecule A     \\
\vtab
\begin{tabular}{|c c c|}
-0.999994	 & 	-0.00286537	 & 	-0.0016974	 \\
0.00283409	 & 	-0.999831	 & 	0.0181489	 \\
-0.00174912	 & 	0.018144	 & 	0.999834
\end{tabular}

\vtab
 EingenValues - Molecule A     \\
\vtab
\begin{tabular}{|c c c|}
103.739	 & 	629.816	 & 	722.221
\end{tabular}
\columnbreak

Molecule B \\ 
s-tmp\_out\_G09
\includegraphics[width=8cm]{Images/ImagesFromVMD/s-tmp_out_G09.jpg}
\\
Inertia Tensor - Molecule B \\
\vtab
\begin{tabular}{|c c c|}
103.707	 & 	1.51142	 & 	1.05403	 \\
1.51142	 & 	630.103	 & 	1.67204	 \\
1.05403	 & 	1.67204	 & 	722.42
\end{tabular}

\vtab
 EingenVectors - Molecule B     \\
\vtab
\begin{tabular}{|c c c|}
-0.999994	 & 	0.00286582	 & 	0.00169581	 \\
-0.0028346	 & 	-0.999832	 & 	0.018136	 \\
0.0017475	 & 	0.0181311	 & 	0.999834
\end{tabular}

\vtab
 EingenValues - Molecule B     \\
\vtab
\begin{tabular}{|c c c|}
103.701	 & 	630.077	 & 	722.452
\end{tabular}

\end{center}
\end{multicols}
\begin{center}
\vtab
\vtab
\textcolor{NavyBlue}{\Large Are Isomer and Stereoisomers}
\end{center}
%%%%%%%%%%%%%%%%%%%%%%%%%%%%%%%%%%%%%%%%%%%%%%%%%%%%%%%%
%%%%%%%%%%%%%%%%%%%%%%%%%%%%%%%%%%%%%%%%%%%%%%%%%%%%%%%%
\newpage
%%%%%%%%%%%%%%%%%%%%%%%%%%%%%%%%%%%%%%%%%%%%%%%%%%%%%%%%
%%%%%%%%%%%%%%%%%%%%%%%%%%%%%%%%%%%%%%%%%%%%%%%%%%%%%%%%

\vtab[-2cm]
\begin{center}
{\large SymmetricEnantiomers}
\end{center}
\begin{multicols}{2}
\begin{center}
Molecule A \\ 
r-tmp\_rotated\_out\_G09
\includegraphics[width=8cm]{Images/ImagesFromVMD/r-tmp_rotated_out_G09.jpg}
\\
Inertia Tensor - Molecule A \\
\vtab
\begin{tabular}{|c c c|}
103.745	 & 	-1.51033	 & 	-1.05456	 \\
-1.51033	 & 	629.842	 & 	1.67376	 \\
-1.05456	 & 	1.67376	 & 	722.189
\end{tabular}

\vtab
 EingenVectors - Molecule A     \\
\vtab
\begin{tabular}{|c c c|}
-0.999994	 & 	-0.00286537	 & 	-0.0016974	 \\
0.00283409	 & 	-0.999831	 & 	0.0181489	 \\
-0.00174912	 & 	0.018144	 & 	0.999834
\end{tabular}

\vtab
 EingenValues - Molecule A     \\
\vtab
\begin{tabular}{|c c c|}
103.739	 & 	629.816	 & 	722.221
\end{tabular}
\columnbreak

Molecule B \\ 
s-tmp\_rotated\_out\_G09
\includegraphics[width=8cm]{Images/ImagesFromVMD/s-tmp_rotated_out_G09.jpg}
\\
Inertia Tensor - Molecule B \\
\vtab
\begin{tabular}{|c c c|}
103.707	 & 	1.51142	 & 	1.05403	 \\
1.51142	 & 	630.103	 & 	1.67204	 \\
1.05403	 & 	1.67204	 & 	722.42
\end{tabular}

\vtab
 EingenVectors - Molecule B     \\
\vtab
\begin{tabular}{|c c c|}
-0.999994	 & 	0.00286582	 & 	0.00169581	 \\
-0.0028346	 & 	-0.999832	 & 	0.018136	 \\
0.0017475	 & 	0.0181311	 & 	0.999834
\end{tabular}

\vtab
 EingenValues - Molecule B     \\
\vtab
\begin{tabular}{|c c c|}
103.701	 & 	630.077	 & 	722.452
\end{tabular}

\end{center}
\end{multicols}
\begin{center}
\vtab
\vtab
\textcolor{NavyBlue}{\Large Are Isomer and Stereoisomers}
\end{center}
%%%%%%%%%%%%%%%%%%%%%%%%%%%%%%%%%%%%%%%%%%%%%%%%%%%%%%%%
%%%%%%%%%%%%%%%%%%%%%%%%%%%%%%%%%%%%%%%%%%%%%%%%%%%%%%%%
\newpage
%%%%%%%%%%%%%%%%%%%%%%%%%%%%%%%%%%%%%%%%%%%%%%%%%%%%%%%%
%%%%%%%%%%%%%%%%%%%%%%%%%%%%%%%%%%%%%%%%%%%%%%%%%%%%%%%%

\vtab[-2cm]
\begin{center}
{\large SymmetricEnantiomers}
\end{center}
\begin{multicols}{2}
\begin{center}
Molecule A \\ 
s-enan\_out\_G09
\includegraphics[width=8cm]{Images/ImagesFromVMD/s-enan_out_G09.jpg}
\\
Inertia Tensor - Molecule A \\
\vtab
\begin{tabular}{|c c c|}
103.707	 & 	-1.51142	 & 	-1.05403	 \\
-1.51142	 & 	630.103	 & 	1.67204	 \\
-1.05403	 & 	1.67204	 & 	722.42
\end{tabular}

\vtab
 EingenVectors - Molecule A     \\
\vtab
\begin{tabular}{|c c c|}
-0.999994	 & 	-0.00286582	 & 	-0.00169581	 \\
0.0028346	 & 	-0.999832	 & 	0.018136	 \\
-0.0017475	 & 	0.0181311	 & 	0.999834
\end{tabular}

\vtab
 EingenValues - Molecule A     \\
\vtab
\begin{tabular}{|c c c|}
103.701	 & 	630.077	 & 	722.452
\end{tabular}
\columnbreak

Molecule B \\ 
s-enan\_rotated\_out\_G09
\includegraphics[width=8cm]{Images/ImagesFromVMD/s-enan_rotated_out_G09.jpg}
\\
Inertia Tensor - Molecule B \\
\vtab
\begin{tabular}{|c c c|}
103.707	 & 	-1.51138	 & 	-1.05403	 \\
-1.51138	 & 	630.103	 & 	1.67204	 \\
-1.05403	 & 	1.67204	 & 	722.42
\end{tabular}

\vtab
 EingenVectors - Molecule B     \\
\vtab
\begin{tabular}{|c c c|}
-0.999994	 & 	-0.00286575	 & 	-0.00169582	 \\
0.00283453	 & 	-0.999832	 & 	0.018136	 \\
-0.0017475	 & 	0.0181311	 & 	0.999834
\end{tabular}

\vtab
 EingenValues - Molecule B     \\
\vtab
\begin{tabular}{|c c c|}
103.701	 & 	630.077	 & 	722.452
\end{tabular}

\end{center}
\end{multicols}
\begin{center}
\vtab
\vtab
\textcolor{NavyBlue}{\Large Are the same molecule}
\end{center}
%%%%%%%%%%%%%%%%%%%%%%%%%%%%%%%%%%%%%%%%%%%%%%%%%%%%%%%%
%%%%%%%%%%%%%%%%%%%%%%%%%%%%%%%%%%%%%%%%%%%%%%%%%%%%%%%%
\newpage
%%%%%%%%%%%%%%%%%%%%%%%%%%%%%%%%%%%%%%%%%%%%%%%%%%%%%%%%
%%%%%%%%%%%%%%%%%%%%%%%%%%%%%%%%%%%%%%%%%%%%%%%%%%%%%%%%

\vtab[-2cm]
\begin{center}
{\large SymmetricEnantiomers}
\end{center}
\begin{multicols}{2}
\begin{center}
Molecule A \\ 
s-enan\_out\_G09
\includegraphics[width=8cm]{Images/ImagesFromVMD/s-enan_out_G09.jpg}
\\
Inertia Tensor - Molecule A \\
\vtab
\begin{tabular}{|c c c|}
103.707	 & 	-1.51142	 & 	-1.05403	 \\
-1.51142	 & 	630.103	 & 	1.67204	 \\
-1.05403	 & 	1.67204	 & 	722.42
\end{tabular}

\vtab
 EingenVectors - Molecule A     \\
\vtab
\begin{tabular}{|c c c|}
-0.999994	 & 	-0.00286582	 & 	-0.00169581	 \\
0.0028346	 & 	-0.999832	 & 	0.018136	 \\
-0.0017475	 & 	0.0181311	 & 	0.999834
\end{tabular}

\vtab
 EingenValues - Molecule A     \\
\vtab
\begin{tabular}{|c c c|}
103.701	 & 	630.077	 & 	722.452
\end{tabular}
\columnbreak

Molecule B \\ 
s-tmp\_out\_G09
\includegraphics[width=8cm]{Images/ImagesFromVMD/s-tmp_out_G09.jpg}
\\
Inertia Tensor - Molecule B \\
\vtab
\begin{tabular}{|c c c|}
103.707	 & 	1.51142	 & 	1.05403	 \\
1.51142	 & 	630.103	 & 	1.67204	 \\
1.05403	 & 	1.67204	 & 	722.42
\end{tabular}

\vtab
 EingenVectors - Molecule B     \\
\vtab
\begin{tabular}{|c c c|}
-0.999994	 & 	0.00286582	 & 	0.00169581	 \\
-0.0028346	 & 	-0.999832	 & 	0.018136	 \\
0.0017475	 & 	0.0181311	 & 	0.999834
\end{tabular}

\vtab
 EingenValues - Molecule B     \\
\vtab
\begin{tabular}{|c c c|}
103.701	 & 	630.077	 & 	722.452
\end{tabular}

\end{center}
\end{multicols}
\begin{center}
\vtab
\vtab
\textcolor{NavyBlue}{\Large Are Isomer and Stereoisomers}
\end{center}
%%%%%%%%%%%%%%%%%%%%%%%%%%%%%%%%%%%%%%%%%%%%%%%%%%%%%%%%
%%%%%%%%%%%%%%%%%%%%%%%%%%%%%%%%%%%%%%%%%%%%%%%%%%%%%%%%
\newpage
%%%%%%%%%%%%%%%%%%%%%%%%%%%%%%%%%%%%%%%%%%%%%%%%%%%%%%%%
%%%%%%%%%%%%%%%%%%%%%%%%%%%%%%%%%%%%%%%%%%%%%%%%%%%%%%%%

\vtab[-2cm]
\begin{center}
{\large SymmetricEnantiomers}
\end{center}
\begin{multicols}{2}
\begin{center}
Molecule A \\ 
s-enan\_out\_G09
\includegraphics[width=8cm]{Images/ImagesFromVMD/s-enan_out_G09.jpg}
\\
Inertia Tensor - Molecule A \\
\vtab
\begin{tabular}{|c c c|}
103.707	 & 	-1.51142	 & 	-1.05403	 \\
-1.51142	 & 	630.103	 & 	1.67204	 \\
-1.05403	 & 	1.67204	 & 	722.42
\end{tabular}

\vtab
 EingenVectors - Molecule A     \\
\vtab
\begin{tabular}{|c c c|}
-0.999994	 & 	-0.00286582	 & 	-0.00169581	 \\
0.0028346	 & 	-0.999832	 & 	0.018136	 \\
-0.0017475	 & 	0.0181311	 & 	0.999834
\end{tabular}

\vtab
 EingenValues - Molecule A     \\
\vtab
\begin{tabular}{|c c c|}
103.701	 & 	630.077	 & 	722.452
\end{tabular}
\columnbreak

Molecule B \\ 
s-tmp\_rotated\_out\_G09
\includegraphics[width=8cm]{Images/ImagesFromVMD/s-tmp_rotated_out_G09.jpg}
\\
Inertia Tensor - Molecule B \\
\vtab
\begin{tabular}{|c c c|}
103.707	 & 	1.51142	 & 	1.05403	 \\
1.51142	 & 	630.103	 & 	1.67204	 \\
1.05403	 & 	1.67204	 & 	722.42
\end{tabular}

\vtab
 EingenVectors - Molecule B     \\
\vtab
\begin{tabular}{|c c c|}
-0.999994	 & 	0.00286582	 & 	0.00169581	 \\
-0.0028346	 & 	-0.999832	 & 	0.018136	 \\
0.0017475	 & 	0.0181311	 & 	0.999834
\end{tabular}

\vtab
 EingenValues - Molecule B     \\
\vtab
\begin{tabular}{|c c c|}
103.701	 & 	630.077	 & 	722.452
\end{tabular}

\end{center}
\end{multicols}
\begin{center}
\vtab
\vtab
\textcolor{NavyBlue}{\Large Are Isomer and Stereoisomers}
\end{center}
%%%%%%%%%%%%%%%%%%%%%%%%%%%%%%%%%%%%%%%%%%%%%%%%%%%%%%%%
%%%%%%%%%%%%%%%%%%%%%%%%%%%%%%%%%%%%%%%%%%%%%%%%%%%%%%%%
\newpage
%%%%%%%%%%%%%%%%%%%%%%%%%%%%%%%%%%%%%%%%%%%%%%%%%%%%%%%%
%%%%%%%%%%%%%%%%%%%%%%%%%%%%%%%%%%%%%%%%%%%%%%%%%%%%%%%%

\vtab[-2cm]
\begin{center}
{\large SymmetricEnantiomers}
\end{center}
\begin{multicols}{2}
\begin{center}
Molecule A \\ 
s-enan\_rotated\_out\_G09
\includegraphics[width=8cm]{Images/ImagesFromVMD/s-enan_rotated_out_G09.jpg}
\\
Inertia Tensor - Molecule A \\
\vtab
\begin{tabular}{|c c c|}
103.707	 & 	-1.51138	 & 	-1.05403	 \\
-1.51138	 & 	630.103	 & 	1.67204	 \\
-1.05403	 & 	1.67204	 & 	722.42
\end{tabular}

\vtab
 EingenVectors - Molecule A     \\
\vtab
\begin{tabular}{|c c c|}
-0.999994	 & 	-0.00286575	 & 	-0.00169582	 \\
0.00283453	 & 	-0.999832	 & 	0.018136	 \\
-0.0017475	 & 	0.0181311	 & 	0.999834
\end{tabular}

\vtab
 EingenValues - Molecule A     \\
\vtab
\begin{tabular}{|c c c|}
103.701	 & 	630.077	 & 	722.452
\end{tabular}
\columnbreak

Molecule B \\ 
s-tmp\_out\_G09
\includegraphics[width=8cm]{Images/ImagesFromVMD/s-tmp_out_G09.jpg}
\\
Inertia Tensor - Molecule B \\
\vtab
\begin{tabular}{|c c c|}
103.707	 & 	1.51142	 & 	1.05403	 \\
1.51142	 & 	630.103	 & 	1.67204	 \\
1.05403	 & 	1.67204	 & 	722.42
\end{tabular}

\vtab
 EingenVectors - Molecule B     \\
\vtab
\begin{tabular}{|c c c|}
-0.999994	 & 	0.00286582	 & 	0.00169581	 \\
-0.0028346	 & 	-0.999832	 & 	0.018136	 \\
0.0017475	 & 	0.0181311	 & 	0.999834
\end{tabular}

\vtab
 EingenValues - Molecule B     \\
\vtab
\begin{tabular}{|c c c|}
103.701	 & 	630.077	 & 	722.452
\end{tabular}

\end{center}
\end{multicols}
\begin{center}
\vtab
\vtab
\textcolor{NavyBlue}{\Large Are Isomer and Stereoisomers}
\end{center}
%%%%%%%%%%%%%%%%%%%%%%%%%%%%%%%%%%%%%%%%%%%%%%%%%%%%%%%%
%%%%%%%%%%%%%%%%%%%%%%%%%%%%%%%%%%%%%%%%%%%%%%%%%%%%%%%%
\newpage
%%%%%%%%%%%%%%%%%%%%%%%%%%%%%%%%%%%%%%%%%%%%%%%%%%%%%%%%
%%%%%%%%%%%%%%%%%%%%%%%%%%%%%%%%%%%%%%%%%%%%%%%%%%%%%%%%

\vtab[-2cm]
\begin{center}
{\large SymmetricEnantiomers}
\end{center}
\begin{multicols}{2}
\begin{center}
Molecule A \\ 
s-enan\_rotated\_out\_G09
\includegraphics[width=8cm]{Images/ImagesFromVMD/s-enan_rotated_out_G09.jpg}
\\
Inertia Tensor - Molecule A \\
\vtab
\begin{tabular}{|c c c|}
103.707	 & 	-1.51138	 & 	-1.05403	 \\
-1.51138	 & 	630.103	 & 	1.67204	 \\
-1.05403	 & 	1.67204	 & 	722.42
\end{tabular}

\vtab
 EingenVectors - Molecule A     \\
\vtab
\begin{tabular}{|c c c|}
-0.999994	 & 	-0.00286575	 & 	-0.00169582	 \\
0.00283453	 & 	-0.999832	 & 	0.018136	 \\
-0.0017475	 & 	0.0181311	 & 	0.999834
\end{tabular}

\vtab
 EingenValues - Molecule A     \\
\vtab
\begin{tabular}{|c c c|}
103.701	 & 	630.077	 & 	722.452
\end{tabular}
\columnbreak

Molecule B \\ 
s-tmp\_rotated\_out\_G09
\includegraphics[width=8cm]{Images/ImagesFromVMD/s-tmp_rotated_out_G09.jpg}
\\
Inertia Tensor - Molecule B \\
\vtab
\begin{tabular}{|c c c|}
103.707	 & 	1.51142	 & 	1.05403	 \\
1.51142	 & 	630.103	 & 	1.67204	 \\
1.05403	 & 	1.67204	 & 	722.42
\end{tabular}

\vtab
 EingenVectors - Molecule B     \\
\vtab
\begin{tabular}{|c c c|}
-0.999994	 & 	0.00286582	 & 	0.00169581	 \\
-0.0028346	 & 	-0.999832	 & 	0.018136	 \\
0.0017475	 & 	0.0181311	 & 	0.999834
\end{tabular}

\vtab
 EingenValues - Molecule B     \\
\vtab
\begin{tabular}{|c c c|}
103.701	 & 	630.077	 & 	722.452
\end{tabular}

\end{center}
\end{multicols}
\begin{center}
\vtab
\vtab
\textcolor{NavyBlue}{\Large Are Isomer and Stereoisomers}
\end{center}
%%%%%%%%%%%%%%%%%%%%%%%%%%%%%%%%%%%%%%%%%%%%%%%%%%%%%%%%
%%%%%%%%%%%%%%%%%%%%%%%%%%%%%%%%%%%%%%%%%%%%%%%%%%%%%%%%
\newpage
%%%%%%%%%%%%%%%%%%%%%%%%%%%%%%%%%%%%%%%%%%%%%%%%%%%%%%%%
%%%%%%%%%%%%%%%%%%%%%%%%%%%%%%%%%%%%%%%%%%%%%%%%%%%%%%%%

\vtab[-2cm]
\begin{center}
{\large SymmetricEnantiomers}
\end{center}
\begin{multicols}{2}
\begin{center}
Molecule A \\ 
s-tmp\_out\_G09
\includegraphics[width=8cm]{Images/ImagesFromVMD/s-tmp_out_G09.jpg}
\\
Inertia Tensor - Molecule A \\
\vtab
\begin{tabular}{|c c c|}
103.707	 & 	1.51142	 & 	1.05403	 \\
1.51142	 & 	630.103	 & 	1.67204	 \\
1.05403	 & 	1.67204	 & 	722.42
\end{tabular}

\vtab
 EingenVectors - Molecule A     \\
\vtab
\begin{tabular}{|c c c|}
-0.999994	 & 	0.00286582	 & 	0.00169581	 \\
-0.0028346	 & 	-0.999832	 & 	0.018136	 \\
0.0017475	 & 	0.0181311	 & 	0.999834
\end{tabular}

\vtab
 EingenValues - Molecule A     \\
\vtab
\begin{tabular}{|c c c|}
103.701	 & 	630.077	 & 	722.452
\end{tabular}
\columnbreak

Molecule B \\ 
s-tmp\_rotated\_out\_G09
\includegraphics[width=8cm]{Images/ImagesFromVMD/s-tmp_rotated_out_G09.jpg}
\\
Inertia Tensor - Molecule B \\
\vtab
\begin{tabular}{|c c c|}
103.707	 & 	1.51142	 & 	1.05403	 \\
1.51142	 & 	630.103	 & 	1.67204	 \\
1.05403	 & 	1.67204	 & 	722.42
\end{tabular}

\vtab
 EingenVectors - Molecule B     \\
\vtab
\begin{tabular}{|c c c|}
-0.999994	 & 	0.00286582	 & 	0.00169581	 \\
-0.0028346	 & 	-0.999832	 & 	0.018136	 \\
0.0017475	 & 	0.0181311	 & 	0.999834
\end{tabular}

\vtab
 EingenValues - Molecule B     \\
\vtab
\begin{tabular}{|c c c|}
103.701	 & 	630.077	 & 	722.452
\end{tabular}

\end{center}
\end{multicols}
\begin{center}
\vtab
\vtab
\textcolor{NavyBlue}{\Large Are the same molecule}
\end{center}
%%%%%%%%%%%%%%%%%%%%%%%%%%%%%%%%%%%%%%%%%%%%%%%%%%%%%%%%
%%%%%%%%%%%%%%%%%%%%%%%%%%%%%%%%%%%%%%%%%%%%%%%%%%%%%%%%
\newpage
%%%%%%%%%%%%%%%%%%%%%%%%%%%%%%%%%%%%%%%%%%%%%%%%%%%%%%%%
%%%%%%%%%%%%%%%%%%%%%%%%%%%%%%%%%%%%%%%%%%%%%%%%%%%%%%%%

 %%%%%%%%%%%%%%%%%%%%

\vtab[-2cm]
\tab[-2cm]
\begin{tabular}{c|m{8cm}|c|c}
\# & Moléculas & Restultado esperado & Resultado programa \\ \hline\hline
\multirow{4}{*}{\tab[2mm] 1 \tab[2mm]} & r-2-chlorobutane\_out\_G09\_inversion &
\multirow{2}{*}{Enantiomers} & \multirow{2}{*}{Enantiomers}
\\
& E = -618.057081944 \tab Freq =115.2921   &    &  \\ \cline{2-2}
& r-2-chlorobutane\_rotated\_out\_G09   & \multicolumn{2}{c}
 {RMS = 0.02719433}
\\
& E = -618.057081944 \tab Freq =115.2921   &    \multicolumn{2}{c} 
{ }
\\ \hline
\multirow{4}{*}{\tab[2mm] 2 \tab[2mm]} & r-2-chlorobutane\_out\_G09\_inversion &
\multirow{2}{*}{Equal} & \multirow{2}{*}{Equal}
\\
& E = -618.057081944 \tab Freq =115.2921   &    &  \\ \cline{2-2}
& s-2-chlorobutane\_out\_G09   & \multicolumn{2}{c}
{ RMS = 4.366941E-07}
\\
& E = -618.057081943 \tab Freq =115.2939   &    \multicolumn{2}{c} 
{ }
\\ \hline
\multirow{4}{*}{\tab[2mm] 3 \tab[2mm]} & r-2-chlorobutane\_out\_G09\_inversion &
\multirow{2}{*}{Enantiomers} & \multirow{2}{*}{Enantiomers}
\\
& E = -618.057081944 \tab Freq =115.2921   &    &  \\ \cline{2-2}
& s-2-chlorobutane\_out\_G09\_inversion   & \multicolumn{2}{c}
 {RMS = 0.2294318}
\\
& E = -618.057081943 \tab Freq =115.2939   &    \multicolumn{2}{c} 
{ }
\\ \hline
\multirow{4}{*}{\tab[2mm] 4 \tab[2mm]} & r-2-chlorobutane\_out\_G09\_inversion &
\multirow{2}{*}{Equal} & \multirow{2}{*}{Equal}
\\
& E = -618.057081944 \tab Freq =115.2921   &    &  \\ \cline{2-2}
& s-2-chlorobutane\_rotated\_out\_G09   & \multicolumn{2}{c}
{ RMS = 5.673033E-07}
\\
& E = -618.057081943 \tab Freq =115.2939   &    \multicolumn{2}{c} 
{ }
\\ \hline
\multirow{4}{*}{\tab[2mm] 5 \tab[2mm]} & r-2-chlorobutane\_out\_G09 &
\multirow{2}{*}{Enantiomers} & \multirow{2}{*}{Enantiomers}
\\
& E = -618.057081944 \tab Freq =115.2921   &    &  \\ \cline{2-2}
& r-2-chlorobutane\_out\_G09\_inversion   & \multicolumn{2}{c}
 {RMS = 0.02719433}
\\
& E = -618.057081944 \tab Freq =115.2921   &    \multicolumn{2}{c} 
{ }
\\ \hline
\multirow{4}{*}{\tab[2mm] 6 \tab[2mm]} & r-2-chlorobutane\_out\_G09 &
\multirow{2}{*}{Equal} & \multirow{2}{*}{Equal}
\\
& E = -618.057081944 \tab Freq =115.2921   &    &  \\ \cline{2-2}
& r-2-chlorobutane\_rotated\_out\_G09   & \multicolumn{2}{c}
{ RMS = 1.236488E-09}
\\
& E = -618.057081944 \tab Freq =115.2921   &    \multicolumn{2}{c} 
{ }
\\ \hline
\multirow{4}{*}{\tab[2mm] 7 \tab[2mm]} & r-2-chlorobutane\_out\_G09 &
\multirow{2}{*}{Enantiomers} & \multirow{2}{*}{Enantiomers}
\\
& E = -618.057081944 \tab Freq =115.2921   &    &  \\ \cline{2-2}
& s-2-chlorobutane\_out\_G09   & \multicolumn{2}{c}
 {RMS = 0.2294273}
\\
& E = -618.057081943 \tab Freq =115.2939   &    \multicolumn{2}{c} 
{ }
\\ \hline
\end{tabular}
\newpage

\vtab[-2cm]
\tab[-2cm]
\begin{tabular}{c|m{8cm}|c|c}
\# & Moléculas & Restultado esperado & Resultado programa \\ \hline\hline
\multirow{4}{*}{\tab[2mm] 8 \tab[2mm]} & r-2-chlorobutane\_out\_G09 &
\multirow{2}{*}{Equal} & \multirow{2}{*}{Equal}
\\
& E = -618.057081944 \tab Freq =115.2921   &    &  \\ \cline{2-2}
& s-2-chlorobutane\_out\_G09\_inversion   & \multicolumn{2}{c}
{ RMS = 4.044324E-06}
\\
& E = -618.057081943 \tab Freq =115.2939   &    \multicolumn{2}{c} 
{ }
\\ \hline
\multirow{4}{*}{\tab[2mm] 9 \tab[2mm]} & r-2-chlorobutane\_out\_G09 &
\multirow{2}{*}{Enantiomers} & \multirow{2}{*}{Enantiomers}
\\
& E = -618.057081944 \tab Freq =115.2921   &    &  \\ \cline{2-2}
& s-2-chlorobutane\_rotated\_out\_G09   & \multicolumn{2}{c}
 {RMS = 0.2294272}
\\
& E = -618.057081943 \tab Freq =115.2939   &    \multicolumn{2}{c} 
{ }
\\ \hline
\multirow{4}{*}{\tab[2mm] 10 \tab[2mm]} & r-2-chlorobutane\_rotated\_out\_G09 &
\multirow{2}{*}{Enantiomers} & \multirow{2}{*}{Enantiomers}
\\
& E = -618.057081944 \tab Freq =115.2921   &    &  \\ \cline{2-2}
& s-2-chlorobutane\_out\_G09   & \multicolumn{2}{c}
 {RMS = 0.2294273}
\\
& E = -618.057081943 \tab Freq =115.2939   &    \multicolumn{2}{c} 
{ }
\\ \hline
\multirow{4}{*}{\tab[2mm] 11 \tab[2mm]} & r-2-chlorobutane\_rotated\_out\_G09 &
\multirow{2}{*}{Equal} & \multirow{2}{*}{Equal}
\\
& E = -618.057081944 \tab Freq =115.2921   &    &  \\ \cline{2-2}
& s-2-chlorobutane\_out\_G09\_inversion   & \multicolumn{2}{c}
{ RMS = 4.043088E-06}
\\
& E = -618.057081943 \tab Freq =115.2939   &    \multicolumn{2}{c} 
{ }
\\ \hline
\multirow{4}{*}{\tab[2mm] 12 \tab[2mm]} & r-2-chlorobutane\_rotated\_out\_G09 &
\multirow{2}{*}{Enantiomers} & \multirow{2}{*}{Enantiomers}
\\
& E = -618.057081944 \tab Freq =115.2921   &    &  \\ \cline{2-2}
& s-2-chlorobutane\_rotated\_out\_G09   & \multicolumn{2}{c}
 {RMS = 0.2294272}
\\
& E = -618.057081943 \tab Freq =115.2939   &    \multicolumn{2}{c} 
{ }
\\ \hline
\multirow{4}{*}{\tab[2mm] 13 \tab[2mm]} & s-2-chlorobutane\_out\_G09\_inversion &
\multirow{2}{*}{Enantiomers} & \multirow{2}{*}{Enantiomers}
\\
& E = -618.057081943 \tab Freq =115.2939   &    &  \\ \cline{2-2}
& s-2-chlorobutane\_rotated\_out\_G09   & \multicolumn{2}{c}
 {RMS = 0.02719894}
\\
& E = -618.057081943 \tab Freq =115.2939   &    \multicolumn{2}{c} 
{ }
\\ \hline
\multirow{4}{*}{\tab[2mm] 14 \tab[2mm]} & s-2-chlorobutane\_out\_G09 &
\multirow{2}{*}{Enantiomers} & \multirow{2}{*}{Enantiomers}
\\
& E = -618.057081943 \tab Freq =115.2939   &    &  \\ \cline{2-2}
& s-2-chlorobutane\_out\_G09\_inversion   & \multicolumn{2}{c}
 {RMS = 0.02719881}
\\
& E = -618.057081943 \tab Freq =115.2939   &    \multicolumn{2}{c} 
{ }
\\ \hline
\end{tabular}
\newpage

\vtab[-2cm]
\tab[-2cm]
\begin{tabular}{c|m{8cm}|c|c}
\# & Moléculas & Restultado esperado & Resultado programa \\ \hline\hline
\multirow{4}{*}{\tab[2mm] 15 \tab[2mm]} & s-2-chlorobutane\_out\_G09 &
\multirow{2}{*}{Equal} & \multirow{2}{*}{Equal}
\\
& E = -618.057081943 \tab Freq =115.2939   &    &  \\ \cline{2-2}
& s-2-chlorobutane\_rotated\_out\_G09   & \multicolumn{2}{c}
{ RMS = 1.306092E-07}
\\
& E = -618.057081943 \tab Freq =115.2939   &    \multicolumn{2}{c} 
{ }
\\ \hline
\multirow{4}{*}{\tab[2mm] 16 \tab[2mm]} & 3NFAACa &
\multirow{2}{*}{\textcolor{Red}{\bf Equal}} & \multirow{2}{*}{\textcolor{Red}{\bf Enantiomers}}
\\
& E = No Data \tab Freq =No Data   &    &  \\ \cline{2-2}
& 3NFAACb   & \multicolumn{2}{c}
 {RMS = 0.02603601}
\\
& E = No Data \tab Freq =No Data   &    \multicolumn{2}{c} 
{ }
\\ \hline
\multirow{4}{*}{\tab[2mm] 17 \tab[2mm]} & 3NFAACa &
\multirow{2}{*}{Different} & \multirow{2}{*}{Different}
\\
& E = No Data \tab Freq =No Data   &    &  \\ \cline{2-2}
& 3NFAACc   & \multicolumn{2}{c}
 {RMS = 0.0334016}
\\
& E = No Data \tab Freq =No Data   &    \multicolumn{2}{c} 
{ }
\\ \hline
\multirow{4}{*}{\tab[2mm] 18 \tab[2mm]} & 3NFAACa &
\multirow{2}{*}{Different} & \multirow{2}{*}{Different}
\\
& E = No Data \tab Freq =No Data   &    &  \\ \cline{2-2}
& 3NFAACd   & \multicolumn{2}{c}
 {RMS = 0.0888511}
\\
& E = No Data \tab Freq =No Data   &    \multicolumn{2}{c} 
{ }
\\ \hline
\multirow{4}{*}{\tab[2mm] 19 \tab[2mm]} & 3NFAACa &
\multirow{2}{*}{Different} & \multirow{2}{*}{Different}
\\
& E = No Data \tab Freq =No Data   &    &  \\ \cline{2-2}
& 3NFAACe   & \multicolumn{2}{c}
 {RMS = 0.162132}
\\
& E = No Data \tab Freq =No Data   &    \multicolumn{2}{c} 
{ }
\\ \hline
\multirow{4}{*}{\tab[2mm] 20 \tab[2mm]} & 3NFAACa &
\multirow{2}{*}{\textcolor{Red}{\bf Equal}} & \multirow{2}{*}{\textcolor{Red}{\bf Enantiomers}}
\\
& E = No Data \tab Freq =No Data   &    &  \\ \cline{2-2}
& 3NFAACf   & \multicolumn{2}{c}
 {RMS = 0.02649945}
\\
& E = No Data \tab Freq =No Data   &    \multicolumn{2}{c} 
{ }
\\ \hline
\multirow{4}{*}{\tab[2mm] 21 \tab[2mm]} & 3NFAACa &
\multirow{2}{*}{Different} & \multirow{2}{*}{Different}
\\
& E = No Data \tab Freq =No Data   &    &  \\ \cline{2-2}
& 3NFAACg   & \multicolumn{2}{c}
 {RMS = 0.0330848}
\\
& E = No Data \tab Freq =No Data   &    \multicolumn{2}{c} 
{ }
\\ \hline
\end{tabular}
\newpage

\vtab[-2cm]
\tab[-2cm]
\begin{tabular}{c|m{8cm}|c|c}
\# & Moléculas & Restultado esperado & Resultado programa \\ \hline\hline
\multirow{4}{*}{\tab[2mm] 22 \tab[2mm]} & 3NFAACa &
\multirow{2}{*}{Different} & \multirow{2}{*}{Different}
\\
& E = No Data \tab Freq =No Data   &    &  \\ \cline{2-2}
& 3NFAACh   & \multicolumn{2}{c}
 {RMS = 0.064417}
\\
& E = No Data \tab Freq =No Data   &    \multicolumn{2}{c} 
{ }
\\ \hline
\multirow{4}{*}{\tab[2mm] 23 \tab[2mm]} & 3NFAACa &
\multirow{2}{*}{Different} & \multirow{2}{*}{Different}
\\
& E = No Data \tab Freq =No Data   &    &  \\ \cline{2-2}
& 3NFAACi   & \multicolumn{2}{c}
 {RMS = 0.0837021}
\\
& E = No Data \tab Freq =No Data   &    \multicolumn{2}{c} 
{ }
\\ \hline
\multirow{4}{*}{\tab[2mm] 24 \tab[2mm]} & 3NFAACa &
\multirow{2}{*}{Different} & \multirow{2}{*}{Different}
\\
& E = No Data \tab Freq =No Data   &    &  \\ \cline{2-2}
& 3NFAACj   & \multicolumn{2}{c}
 {RMS = 0.0903201}
\\
& E = No Data \tab Freq =No Data   &    \multicolumn{2}{c} 
{ }
\\ \hline
\multirow{4}{*}{\tab[2mm] 25 \tab[2mm]} & 3NFAACa &
\multirow{2}{*}{Different} & \multirow{2}{*}{Different}
\\
& E = No Data \tab Freq =No Data   &    &  \\ \cline{2-2}
& 3NFAACk   & \multicolumn{2}{c}
 {RMS = 0.0114116}
\\
& E = No Data \tab Freq =No Data   &    \multicolumn{2}{c} 
{ }
\\ \hline
\multirow{4}{*}{\tab[2mm] 26 \tab[2mm]} & 3NFAACa &
\multirow{2}{*}{Different} & \multirow{2}{*}{Different}
\\
& E = No Data \tab Freq =No Data   &    &  \\ \cline{2-2}
& 3NFAACl   & \multicolumn{2}{c}
 {RMS = 0.0362294}
\\
& E = No Data \tab Freq =No Data   &    \multicolumn{2}{c} 
{ }
\\ \hline
\multirow{4}{*}{\tab[2mm] 27 \tab[2mm]} & 3NFAACa &
\multirow{2}{*}{Different} & \multirow{2}{*}{Different}
\\
& E = No Data \tab Freq =No Data   &    &  \\ \cline{2-2}
& 3NFAACm   & \multicolumn{2}{c}
 {RMS = 0.152833}
\\
& E = No Data \tab Freq =No Data   &    \multicolumn{2}{c} 
{ }
\\ \hline
\multirow{4}{*}{\tab[2mm] 28 \tab[2mm]} & 3NFAACa &
\multirow{2}{*}{Different} & \multirow{2}{*}{Different}
\\
& E = No Data \tab Freq =No Data   &    &  \\ \cline{2-2}
& 3NFAACn   & \multicolumn{2}{c}
 {RMS = 0.157182}
\\
& E = No Data \tab Freq =No Data   &    \multicolumn{2}{c} 
{ }
\\ \hline
\end{tabular}
\newpage

\vtab[-2cm]
\tab[-2cm]
\begin{tabular}{c|m{8cm}|c|c}
\# & Moléculas & Restultado esperado & Resultado programa \\ \hline\hline
\multirow{4}{*}{\tab[2mm] 29 \tab[2mm]} & 3NFAACa &
\multirow{2}{*}{Different} & \multirow{2}{*}{Different}
\\
& E = No Data \tab Freq =No Data   &    &  \\ \cline{2-2}
& 4NFAACa   & \multicolumn{2}{c}
 {RMS = 0.219778}
\\
& E = No Data \tab Freq =No Data   &    \multicolumn{2}{c} 
{ }
\\ \hline
\multirow{4}{*}{\tab[2mm] 30 \tab[2mm]} & 3NFAACa &
\multirow{2}{*}{Different} & \multirow{2}{*}{Different}
\\
& E = No Data \tab Freq =No Data   &    &  \\ \cline{2-2}
& 4NFAACb   & \multicolumn{2}{c}
 {RMS = 0.258442}
\\
& E = No Data \tab Freq =No Data   &    \multicolumn{2}{c} 
{ }
\\ \hline
\multirow{4}{*}{\tab[2mm] 31 \tab[2mm]} & 3NFAACa &
\multirow{2}{*}{Different} & \multirow{2}{*}{Different}
\\
& E = No Data \tab Freq =No Data   &    &  \\ \cline{2-2}
& 4NFAACc   & \multicolumn{2}{c}
 {RMS = 0.0186255}
\\
& E = No Data \tab Freq =No Data   &    \multicolumn{2}{c} 
{ }
\\ \hline
\multirow{4}{*}{\tab[2mm] 32 \tab[2mm]} & 3NFAACa &
\multirow{2}{*}{Different} & \multirow{2}{*}{Different}
\\
& E = No Data \tab Freq =No Data   &    &  \\ \cline{2-2}
& 4NFAACd   & \multicolumn{2}{c}
 {RMS = 0.0149078}
\\
& E = No Data \tab Freq =No Data   &    \multicolumn{2}{c} 
{ }
\\ \hline
\multirow{4}{*}{\tab[2mm] 33 \tab[2mm]} & 3NFAACa &
\multirow{2}{*}{Different} & \multirow{2}{*}{Different}
\\
& E = No Data \tab Freq =No Data   &    &  \\ \cline{2-2}
& 4NFAACe   & \multicolumn{2}{c}
 {RMS = 0.0228438}
\\
& E = No Data \tab Freq =No Data   &    \multicolumn{2}{c} 
{ }
\\ \hline
\multirow{4}{*}{\tab[2mm] 34 \tab[2mm]} & 3NFAACa &
\multirow{2}{*}{Different} & \multirow{2}{*}{Different}
\\
& E = No Data \tab Freq =No Data   &    &  \\ \cline{2-2}
& 4NFAACf   & \multicolumn{2}{c}
 {RMS = 0.118302}
\\
& E = No Data \tab Freq =No Data   &    \multicolumn{2}{c} 
{ }
\\ \hline
\multirow{4}{*}{\tab[2mm] 35 \tab[2mm]} & 3NFAACa &
\multirow{2}{*}{Different} & \multirow{2}{*}{Different}
\\
& E = No Data \tab Freq =No Data   &    &  \\ \cline{2-2}
& 4NFAACg   & \multicolumn{2}{c}
 {RMS = 0.0338585}
\\
& E = No Data \tab Freq =No Data   &    \multicolumn{2}{c} 
{ }
\\ \hline
\end{tabular}
\newpage

\vtab[-2cm]
\tab[-2cm]
\begin{tabular}{c|m{8cm}|c|c}
\# & Moléculas & Restultado esperado & Resultado programa \\ \hline\hline
\multirow{4}{*}{\tab[2mm] 36 \tab[2mm]} & 3NFAACa &
\multirow{2}{*}{Different} & \multirow{2}{*}{Different}
\\
& E = No Data \tab Freq =No Data   &    &  \\ \cline{2-2}
& 4NFAACi   & \multicolumn{2}{c}
 {RMS = 0.00146242}
\\
& E = No Data \tab Freq =No Data   &    \multicolumn{2}{c} 
{ }
\\ \hline
\multirow{4}{*}{\tab[2mm] 37 \tab[2mm]} & 3NFAACa &
\multirow{2}{*}{Different} & \multirow{2}{*}{Different}
\\
& E = No Data \tab Freq =No Data   &    &  \\ \cline{2-2}
& 4NFAACj   & \multicolumn{2}{c}
 {RMS = 0.0868356}
\\
& E = No Data \tab Freq =No Data   &    \multicolumn{2}{c} 
{ }
\\ \hline
\multirow{4}{*}{\tab[2mm] 38 \tab[2mm]} & 3NFAACa &
\multirow{2}{*}{Different} & \multirow{2}{*}{Different}
\\
& E = No Data \tab Freq =No Data   &    &  \\ \cline{2-2}
& 4NFAACl-3   & \multicolumn{2}{c}
 {RMS = 0.0414826}
\\
& E = No Data \tab Freq =No Data   &    \multicolumn{2}{c} 
{ }
\\ \hline
\multirow{4}{*}{\tab[2mm] 39 \tab[2mm]} & 3NFAACb &
\multirow{2}{*}{Different} & \multirow{2}{*}{Different}
\\
& E = No Data \tab Freq =No Data   &    &  \\ \cline{2-2}
& 3NFAACc   & \multicolumn{2}{c}
 {RMS = 0.0881291}
\\
& E = No Data \tab Freq =No Data   &    \multicolumn{2}{c} 
{ }
\\ \hline
\multirow{4}{*}{\tab[2mm] 40 \tab[2mm]} & 3NFAACb &
\multirow{2}{*}{Different} & \multirow{2}{*}{Different}
\\
& E = No Data \tab Freq =No Data   &    &  \\ \cline{2-2}
& 3NFAACd   & \multicolumn{2}{c}
 {RMS = 0.143579}
\\
& E = No Data \tab Freq =No Data   &    \multicolumn{2}{c} 
{ }
\\ \hline
\multirow{4}{*}{\tab[2mm] 41 \tab[2mm]} & 3NFAACb &
\multirow{2}{*}{Different} & \multirow{2}{*}{Different}
\\
& E = No Data \tab Freq =No Data   &    &  \\ \cline{2-2}
& 3NFAACe   & \multicolumn{2}{c}
 {RMS = 0.107404}
\\
& E = No Data \tab Freq =No Data   &    \multicolumn{2}{c} 
{ }
\\ \hline
\multirow{4}{*}{\tab[2mm] 42 \tab[2mm]} & 3NFAACb &
\multirow{2}{*}{\textcolor{Red}{\bf Equal}} & \multirow{2}{*}{\textcolor{Red}{\bf Enantiomers}}
\\
& E = No Data \tab Freq =No Data   &    &  \\ \cline{2-2}
& 3NFAACf   & \multicolumn{2}{c}
{\textcolor{Red}{ RMS = 0.0004634431}}
\\
& E = No Data \tab Freq =No Data   &    \multicolumn{2}{c} 
{ }
\\ \hline
\end{tabular}
\newpage

\vtab[-2cm]
\tab[-2cm]
\begin{tabular}{c|m{8cm}|c|c}
\# & Moléculas & Restultado esperado & Resultado programa \\ \hline\hline
\multirow{4}{*}{\tab[2mm] 43 \tab[2mm]} & 3NFAACb &
\multirow{2}{*}{Different} & \multirow{2}{*}{Different}
\\
& E = No Data \tab Freq =No Data   &    &  \\ \cline{2-2}
& 3NFAACg   & \multicolumn{2}{c}
 {RMS = 0.0878122}
\\
& E = No Data \tab Freq =No Data   &    \multicolumn{2}{c} 
{ }
\\ \hline
\multirow{4}{*}{\tab[2mm] 44 \tab[2mm]} & 3NFAACb &
\multirow{2}{*}{Different} & \multirow{2}{*}{Different}
\\
& E = No Data \tab Freq =No Data   &    &  \\ \cline{2-2}
& 3NFAACh   & \multicolumn{2}{c}
 {RMS = 0.00968954}
\\
& E = No Data \tab Freq =No Data   &    \multicolumn{2}{c} 
{ }
\\ \hline
\multirow{4}{*}{\tab[2mm] 45 \tab[2mm]} & 3NFAACb &
\multirow{2}{*}{Different} & \multirow{2}{*}{Different}
\\
& E = No Data \tab Freq =No Data   &    &  \\ \cline{2-2}
& 3NFAACi   & \multicolumn{2}{c}
 {RMS = 0.13843}
\\
& E = No Data \tab Freq =No Data   &    \multicolumn{2}{c} 
{ }
\\ \hline
\multirow{4}{*}{\tab[2mm] 46 \tab[2mm]} & 3NFAACb &
\multirow{2}{*}{Different} & \multirow{2}{*}{Different}
\\
& E = No Data \tab Freq =No Data   &    &  \\ \cline{2-2}
& 3NFAACj   & \multicolumn{2}{c}
 {RMS = 0.145048}
\\
& E = No Data \tab Freq =No Data   &    \multicolumn{2}{c} 
{ }
\\ \hline
\multirow{4}{*}{\tab[2mm] 47 \tab[2mm]} & 3NFAACb &
\multirow{2}{*}{Different} & \multirow{2}{*}{Different}
\\
& E = No Data \tab Freq =No Data   &    &  \\ \cline{2-2}
& 3NFAACk   & \multicolumn{2}{c}
 {RMS = 0.0661391}
\\
& E = No Data \tab Freq =No Data   &    \multicolumn{2}{c} 
{ }
\\ \hline
\multirow{4}{*}{\tab[2mm] 48 \tab[2mm]} & 3NFAACb &
\multirow{2}{*}{Different} & \multirow{2}{*}{Different}
\\
& E = No Data \tab Freq =No Data   &    &  \\ \cline{2-2}
& 3NFAACl   & \multicolumn{2}{c}
 {RMS = 0.018498}
\\
& E = No Data \tab Freq =No Data   &    \multicolumn{2}{c} 
{ }
\\ \hline
\multirow{4}{*}{\tab[2mm] 49 \tab[2mm]} & 3NFAACb &
\multirow{2}{*}{Different} & \multirow{2}{*}{Different}
\\
& E = No Data \tab Freq =No Data   &    &  \\ \cline{2-2}
& 3NFAACm   & \multicolumn{2}{c}
 {RMS = 0.0981054}
\\
& E = No Data \tab Freq =No Data   &    \multicolumn{2}{c} 
{ }
\\ \hline
\end{tabular}
\newpage

\vtab[-2cm]
\tab[-2cm]
\begin{tabular}{c|m{8cm}|c|c}
\# & Moléculas & Restultado esperado & Resultado programa \\ \hline\hline
\multirow{4}{*}{\tab[2mm] 50 \tab[2mm]} & 3NFAACb &
\multirow{2}{*}{Different} & \multirow{2}{*}{Different}
\\
& E = No Data \tab Freq =No Data   &    &  \\ \cline{2-2}
& 3NFAACn   & \multicolumn{2}{c}
 {RMS = 0.102455}
\\
& E = No Data \tab Freq =No Data   &    \multicolumn{2}{c} 
{ }
\\ \hline
\multirow{4}{*}{\tab[2mm] 51 \tab[2mm]} & 3NFAACb &
\multirow{2}{*}{Different} & \multirow{2}{*}{Different}
\\
& E = No Data \tab Freq =No Data   &    &  \\ \cline{2-2}
& 4NFAACa   & \multicolumn{2}{c}
 {RMS = 0.274505}
\\
& E = No Data \tab Freq =No Data   &    \multicolumn{2}{c} 
{ }
\\ \hline
\multirow{4}{*}{\tab[2mm] 52 \tab[2mm]} & 3NFAACb &
\multirow{2}{*}{Different} & \multirow{2}{*}{Different}
\\
& E = No Data \tab Freq =No Data   &    &  \\ \cline{2-2}
& 4NFAACb   & \multicolumn{2}{c}
 {RMS = 0.203714}
\\
& E = No Data \tab Freq =No Data   &    \multicolumn{2}{c} 
{ }
\\ \hline
\multirow{4}{*}{\tab[2mm] 53 \tab[2mm]} & 3NFAACb &
\multirow{2}{*}{Different} & \multirow{2}{*}{Different}
\\
& E = No Data \tab Freq =No Data   &    &  \\ \cline{2-2}
& 4NFAACc   & \multicolumn{2}{c}
 {RMS = 0.036102}
\\
& E = No Data \tab Freq =No Data   &    \multicolumn{2}{c} 
{ }
\\ \hline
\multirow{4}{*}{\tab[2mm] 54 \tab[2mm]} & 3NFAACb &
\multirow{2}{*}{Different} & \multirow{2}{*}{Different}
\\
& E = No Data \tab Freq =No Data   &    &  \\ \cline{2-2}
& 4NFAACd   & \multicolumn{2}{c}
 {RMS = 0.0398197}
\\
& E = No Data \tab Freq =No Data   &    \multicolumn{2}{c} 
{ }
\\ \hline
\multirow{4}{*}{\tab[2mm] 55 \tab[2mm]} & 3NFAACb &
\multirow{2}{*}{Different} & \multirow{2}{*}{Different}
\\
& E = No Data \tab Freq =No Data   &    &  \\ \cline{2-2}
& 4NFAACe   & \multicolumn{2}{c}
 {RMS = 0.0318836}
\\
& E = No Data \tab Freq =No Data   &    \multicolumn{2}{c} 
{ }
\\ \hline
\multirow{4}{*}{\tab[2mm] 56 \tab[2mm]} & 3NFAACb &
\multirow{2}{*}{Different} & \multirow{2}{*}{Different}
\\
& E = No Data \tab Freq =No Data   &    &  \\ \cline{2-2}
& 4NFAACf   & \multicolumn{2}{c}
 {RMS = 0.0635749}
\\
& E = No Data \tab Freq =No Data   &    \multicolumn{2}{c} 
{ }
\\ \hline
\end{tabular}
\newpage

\vtab[-2cm]
\tab[-2cm]
\begin{tabular}{c|m{8cm}|c|c}
\# & Moléculas & Restultado esperado & Resultado programa \\ \hline\hline
\multirow{4}{*}{\tab[2mm] 57 \tab[2mm]} & 3NFAACb &
\multirow{2}{*}{Different} & \multirow{2}{*}{Different}
\\
& E = No Data \tab Freq =No Data   &    &  \\ \cline{2-2}
& 4NFAACg   & \multicolumn{2}{c}
 {RMS = 0.020869}
\\
& E = No Data \tab Freq =No Data   &    \multicolumn{2}{c} 
{ }
\\ \hline
\multirow{4}{*}{\tab[2mm] 58 \tab[2mm]} & 3NFAACb &
\multirow{2}{*}{Different} & \multirow{2}{*}{Different}
\\
& E = No Data \tab Freq =No Data   &    &  \\ \cline{2-2}
& 4NFAACi   & \multicolumn{2}{c}
 {RMS = 0.0561899}
\\
& E = No Data \tab Freq =No Data   &    \multicolumn{2}{c} 
{ }
\\ \hline
\multirow{4}{*}{\tab[2mm] 59 \tab[2mm]} & 3NFAACb &
\multirow{2}{*}{Different} & \multirow{2}{*}{Different}
\\
& E = No Data \tab Freq =No Data   &    &  \\ \cline{2-2}
& 4NFAACj   & \multicolumn{2}{c}
 {RMS = 0.0321081}
\\
& E = No Data \tab Freq =No Data   &    \multicolumn{2}{c} 
{ }
\\ \hline
\multirow{4}{*}{\tab[2mm] 60 \tab[2mm]} & 3NFAACb &
\multirow{2}{*}{Different} & \multirow{2}{*}{Different}
\\
& E = No Data \tab Freq =No Data   &    &  \\ \cline{2-2}
& 4NFAACl-3   & \multicolumn{2}{c}
 {RMS = 0.0132449}
\\
& E = No Data \tab Freq =No Data   &    \multicolumn{2}{c} 
{ }
\\ \hline
\multirow{4}{*}{\tab[2mm] 61 \tab[2mm]} & 3NFAACc &
\multirow{2}{*}{Different} & \multirow{2}{*}{Different}
\\
& E = No Data \tab Freq =No Data   &    &  \\ \cline{2-2}
& 3NFAACd   & \multicolumn{2}{c}
 {RMS = 0.0554495}
\\
& E = No Data \tab Freq =No Data   &    \multicolumn{2}{c} 
{ }
\\ \hline
\multirow{4}{*}{\tab[2mm] 62 \tab[2mm]} & 3NFAACc &
\multirow{2}{*}{Different} & \multirow{2}{*}{Different}
\\
& E = No Data \tab Freq =No Data   &    &  \\ \cline{2-2}
& 3NFAACe   & \multicolumn{2}{c}
 {RMS = 0.195533}
\\
& E = No Data \tab Freq =No Data   &    \multicolumn{2}{c} 
{ }
\\ \hline
\multirow{4}{*}{\tab[2mm] 63 \tab[2mm]} & 3NFAACc &
\multirow{2}{*}{Different} & \multirow{2}{*}{Different}
\\
& E = No Data \tab Freq =No Data   &    &  \\ \cline{2-2}
& 3NFAACf   & \multicolumn{2}{c}
 {RMS = 0.0879682}
\\
& E = No Data \tab Freq =No Data   &    \multicolumn{2}{c} 
{ }
\\ \hline
\end{tabular}
\newpage

\vtab[-2cm]
\tab[-2cm]
\begin{tabular}{c|m{8cm}|c|c}
\# & Moléculas & Restultado esperado & Resultado programa \\ \hline\hline
\multirow{4}{*}{\tab[2mm] 64 \tab[2mm]} & 3NFAACc &
\multirow{2}{*}{Different} & \multirow{2}{*}{Different}
\\
& E = No Data \tab Freq =No Data   &    &  \\ \cline{2-2}
& 3NFAACg   & \multicolumn{2}{c}
 {RMS = 0.000316817}
\\
& E = No Data \tab Freq =No Data   &    \multicolumn{2}{c} 
{ }
\\ \hline
\multirow{4}{*}{\tab[2mm] 65 \tab[2mm]} & 3NFAACc &
\multirow{2}{*}{Different} & \multirow{2}{*}{Different}
\\
& E = No Data \tab Freq =No Data   &    &  \\ \cline{2-2}
& 3NFAACh   & \multicolumn{2}{c}
 {RMS = 0.0978186}
\\
& E = No Data \tab Freq =No Data   &    \multicolumn{2}{c} 
{ }
\\ \hline
\multirow{4}{*}{\tab[2mm] 66 \tab[2mm]} & 3NFAACc &
\multirow{2}{*}{Different} & \multirow{2}{*}{Different}
\\
& E = No Data \tab Freq =No Data   &    &  \\ \cline{2-2}
& 3NFAACi   & \multicolumn{2}{c}
 {RMS = 0.0503005}
\\
& E = No Data \tab Freq =No Data   &    \multicolumn{2}{c} 
{ }
\\ \hline
\multirow{4}{*}{\tab[2mm] 67 \tab[2mm]} & 3NFAACc &
\multirow{2}{*}{Different} & \multirow{2}{*}{Different}
\\
& E = No Data \tab Freq =No Data   &    &  \\ \cline{2-2}
& 3NFAACj   & \multicolumn{2}{c}
 {RMS = 0.0569185}
\\
& E = No Data \tab Freq =No Data   &    \multicolumn{2}{c} 
{ }
\\ \hline
\multirow{4}{*}{\tab[2mm] 68 \tab[2mm]} & 3NFAACc &
\multirow{2}{*}{Different} & \multirow{2}{*}{Different}
\\
& E = No Data \tab Freq =No Data   &    &  \\ \cline{2-2}
& 3NFAACk   & \multicolumn{2}{c}
 {RMS = 0.02199}
\\
& E = No Data \tab Freq =No Data   &    \multicolumn{2}{c} 
{ }
\\ \hline
\multirow{4}{*}{\tab[2mm] 69 \tab[2mm]} & 3NFAACc &
\multirow{2}{*}{Different} & \multirow{2}{*}{Different}
\\
& E = No Data \tab Freq =No Data   &    &  \\ \cline{2-2}
& 3NFAACl   & \multicolumn{2}{c}
 {RMS = 0.069631}
\\
& E = No Data \tab Freq =No Data   &    \multicolumn{2}{c} 
{ }
\\ \hline
\multirow{4}{*}{\tab[2mm] 70 \tab[2mm]} & 3NFAACc &
\multirow{2}{*}{Different} & \multirow{2}{*}{Different}
\\
& E = No Data \tab Freq =No Data   &    &  \\ \cline{2-2}
& 3NFAACm   & \multicolumn{2}{c}
 {RMS = 0.186234}
\\
& E = No Data \tab Freq =No Data   &    \multicolumn{2}{c} 
{ }
\\ \hline
\end{tabular}
\newpage

\vtab[-2cm]
\tab[-2cm]
\begin{tabular}{c|m{8cm}|c|c}
\# & Moléculas & Restultado esperado & Resultado programa \\ \hline\hline
\multirow{4}{*}{\tab[2mm] 71 \tab[2mm]} & 3NFAACc &
\multirow{2}{*}{Enantiomers} & \multirow{2}{*}{Enantiomers}
\\
& E = No Data \tab Freq =No Data   &    &  \\ \cline{2-2}
& 3NFAACn   & \multicolumn{2}{c}
 {RMS = 0.01522269}
\\
& E = No Data \tab Freq =No Data   &    \multicolumn{2}{c} 
{ }
\\ \hline
\multirow{4}{*}{\tab[2mm] 72 \tab[2mm]} & 3NFAACc &
\multirow{2}{*}{Different} & \multirow{2}{*}{Different}
\\
& E = No Data \tab Freq =No Data   &    &  \\ \cline{2-2}
& 4NFAACa   & \multicolumn{2}{c}
 {RMS = 0.186376}
\\
& E = No Data \tab Freq =No Data   &    \multicolumn{2}{c} 
{ }
\\ \hline
\multirow{4}{*}{\tab[2mm] 73 \tab[2mm]} & 3NFAACc &
\multirow{2}{*}{Different} & \multirow{2}{*}{Different}
\\
& E = No Data \tab Freq =No Data   &    &  \\ \cline{2-2}
& 4NFAACb   & \multicolumn{2}{c}
 {RMS = 0.291844}
\\
& E = No Data \tab Freq =No Data   &    \multicolumn{2}{c} 
{ }
\\ \hline
\multirow{4}{*}{\tab[2mm] 74 \tab[2mm]} & 3NFAACc &
\multirow{2}{*}{Different} & \multirow{2}{*}{Different}
\\
& E = No Data \tab Freq =No Data   &    &  \\ \cline{2-2}
& 4NFAACc   & \multicolumn{2}{c}
 {RMS = 0.0520271}
\\
& E = No Data \tab Freq =No Data   &    \multicolumn{2}{c} 
{ }
\\ \hline
\multirow{4}{*}{\tab[2mm] 75 \tab[2mm]} & 3NFAACc &
\multirow{2}{*}{Different} & \multirow{2}{*}{Different}
\\
& E = No Data \tab Freq =No Data   &    &  \\ \cline{2-2}
& 4NFAACd   & \multicolumn{2}{c}
 {RMS = 0.0483094}
\\
& E = No Data \tab Freq =No Data   &    \multicolumn{2}{c} 
{ }
\\ \hline
\multirow{4}{*}{\tab[2mm] 76 \tab[2mm]} & 3NFAACc &
\multirow{2}{*}{Different} & \multirow{2}{*}{Different}
\\
& E = No Data \tab Freq =No Data   &    &  \\ \cline{2-2}
& 4NFAACe   & \multicolumn{2}{c}
 {RMS = 0.0562454}
\\
& E = No Data \tab Freq =No Data   &    \multicolumn{2}{c} 
{ }
\\ \hline
\multirow{4}{*}{\tab[2mm] 77 \tab[2mm]} & 3NFAACc &
\multirow{2}{*}{Different} & \multirow{2}{*}{Different}
\\
& E = No Data \tab Freq =No Data   &    &  \\ \cline{2-2}
& 4NFAACf   & \multicolumn{2}{c}
 {RMS = 0.151704}
\\
& E = No Data \tab Freq =No Data   &    \multicolumn{2}{c} 
{ }
\\ \hline
\end{tabular}
\newpage

\vtab[-2cm]
\tab[-2cm]
\begin{tabular}{c|m{8cm}|c|c}
\# & Moléculas & Restultado esperado & Resultado programa \\ \hline\hline
\multirow{4}{*}{\tab[2mm] 78 \tab[2mm]} & 3NFAACc &
\multirow{2}{*}{Different} & \multirow{2}{*}{Different}
\\
& E = No Data \tab Freq =No Data   &    &  \\ \cline{2-2}
& 4NFAACg   & \multicolumn{2}{c}
 {RMS = 0.0672601}
\\
& E = No Data \tab Freq =No Data   &    \multicolumn{2}{c} 
{ }
\\ \hline
\multirow{4}{*}{\tab[2mm] 79 \tab[2mm]} & 3NFAACc &
\multirow{2}{*}{Different} & \multirow{2}{*}{Different}
\\
& E = No Data \tab Freq =No Data   &    &  \\ \cline{2-2}
& 4NFAACi   & \multicolumn{2}{c}
 {RMS = 0.0319392}
\\
& E = No Data \tab Freq =No Data   &    \multicolumn{2}{c} 
{ }
\\ \hline
\multirow{4}{*}{\tab[2mm] 80 \tab[2mm]} & 3NFAACc &
\multirow{2}{*}{Different} & \multirow{2}{*}{Different}
\\
& E = No Data \tab Freq =No Data   &    &  \\ \cline{2-2}
& 4NFAACj   & \multicolumn{2}{c}
 {RMS = 0.120237}
\\
& E = No Data \tab Freq =No Data   &    \multicolumn{2}{c} 
{ }
\\ \hline
\multirow{4}{*}{\tab[2mm] 81 \tab[2mm]} & 3NFAACc &
\multirow{2}{*}{Different} & \multirow{2}{*}{Different}
\\
& E = No Data \tab Freq =No Data   &    &  \\ \cline{2-2}
& 4NFAACl-3   & \multicolumn{2}{c}
 {RMS = 0.0748842}
\\
& E = No Data \tab Freq =No Data   &    \multicolumn{2}{c} 
{ }
\\ \hline
\multirow{4}{*}{\tab[2mm] 82 \tab[2mm]} & 3NFAACd &
\multirow{2}{*}{\textcolor{Red}{\bf Equal}} & \multirow{2}{*}{\textcolor{Red}{\bf Enantiomers}}
\\
& E = No Data \tab Freq =No Data   &    &  \\ \cline{2-2}
& 3NFAACe   & \multicolumn{2}{c}
{\textcolor{Red}{ RMS = 0.0004675477}}
\\
& E = No Data \tab Freq =No Data   &    \multicolumn{2}{c} 
{ }
\\ \hline
\multirow{4}{*}{\tab[2mm] 83 \tab[2mm]} & 3NFAACd &
\multirow{2}{*}{Different} & \multirow{2}{*}{Different}
\\
& E = No Data \tab Freq =No Data   &    &  \\ \cline{2-2}
& 3NFAACf   & \multicolumn{2}{c}
 {RMS = 0.143418}
\\
& E = No Data \tab Freq =No Data   &    \multicolumn{2}{c} 
{ }
\\ \hline
\multirow{4}{*}{\tab[2mm] 84 \tab[2mm]} & 3NFAACd &
\multirow{2}{*}{Different} & \multirow{2}{*}{Different}
\\
& E = No Data \tab Freq =No Data   &    &  \\ \cline{2-2}
& 3NFAACg   & \multicolumn{2}{c}
 {RMS = 0.0557663}
\\
& E = No Data \tab Freq =No Data   &    \multicolumn{2}{c} 
{ }
\\ \hline
\end{tabular}
\newpage

\vtab[-2cm]
\tab[-2cm]
\begin{tabular}{c|m{8cm}|c|c}
\# & Moléculas & Restultado esperado & Resultado programa \\ \hline\hline
\multirow{4}{*}{\tab[2mm] 85 \tab[2mm]} & 3NFAACd &
\multirow{2}{*}{Different} & \multirow{2}{*}{Different}
\\
& E = No Data \tab Freq =No Data   &    &  \\ \cline{2-2}
& 3NFAACh   & \multicolumn{2}{c}
 {RMS = 0.153268}
\\
& E = No Data \tab Freq =No Data   &    \multicolumn{2}{c} 
{ }
\\ \hline
\multirow{4}{*}{\tab[2mm] 86 \tab[2mm]} & 3NFAACd &
\multirow{2}{*}{Different} & \multirow{2}{*}{Different}
\\
& E = No Data \tab Freq =No Data   &    &  \\ \cline{2-2}
& 3NFAACi   & \multicolumn{2}{c}
 {RMS = 0.00514902}
\\
& E = No Data \tab Freq =No Data   &    \multicolumn{2}{c} 
{ }
\\ \hline
\multirow{4}{*}{\tab[2mm] 87 \tab[2mm]} & 3NFAACd &
\multirow{2}{*}{Different} & \multirow{2}{*}{Different}
\\
& E = No Data \tab Freq =No Data   &    &  \\ \cline{2-2}
& 3NFAACj   & \multicolumn{2}{c}
 {RMS = 0.00146894}
\\
& E = No Data \tab Freq =No Data   &    \multicolumn{2}{c} 
{ }
\\ \hline
\multirow{4}{*}{\tab[2mm] 88 \tab[2mm]} & 3NFAACd &
\multirow{2}{*}{Different} & \multirow{2}{*}{Different}
\\
& E = No Data \tab Freq =No Data   &    &  \\ \cline{2-2}
& 3NFAACk   & \multicolumn{2}{c}
 {RMS = 0.0774395}
\\
& E = No Data \tab Freq =No Data   &    \multicolumn{2}{c} 
{ }
\\ \hline
\multirow{4}{*}{\tab[2mm] 89 \tab[2mm]} & 3NFAACd &
\multirow{2}{*}{Different} & \multirow{2}{*}{Different}
\\
& E = No Data \tab Freq =No Data   &    &  \\ \cline{2-2}
& 3NFAACl   & \multicolumn{2}{c}
 {RMS = 0.125081}
\\
& E = No Data \tab Freq =No Data   &    \multicolumn{2}{c} 
{ }
\\ \hline
\multirow{4}{*}{\tab[2mm] 90 \tab[2mm]} & 3NFAACd &
\multirow{2}{*}{Different} & \multirow{2}{*}{Different}
\\
& E = No Data \tab Freq =No Data   &    &  \\ \cline{2-2}
& 3NFAACm   & \multicolumn{2}{c}
 {RMS = 0.241684}
\\
& E = No Data \tab Freq =No Data   &    \multicolumn{2}{c} 
{ }
\\ \hline
\multirow{4}{*}{\tab[2mm] 91 \tab[2mm]} & 3NFAACd &
\multirow{2}{*}{Different} & \multirow{2}{*}{Different}
\\
& E = No Data \tab Freq =No Data   &    &  \\ \cline{2-2}
& 3NFAACn   & \multicolumn{2}{c}
 {RMS = 0.246033}
\\
& E = No Data \tab Freq =No Data   &    \multicolumn{2}{c} 
{ }
\\ \hline
\end{tabular}
\newpage

\vtab[-2cm]
\tab[-2cm]
\begin{tabular}{c|m{8cm}|c|c}
\# & Moléculas & Restultado esperado & Resultado programa \\ \hline\hline
\multirow{4}{*}{\tab[2mm] 92 \tab[2mm]} & 3NFAACd &
\multirow{2}{*}{Different} & \multirow{2}{*}{Different}
\\
& E = No Data \tab Freq =No Data   &    &  \\ \cline{2-2}
& 4NFAACa   & \multicolumn{2}{c}
 {RMS = 0.130927}
\\
& E = No Data \tab Freq =No Data   &    \multicolumn{2}{c} 
{ }
\\ \hline
\multirow{4}{*}{\tab[2mm] 93 \tab[2mm]} & 3NFAACd &
\multirow{2}{*}{Different} & \multirow{2}{*}{Different}
\\
& E = No Data \tab Freq =No Data   &    &  \\ \cline{2-2}
& 4NFAACb   & \multicolumn{2}{c}
 {RMS = 0.347293}
\\
& E = No Data \tab Freq =No Data   &    \multicolumn{2}{c} 
{ }
\\ \hline
\multirow{4}{*}{\tab[2mm] 94 \tab[2mm]} & 3NFAACd &
\multirow{2}{*}{Different} & \multirow{2}{*}{Different}
\\
& E = No Data \tab Freq =No Data   &    &  \\ \cline{2-2}
& 4NFAACc   & \multicolumn{2}{c}
 {RMS = 0.107477}
\\
& E = No Data \tab Freq =No Data   &    \multicolumn{2}{c} 
{ }
\\ \hline
\multirow{4}{*}{\tab[2mm] 95 \tab[2mm]} & 3NFAACd &
\multirow{2}{*}{Different} & \multirow{2}{*}{Different}
\\
& E = No Data \tab Freq =No Data   &    &  \\ \cline{2-2}
& 4NFAACd   & \multicolumn{2}{c}
 {RMS = 0.103759}
\\
& E = No Data \tab Freq =No Data   &    \multicolumn{2}{c} 
{ }
\\ \hline
\multirow{4}{*}{\tab[2mm] 96 \tab[2mm]} & 3NFAACd &
\multirow{2}{*}{Different} & \multirow{2}{*}{Different}
\\
& E = No Data \tab Freq =No Data   &    &  \\ \cline{2-2}
& 4NFAACe   & \multicolumn{2}{c}
 {RMS = 0.111695}
\\
& E = No Data \tab Freq =No Data   &    \multicolumn{2}{c} 
{ }
\\ \hline
\multirow{4}{*}{\tab[2mm] 97 \tab[2mm]} & 3NFAACd &
\multirow{2}{*}{Different} & \multirow{2}{*}{Different}
\\
& E = No Data \tab Freq =No Data   &    &  \\ \cline{2-2}
& 4NFAACf   & \multicolumn{2}{c}
 {RMS = 0.207153}
\\
& E = No Data \tab Freq =No Data   &    \multicolumn{2}{c} 
{ }
\\ \hline
\multirow{4}{*}{\tab[2mm] 98 \tab[2mm]} & 3NFAACd &
\multirow{2}{*}{Different} & \multirow{2}{*}{Different}
\\
& E = No Data \tab Freq =No Data   &    &  \\ \cline{2-2}
& 4NFAACg   & \multicolumn{2}{c}
 {RMS = 0.12271}
\\
& E = No Data \tab Freq =No Data   &    \multicolumn{2}{c} 
{ }
\\ \hline
\end{tabular}
\newpage

\vtab[-2cm]
\tab[-2cm]
\begin{tabular}{c|m{8cm}|c|c}
\# & Moléculas & Restultado esperado & Resultado programa \\ \hline\hline
\multirow{4}{*}{\tab[2mm] 99 \tab[2mm]} & 3NFAACd &
\multirow{2}{*}{Different} & \multirow{2}{*}{Different}
\\
& E = No Data \tab Freq =No Data   &    &  \\ \cline{2-2}
& 4NFAACi   & \multicolumn{2}{c}
 {RMS = 0.0873887}
\\
& E = No Data \tab Freq =No Data   &    \multicolumn{2}{c} 
{ }
\\ \hline
\multirow{4}{*}{\tab[2mm] 100 \tab[2mm]} & 3NFAACd &
\multirow{2}{*}{Different} & \multirow{2}{*}{Different}
\\
& E = No Data \tab Freq =No Data   &    &  \\ \cline{2-2}
& 4NFAACj   & \multicolumn{2}{c}
 {RMS = 0.175687}
\\
& E = No Data \tab Freq =No Data   &    \multicolumn{2}{c} 
{ }
\\ \hline
\multirow{4}{*}{\tab[2mm] 101 \tab[2mm]} & 3NFAACd &
\multirow{2}{*}{Different} & \multirow{2}{*}{Different}
\\
& E = No Data \tab Freq =No Data   &    &  \\ \cline{2-2}
& 4NFAACl-3   & \multicolumn{2}{c}
 {RMS = 0.130334}
\\
& E = No Data \tab Freq =No Data   &    \multicolumn{2}{c} 
{ }
\\ \hline
\multirow{4}{*}{\tab[2mm] 102 \tab[2mm]} & 3NFAACe &
\multirow{2}{*}{Different} & \multirow{2}{*}{Different}
\\
& E = No Data \tab Freq =No Data   &    &  \\ \cline{2-2}
& 3NFAACf   & \multicolumn{2}{c}
 {RMS = 0.107565}
\\
& E = No Data \tab Freq =No Data   &    \multicolumn{2}{c} 
{ }
\\ \hline
\multirow{4}{*}{\tab[2mm] 103 \tab[2mm]} & 3NFAACe &
\multirow{2}{*}{Different} & \multirow{2}{*}{Different}
\\
& E = No Data \tab Freq =No Data   &    &  \\ \cline{2-2}
& 3NFAACg   & \multicolumn{2}{c}
 {RMS = 0.195216}
\\
& E = No Data \tab Freq =No Data   &    \multicolumn{2}{c} 
{ }
\\ \hline
\multirow{4}{*}{\tab[2mm] 104 \tab[2mm]} & 3NFAACe &
\multirow{2}{*}{Different} & \multirow{2}{*}{Different}
\\
& E = No Data \tab Freq =No Data   &    &  \\ \cline{2-2}
& 3NFAACh   & \multicolumn{2}{c}
 {RMS = 0.0977146}
\\
& E = No Data \tab Freq =No Data   &    \multicolumn{2}{c} 
{ }
\\ \hline
\multirow{4}{*}{\tab[2mm] 105 \tab[2mm]} & 3NFAACe &
\multirow{2}{*}{Different} & \multirow{2}{*}{Different}
\\
& E = No Data \tab Freq =No Data   &    &  \\ \cline{2-2}
& 3NFAACi   & \multicolumn{2}{c}
 {RMS = 0.245834}
\\
& E = No Data \tab Freq =No Data   &    \multicolumn{2}{c} 
{ }
\\ \hline
\end{tabular}
\newpage

\vtab[-2cm]
\tab[-2cm]
\begin{tabular}{c|m{8cm}|c|c}
\# & Moléculas & Restultado esperado & Resultado programa \\ \hline\hline
\multirow{4}{*}{\tab[2mm] 106 \tab[2mm]} & 3NFAACe &
\multirow{2}{*}{Different} & \multirow{2}{*}{Different}
\\
& E = No Data \tab Freq =No Data   &    &  \\ \cline{2-2}
& 3NFAACj   & \multicolumn{2}{c}
 {RMS = 0.252452}
\\
& E = No Data \tab Freq =No Data   &    \multicolumn{2}{c} 
{ }
\\ \hline
\multirow{4}{*}{\tab[2mm] 107 \tab[2mm]} & 3NFAACe &
\multirow{2}{*}{Different} & \multirow{2}{*}{Different}
\\
& E = No Data \tab Freq =No Data   &    &  \\ \cline{2-2}
& 3NFAACk   & \multicolumn{2}{c}
 {RMS = 0.173543}
\\
& E = No Data \tab Freq =No Data   &    \multicolumn{2}{c} 
{ }
\\ \hline
\multirow{4}{*}{\tab[2mm] 108 \tab[2mm]} & 3NFAACe &
\multirow{2}{*}{Different} & \multirow{2}{*}{Different}
\\
& E = No Data \tab Freq =No Data   &    &  \\ \cline{2-2}
& 3NFAACl   & \multicolumn{2}{c}
 {RMS = 0.125902}
\\
& E = No Data \tab Freq =No Data   &    \multicolumn{2}{c} 
{ }
\\ \hline
\multirow{4}{*}{\tab[2mm] 109 \tab[2mm]} & 3NFAACe &
\multirow{2}{*}{Different} & \multirow{2}{*}{Different}
\\
& E = No Data \tab Freq =No Data   &    &  \\ \cline{2-2}
& 3NFAACm   & \multicolumn{2}{c}
 {RMS = 0.00929873}
\\
& E = No Data \tab Freq =No Data   &    \multicolumn{2}{c} 
{ }
\\ \hline
\multirow{4}{*}{\tab[2mm] 110 \tab[2mm]} & 3NFAACe &
\multirow{2}{*}{Different} & \multirow{2}{*}{Different}
\\
& E = No Data \tab Freq =No Data   &    &  \\ \cline{2-2}
& 3NFAACn   & \multicolumn{2}{c}
 {RMS = 0.00494937}
\\
& E = No Data \tab Freq =No Data   &    \multicolumn{2}{c} 
{ }
\\ \hline
\multirow{4}{*}{\tab[2mm] 111 \tab[2mm]} & 3NFAACe &
\multirow{2}{*}{Different} & \multirow{2}{*}{Different}
\\
& E = No Data \tab Freq =No Data   &    &  \\ \cline{2-2}
& 4NFAACa   & \multicolumn{2}{c}
 {RMS = 0.381909}
\\
& E = No Data \tab Freq =No Data   &    \multicolumn{2}{c} 
{ }
\\ \hline
\multirow{4}{*}{\tab[2mm] 112 \tab[2mm]} & 3NFAACe &
\multirow{2}{*}{Different} & \multirow{2}{*}{Different}
\\
& E = No Data \tab Freq =No Data   &    &  \\ \cline{2-2}
& 4NFAACb   & \multicolumn{2}{c}
 {RMS = 0.0963103}
\\
& E = No Data \tab Freq =No Data   &    \multicolumn{2}{c} 
{ }
\\ \hline
\end{tabular}
\newpage

\vtab[-2cm]
\tab[-2cm]
\begin{tabular}{c|m{8cm}|c|c}
\# & Moléculas & Restultado esperado & Resultado programa \\ \hline\hline
\multirow{4}{*}{\tab[2mm] 113 \tab[2mm]} & 3NFAACe &
\multirow{2}{*}{Different} & \multirow{2}{*}{Different}
\\
& E = No Data \tab Freq =No Data   &    &  \\ \cline{2-2}
& 4NFAACc   & \multicolumn{2}{c}
 {RMS = 0.143506}
\\
& E = No Data \tab Freq =No Data   &    \multicolumn{2}{c} 
{ }
\\ \hline
\multirow{4}{*}{\tab[2mm] 114 \tab[2mm]} & 3NFAACe &
\multirow{2}{*}{Different} & \multirow{2}{*}{Different}
\\
& E = No Data \tab Freq =No Data   &    &  \\ \cline{2-2}
& 4NFAACd   & \multicolumn{2}{c}
 {RMS = 0.147224}
\\
& E = No Data \tab Freq =No Data   &    \multicolumn{2}{c} 
{ }
\\ \hline
\multirow{4}{*}{\tab[2mm] 115 \tab[2mm]} & 3NFAACe &
\multirow{2}{*}{Different} & \multirow{2}{*}{Different}
\\
& E = No Data \tab Freq =No Data   &    &  \\ \cline{2-2}
& 4NFAACe   & \multicolumn{2}{c}
 {RMS = 0.139288}
\\
& E = No Data \tab Freq =No Data   &    \multicolumn{2}{c} 
{ }
\\ \hline
\multirow{4}{*}{\tab[2mm] 116 \tab[2mm]} & 3NFAACe &
\multirow{2}{*}{Different} & \multirow{2}{*}{Different}
\\
& E = No Data \tab Freq =No Data   &    &  \\ \cline{2-2}
& 4NFAACf   & \multicolumn{2}{c}
 {RMS = 0.0438293}
\\
& E = No Data \tab Freq =No Data   &    \multicolumn{2}{c} 
{ }
\\ \hline
\multirow{4}{*}{\tab[2mm] 117 \tab[2mm]} & 3NFAACe &
\multirow{2}{*}{Different} & \multirow{2}{*}{Different}
\\
& E = No Data \tab Freq =No Data   &    &  \\ \cline{2-2}
& 4NFAACg   & \multicolumn{2}{c}
 {RMS = 0.128273}
\\
& E = No Data \tab Freq =No Data   &    \multicolumn{2}{c} 
{ }
\\ \hline
\multirow{4}{*}{\tab[2mm] 118 \tab[2mm]} & 3NFAACe &
\multirow{2}{*}{Different} & \multirow{2}{*}{Different}
\\
& E = No Data \tab Freq =No Data   &    &  \\ \cline{2-2}
& 4NFAACi   & \multicolumn{2}{c}
 {RMS = 0.163594}
\\
& E = No Data \tab Freq =No Data   &    \multicolumn{2}{c} 
{ }
\\ \hline
\multirow{4}{*}{\tab[2mm] 119 \tab[2mm]} & 3NFAACe &
\multirow{2}{*}{Different} & \multirow{2}{*}{Different}
\\
& E = No Data \tab Freq =No Data   &    &  \\ \cline{2-2}
& 4NFAACj   & \multicolumn{2}{c}
 {RMS = 0.075296}
\\
& E = No Data \tab Freq =No Data   &    \multicolumn{2}{c} 
{ }
\\ \hline
\end{tabular}
\newpage

\vtab[-2cm]
\tab[-2cm]
\begin{tabular}{c|m{8cm}|c|c}
\# & Moléculas & Restultado esperado & Resultado programa \\ \hline\hline
\multirow{4}{*}{\tab[2mm] 120 \tab[2mm]} & 3NFAACe &
\multirow{2}{*}{Different} & \multirow{2}{*}{Different}
\\
& E = No Data \tab Freq =No Data   &    &  \\ \cline{2-2}
& 4NFAACl-3   & \multicolumn{2}{c}
 {RMS = 0.120649}
\\
& E = No Data \tab Freq =No Data   &    \multicolumn{2}{c} 
{ }
\\ \hline
\multirow{4}{*}{\tab[2mm] 121 \tab[2mm]} & 3NFAACf &
\multirow{2}{*}{Different} & \multirow{2}{*}{Different}
\\
& E = No Data \tab Freq =No Data   &    &  \\ \cline{2-2}
& 3NFAACg   & \multicolumn{2}{c}
 {RMS = 0.0876514}
\\
& E = No Data \tab Freq =No Data   &    \multicolumn{2}{c} 
{ }
\\ \hline
\multirow{4}{*}{\tab[2mm] 122 \tab[2mm]} & 3NFAACf &
\multirow{2}{*}{Different} & \multirow{2}{*}{Different}
\\
& E = No Data \tab Freq =No Data   &    &  \\ \cline{2-2}
& 3NFAACh   & \multicolumn{2}{c}
 {RMS = 0.00985042}
\\
& E = No Data \tab Freq =No Data   &    \multicolumn{2}{c} 
{ }
\\ \hline
\multirow{4}{*}{\tab[2mm] 123 \tab[2mm]} & 3NFAACf &
\multirow{2}{*}{Different} & \multirow{2}{*}{Different}
\\
& E = No Data \tab Freq =No Data   &    &  \\ \cline{2-2}
& 3NFAACi   & \multicolumn{2}{c}
 {RMS = 0.138269}
\\
& E = No Data \tab Freq =No Data   &    \multicolumn{2}{c} 
{ }
\\ \hline
\multirow{4}{*}{\tab[2mm] 124 \tab[2mm]} & 3NFAACf &
\multirow{2}{*}{Different} & \multirow{2}{*}{Different}
\\
& E = No Data \tab Freq =No Data   &    &  \\ \cline{2-2}
& 3NFAACj   & \multicolumn{2}{c}
 {RMS = 0.144887}
\\
& E = No Data \tab Freq =No Data   &    \multicolumn{2}{c} 
{ }
\\ \hline
\multirow{4}{*}{\tab[2mm] 125 \tab[2mm]} & 3NFAACf &
\multirow{2}{*}{Different} & \multirow{2}{*}{Different}
\\
& E = No Data \tab Freq =No Data   &    &  \\ \cline{2-2}
& 3NFAACk   & \multicolumn{2}{c}
 {RMS = 0.0659782}
\\
& E = No Data \tab Freq =No Data   &    \multicolumn{2}{c} 
{ }
\\ \hline
\multirow{4}{*}{\tab[2mm] 126 \tab[2mm]} & 3NFAACf &
\multirow{2}{*}{Different} & \multirow{2}{*}{Different}
\\
& E = No Data \tab Freq =No Data   &    &  \\ \cline{2-2}
& 3NFAACl   & \multicolumn{2}{c}
 {RMS = 0.0183372}
\\
& E = No Data \tab Freq =No Data   &    \multicolumn{2}{c} 
{ }
\\ \hline
\end{tabular}
\newpage

\vtab[-2cm]
\tab[-2cm]
\begin{tabular}{c|m{8cm}|c|c}
\# & Moléculas & Restultado esperado & Resultado programa \\ \hline\hline
\multirow{4}{*}{\tab[2mm] 127 \tab[2mm]} & 3NFAACf &
\multirow{2}{*}{Different} & \multirow{2}{*}{Different}
\\
& E = No Data \tab Freq =No Data   &    &  \\ \cline{2-2}
& 3NFAACm   & \multicolumn{2}{c}
 {RMS = 0.0982663}
\\
& E = No Data \tab Freq =No Data   &    \multicolumn{2}{c} 
{ }
\\ \hline
\multirow{4}{*}{\tab[2mm] 128 \tab[2mm]} & 3NFAACf &
\multirow{2}{*}{Different} & \multirow{2}{*}{Different}
\\
& E = No Data \tab Freq =No Data   &    &  \\ \cline{2-2}
& 3NFAACn   & \multicolumn{2}{c}
 {RMS = 0.102616}
\\
& E = No Data \tab Freq =No Data   &    \multicolumn{2}{c} 
{ }
\\ \hline
\multirow{4}{*}{\tab[2mm] 129 \tab[2mm]} & 3NFAACf &
\multirow{2}{*}{Different} & \multirow{2}{*}{Different}
\\
& E = No Data \tab Freq =No Data   &    &  \\ \cline{2-2}
& 4NFAACa   & \multicolumn{2}{c}
 {RMS = 0.274344}
\\
& E = No Data \tab Freq =No Data   &    \multicolumn{2}{c} 
{ }
\\ \hline
\multirow{4}{*}{\tab[2mm] 130 \tab[2mm]} & 3NFAACf &
\multirow{2}{*}{Different} & \multirow{2}{*}{Different}
\\
& E = No Data \tab Freq =No Data   &    &  \\ \cline{2-2}
& 4NFAACb   & \multicolumn{2}{c}
 {RMS = 0.203875}
\\
& E = No Data \tab Freq =No Data   &    \multicolumn{2}{c} 
{ }
\\ \hline
\multirow{4}{*}{\tab[2mm] 131 \tab[2mm]} & 3NFAACf &
\multirow{2}{*}{Different} & \multirow{2}{*}{Different}
\\
& E = No Data \tab Freq =No Data   &    &  \\ \cline{2-2}
& 4NFAACc   & \multicolumn{2}{c}
 {RMS = 0.0359411}
\\
& E = No Data \tab Freq =No Data   &    \multicolumn{2}{c} 
{ }
\\ \hline
\multirow{4}{*}{\tab[2mm] 132 \tab[2mm]} & 3NFAACf &
\multirow{2}{*}{Different} & \multirow{2}{*}{Different}
\\
& E = No Data \tab Freq =No Data   &    &  \\ \cline{2-2}
& 4NFAACd   & \multicolumn{2}{c}
 {RMS = 0.0396588}
\\
& E = No Data \tab Freq =No Data   &    \multicolumn{2}{c} 
{ }
\\ \hline
\multirow{4}{*}{\tab[2mm] 133 \tab[2mm]} & 3NFAACf &
\multirow{2}{*}{Different} & \multirow{2}{*}{Different}
\\
& E = No Data \tab Freq =No Data   &    &  \\ \cline{2-2}
& 4NFAACe   & \multicolumn{2}{c}
 {RMS = 0.0317228}
\\
& E = No Data \tab Freq =No Data   &    \multicolumn{2}{c} 
{ }
\\ \hline
\end{tabular}
\newpage

\vtab[-2cm]
\tab[-2cm]
\begin{tabular}{c|m{8cm}|c|c}
\# & Moléculas & Restultado esperado & Resultado programa \\ \hline\hline
\multirow{4}{*}{\tab[2mm] 134 \tab[2mm]} & 3NFAACf &
\multirow{2}{*}{Different} & \multirow{2}{*}{Different}
\\
& E = No Data \tab Freq =No Data   &    &  \\ \cline{2-2}
& 4NFAACf   & \multicolumn{2}{c}
 {RMS = 0.0637358}
\\
& E = No Data \tab Freq =No Data   &    \multicolumn{2}{c} 
{ }
\\ \hline
\multirow{4}{*}{\tab[2mm] 135 \tab[2mm]} & 3NFAACf &
\multirow{2}{*}{Different} & \multirow{2}{*}{Different}
\\
& E = No Data \tab Freq =No Data   &    &  \\ \cline{2-2}
& 4NFAACg   & \multicolumn{2}{c}
 {RMS = 0.0207081}
\\
& E = No Data \tab Freq =No Data   &    \multicolumn{2}{c} 
{ }
\\ \hline
\multirow{4}{*}{\tab[2mm] 136 \tab[2mm]} & 3NFAACf &
\multirow{2}{*}{Different} & \multirow{2}{*}{Different}
\\
& E = No Data \tab Freq =No Data   &    &  \\ \cline{2-2}
& 4NFAACi   & \multicolumn{2}{c}
 {RMS = 0.056029}
\\
& E = No Data \tab Freq =No Data   &    \multicolumn{2}{c} 
{ }
\\ \hline
\multirow{4}{*}{\tab[2mm] 137 \tab[2mm]} & 3NFAACf &
\multirow{2}{*}{Different} & \multirow{2}{*}{Different}
\\
& E = No Data \tab Freq =No Data   &    &  \\ \cline{2-2}
& 4NFAACj   & \multicolumn{2}{c}
 {RMS = 0.032269}
\\
& E = No Data \tab Freq =No Data   &    \multicolumn{2}{c} 
{ }
\\ \hline
\multirow{4}{*}{\tab[2mm] 138 \tab[2mm]} & 3NFAACf &
\multirow{2}{*}{Different} & \multirow{2}{*}{Different}
\\
& E = No Data \tab Freq =No Data   &    &  \\ \cline{2-2}
& 4NFAACl-3   & \multicolumn{2}{c}
 {RMS = 0.013084}
\\
& E = No Data \tab Freq =No Data   &    \multicolumn{2}{c} 
{ }
\\ \hline
\multirow{4}{*}{\tab[2mm] 139 \tab[2mm]} & 3NFAACg &
\multirow{2}{*}{Different} & \multirow{2}{*}{Different}
\\
& E = No Data \tab Freq =No Data   &    &  \\ \cline{2-2}
& 3NFAACh   & \multicolumn{2}{c}
 {RMS = 0.0975018}
\\
& E = No Data \tab Freq =No Data   &    \multicolumn{2}{c} 
{ }
\\ \hline
\multirow{4}{*}{\tab[2mm] 140 \tab[2mm]} & 3NFAACg &
\multirow{2}{*}{Different} & \multirow{2}{*}{Different}
\\
& E = No Data \tab Freq =No Data   &    &  \\ \cline{2-2}
& 3NFAACi   & \multicolumn{2}{c}
 {RMS = 0.0506173}
\\
& E = No Data \tab Freq =No Data   &    \multicolumn{2}{c} 
{ }
\\ \hline
\end{tabular}
\newpage

\vtab[-2cm]
\tab[-2cm]
\begin{tabular}{c|m{8cm}|c|c}
\# & Moléculas & Restultado esperado & Resultado programa \\ \hline\hline
\multirow{4}{*}{\tab[2mm] 141 \tab[2mm]} & 3NFAACg &
\multirow{2}{*}{Different} & \multirow{2}{*}{Different}
\\
& E = No Data \tab Freq =No Data   &    &  \\ \cline{2-2}
& 3NFAACj   & \multicolumn{2}{c}
 {RMS = 0.0572353}
\\
& E = No Data \tab Freq =No Data   &    \multicolumn{2}{c} 
{ }
\\ \hline
\multirow{4}{*}{\tab[2mm] 142 \tab[2mm]} & 3NFAACg &
\multirow{2}{*}{Different} & \multirow{2}{*}{Different}
\\
& E = No Data \tab Freq =No Data   &    &  \\ \cline{2-2}
& 3NFAACk   & \multicolumn{2}{c}
 {RMS = 0.0216732}
\\
& E = No Data \tab Freq =No Data   &    \multicolumn{2}{c} 
{ }
\\ \hline
\multirow{4}{*}{\tab[2mm] 143 \tab[2mm]} & 3NFAACg &
\multirow{2}{*}{Different} & \multirow{2}{*}{Different}
\\
& E = No Data \tab Freq =No Data   &    &  \\ \cline{2-2}
& 3NFAACl   & \multicolumn{2}{c}
 {RMS = 0.0693142}
\\
& E = No Data \tab Freq =No Data   &    \multicolumn{2}{c} 
{ }
\\ \hline
\multirow{4}{*}{\tab[2mm] 144 \tab[2mm]} & 3NFAACg &
\multirow{2}{*}{Different} & \multirow{2}{*}{Different}
\\
& E = No Data \tab Freq =No Data   &    &  \\ \cline{2-2}
& 3NFAACm   & \multicolumn{2}{c}
 {RMS = 0.185918}
\\
& E = No Data \tab Freq =No Data   &    \multicolumn{2}{c} 
{ }
\\ \hline
\multirow{4}{*}{\tab[2mm] 145 \tab[2mm]} & 3NFAACg &
\multirow{2}{*}{Different} & \multirow{2}{*}{Different}
\\
& E = No Data \tab Freq =No Data   &    &  \\ \cline{2-2}
& 3NFAACn   & \multicolumn{2}{c}
 {RMS = 0.190267}
\\
& E = No Data \tab Freq =No Data   &    \multicolumn{2}{c} 
{ }
\\ \hline
\multirow{4}{*}{\tab[2mm] 146 \tab[2mm]} & 3NFAACg &
\multirow{2}{*}{Different} & \multirow{2}{*}{Different}
\\
& E = No Data \tab Freq =No Data   &    &  \\ \cline{2-2}
& 4NFAACa   & \multicolumn{2}{c}
 {RMS = 0.186693}
\\
& E = No Data \tab Freq =No Data   &    \multicolumn{2}{c} 
{ }
\\ \hline
\multirow{4}{*}{\tab[2mm] 147 \tab[2mm]} & 3NFAACg &
\multirow{2}{*}{Different} & \multirow{2}{*}{Different}
\\
& E = No Data \tab Freq =No Data   &    &  \\ \cline{2-2}
& 4NFAACb   & \multicolumn{2}{c}
 {RMS = 0.291527}
\\
& E = No Data \tab Freq =No Data   &    \multicolumn{2}{c} 
{ }
\\ \hline
\end{tabular}
\newpage

\vtab[-2cm]
\tab[-2cm]
\begin{tabular}{c|m{8cm}|c|c}
\# & Moléculas & Restultado esperado & Resultado programa \\ \hline\hline
\multirow{4}{*}{\tab[2mm] 148 \tab[2mm]} & 3NFAACg &
\multirow{2}{*}{Different} & \multirow{2}{*}{Different}
\\
& E = No Data \tab Freq =No Data   &    &  \\ \cline{2-2}
& 4NFAACc   & \multicolumn{2}{c}
 {RMS = 0.0517103}
\\
& E = No Data \tab Freq =No Data   &    \multicolumn{2}{c} 
{ }
\\ \hline
\multirow{4}{*}{\tab[2mm] 149 \tab[2mm]} & 3NFAACg &
\multirow{2}{*}{Different} & \multirow{2}{*}{Different}
\\
& E = No Data \tab Freq =No Data   &    &  \\ \cline{2-2}
& 4NFAACd   & \multicolumn{2}{c}
 {RMS = 0.0479926}
\\
& E = No Data \tab Freq =No Data   &    \multicolumn{2}{c} 
{ }
\\ \hline
\multirow{4}{*}{\tab[2mm] 150 \tab[2mm]} & 3NFAACg &
\multirow{2}{*}{Different} & \multirow{2}{*}{Different}
\\
& E = No Data \tab Freq =No Data   &    &  \\ \cline{2-2}
& 4NFAACe   & \multicolumn{2}{c}
 {RMS = 0.0559286}
\\
& E = No Data \tab Freq =No Data   &    \multicolumn{2}{c} 
{ }
\\ \hline
\multirow{4}{*}{\tab[2mm] 151 \tab[2mm]} & 3NFAACg &
\multirow{2}{*}{Different} & \multirow{2}{*}{Different}
\\
& E = No Data \tab Freq =No Data   &    &  \\ \cline{2-2}
& 4NFAACf   & \multicolumn{2}{c}
 {RMS = 0.151387}
\\
& E = No Data \tab Freq =No Data   &    \multicolumn{2}{c} 
{ }
\\ \hline
\multirow{4}{*}{\tab[2mm] 152 \tab[2mm]} & 3NFAACg &
\multirow{2}{*}{Different} & \multirow{2}{*}{Different}
\\
& E = No Data \tab Freq =No Data   &    &  \\ \cline{2-2}
& 4NFAACg   & \multicolumn{2}{c}
 {RMS = 0.0669433}
\\
& E = No Data \tab Freq =No Data   &    \multicolumn{2}{c} 
{ }
\\ \hline
\multirow{4}{*}{\tab[2mm] 153 \tab[2mm]} & 3NFAACg &
\multirow{2}{*}{Different} & \multirow{2}{*}{Different}
\\
& E = No Data \tab Freq =No Data   &    &  \\ \cline{2-2}
& 4NFAACi   & \multicolumn{2}{c}
 {RMS = 0.0316224}
\\
& E = No Data \tab Freq =No Data   &    \multicolumn{2}{c} 
{ }
\\ \hline
\multirow{4}{*}{\tab[2mm] 154 \tab[2mm]} & 3NFAACg &
\multirow{2}{*}{Different} & \multirow{2}{*}{Different}
\\
& E = No Data \tab Freq =No Data   &    &  \\ \cline{2-2}
& 4NFAACj   & \multicolumn{2}{c}
 {RMS = 0.11992}
\\
& E = No Data \tab Freq =No Data   &    \multicolumn{2}{c} 
{ }
\\ \hline
\end{tabular}
\newpage

\vtab[-2cm]
\tab[-2cm]
\begin{tabular}{c|m{8cm}|c|c}
\# & Moléculas & Restultado esperado & Resultado programa \\ \hline\hline
\multirow{4}{*}{\tab[2mm] 155 \tab[2mm]} & 3NFAACg &
\multirow{2}{*}{Different} & \multirow{2}{*}{Different}
\\
& E = No Data \tab Freq =No Data   &    &  \\ \cline{2-2}
& 4NFAACl-3   & \multicolumn{2}{c}
 {RMS = 0.0745674}
\\
& E = No Data \tab Freq =No Data   &    \multicolumn{2}{c} 
{ }
\\ \hline
\multirow{4}{*}{\tab[2mm] 156 \tab[2mm]} & 3NFAACh &
\multirow{2}{*}{Different} & \multirow{2}{*}{Different}
\\
& E = No Data \tab Freq =No Data   &    &  \\ \cline{2-2}
& 3NFAACi   & \multicolumn{2}{c}
 {RMS = 0.148119}
\\
& E = No Data \tab Freq =No Data   &    \multicolumn{2}{c} 
{ }
\\ \hline
\multirow{4}{*}{\tab[2mm] 157 \tab[2mm]} & 3NFAACh &
\multirow{2}{*}{Different} & \multirow{2}{*}{Different}
\\
& E = No Data \tab Freq =No Data   &    &  \\ \cline{2-2}
& 3NFAACj   & \multicolumn{2}{c}
 {RMS = 0.154737}
\\
& E = No Data \tab Freq =No Data   &    \multicolumn{2}{c} 
{ }
\\ \hline
\multirow{4}{*}{\tab[2mm] 158 \tab[2mm]} & 3NFAACh &
\multirow{2}{*}{Different} & \multirow{2}{*}{Different}
\\
& E = No Data \tab Freq =No Data   &    &  \\ \cline{2-2}
& 3NFAACk   & \multicolumn{2}{c}
 {RMS = 0.0758286}
\\
& E = No Data \tab Freq =No Data   &    \multicolumn{2}{c} 
{ }
\\ \hline
\multirow{4}{*}{\tab[2mm] 159 \tab[2mm]} & 3NFAACh &
\multirow{2}{*}{Different} & \multirow{2}{*}{Different}
\\
& E = No Data \tab Freq =No Data   &    &  \\ \cline{2-2}
& 3NFAACl   & \multicolumn{2}{c}
 {RMS = 0.0281876}
\\
& E = No Data \tab Freq =No Data   &    \multicolumn{2}{c} 
{ }
\\ \hline
\multirow{4}{*}{\tab[2mm] 160 \tab[2mm]} & 3NFAACh &
\multirow{2}{*}{Different} & \multirow{2}{*}{Different}
\\
& E = No Data \tab Freq =No Data   &    &  \\ \cline{2-2}
& 3NFAACm   & \multicolumn{2}{c}
 {RMS = 0.0884159}
\\
& E = No Data \tab Freq =No Data   &    \multicolumn{2}{c} 
{ }
\\ \hline
\multirow{4}{*}{\tab[2mm] 161 \tab[2mm]} & 3NFAACh &
\multirow{2}{*}{Different} & \multirow{2}{*}{Different}
\\
& E = No Data \tab Freq =No Data   &    &  \\ \cline{2-2}
& 3NFAACn   & \multicolumn{2}{c}
 {RMS = 0.0927652}
\\
& E = No Data \tab Freq =No Data   &    \multicolumn{2}{c} 
{ }
\\ \hline
\end{tabular}
\newpage

\vtab[-2cm]
\tab[-2cm]
\begin{tabular}{c|m{8cm}|c|c}
\# & Moléculas & Restultado esperado & Resultado programa \\ \hline\hline
\multirow{4}{*}{\tab[2mm] 162 \tab[2mm]} & 3NFAACh &
\multirow{2}{*}{Different} & \multirow{2}{*}{Different}
\\
& E = No Data \tab Freq =No Data   &    &  \\ \cline{2-2}
& 4NFAACa   & \multicolumn{2}{c}
 {RMS = 0.284195}
\\
& E = No Data \tab Freq =No Data   &    \multicolumn{2}{c} 
{ }
\\ \hline
\multirow{4}{*}{\tab[2mm] 163 \tab[2mm]} & 3NFAACh &
\multirow{2}{*}{Different} & \multirow{2}{*}{Different}
\\
& E = No Data \tab Freq =No Data   &    &  \\ \cline{2-2}
& 4NFAACb   & \multicolumn{2}{c}
 {RMS = 0.194025}
\\
& E = No Data \tab Freq =No Data   &    \multicolumn{2}{c} 
{ }
\\ \hline
\multirow{4}{*}{\tab[2mm] 164 \tab[2mm]} & 3NFAACh &
\multirow{2}{*}{Different} & \multirow{2}{*}{Different}
\\
& E = No Data \tab Freq =No Data   &    &  \\ \cline{2-2}
& 4NFAACc   & \multicolumn{2}{c}
 {RMS = 0.0457915}
\\
& E = No Data \tab Freq =No Data   &    \multicolumn{2}{c} 
{ }
\\ \hline
\multirow{4}{*}{\tab[2mm] 165 \tab[2mm]} & 3NFAACh &
\multirow{2}{*}{Different} & \multirow{2}{*}{Different}
\\
& E = No Data \tab Freq =No Data   &    &  \\ \cline{2-2}
& 4NFAACd   & \multicolumn{2}{c}
 {RMS = 0.0495092}
\\
& E = No Data \tab Freq =No Data   &    \multicolumn{2}{c} 
{ }
\\ \hline
\multirow{4}{*}{\tab[2mm] 166 \tab[2mm]} & 3NFAACh &
\multirow{2}{*}{Different} & \multirow{2}{*}{Different}
\\
& E = No Data \tab Freq =No Data   &    &  \\ \cline{2-2}
& 4NFAACe   & \multicolumn{2}{c}
 {RMS = 0.0415732}
\\
& E = No Data \tab Freq =No Data   &    \multicolumn{2}{c} 
{ }
\\ \hline
\multirow{4}{*}{\tab[2mm] 167 \tab[2mm]} & 3NFAACh &
\multirow{2}{*}{Different} & \multirow{2}{*}{Different}
\\
& E = No Data \tab Freq =No Data   &    &  \\ \cline{2-2}
& 4NFAACf   & \multicolumn{2}{c}
 {RMS = 0.0538853}
\\
& E = No Data \tab Freq =No Data   &    \multicolumn{2}{c} 
{ }
\\ \hline
\multirow{4}{*}{\tab[2mm] 168 \tab[2mm]} & 3NFAACh &
\multirow{2}{*}{Different} & \multirow{2}{*}{Different}
\\
& E = No Data \tab Freq =No Data   &    &  \\ \cline{2-2}
& 4NFAACg   & \multicolumn{2}{c}
 {RMS = 0.0305585}
\\
& E = No Data \tab Freq =No Data   &    \multicolumn{2}{c} 
{ }
\\ \hline
\end{tabular}
\newpage

\vtab[-2cm]
\tab[-2cm]
\begin{tabular}{c|m{8cm}|c|c}
\# & Moléculas & Restultado esperado & Resultado programa \\ \hline\hline
\multirow{4}{*}{\tab[2mm] 169 \tab[2mm]} & 3NFAACh &
\multirow{2}{*}{Different} & \multirow{2}{*}{Different}
\\
& E = No Data \tab Freq =No Data   &    &  \\ \cline{2-2}
& 4NFAACi   & \multicolumn{2}{c}
 {RMS = 0.0658794}
\\
& E = No Data \tab Freq =No Data   &    \multicolumn{2}{c} 
{ }
\\ \hline
\multirow{4}{*}{\tab[2mm] 170 \tab[2mm]} & 3NFAACh &
\multirow{2}{*}{Different} & \multirow{2}{*}{Different}
\\
& E = No Data \tab Freq =No Data   &    &  \\ \cline{2-2}
& 4NFAACj   & \multicolumn{2}{c}
 {RMS = 0.0224186}
\\
& E = No Data \tab Freq =No Data   &    \multicolumn{2}{c} 
{ }
\\ \hline
\multirow{4}{*}{\tab[2mm] 171 \tab[2mm]} & 3NFAACh &
\multirow{2}{*}{Different} & \multirow{2}{*}{Different}
\\
& E = No Data \tab Freq =No Data   &    &  \\ \cline{2-2}
& 4NFAACl-3   & \multicolumn{2}{c}
 {RMS = 0.0229344}
\\
& E = No Data \tab Freq =No Data   &    \multicolumn{2}{c} 
{ }
\\ \hline
\multirow{4}{*}{\tab[2mm] 172 \tab[2mm]} & 3NFAACi &
\multirow{2}{*}{Different} & \multirow{2}{*}{Different}
\\
& E = No Data \tab Freq =No Data   &    &  \\ \cline{2-2}
& 3NFAACj   & \multicolumn{2}{c}
 {RMS = 0.00661796}
\\
& E = No Data \tab Freq =No Data   &    \multicolumn{2}{c} 
{ }
\\ \hline
\multirow{4}{*}{\tab[2mm] 173 \tab[2mm]} & 3NFAACi &
\multirow{2}{*}{Different} & \multirow{2}{*}{Different}
\\
& E = No Data \tab Freq =No Data   &    &  \\ \cline{2-2}
& 3NFAACk   & \multicolumn{2}{c}
 {RMS = 0.0722905}
\\
& E = No Data \tab Freq =No Data   &    \multicolumn{2}{c} 
{ }
\\ \hline
\multirow{4}{*}{\tab[2mm] 174 \tab[2mm]} & 3NFAACi &
\multirow{2}{*}{Different} & \multirow{2}{*}{Different}
\\
& E = No Data \tab Freq =No Data   &    &  \\ \cline{2-2}
& 3NFAACl   & \multicolumn{2}{c}
 {RMS = 0.119932}
\\
& E = No Data \tab Freq =No Data   &    \multicolumn{2}{c} 
{ }
\\ \hline
\multirow{4}{*}{\tab[2mm] 175 \tab[2mm]} & 3NFAACi &
\multirow{2}{*}{Different} & \multirow{2}{*}{Different}
\\
& E = No Data \tab Freq =No Data   &    &  \\ \cline{2-2}
& 3NFAACm   & \multicolumn{2}{c}
 {RMS = 0.236535}
\\
& E = No Data \tab Freq =No Data   &    \multicolumn{2}{c} 
{ }
\\ \hline
\end{tabular}
\newpage

\vtab[-2cm]
\tab[-2cm]
\begin{tabular}{c|m{8cm}|c|c}
\# & Moléculas & Restultado esperado & Resultado programa \\ \hline\hline
\multirow{4}{*}{\tab[2mm] 176 \tab[2mm]} & 3NFAACi &
\multirow{2}{*}{Different} & \multirow{2}{*}{Different}
\\
& E = No Data \tab Freq =No Data   &    &  \\ \cline{2-2}
& 3NFAACn   & \multicolumn{2}{c}
 {RMS = 0.240884}
\\
& E = No Data \tab Freq =No Data   &    \multicolumn{2}{c} 
{ }
\\ \hline
\multirow{4}{*}{\tab[2mm] 177 \tab[2mm]} & 3NFAACi &
\multirow{2}{*}{Different} & \multirow{2}{*}{Different}
\\
& E = No Data \tab Freq =No Data   &    &  \\ \cline{2-2}
& 4NFAACa   & \multicolumn{2}{c}
 {RMS = 0.136076}
\\
& E = No Data \tab Freq =No Data   &    \multicolumn{2}{c} 
{ }
\\ \hline
\multirow{4}{*}{\tab[2mm] 178 \tab[2mm]} & 3NFAACi &
\multirow{2}{*}{Different} & \multirow{2}{*}{Different}
\\
& E = No Data \tab Freq =No Data   &    &  \\ \cline{2-2}
& 4NFAACb   & \multicolumn{2}{c}
 {RMS = 0.342144}
\\
& E = No Data \tab Freq =No Data   &    \multicolumn{2}{c} 
{ }
\\ \hline
\multirow{4}{*}{\tab[2mm] 179 \tab[2mm]} & 3NFAACi &
\multirow{2}{*}{Different} & \multirow{2}{*}{Different}
\\
& E = No Data \tab Freq =No Data   &    &  \\ \cline{2-2}
& 4NFAACc   & \multicolumn{2}{c}
 {RMS = 0.102328}
\\
& E = No Data \tab Freq =No Data   &    \multicolumn{2}{c} 
{ }
\\ \hline
\multirow{4}{*}{\tab[2mm] 180 \tab[2mm]} & 3NFAACi &
\multirow{2}{*}{Different} & \multirow{2}{*}{Different}
\\
& E = No Data \tab Freq =No Data   &    &  \\ \cline{2-2}
& 4NFAACd   & \multicolumn{2}{c}
 {RMS = 0.0986099}
\\
& E = No Data \tab Freq =No Data   &    \multicolumn{2}{c} 
{ }
\\ \hline
\multirow{4}{*}{\tab[2mm] 181 \tab[2mm]} & 3NFAACi &
\multirow{2}{*}{Different} & \multirow{2}{*}{Different}
\\
& E = No Data \tab Freq =No Data   &    &  \\ \cline{2-2}
& 4NFAACe   & \multicolumn{2}{c}
 {RMS = 0.106546}
\\
& E = No Data \tab Freq =No Data   &    \multicolumn{2}{c} 
{ }
\\ \hline
\multirow{4}{*}{\tab[2mm] 182 \tab[2mm]} & 3NFAACi &
\multirow{2}{*}{Different} & \multirow{2}{*}{Different}
\\
& E = No Data \tab Freq =No Data   &    &  \\ \cline{2-2}
& 4NFAACf   & \multicolumn{2}{c}
 {RMS = 0.202004}
\\
& E = No Data \tab Freq =No Data   &    \multicolumn{2}{c} 
{ }
\\ \hline
\end{tabular}
\newpage

\vtab[-2cm]
\tab[-2cm]
\begin{tabular}{c|m{8cm}|c|c}
\# & Moléculas & Restultado esperado & Resultado programa \\ \hline\hline
\multirow{4}{*}{\tab[2mm] 183 \tab[2mm]} & 3NFAACi &
\multirow{2}{*}{Different} & \multirow{2}{*}{Different}
\\
& E = No Data \tab Freq =No Data   &    &  \\ \cline{2-2}
& 4NFAACg   & \multicolumn{2}{c}
 {RMS = 0.117561}
\\
& E = No Data \tab Freq =No Data   &    \multicolumn{2}{c} 
{ }
\\ \hline
\multirow{4}{*}{\tab[2mm] 184 \tab[2mm]} & 3NFAACi &
\multirow{2}{*}{Different} & \multirow{2}{*}{Different}
\\
& E = No Data \tab Freq =No Data   &    &  \\ \cline{2-2}
& 4NFAACi   & \multicolumn{2}{c}
 {RMS = 0.0822397}
\\
& E = No Data \tab Freq =No Data   &    \multicolumn{2}{c} 
{ }
\\ \hline
\multirow{4}{*}{\tab[2mm] 185 \tab[2mm]} & 3NFAACi &
\multirow{2}{*}{Different} & \multirow{2}{*}{Different}
\\
& E = No Data \tab Freq =No Data   &    &  \\ \cline{2-2}
& 4NFAACj   & \multicolumn{2}{c}
 {RMS = 0.170538}
\\
& E = No Data \tab Freq =No Data   &    \multicolumn{2}{c} 
{ }
\\ \hline
\multirow{4}{*}{\tab[2mm] 186 \tab[2mm]} & 3NFAACi &
\multirow{2}{*}{Different} & \multirow{2}{*}{Different}
\\
& E = No Data \tab Freq =No Data   &    &  \\ \cline{2-2}
& 4NFAACl-3   & \multicolumn{2}{c}
 {RMS = 0.125185}
\\
& E = No Data \tab Freq =No Data   &    \multicolumn{2}{c} 
{ }
\\ \hline
\multirow{4}{*}{\tab[2mm] 187 \tab[2mm]} & 3NFAACj &
\multirow{2}{*}{Different} & \multirow{2}{*}{Different}
\\
& E = No Data \tab Freq =No Data   &    &  \\ \cline{2-2}
& 3NFAACk   & \multicolumn{2}{c}
 {RMS = 0.0789084}
\\
& E = No Data \tab Freq =No Data   &    \multicolumn{2}{c} 
{ }
\\ \hline
\multirow{4}{*}{\tab[2mm] 188 \tab[2mm]} & 3NFAACj &
\multirow{2}{*}{Different} & \multirow{2}{*}{Different}
\\
& E = No Data \tab Freq =No Data   &    &  \\ \cline{2-2}
& 3NFAACl   & \multicolumn{2}{c}
 {RMS = 0.126549}
\\
& E = No Data \tab Freq =No Data   &    \multicolumn{2}{c} 
{ }
\\ \hline
\multirow{4}{*}{\tab[2mm] 189 \tab[2mm]} & 3NFAACj &
\multirow{2}{*}{Different} & \multirow{2}{*}{Different}
\\
& E = No Data \tab Freq =No Data   &    &  \\ \cline{2-2}
& 3NFAACm   & \multicolumn{2}{c}
 {RMS = 0.243153}
\\
& E = No Data \tab Freq =No Data   &    \multicolumn{2}{c} 
{ }
\\ \hline
\end{tabular}
\newpage

\vtab[-2cm]
\tab[-2cm]
\begin{tabular}{c|m{8cm}|c|c}
\# & Moléculas & Restultado esperado & Resultado programa \\ \hline\hline
\multirow{4}{*}{\tab[2mm] 190 \tab[2mm]} & 3NFAACj &
\multirow{2}{*}{Different} & \multirow{2}{*}{Different}
\\
& E = No Data \tab Freq =No Data   &    &  \\ \cline{2-2}
& 3NFAACn   & \multicolumn{2}{c}
 {RMS = 0.247502}
\\
& E = No Data \tab Freq =No Data   &    \multicolumn{2}{c} 
{ }
\\ \hline
\multirow{4}{*}{\tab[2mm] 191 \tab[2mm]} & 3NFAACj &
\multirow{2}{*}{Different} & \multirow{2}{*}{Different}
\\
& E = No Data \tab Freq =No Data   &    &  \\ \cline{2-2}
& 4NFAACa   & \multicolumn{2}{c}
 {RMS = 0.129458}
\\
& E = No Data \tab Freq =No Data   &    \multicolumn{2}{c} 
{ }
\\ \hline
\multirow{4}{*}{\tab[2mm] 192 \tab[2mm]} & 3NFAACj &
\multirow{2}{*}{Different} & \multirow{2}{*}{Different}
\\
& E = No Data \tab Freq =No Data   &    &  \\ \cline{2-2}
& 4NFAACb   & \multicolumn{2}{c}
 {RMS = 0.348762}
\\
& E = No Data \tab Freq =No Data   &    \multicolumn{2}{c} 
{ }
\\ \hline
\multirow{4}{*}{\tab[2mm] 193 \tab[2mm]} & 3NFAACj &
\multirow{2}{*}{Different} & \multirow{2}{*}{Different}
\\
& E = No Data \tab Freq =No Data   &    &  \\ \cline{2-2}
& 4NFAACc   & \multicolumn{2}{c}
 {RMS = 0.108946}
\\
& E = No Data \tab Freq =No Data   &    \multicolumn{2}{c} 
{ }
\\ \hline
\multirow{4}{*}{\tab[2mm] 194 \tab[2mm]} & 3NFAACj &
\multirow{2}{*}{Different} & \multirow{2}{*}{Different}
\\
& E = No Data \tab Freq =No Data   &    &  \\ \cline{2-2}
& 4NFAACd   & \multicolumn{2}{c}
 {RMS = 0.105228}
\\
& E = No Data \tab Freq =No Data   &    \multicolumn{2}{c} 
{ }
\\ \hline
\multirow{4}{*}{\tab[2mm] 195 \tab[2mm]} & 3NFAACj &
\multirow{2}{*}{Different} & \multirow{2}{*}{Different}
\\
& E = No Data \tab Freq =No Data   &    &  \\ \cline{2-2}
& 4NFAACe   & \multicolumn{2}{c}
 {RMS = 0.113164}
\\
& E = No Data \tab Freq =No Data   &    \multicolumn{2}{c} 
{ }
\\ \hline
\multirow{4}{*}{\tab[2mm] 196 \tab[2mm]} & 3NFAACj &
\multirow{2}{*}{Different} & \multirow{2}{*}{Different}
\\
& E = No Data \tab Freq =No Data   &    &  \\ \cline{2-2}
& 4NFAACf   & \multicolumn{2}{c}
 {RMS = 0.208622}
\\
& E = No Data \tab Freq =No Data   &    \multicolumn{2}{c} 
{ }
\\ \hline
\end{tabular}
\newpage

\vtab[-2cm]
\tab[-2cm]
\begin{tabular}{c|m{8cm}|c|c}
\# & Moléculas & Restultado esperado & Resultado programa \\ \hline\hline
\multirow{4}{*}{\tab[2mm] 197 \tab[2mm]} & 3NFAACj &
\multirow{2}{*}{Different} & \multirow{2}{*}{Different}
\\
& E = No Data \tab Freq =No Data   &    &  \\ \cline{2-2}
& 4NFAACg   & \multicolumn{2}{c}
 {RMS = 0.124179}
\\
& E = No Data \tab Freq =No Data   &    \multicolumn{2}{c} 
{ }
\\ \hline
\multirow{4}{*}{\tab[2mm] 198 \tab[2mm]} & 3NFAACj &
\multirow{2}{*}{Different} & \multirow{2}{*}{Different}
\\
& E = No Data \tab Freq =No Data   &    &  \\ \cline{2-2}
& 4NFAACi   & \multicolumn{2}{c}
 {RMS = 0.0888577}
\\
& E = No Data \tab Freq =No Data   &    \multicolumn{2}{c} 
{ }
\\ \hline
\multirow{4}{*}{\tab[2mm] 199 \tab[2mm]} & 3NFAACj &
\multirow{2}{*}{Different} & \multirow{2}{*}{Different}
\\
& E = No Data \tab Freq =No Data   &    &  \\ \cline{2-2}
& 4NFAACj   & \multicolumn{2}{c}
 {RMS = 0.177156}
\\
& E = No Data \tab Freq =No Data   &    \multicolumn{2}{c} 
{ }
\\ \hline
\multirow{4}{*}{\tab[2mm] 200 \tab[2mm]} & 3NFAACj &
\multirow{2}{*}{Different} & \multirow{2}{*}{Different}
\\
& E = No Data \tab Freq =No Data   &    &  \\ \cline{2-2}
& 4NFAACl-3   & \multicolumn{2}{c}
 {RMS = 0.131803}
\\
& E = No Data \tab Freq =No Data   &    \multicolumn{2}{c} 
{ }
\\ \hline
\multirow{4}{*}{\tab[2mm] 201 \tab[2mm]} & 3NFAACk &
\multirow{2}{*}{Different} & \multirow{2}{*}{Different}
\\
& E = No Data \tab Freq =No Data   &    &  \\ \cline{2-2}
& 3NFAACl   & \multicolumn{2}{c}
 {RMS = 0.047641}
\\
& E = No Data \tab Freq =No Data   &    \multicolumn{2}{c} 
{ }
\\ \hline
\multirow{4}{*}{\tab[2mm] 202 \tab[2mm]} & 3NFAACk &
\multirow{2}{*}{Different} & \multirow{2}{*}{Different}
\\
& E = No Data \tab Freq =No Data   &    &  \\ \cline{2-2}
& 3NFAACm   & \multicolumn{2}{c}
 {RMS = 0.164244}
\\
& E = No Data \tab Freq =No Data   &    \multicolumn{2}{c} 
{ }
\\ \hline
\multirow{4}{*}{\tab[2mm] 203 \tab[2mm]} & 3NFAACk &
\multirow{2}{*}{Different} & \multirow{2}{*}{Different}
\\
& E = No Data \tab Freq =No Data   &    &  \\ \cline{2-2}
& 3NFAACn   & \multicolumn{2}{c}
 {RMS = 0.168594}
\\
& E = No Data \tab Freq =No Data   &    \multicolumn{2}{c} 
{ }
\\ \hline
\end{tabular}
\newpage

\vtab[-2cm]
\tab[-2cm]
\begin{tabular}{c|m{8cm}|c|c}
\# & Moléculas & Restultado esperado & Resultado programa \\ \hline\hline
\multirow{4}{*}{\tab[2mm] 204 \tab[2mm]} & 3NFAACk &
\multirow{2}{*}{Different} & \multirow{2}{*}{Different}
\\
& E = No Data \tab Freq =No Data   &    &  \\ \cline{2-2}
& 4NFAACa   & \multicolumn{2}{c}
 {RMS = 0.208366}
\\
& E = No Data \tab Freq =No Data   &    \multicolumn{2}{c} 
{ }
\\ \hline
\multirow{4}{*}{\tab[2mm] 205 \tab[2mm]} & 3NFAACk &
\multirow{2}{*}{Different} & \multirow{2}{*}{Different}
\\
& E = No Data \tab Freq =No Data   &    &  \\ \cline{2-2}
& 4NFAACb   & \multicolumn{2}{c}
 {RMS = 0.269854}
\\
& E = No Data \tab Freq =No Data   &    \multicolumn{2}{c} 
{ }
\\ \hline
\multirow{4}{*}{\tab[2mm] 206 \tab[2mm]} & 3NFAACk &
\multirow{2}{*}{Different} & \multirow{2}{*}{Different}
\\
& E = No Data \tab Freq =No Data   &    &  \\ \cline{2-2}
& 4NFAACc   & \multicolumn{2}{c}
 {RMS = 0.0300371}
\\
& E = No Data \tab Freq =No Data   &    \multicolumn{2}{c} 
{ }
\\ \hline
\multirow{4}{*}{\tab[2mm] 207 \tab[2mm]} & 3NFAACk &
\multirow{2}{*}{Different} & \multirow{2}{*}{Different}
\\
& E = No Data \tab Freq =No Data   &    &  \\ \cline{2-2}
& 4NFAACd   & \multicolumn{2}{c}
 {RMS = 0.0263194}
\\
& E = No Data \tab Freq =No Data   &    \multicolumn{2}{c} 
{ }
\\ \hline
\multirow{4}{*}{\tab[2mm] 208 \tab[2mm]} & 3NFAACk &
\multirow{2}{*}{Different} & \multirow{2}{*}{Different}
\\
& E = No Data \tab Freq =No Data   &    &  \\ \cline{2-2}
& 4NFAACe   & \multicolumn{2}{c}
 {RMS = 0.0342554}
\\
& E = No Data \tab Freq =No Data   &    \multicolumn{2}{c} 
{ }
\\ \hline
\multirow{4}{*}{\tab[2mm] 209 \tab[2mm]} & 3NFAACk &
\multirow{2}{*}{Different} & \multirow{2}{*}{Different}
\\
& E = No Data \tab Freq =No Data   &    &  \\ \cline{2-2}
& 4NFAACf   & \multicolumn{2}{c}
 {RMS = 0.129714}
\\
& E = No Data \tab Freq =No Data   &    \multicolumn{2}{c} 
{ }
\\ \hline
\multirow{4}{*}{\tab[2mm] 210 \tab[2mm]} & 3NFAACk &
\multirow{2}{*}{Different} & \multirow{2}{*}{Different}
\\
& E = No Data \tab Freq =No Data   &    &  \\ \cline{2-2}
& 4NFAACg   & \multicolumn{2}{c}
 {RMS = 0.0452701}
\\
& E = No Data \tab Freq =No Data   &    \multicolumn{2}{c} 
{ }
\\ \hline
\end{tabular}
\newpage

\vtab[-2cm]
\tab[-2cm]
\begin{tabular}{c|m{8cm}|c|c}
\# & Moléculas & Restultado esperado & Resultado programa \\ \hline\hline
\multirow{4}{*}{\tab[2mm] 211 \tab[2mm]} & 3NFAACk &
\multirow{2}{*}{Different} & \multirow{2}{*}{Different}
\\
& E = No Data \tab Freq =No Data   &    &  \\ \cline{2-2}
& 4NFAACi   & \multicolumn{2}{c}
 {RMS = 0.00994921}
\\
& E = No Data \tab Freq =No Data   &    \multicolumn{2}{c} 
{ }
\\ \hline
\multirow{4}{*}{\tab[2mm] 212 \tab[2mm]} & 3NFAACk &
\multirow{2}{*}{Different} & \multirow{2}{*}{Different}
\\
& E = No Data \tab Freq =No Data   &    &  \\ \cline{2-2}
& 4NFAACj   & \multicolumn{2}{c}
 {RMS = 0.0982472}
\\
& E = No Data \tab Freq =No Data   &    \multicolumn{2}{c} 
{ }
\\ \hline
\multirow{4}{*}{\tab[2mm] 213 \tab[2mm]} & 3NFAACk &
\multirow{2}{*}{Different} & \multirow{2}{*}{Different}
\\
& E = No Data \tab Freq =No Data   &    &  \\ \cline{2-2}
& 4NFAACl-3   & \multicolumn{2}{c}
 {RMS = 0.0528942}
\\
& E = No Data \tab Freq =No Data   &    \multicolumn{2}{c} 
{ }
\\ \hline
\multirow{4}{*}{\tab[2mm] 214 \tab[2mm]} & 3NFAACl &
\multirow{2}{*}{Enantiomers} & \multirow{2}{*}{Enantiomers}
\\
& E = No Data \tab Freq =No Data   &    &  \\ \cline{2-2}
& 3NFAACm   & \multicolumn{2}{c}
 {RMS = 0.02978757}
\\
& E = No Data \tab Freq =No Data   &    \multicolumn{2}{c} 
{ }
\\ \hline
\multirow{4}{*}{\tab[2mm] 215 \tab[2mm]} & 3NFAACl &
\multirow{2}{*}{Different} & \multirow{2}{*}{Different}
\\
& E = No Data \tab Freq =No Data   &    &  \\ \cline{2-2}
& 3NFAACn   & \multicolumn{2}{c}
 {RMS = 0.120953}
\\
& E = No Data \tab Freq =No Data   &    \multicolumn{2}{c} 
{ }
\\ \hline
\multirow{4}{*}{\tab[2mm] 216 \tab[2mm]} & 3NFAACl &
\multirow{2}{*}{Different} & \multirow{2}{*}{Different}
\\
& E = No Data \tab Freq =No Data   &    &  \\ \cline{2-2}
& 4NFAACa   & \multicolumn{2}{c}
 {RMS = 0.256007}
\\
& E = No Data \tab Freq =No Data   &    \multicolumn{2}{c} 
{ }
\\ \hline
\multirow{4}{*}{\tab[2mm] 217 \tab[2mm]} & 3NFAACl &
\multirow{2}{*}{Different} & \multirow{2}{*}{Different}
\\
& E = No Data \tab Freq =No Data   &    &  \\ \cline{2-2}
& 4NFAACb   & \multicolumn{2}{c}
 {RMS = 0.222213}
\\
& E = No Data \tab Freq =No Data   &    \multicolumn{2}{c} 
{ }
\\ \hline
\end{tabular}
\newpage

\vtab[-2cm]
\tab[-2cm]
\begin{tabular}{c|m{8cm}|c|c}
\# & Moléculas & Restultado esperado & Resultado programa \\ \hline\hline
\multirow{4}{*}{\tab[2mm] 218 \tab[2mm]} & 3NFAACl &
\multirow{2}{*}{Different} & \multirow{2}{*}{Different}
\\
& E = No Data \tab Freq =No Data   &    &  \\ \cline{2-2}
& 4NFAACc   & \multicolumn{2}{c}
 {RMS = 0.0176039}
\\
& E = No Data \tab Freq =No Data   &    \multicolumn{2}{c} 
{ }
\\ \hline
\multirow{4}{*}{\tab[2mm] 219 \tab[2mm]} & 3NFAACl &
\multirow{2}{*}{Different} & \multirow{2}{*}{Different}
\\
& E = No Data \tab Freq =No Data   &    &  \\ \cline{2-2}
& 4NFAACd   & \multicolumn{2}{c}
 {RMS = 0.0213216}
\\
& E = No Data \tab Freq =No Data   &    \multicolumn{2}{c} 
{ }
\\ \hline
\multirow{4}{*}{\tab[2mm] 220 \tab[2mm]} & 3NFAACl &
\multirow{2}{*}{Different} & \multirow{2}{*}{Different}
\\
& E = No Data \tab Freq =No Data   &    &  \\ \cline{2-2}
& 4NFAACe   & \multicolumn{2}{c}
 {RMS = 0.0133856}
\\
& E = No Data \tab Freq =No Data   &    \multicolumn{2}{c} 
{ }
\\ \hline
\multirow{4}{*}{\tab[2mm] 221 \tab[2mm]} & 3NFAACl &
\multirow{2}{*}{Different} & \multirow{2}{*}{Different}
\\
& E = No Data \tab Freq =No Data   &    &  \\ \cline{2-2}
& 4NFAACf   & \multicolumn{2}{c}
 {RMS = 0.0820729}
\\
& E = No Data \tab Freq =No Data   &    \multicolumn{2}{c} 
{ }
\\ \hline
\multirow{4}{*}{\tab[2mm] 222 \tab[2mm]} & 3NFAACl &
\multirow{2}{*}{Different} & \multirow{2}{*}{Different}
\\
& E = No Data \tab Freq =No Data   &    &  \\ \cline{2-2}
& 4NFAACg   & \multicolumn{2}{c}
 {RMS = 0.00237094}
\\
& E = No Data \tab Freq =No Data   &    \multicolumn{2}{c} 
{ }
\\ \hline
\multirow{4}{*}{\tab[2mm] 223 \tab[2mm]} & 3NFAACl &
\multirow{2}{*}{Different} & \multirow{2}{*}{Different}
\\
& E = No Data \tab Freq =No Data   &    &  \\ \cline{2-2}
& 4NFAACi   & \multicolumn{2}{c}
 {RMS = 0.0376918}
\\
& E = No Data \tab Freq =No Data   &    \multicolumn{2}{c} 
{ }
\\ \hline
\multirow{4}{*}{\tab[2mm] 224 \tab[2mm]} & 3NFAACl &
\multirow{2}{*}{Different} & \multirow{2}{*}{Different}
\\
& E = No Data \tab Freq =No Data   &    &  \\ \cline{2-2}
& 4NFAACj   & \multicolumn{2}{c}
 {RMS = 0.0506062}
\\
& E = No Data \tab Freq =No Data   &    \multicolumn{2}{c} 
{ }
\\ \hline
\end{tabular}
\newpage

\vtab[-2cm]
\tab[-2cm]
\begin{tabular}{c|m{8cm}|c|c}
\# & Moléculas & Restultado esperado & Resultado programa \\ \hline\hline
\multirow{4}{*}{\tab[2mm] 225 \tab[2mm]} & 3NFAACl &
\multirow{2}{*}{Different} & \multirow{2}{*}{Different}
\\
& E = No Data \tab Freq =No Data   &    &  \\ \cline{2-2}
& 4NFAACl-3   & \multicolumn{2}{c}
 {RMS = 0.0052532}
\\
& E = No Data \tab Freq =No Data   &    \multicolumn{2}{c} 
{ }
\\ \hline
\multirow{4}{*}{\tab[2mm] 226 \tab[2mm]} & 3NFAACm &
\multirow{2}{*}{Different} & \multirow{2}{*}{Different}
\\
& E = No Data \tab Freq =No Data   &    &  \\ \cline{2-2}
& 3NFAACn   & \multicolumn{2}{c}
 {RMS = 0.00434937}
\\
& E = No Data \tab Freq =No Data   &    \multicolumn{2}{c} 
{ }
\\ \hline
\multirow{4}{*}{\tab[2mm] 227 \tab[2mm]} & 3NFAACm &
\multirow{2}{*}{Different} & \multirow{2}{*}{Different}
\\
& E = No Data \tab Freq =No Data   &    &  \\ \cline{2-2}
& 4NFAACa   & \multicolumn{2}{c}
 {RMS = 0.372611}
\\
& E = No Data \tab Freq =No Data   &    \multicolumn{2}{c} 
{ }
\\ \hline
\multirow{4}{*}{\tab[2mm] 228 \tab[2mm]} & 3NFAACm &
\multirow{2}{*}{Different} & \multirow{2}{*}{Different}
\\
& E = No Data \tab Freq =No Data   &    &  \\ \cline{2-2}
& 4NFAACb   & \multicolumn{2}{c}
 {RMS = 0.105609}
\\
& E = No Data \tab Freq =No Data   &    \multicolumn{2}{c} 
{ }
\\ \hline
\multirow{4}{*}{\tab[2mm] 229 \tab[2mm]} & 3NFAACm &
\multirow{2}{*}{Different} & \multirow{2}{*}{Different}
\\
& E = No Data \tab Freq =No Data   &    &  \\ \cline{2-2}
& 4NFAACc   & \multicolumn{2}{c}
 {RMS = 0.134207}
\\
& E = No Data \tab Freq =No Data   &    \multicolumn{2}{c} 
{ }
\\ \hline
\multirow{4}{*}{\tab[2mm] 230 \tab[2mm]} & 3NFAACm &
\multirow{2}{*}{Different} & \multirow{2}{*}{Different}
\\
& E = No Data \tab Freq =No Data   &    &  \\ \cline{2-2}
& 4NFAACd   & \multicolumn{2}{c}
 {RMS = 0.137925}
\\
& E = No Data \tab Freq =No Data   &    \multicolumn{2}{c} 
{ }
\\ \hline
\multirow{4}{*}{\tab[2mm] 231 \tab[2mm]} & 3NFAACm &
\multirow{2}{*}{Different} & \multirow{2}{*}{Different}
\\
& E = No Data \tab Freq =No Data   &    &  \\ \cline{2-2}
& 4NFAACe   & \multicolumn{2}{c}
 {RMS = 0.129989}
\\
& E = No Data \tab Freq =No Data   &    \multicolumn{2}{c} 
{ }
\\ \hline
\end{tabular}
\newpage

\vtab[-2cm]
\tab[-2cm]
\begin{tabular}{c|m{8cm}|c|c}
\# & Moléculas & Restultado esperado & Resultado programa \\ \hline\hline
\multirow{4}{*}{\tab[2mm] 232 \tab[2mm]} & 3NFAACm &
\multirow{2}{*}{Different} & \multirow{2}{*}{Different}
\\
& E = No Data \tab Freq =No Data   &    &  \\ \cline{2-2}
& 4NFAACf   & \multicolumn{2}{c}
 {RMS = 0.0345305}
\\
& E = No Data \tab Freq =No Data   &    \multicolumn{2}{c} 
{ }
\\ \hline
\multirow{4}{*}{\tab[2mm] 233 \tab[2mm]} & 3NFAACm &
\multirow{2}{*}{Different} & \multirow{2}{*}{Different}
\\
& E = No Data \tab Freq =No Data   &    &  \\ \cline{2-2}
& 4NFAACg   & \multicolumn{2}{c}
 {RMS = 0.118974}
\\
& E = No Data \tab Freq =No Data   &    \multicolumn{2}{c} 
{ }
\\ \hline
\multirow{4}{*}{\tab[2mm] 234 \tab[2mm]} & 3NFAACm &
\multirow{2}{*}{Different} & \multirow{2}{*}{Different}
\\
& E = No Data \tab Freq =No Data   &    &  \\ \cline{2-2}
& 4NFAACi   & \multicolumn{2}{c}
 {RMS = 0.154295}
\\
& E = No Data \tab Freq =No Data   &    \multicolumn{2}{c} 
{ }
\\ \hline
\multirow{4}{*}{\tab[2mm] 235 \tab[2mm]} & 3NFAACm &
\multirow{2}{*}{Different} & \multirow{2}{*}{Different}
\\
& E = No Data \tab Freq =No Data   &    &  \\ \cline{2-2}
& 4NFAACj   & \multicolumn{2}{c}
 {RMS = 0.0659973}
\\
& E = No Data \tab Freq =No Data   &    \multicolumn{2}{c} 
{ }
\\ \hline
\multirow{4}{*}{\tab[2mm] 236 \tab[2mm]} & 3NFAACm &
\multirow{2}{*}{Different} & \multirow{2}{*}{Different}
\\
& E = No Data \tab Freq =No Data   &    &  \\ \cline{2-2}
& 4NFAACl-3   & \multicolumn{2}{c}
 {RMS = 0.11135}
\\
& E = No Data \tab Freq =No Data   &    \multicolumn{2}{c} 
{ }
\\ \hline
\multirow{4}{*}{\tab[2mm] 237 \tab[2mm]} & 3NFAACn &
\multirow{2}{*}{Different} & \multirow{2}{*}{Different}
\\
& E = No Data \tab Freq =No Data   &    &  \\ \cline{2-2}
& 4NFAACa   & \multicolumn{2}{c}
 {RMS = 0.37696}
\\
& E = No Data \tab Freq =No Data   &    \multicolumn{2}{c} 
{ }
\\ \hline
\multirow{4}{*}{\tab[2mm] 238 \tab[2mm]} & 3NFAACn &
\multirow{2}{*}{Different} & \multirow{2}{*}{Different}
\\
& E = No Data \tab Freq =No Data   &    &  \\ \cline{2-2}
& 4NFAACb   & \multicolumn{2}{c}
 {RMS = 0.10126}
\\
& E = No Data \tab Freq =No Data   &    \multicolumn{2}{c} 
{ }
\\ \hline
\end{tabular}
\newpage

\vtab[-2cm]
\tab[-2cm]
\begin{tabular}{c|m{8cm}|c|c}
\# & Moléculas & Restultado esperado & Resultado programa \\ \hline\hline
\multirow{4}{*}{\tab[2mm] 239 \tab[2mm]} & 3NFAACn &
\multirow{2}{*}{Different} & \multirow{2}{*}{Different}
\\
& E = No Data \tab Freq =No Data   &    &  \\ \cline{2-2}
& 4NFAACc   & \multicolumn{2}{c}
 {RMS = 0.138557}
\\
& E = No Data \tab Freq =No Data   &    \multicolumn{2}{c} 
{ }
\\ \hline
\multirow{4}{*}{\tab[2mm] 240 \tab[2mm]} & 3NFAACn &
\multirow{2}{*}{Different} & \multirow{2}{*}{Different}
\\
& E = No Data \tab Freq =No Data   &    &  \\ \cline{2-2}
& 4NFAACd   & \multicolumn{2}{c}
 {RMS = 0.142274}
\\
& E = No Data \tab Freq =No Data   &    \multicolumn{2}{c} 
{ }
\\ \hline
\multirow{4}{*}{\tab[2mm] 241 \tab[2mm]} & 3NFAACn &
\multirow{2}{*}{Different} & \multirow{2}{*}{Different}
\\
& E = No Data \tab Freq =No Data   &    &  \\ \cline{2-2}
& 4NFAACe   & \multicolumn{2}{c}
 {RMS = 0.134338}
\\
& E = No Data \tab Freq =No Data   &    \multicolumn{2}{c} 
{ }
\\ \hline
\multirow{4}{*}{\tab[2mm] 242 \tab[2mm]} & 3NFAACn &
\multirow{2}{*}{Different} & \multirow{2}{*}{Different}
\\
& E = No Data \tab Freq =No Data   &    &  \\ \cline{2-2}
& 4NFAACf   & \multicolumn{2}{c}
 {RMS = 0.0388799}
\\
& E = No Data \tab Freq =No Data   &    \multicolumn{2}{c} 
{ }
\\ \hline
\multirow{4}{*}{\tab[2mm] 243 \tab[2mm]} & 3NFAACn &
\multirow{2}{*}{Different} & \multirow{2}{*}{Different}
\\
& E = No Data \tab Freq =No Data   &    &  \\ \cline{2-2}
& 4NFAACg   & \multicolumn{2}{c}
 {RMS = 0.123324}
\\
& E = No Data \tab Freq =No Data   &    \multicolumn{2}{c} 
{ }
\\ \hline
\multirow{4}{*}{\tab[2mm] 244 \tab[2mm]} & 3NFAACn &
\multirow{2}{*}{Different} & \multirow{2}{*}{Different}
\\
& E = No Data \tab Freq =No Data   &    &  \\ \cline{2-2}
& 4NFAACi   & \multicolumn{2}{c}
 {RMS = 0.158645}
\\
& E = No Data \tab Freq =No Data   &    \multicolumn{2}{c} 
{ }
\\ \hline
\multirow{4}{*}{\tab[2mm] 245 \tab[2mm]} & 3NFAACn &
\multirow{2}{*}{Different} & \multirow{2}{*}{Different}
\\
& E = No Data \tab Freq =No Data   &    &  \\ \cline{2-2}
& 4NFAACj   & \multicolumn{2}{c}
 {RMS = 0.0703467}
\\
& E = No Data \tab Freq =No Data   &    \multicolumn{2}{c} 
{ }
\\ \hline
\end{tabular}
\newpage

\vtab[-2cm]
\tab[-2cm]
\begin{tabular}{c|m{8cm}|c|c}
\# & Moléculas & Restultado esperado & Resultado programa \\ \hline\hline
\multirow{4}{*}{\tab[2mm] 246 \tab[2mm]} & 3NFAACn &
\multirow{2}{*}{Different} & \multirow{2}{*}{Different}
\\
& E = No Data \tab Freq =No Data   &    &  \\ \cline{2-2}
& 4NFAACl-3   & \multicolumn{2}{c}
 {RMS = 0.1157}
\\
& E = No Data \tab Freq =No Data   &    \multicolumn{2}{c} 
{ }
\\ \hline
\multirow{4}{*}{\tab[2mm] 247 \tab[2mm]} & 4NFAACa &
\multirow{2}{*}{Enantiomers} & \multirow{2}{*}{Enantiomers}
\\
& E = No Data \tab Freq =No Data   &    &  \\ \cline{2-2}
& 4NFAACb   & \multicolumn{2}{c}
 {RMS = 0.485443}
\\
& E = No Data \tab Freq =No Data   &    \multicolumn{2}{c} 
{ }
\\ \hline
\multirow{4}{*}{\tab[2mm] 248 \tab[2mm]} & 4NFAACa &
\multirow{2}{*}{Different} & \multirow{2}{*}{Different}
\\
& E = No Data \tab Freq =No Data   &    &  \\ \cline{2-2}
& 4NFAACc   & \multicolumn{2}{c}
 {RMS = 0.238403}
\\
& E = No Data \tab Freq =No Data   &    \multicolumn{2}{c} 
{ }
\\ \hline
\multirow{4}{*}{\tab[2mm] 249 \tab[2mm]} & 4NFAACa &
\multirow{2}{*}{Different} & \multirow{2}{*}{Different}
\\
& E = No Data \tab Freq =No Data   &    &  \\ \cline{2-2}
& 4NFAACd   & \multicolumn{2}{c}
 {RMS = 0.234686}
\\
& E = No Data \tab Freq =No Data   &    \multicolumn{2}{c} 
{ }
\\ \hline
\multirow{4}{*}{\tab[2mm] 250 \tab[2mm]} & 4NFAACa &
\multirow{2}{*}{Different} & \multirow{2}{*}{Different}
\\
& E = No Data \tab Freq =No Data   &    &  \\ \cline{2-2}
& 4NFAACe   & \multicolumn{2}{c}
 {RMS = 0.242622}
\\
& E = No Data \tab Freq =No Data   &    \multicolumn{2}{c} 
{ }
\\ \hline
\multirow{4}{*}{\tab[2mm] 251 \tab[2mm]} & 4NFAACa &
\multirow{2}{*}{Different} & \multirow{2}{*}{Different}
\\
& E = No Data \tab Freq =No Data   &    &  \\ \cline{2-2}
& 4NFAACf   & \multicolumn{2}{c}
 {RMS = 0.33808}
\\
& E = No Data \tab Freq =No Data   &    \multicolumn{2}{c} 
{ }
\\ \hline
\multirow{4}{*}{\tab[2mm] 252 \tab[2mm]} & 4NFAACa &
\multirow{2}{*}{Different} & \multirow{2}{*}{Different}
\\
& E = No Data \tab Freq =No Data   &    &  \\ \cline{2-2}
& 4NFAACg   & \multicolumn{2}{c}
 {RMS = 0.253636}
\\
& E = No Data \tab Freq =No Data   &    \multicolumn{2}{c} 
{ }
\\ \hline
\end{tabular}
\newpage

\vtab[-2cm]
\tab[-2cm]
\begin{tabular}{c|m{8cm}|c|c}
\# & Moléculas & Restultado esperado & Resultado programa \\ \hline\hline
\multirow{4}{*}{\tab[2mm] 253 \tab[2mm]} & 4NFAACa &
\multirow{2}{*}{Different} & \multirow{2}{*}{Different}
\\
& E = No Data \tab Freq =No Data   &    &  \\ \cline{2-2}
& 4NFAACi   & \multicolumn{2}{c}
 {RMS = 0.218315}
\\
& E = No Data \tab Freq =No Data   &    \multicolumn{2}{c} 
{ }
\\ \hline
\multirow{4}{*}{\tab[2mm] 254 \tab[2mm]} & 4NFAACa &
\multirow{2}{*}{Different} & \multirow{2}{*}{Different}
\\
& E = No Data \tab Freq =No Data   &    &  \\ \cline{2-2}
& 4NFAACj   & \multicolumn{2}{c}
 {RMS = 0.306613}
\\
& E = No Data \tab Freq =No Data   &    \multicolumn{2}{c} 
{ }
\\ \hline
\multirow{4}{*}{\tab[2mm] 255 \tab[2mm]} & 4NFAACa &
\multirow{2}{*}{Different} & \multirow{2}{*}{Different}
\\
& E = No Data \tab Freq =No Data   &    &  \\ \cline{2-2}
& 4NFAACl-3   & \multicolumn{2}{c}
 {RMS = 0.26126}
\\
& E = No Data \tab Freq =No Data   &    \multicolumn{2}{c} 
{ }
\\ \hline
\multirow{4}{*}{\tab[2mm] 256 \tab[2mm]} & 4NFAACb &
\multirow{2}{*}{Different} & \multirow{2}{*}{Different}
\\
& E = No Data \tab Freq =No Data   &    &  \\ \cline{2-2}
& 4NFAACc   & \multicolumn{2}{c}
 {RMS = 0.239816}
\\
& E = No Data \tab Freq =No Data   &    \multicolumn{2}{c} 
{ }
\\ \hline
\multirow{4}{*}{\tab[2mm] 257 \tab[2mm]} & 4NFAACb &
\multirow{2}{*}{Different} & \multirow{2}{*}{Different}
\\
& E = No Data \tab Freq =No Data   &    &  \\ \cline{2-2}
& 4NFAACd   & \multicolumn{2}{c}
 {RMS = 0.243534}
\\
& E = No Data \tab Freq =No Data   &    \multicolumn{2}{c} 
{ }
\\ \hline
\multirow{4}{*}{\tab[2mm] 258 \tab[2mm]} & 4NFAACb &
\multirow{2}{*}{Different} & \multirow{2}{*}{Different}
\\
& E = No Data \tab Freq =No Data   &    &  \\ \cline{2-2}
& 4NFAACe   & \multicolumn{2}{c}
 {RMS = 0.235598}
\\
& E = No Data \tab Freq =No Data   &    \multicolumn{2}{c} 
{ }
\\ \hline
\multirow{4}{*}{\tab[2mm] 259 \tab[2mm]} & 4NFAACb &
\multirow{2}{*}{Different} & \multirow{2}{*}{Different}
\\
& E = No Data \tab Freq =No Data   &    &  \\ \cline{2-2}
& 4NFAACf   & \multicolumn{2}{c}
 {RMS = 0.14014}
\\
& E = No Data \tab Freq =No Data   &    \multicolumn{2}{c} 
{ }
\\ \hline
\end{tabular}
\newpage

\vtab[-2cm]
\tab[-2cm]
\begin{tabular}{c|m{8cm}|c|c}
\# & Moléculas & Restultado esperado & Resultado programa \\ \hline\hline
\multirow{4}{*}{\tab[2mm] 260 \tab[2mm]} & 4NFAACb &
\multirow{2}{*}{Different} & \multirow{2}{*}{Different}
\\
& E = No Data \tab Freq =No Data   &    &  \\ \cline{2-2}
& 4NFAACg   & \multicolumn{2}{c}
 {RMS = 0.224583}
\\
& E = No Data \tab Freq =No Data   &    \multicolumn{2}{c} 
{ }
\\ \hline
\multirow{4}{*}{\tab[2mm] 261 \tab[2mm]} & 4NFAACb &
\multirow{2}{*}{Different} & \multirow{2}{*}{Different}
\\
& E = No Data \tab Freq =No Data   &    &  \\ \cline{2-2}
& 4NFAACi   & \multicolumn{2}{c}
 {RMS = 0.259904}
\\
& E = No Data \tab Freq =No Data   &    \multicolumn{2}{c} 
{ }
\\ \hline
\multirow{4}{*}{\tab[2mm] 262 \tab[2mm]} & 4NFAACb &
\multirow{2}{*}{Different} & \multirow{2}{*}{Different}
\\
& E = No Data \tab Freq =No Data   &    &  \\ \cline{2-2}
& 4NFAACj   & \multicolumn{2}{c}
 {RMS = 0.171606}
\\
& E = No Data \tab Freq =No Data   &    \multicolumn{2}{c} 
{ }
\\ \hline
\multirow{4}{*}{\tab[2mm] 263 \tab[2mm]} & 4NFAACb &
\multirow{2}{*}{Different} & \multirow{2}{*}{Different}
\\
& E = No Data \tab Freq =No Data   &    &  \\ \cline{2-2}
& 4NFAACl-3   & \multicolumn{2}{c}
 {RMS = 0.216959}
\\
& E = No Data \tab Freq =No Data   &    \multicolumn{2}{c} 
{ }
\\ \hline
\multirow{4}{*}{\tab[2mm] 264 \tab[2mm]} & 4NFAACc &
\multirow{2}{*}{Different} & \multirow{2}{*}{Different}
\\
& E = No Data \tab Freq =No Data   &    &  \\ \cline{2-2}
& 4NFAACd   & \multicolumn{2}{c}
 {RMS = 0.00371771}
\\
& E = No Data \tab Freq =No Data   &    \multicolumn{2}{c} 
{ }
\\ \hline
\multirow{4}{*}{\tab[2mm] 265 \tab[2mm]} & 4NFAACc &
\multirow{2}{*}{Different} & \multirow{2}{*}{Different}
\\
& E = No Data \tab Freq =No Data   &    &  \\ \cline{2-2}
& 4NFAACe   & \multicolumn{2}{c}
 {RMS = 0.00421832}
\\
& E = No Data \tab Freq =No Data   &    \multicolumn{2}{c} 
{ }
\\ \hline
\multirow{4}{*}{\tab[2mm] 266 \tab[2mm]} & 4NFAACc &
\multirow{2}{*}{Different} & \multirow{2}{*}{Different}
\\
& E = No Data \tab Freq =No Data   &    &  \\ \cline{2-2}
& 4NFAACf   & \multicolumn{2}{c}
 {RMS = 0.0996768}
\\
& E = No Data \tab Freq =No Data   &    \multicolumn{2}{c} 
{ }
\\ \hline
\end{tabular}
\newpage

\vtab[-2cm]
\tab[-2cm]
\begin{tabular}{c|m{8cm}|c|c}
\# & Moléculas & Restultado esperado & Resultado programa \\ \hline\hline
\multirow{4}{*}{\tab[2mm] 267 \tab[2mm]} & 4NFAACc &
\multirow{2}{*}{Different} & \multirow{2}{*}{Different}
\\
& E = No Data \tab Freq =No Data   &    &  \\ \cline{2-2}
& 4NFAACg   & \multicolumn{2}{c}
 {RMS = 0.015233}
\\
& E = No Data \tab Freq =No Data   &    \multicolumn{2}{c} 
{ }
\\ \hline
\multirow{4}{*}{\tab[2mm] 268 \tab[2mm]} & 4NFAACc &
\multirow{2}{*}{Different} & \multirow{2}{*}{Different}
\\
& E = No Data \tab Freq =No Data   &    &  \\ \cline{2-2}
& 4NFAACi   & \multicolumn{2}{c}
 {RMS = 0.0200879}
\\
& E = No Data \tab Freq =No Data   &    \multicolumn{2}{c} 
{ }
\\ \hline
\multirow{4}{*}{\tab[2mm] 269 \tab[2mm]} & 4NFAACc &
\multirow{2}{*}{Different} & \multirow{2}{*}{Different}
\\
& E = No Data \tab Freq =No Data   &    &  \\ \cline{2-2}
& 4NFAACj   & \multicolumn{2}{c}
 {RMS = 0.0682101}
\\
& E = No Data \tab Freq =No Data   &    \multicolumn{2}{c} 
{ }
\\ \hline
\multirow{4}{*}{\tab[2mm] 270 \tab[2mm]} & 4NFAACc &
\multirow{2}{*}{Different} & \multirow{2}{*}{Different}
\\
& E = No Data \tab Freq =No Data   &    &  \\ \cline{2-2}
& 4NFAACl-3   & \multicolumn{2}{c}
 {RMS = 0.0228571}
\\
& E = No Data \tab Freq =No Data   &    \multicolumn{2}{c} 
{ }
\\ \hline
\multirow{4}{*}{\tab[2mm] 271 \tab[2mm]} & 4NFAACd &
\multirow{2}{*}{Different} & \multirow{2}{*}{Different}
\\
& E = No Data \tab Freq =No Data   &    &  \\ \cline{2-2}
& 4NFAACe   & \multicolumn{2}{c}
 {RMS = 0.00793602}
\\
& E = No Data \tab Freq =No Data   &    \multicolumn{2}{c} 
{ }
\\ \hline
\multirow{4}{*}{\tab[2mm] 272 \tab[2mm]} & 4NFAACd &
\multirow{2}{*}{Different} & \multirow{2}{*}{Different}
\\
& E = No Data \tab Freq =No Data   &    &  \\ \cline{2-2}
& 4NFAACf   & \multicolumn{2}{c}
 {RMS = 0.103395}
\\
& E = No Data \tab Freq =No Data   &    \multicolumn{2}{c} 
{ }
\\ \hline
\multirow{4}{*}{\tab[2mm] 273 \tab[2mm]} & 4NFAACd &
\multirow{2}{*}{Different} & \multirow{2}{*}{Different}
\\
& E = No Data \tab Freq =No Data   &    &  \\ \cline{2-2}
& 4NFAACg   & \multicolumn{2}{c}
 {RMS = 0.0189507}
\\
& E = No Data \tab Freq =No Data   &    \multicolumn{2}{c} 
{ }
\\ \hline
\end{tabular}
\newpage

\vtab[-2cm]
\tab[-2cm]
\begin{tabular}{c|m{8cm}|c|c}
\# & Moléculas & Restultado esperado & Resultado programa \\ \hline\hline
\multirow{4}{*}{\tab[2mm] 274 \tab[2mm]} & 4NFAACd &
\multirow{2}{*}{Different} & \multirow{2}{*}{Different}
\\
& E = No Data \tab Freq =No Data   &    &  \\ \cline{2-2}
& 4NFAACi   & \multicolumn{2}{c}
 {RMS = 0.0163702}
\\
& E = No Data \tab Freq =No Data   &    \multicolumn{2}{c} 
{ }
\\ \hline
\multirow{4}{*}{\tab[2mm] 275 \tab[2mm]} & 4NFAACd &
\multirow{2}{*}{Different} & \multirow{2}{*}{Different}
\\
& E = No Data \tab Freq =No Data   &    &  \\ \cline{2-2}
& 4NFAACj   & \multicolumn{2}{c}
 {RMS = 0.0719278}
\\
& E = No Data \tab Freq =No Data   &    \multicolumn{2}{c} 
{ }
\\ \hline
\multirow{4}{*}{\tab[2mm] 276 \tab[2mm]} & 4NFAACd &
\multirow{2}{*}{Different} & \multirow{2}{*}{Different}
\\
& E = No Data \tab Freq =No Data   &    &  \\ \cline{2-2}
& 4NFAACl-3   & \multicolumn{2}{c}
 {RMS = 0.0265748}
\\
& E = No Data \tab Freq =No Data   &    \multicolumn{2}{c} 
{ }
\\ \hline
\multirow{4}{*}{\tab[2mm] 277 \tab[2mm]} & 4NFAACe &
\multirow{2}{*}{Different} & \multirow{2}{*}{Different}
\\
& E = No Data \tab Freq =No Data   &    &  \\ \cline{2-2}
& 4NFAACf   & \multicolumn{2}{c}
 {RMS = 0.0954585}
\\
& E = No Data \tab Freq =No Data   &    \multicolumn{2}{c} 
{ }
\\ \hline
\multirow{4}{*}{\tab[2mm] 278 \tab[2mm]} & 4NFAACe &
\multirow{2}{*}{Different} & \multirow{2}{*}{Different}
\\
& E = No Data \tab Freq =No Data   &    &  \\ \cline{2-2}
& 4NFAACg   & \multicolumn{2}{c}
 {RMS = 0.0110146}
\\
& E = No Data \tab Freq =No Data   &    \multicolumn{2}{c} 
{ }
\\ \hline
\multirow{4}{*}{\tab[2mm] 279 \tab[2mm]} & 4NFAACe &
\multirow{2}{*}{Different} & \multirow{2}{*}{Different}
\\
& E = No Data \tab Freq =No Data   &    &  \\ \cline{2-2}
& 4NFAACi   & \multicolumn{2}{c}
 {RMS = 0.0243062}
\\
& E = No Data \tab Freq =No Data   &    \multicolumn{2}{c} 
{ }
\\ \hline
\multirow{4}{*}{\tab[2mm] 280 \tab[2mm]} & 4NFAACe &
\multirow{2}{*}{Different} & \multirow{2}{*}{Different}
\\
& E = No Data \tab Freq =No Data   &    &  \\ \cline{2-2}
& 4NFAACj   & \multicolumn{2}{c}
 {RMS = 0.0639918}
\\
& E = No Data \tab Freq =No Data   &    \multicolumn{2}{c} 
{ }
\\ \hline
\end{tabular}
\newpage

\vtab[-2cm]
\tab[-2cm]
\begin{tabular}{c|m{8cm}|c|c}
\# & Moléculas & Restultado esperado & Resultado programa \\ \hline\hline
\multirow{4}{*}{\tab[2mm] 281 \tab[2mm]} & 4NFAACe &
\multirow{2}{*}{Different} & \multirow{2}{*}{Different}
\\
& E = No Data \tab Freq =No Data   &    &  \\ \cline{2-2}
& 4NFAACl-3   & \multicolumn{2}{c}
 {RMS = 0.0186388}
\\
& E = No Data \tab Freq =No Data   &    \multicolumn{2}{c} 
{ }
\\ \hline
\multirow{4}{*}{\tab[2mm] 282 \tab[2mm]} & 4NFAACf &
\multirow{2}{*}{Different} & \multirow{2}{*}{Different}
\\
& E = No Data \tab Freq =No Data   &    &  \\ \cline{2-2}
& 4NFAACg   & \multicolumn{2}{c}
 {RMS = 0.0844439}
\\
& E = No Data \tab Freq =No Data   &    \multicolumn{2}{c} 
{ }
\\ \hline
\multirow{4}{*}{\tab[2mm] 283 \tab[2mm]} & 4NFAACf &
\multirow{2}{*}{Different} & \multirow{2}{*}{Different}
\\
& E = No Data \tab Freq =No Data   &    &  \\ \cline{2-2}
& 4NFAACi   & \multicolumn{2}{c}
 {RMS = 0.119765}
\\
& E = No Data \tab Freq =No Data   &    \multicolumn{2}{c} 
{ }
\\ \hline
\multirow{4}{*}{\tab[2mm] 284 \tab[2mm]} & 4NFAACf &
\multirow{2}{*}{Different} & \multirow{2}{*}{Different}
\\
& E = No Data \tab Freq =No Data   &    &  \\ \cline{2-2}
& 4NFAACj   & \multicolumn{2}{c}
 {RMS = 0.0314668}
\\
& E = No Data \tab Freq =No Data   &    \multicolumn{2}{c} 
{ }
\\ \hline
\multirow{4}{*}{\tab[2mm] 285 \tab[2mm]} & 4NFAACf &
\multirow{2}{*}{Different} & \multirow{2}{*}{Different}
\\
& E = No Data \tab Freq =No Data   &    &  \\ \cline{2-2}
& 4NFAACl-3   & \multicolumn{2}{c}
 {RMS = 0.0768197}
\\
& E = No Data \tab Freq =No Data   &    \multicolumn{2}{c} 
{ }
\\ \hline
\multirow{4}{*}{\tab[2mm] 286 \tab[2mm]} & 4NFAACg &
\multirow{2}{*}{Different} & \multirow{2}{*}{Different}
\\
& E = No Data \tab Freq =No Data   &    &  \\ \cline{2-2}
& 4NFAACi   & \multicolumn{2}{c}
 {RMS = 0.0353209}
\\
& E = No Data \tab Freq =No Data   &    \multicolumn{2}{c} 
{ }
\\ \hline
\multirow{4}{*}{\tab[2mm] 287 \tab[2mm]} & 4NFAACg &
\multirow{2}{*}{Different} & \multirow{2}{*}{Different}
\\
& E = No Data \tab Freq =No Data   &    &  \\ \cline{2-2}
& 4NFAACj   & \multicolumn{2}{c}
 {RMS = 0.0529771}
\\
& E = No Data \tab Freq =No Data   &    \multicolumn{2}{c} 
{ }
\\ \hline
\end{tabular}
\newpage

\vtab[-2cm]
\tab[-2cm]
\begin{tabular}{c|m{8cm}|c|c}
\# & Moléculas & Restultado esperado & Resultado programa \\ \hline\hline
\multirow{4}{*}{\tab[2mm] 288 \tab[2mm]} & 4NFAACg &
\multirow{2}{*}{Different} & \multirow{2}{*}{Different}
\\
& E = No Data \tab Freq =No Data   &    &  \\ \cline{2-2}
& 4NFAACl-3   & \multicolumn{2}{c}
 {RMS = 0.00762414}
\\
& E = No Data \tab Freq =No Data   &    \multicolumn{2}{c} 
{ }
\\ \hline
\multirow{4}{*}{\tab[2mm] 289 \tab[2mm]} & 4NFAACi &
\multirow{2}{*}{Different} & \multirow{2}{*}{Different}
\\
& E = No Data \tab Freq =No Data   &    &  \\ \cline{2-2}
& 4NFAACj   & \multicolumn{2}{c}
 {RMS = 0.088298}
\\
& E = No Data \tab Freq =No Data   &    \multicolumn{2}{c} 
{ }
\\ \hline
\multirow{4}{*}{\tab[2mm] 290 \tab[2mm]} & 4NFAACi &
\multirow{2}{*}{Different} & \multirow{2}{*}{Different}
\\
& E = No Data \tab Freq =No Data   &    &  \\ \cline{2-2}
& 4NFAACl-3   & \multicolumn{2}{c}
 {RMS = 0.042945}
\\
& E = No Data \tab Freq =No Data   &    \multicolumn{2}{c} 
{ }
\\ \hline
\multirow{4}{*}{\tab[2mm] 291 \tab[2mm]} & 4NFAACj &
\multirow{2}{*}{Different} & \multirow{2}{*}{Different}
\\
& E = No Data \tab Freq =No Data   &    &  \\ \cline{2-2}
& 4NFAACl-3   & \multicolumn{2}{c}
 {RMS = 0.045353}
\\
& E = No Data \tab Freq =No Data   &    \multicolumn{2}{c} 
{ }
\\ \hline
\multirow{4}{*}{\tab[2mm] 292 \tab[2mm]} & 2,2-dimethyl-butane\_out\_G09\_invertion &
\multirow{2}{*}{Different} & \multirow{2}{*}{Different}
\\
& E = -237.083690318 \tab Freq =92.1843   &    &  \\ \cline{2-2}
& 2,3-dimethyl-butane\_out\_G09   & \multicolumn{2}{c}
 {RMS = 0.269312}
\\
& E = -237.082490952 \tab Freq =58.6018   &    \multicolumn{2}{c} 
{ }
\\ \hline
\multirow{4}{*}{\tab[2mm] 293 \tab[2mm]} & 2,2-dimethyl-butane\_out\_G09\_invertion &
\multirow{2}{*}{Different} & \multirow{2}{*}{Different}
\\
& E = -237.083690318 \tab Freq =92.1843   &    &  \\ \cline{2-2}
& 2,3-dimethyl-butane\_out\_G09\_invertion   & \multicolumn{2}{c}
 {RMS = 0.269312}
\\
& E = -237.082490952 \tab Freq =58.6018   &    \multicolumn{2}{c} 
{ }
\\ \hline
\multirow{4}{*}{\tab[2mm] 294 \tab[2mm]} & 2,2-dimethyl-butane\_out\_G09\_invertion &
\multirow{2}{*}{Different} & \multirow{2}{*}{Different}
\\
& E = -237.083690318 \tab Freq =92.1843   &    &  \\ \cline{2-2}
& 2-methyl-pentane\_out\_G09   & \multicolumn{2}{c}
 {RMS = 0.0426031}
\\
& E = -237.084759065 \tab Freq =80.9103   &    \multicolumn{2}{c} 
{ }
\\ \hline
\end{tabular}
\newpage

\vtab[-2cm]
\tab[-2cm]
\begin{tabular}{c|m{8cm}|c|c}
\# & Moléculas & Restultado esperado & Resultado programa \\ \hline\hline
\multirow{4}{*}{\tab[2mm] 295 \tab[2mm]} & 2,2-dimethyl-butane\_out\_G09\_invertion &
\multirow{2}{*}{Different} & \multirow{2}{*}{Different}
\\
& E = -237.083690318 \tab Freq =92.1843   &    &  \\ \cline{2-2}
& 2-methyl-pentane\_out\_G09\_invertion   & \multicolumn{2}{c}
 {RMS = 0.0778147}
\\
& E = -237.084759065 \tab Freq =80.9103   &    \multicolumn{2}{c} 
{ }
\\ \hline
\multirow{4}{*}{\tab[2mm] 296 \tab[2mm]} & 2,2-dimethyl-butane\_out\_G09\_invertion &
\multirow{2}{*}{Different} & \multirow{2}{*}{Different}
\\
& E = -237.083690318 \tab Freq =92.1843   &    &  \\ \cline{2-2}
& 3-methyl-pentane\_out\_G09   & \multicolumn{2}{c}
 {RMS = 0.766027}
\\
& E = -237.083213438 \tab Freq =81.4725   &    \multicolumn{2}{c} 
{ }
\\ \hline
\multirow{4}{*}{\tab[2mm] 297 \tab[2mm]} & 2,2-dimethyl-butane\_out\_G09\_invertion &
\multirow{2}{*}{Different} & \multirow{2}{*}{Different}
\\
& E = -237.083690318 \tab Freq =92.1843   &    &  \\ \cline{2-2}
& 3-methyl-pentane\_out\_G09\_invertion   & \multicolumn{2}{c}
 {RMS = 0.766028}
\\
& E = -237.083213438 \tab Freq =81.4725   &    \multicolumn{2}{c} 
{ }
\\ \hline
\multirow{4}{*}{\tab[2mm] 298 \tab[2mm]} & 2,2-dimethyl-butane\_out\_G09\_invertion &
\multirow{2}{*}{Different} & \multirow{2}{*}{Different}
\\
& E = -237.083690318 \tab Freq =92.1843   &    &  \\ \cline{2-2}
& hexane\_out\_G09   & \multicolumn{2}{c}
 {RMS = 0}
\\
& E = -237.085500432 \tab Freq =72.0752   &    \multicolumn{2}{c} 
{ }
\\ \hline
\multirow{4}{*}{\tab[2mm] 299 \tab[2mm]} & 2,2-dimethyl-butane\_out\_G09\_invertion &
\multirow{2}{*}{Different} & \multirow{2}{*}{Different}
\\
& E = -237.083690318 \tab Freq =92.1843   &    &  \\ \cline{2-2}
& hexane\_out\_G09\_invertion   & \multicolumn{2}{c}
 {RMS = 0}
\\
& E = -237.085500432 \tab Freq =72.0752   &    \multicolumn{2}{c} 
{ }
\\ \hline
\multirow{4}{*}{\tab[2mm] 300 \tab[2mm]} & 2,2-dimethyl-butane\_out\_G09 &
\multirow{2}{*}{Equal} & \multirow{2}{*}{Equal}
\\
& E = -237.083690318 \tab Freq =92.1843   &    &  \\ \cline{2-2}
& 2,2-dimethyl-butane\_out\_G09\_invertion   & \multicolumn{2}{c}
{ RMS = 0}
\\
& E = -237.083690318 \tab Freq =92.1843   &    \multicolumn{2}{c} 
{ }
\\ \hline
\multirow{4}{*}{\tab[2mm] 301 \tab[2mm]} & 2,2-dimethyl-butane\_out\_G09 &
\multirow{2}{*}{Different} & \multirow{2}{*}{Different}
\\
& E = -237.083690318 \tab Freq =92.1843   &    &  \\ \cline{2-2}
& 2,3-dimethyl-butane\_out\_G09   & \multicolumn{2}{c}
 {RMS = 0.269312}
\\
& E = -237.082490952 \tab Freq =58.6018   &    \multicolumn{2}{c} 
{ }
\\ \hline
\end{tabular}
\newpage

\vtab[-2cm]
\tab[-2cm]
\begin{tabular}{c|m{8cm}|c|c}
\# & Moléculas & Restultado esperado & Resultado programa \\ \hline\hline
\multirow{4}{*}{\tab[2mm] 302 \tab[2mm]} & 2,2-dimethyl-butane\_out\_G09 &
\multirow{2}{*}{Different} & \multirow{2}{*}{Different}
\\
& E = -237.083690318 \tab Freq =92.1843   &    &  \\ \cline{2-2}
& 2,3-dimethyl-butane\_out\_G09\_invertion   & \multicolumn{2}{c}
 {RMS = 0.269312}
\\
& E = -237.082490952 \tab Freq =58.6018   &    \multicolumn{2}{c} 
{ }
\\ \hline
\multirow{4}{*}{\tab[2mm] 303 \tab[2mm]} & 2,2-dimethyl-butane\_out\_G09 &
\multirow{2}{*}{Different} & \multirow{2}{*}{Different}
\\
& E = -237.083690318 \tab Freq =92.1843   &    &  \\ \cline{2-2}
& 2-methyl-pentane\_out\_G09   & \multicolumn{2}{c}
 {RMS = 0.0426031}
\\
& E = -237.084759065 \tab Freq =80.9103   &    \multicolumn{2}{c} 
{ }
\\ \hline
\multirow{4}{*}{\tab[2mm] 304 \tab[2mm]} & 2,2-dimethyl-butane\_out\_G09 &
\multirow{2}{*}{Different} & \multirow{2}{*}{Different}
\\
& E = -237.083690318 \tab Freq =92.1843   &    &  \\ \cline{2-2}
& 2-methyl-pentane\_out\_G09\_invertion   & \multicolumn{2}{c}
 {RMS = 0.0778147}
\\
& E = -237.084759065 \tab Freq =80.9103   &    \multicolumn{2}{c} 
{ }
\\ \hline
\multirow{4}{*}{\tab[2mm] 305 \tab[2mm]} & 2,2-dimethyl-butane\_out\_G09 &
\multirow{2}{*}{Different} & \multirow{2}{*}{Different}
\\
& E = -237.083690318 \tab Freq =92.1843   &    &  \\ \cline{2-2}
& 3-methyl-pentane\_out\_G09   & \multicolumn{2}{c}
 {RMS = 0.766027}
\\
& E = -237.083213438 \tab Freq =81.4725   &    \multicolumn{2}{c} 
{ }
\\ \hline
\multirow{4}{*}{\tab[2mm] 306 \tab[2mm]} & 2,2-dimethyl-butane\_out\_G09 &
\multirow{2}{*}{Different} & \multirow{2}{*}{Different}
\\
& E = -237.083690318 \tab Freq =92.1843   &    &  \\ \cline{2-2}
& 3-methyl-pentane\_out\_G09\_invertion   & \multicolumn{2}{c}
 {RMS = 0.766028}
\\
& E = -237.083213438 \tab Freq =81.4725   &    \multicolumn{2}{c} 
{ }
\\ \hline
\multirow{4}{*}{\tab[2mm] 307 \tab[2mm]} & 2,2-dimethyl-butane\_out\_G09 &
\multirow{2}{*}{Different} & \multirow{2}{*}{Different}
\\
& E = -237.083690318 \tab Freq =92.1843   &    &  \\ \cline{2-2}
& hexane\_out\_G09   & \multicolumn{2}{c}
 {RMS = 0}
\\
& E = -237.085500432 \tab Freq =72.0752   &    \multicolumn{2}{c} 
{ }
\\ \hline
\multirow{4}{*}{\tab[2mm] 308 \tab[2mm]} & 2,2-dimethyl-butane\_out\_G09 &
\multirow{2}{*}{Different} & \multirow{2}{*}{Different}
\\
& E = -237.083690318 \tab Freq =92.1843   &    &  \\ \cline{2-2}
& hexane\_out\_G09\_invertion   & \multicolumn{2}{c}
 {RMS = 0}
\\
& E = -237.085500432 \tab Freq =72.0752   &    \multicolumn{2}{c} 
{ }
\\ \hline
\end{tabular}
\newpage

\vtab[-2cm]
\tab[-2cm]
\begin{tabular}{c|m{8cm}|c|c}
\# & Moléculas & Restultado esperado & Resultado programa \\ \hline\hline
\multirow{4}{*}{\tab[2mm] 309 \tab[2mm]} & 2,3-dimethyl-butane\_out\_G09\_invertion &
\multirow{2}{*}{Different} & \multirow{2}{*}{Different}
\\
& E = -237.082490952 \tab Freq =58.6018   &    &  \\ \cline{2-2}
& 2-methyl-pentane\_out\_G09   & \multicolumn{2}{c}
 {RMS = 0.311915}
\\
& E = -237.084759065 \tab Freq =80.9103   &    \multicolumn{2}{c} 
{ }
\\ \hline
\multirow{4}{*}{\tab[2mm] 310 \tab[2mm]} & 2,3-dimethyl-butane\_out\_G09\_invertion &
\multirow{2}{*}{Different} & \multirow{2}{*}{Different}
\\
& E = -237.082490952 \tab Freq =58.6018   &    &  \\ \cline{2-2}
& 2-methyl-pentane\_out\_G09\_invertion   & \multicolumn{2}{c}
 {RMS = 0.191497}
\\
& E = -237.084759065 \tab Freq =80.9103   &    \multicolumn{2}{c} 
{ }
\\ \hline
\multirow{4}{*}{\tab[2mm] 311 \tab[2mm]} & 2,3-dimethyl-butane\_out\_G09\_invertion &
\multirow{2}{*}{Different} & \multirow{2}{*}{Different}
\\
& E = -237.082490952 \tab Freq =58.6018   &    &  \\ \cline{2-2}
& 3-methyl-pentane\_out\_G09   & \multicolumn{2}{c}
 {RMS = 0.496715}
\\
& E = -237.083213438 \tab Freq =81.4725   &    \multicolumn{2}{c} 
{ }
\\ \hline
\multirow{4}{*}{\tab[2mm] 312 \tab[2mm]} & 2,3-dimethyl-butane\_out\_G09\_invertion &
\multirow{2}{*}{Different} & \multirow{2}{*}{Different}
\\
& E = -237.082490952 \tab Freq =58.6018   &    &  \\ \cline{2-2}
& 3-methyl-pentane\_out\_G09\_invertion   & \multicolumn{2}{c}
 {RMS = 0.496716}
\\
& E = -237.083213438 \tab Freq =81.4725   &    \multicolumn{2}{c} 
{ }
\\ \hline
\multirow{4}{*}{\tab[2mm] 313 \tab[2mm]} & 2,3-dimethyl-butane\_out\_G09\_invertion &
\multirow{2}{*}{Different} & \multirow{2}{*}{Different}
\\
& E = -237.082490952 \tab Freq =58.6018   &    &  \\ \cline{2-2}
& hexane\_out\_G09   & \multicolumn{2}{c}
 {RMS = 0.269312}
\\
& E = -237.085500432 \tab Freq =72.0752   &    \multicolumn{2}{c} 
{ }
\\ \hline
\multirow{4}{*}{\tab[2mm] 314 \tab[2mm]} & 2,3-dimethyl-butane\_out\_G09\_invertion &
\multirow{2}{*}{Different} & \multirow{2}{*}{Different}
\\
& E = -237.082490952 \tab Freq =58.6018   &    &  \\ \cline{2-2}
& hexane\_out\_G09\_invertion   & \multicolumn{2}{c}
 {RMS = 0.269312}
\\
& E = -237.085500432 \tab Freq =72.0752   &    \multicolumn{2}{c} 
{ }
\\ \hline
\multirow{4}{*}{\tab[2mm] 315 \tab[2mm]} & 2,3-dimethyl-butane\_out\_G09 &
\multirow{2}{*}{Equal} & \multirow{2}{*}{Equal}
\\
& E = -237.082490952 \tab Freq =58.6018   &    &  \\ \cline{2-2}
& 2,3-dimethyl-butane\_out\_G09\_invertion   & \multicolumn{2}{c}
{ RMS = 1.773561E-07}
\\
& E = -237.082490952 \tab Freq =58.6018   &    \multicolumn{2}{c} 
{ }
\\ \hline
\end{tabular}
\newpage

\vtab[-2cm]
\tab[-2cm]
\begin{tabular}{c|m{8cm}|c|c}
\# & Moléculas & Restultado esperado & Resultado programa \\ \hline\hline
\multirow{4}{*}{\tab[2mm] 316 \tab[2mm]} & 2,3-dimethyl-butane\_out\_G09 &
\multirow{2}{*}{Different} & \multirow{2}{*}{Different}
\\
& E = -237.082490952 \tab Freq =58.6018   &    &  \\ \cline{2-2}
& 2-methyl-pentane\_out\_G09   & \multicolumn{2}{c}
 {RMS = 0.311916}
\\
& E = -237.084759065 \tab Freq =80.9103   &    \multicolumn{2}{c} 
{ }
\\ \hline
\multirow{4}{*}{\tab[2mm] 317 \tab[2mm]} & 2,3-dimethyl-butane\_out\_G09 &
\multirow{2}{*}{Different} & \multirow{2}{*}{Different}
\\
& E = -237.082490952 \tab Freq =58.6018   &    &  \\ \cline{2-2}
& 2-methyl-pentane\_out\_G09\_invertion   & \multicolumn{2}{c}
 {RMS = 0.191498}
\\
& E = -237.084759065 \tab Freq =80.9103   &    \multicolumn{2}{c} 
{ }
\\ \hline
\multirow{4}{*}{\tab[2mm] 318 \tab[2mm]} & 2,3-dimethyl-butane\_out\_G09 &
\multirow{2}{*}{Different} & \multirow{2}{*}{Different}
\\
& E = -237.082490952 \tab Freq =58.6018   &    &  \\ \cline{2-2}
& 3-methyl-pentane\_out\_G09   & \multicolumn{2}{c}
 {RMS = 0.496715}
\\
& E = -237.083213438 \tab Freq =81.4725   &    \multicolumn{2}{c} 
{ }
\\ \hline
\multirow{4}{*}{\tab[2mm] 319 \tab[2mm]} & 2,3-dimethyl-butane\_out\_G09 &
\multirow{2}{*}{Different} & \multirow{2}{*}{Different}
\\
& E = -237.082490952 \tab Freq =58.6018   &    &  \\ \cline{2-2}
& 3-methyl-pentane\_out\_G09\_invertion   & \multicolumn{2}{c}
 {RMS = 0.496715}
\\
& E = -237.083213438 \tab Freq =81.4725   &    \multicolumn{2}{c} 
{ }
\\ \hline
\multirow{4}{*}{\tab[2mm] 320 \tab[2mm]} & 2,3-dimethyl-butane\_out\_G09 &
\multirow{2}{*}{Different} & \multirow{2}{*}{Different}
\\
& E = -237.082490952 \tab Freq =58.6018   &    &  \\ \cline{2-2}
& hexane\_out\_G09   & \multicolumn{2}{c}
 {RMS = 0.269312}
\\
& E = -237.085500432 \tab Freq =72.0752   &    \multicolumn{2}{c} 
{ }
\\ \hline
\multirow{4}{*}{\tab[2mm] 321 \tab[2mm]} & 2,3-dimethyl-butane\_out\_G09 &
\multirow{2}{*}{Different} & \multirow{2}{*}{Different}
\\
& E = -237.082490952 \tab Freq =58.6018   &    &  \\ \cline{2-2}
& hexane\_out\_G09\_invertion   & \multicolumn{2}{c}
 {RMS = 0.269312}
\\
& E = -237.085500432 \tab Freq =72.0752   &    \multicolumn{2}{c} 
{ }
\\ \hline
\multirow{4}{*}{\tab[2mm] 322 \tab[2mm]} & 2-methyl-pentane\_out\_G09\_invertion &
\multirow{2}{*}{Different} & \multirow{2}{*}{Different}
\\
& E = -237.084759065 \tab Freq =80.9103   &    &  \\ \cline{2-2}
& 3-methyl-pentane\_out\_G09   & \multicolumn{2}{c}
 {RMS = 0.688213}
\\
& E = -237.083213438 \tab Freq =81.4725   &    \multicolumn{2}{c} 
{ }
\\ \hline
\end{tabular}
\newpage

\vtab[-2cm]
\tab[-2cm]
\begin{tabular}{c|m{8cm}|c|c}
\# & Moléculas & Restultado esperado & Resultado programa \\ \hline\hline
\multirow{4}{*}{\tab[2mm] 323 \tab[2mm]} & 2-methyl-pentane\_out\_G09\_invertion &
\multirow{2}{*}{Different} & \multirow{2}{*}{Different}
\\
& E = -237.084759065 \tab Freq =80.9103   &    &  \\ \cline{2-2}
& 3-methyl-pentane\_out\_G09\_invertion   & \multicolumn{2}{c}
 {RMS = 0.688213}
\\
& E = -237.083213438 \tab Freq =81.4725   &    \multicolumn{2}{c} 
{ }
\\ \hline
\multirow{4}{*}{\tab[2mm] 324 \tab[2mm]} & 2-methyl-pentane\_out\_G09\_invertion &
\multirow{2}{*}{Different} & \multirow{2}{*}{Different}
\\
& E = -237.084759065 \tab Freq =80.9103   &    &  \\ \cline{2-2}
& hexane\_out\_G09   & \multicolumn{2}{c}
 {RMS = 0.0778147}
\\
& E = -237.085500432 \tab Freq =72.0752   &    \multicolumn{2}{c} 
{ }
\\ \hline
\multirow{4}{*}{\tab[2mm] 325 \tab[2mm]} & 2-methyl-pentane\_out\_G09\_invertion &
\multirow{2}{*}{Different} & \multirow{2}{*}{Different}
\\
& E = -237.084759065 \tab Freq =80.9103   &    &  \\ \cline{2-2}
& hexane\_out\_G09\_invertion   & \multicolumn{2}{c}
 {RMS = 0.0778147}
\\
& E = -237.085500432 \tab Freq =72.0752   &    \multicolumn{2}{c} 
{ }
\\ \hline
\multirow{4}{*}{\tab[2mm] 326 \tab[2mm]} & 2-methyl-pentane\_out\_G09 &
\multirow{2}{*}{Enantiomers} & \multirow{2}{*}{Enantiomers}
\\
& E = -237.084759065 \tab Freq =80.9103   &    &  \\ \cline{2-2}
& 2-methyl-pentane\_out\_G09\_invertion   & \multicolumn{2}{c}
 {RMS = 0.1203334}
\\
& E = -237.084759065 \tab Freq =80.9103   &    \multicolumn{2}{c} 
{ }
\\ \hline
\multirow{4}{*}{\tab[2mm] 327 \tab[2mm]} & 2-methyl-pentane\_out\_G09 &
\multirow{2}{*}{Different} & \multirow{2}{*}{Different}
\\
& E = -237.084759065 \tab Freq =80.9103   &    &  \\ \cline{2-2}
& 3-methyl-pentane\_out\_G09   & \multicolumn{2}{c}
 {RMS = 0.80863}
\\
& E = -237.083213438 \tab Freq =81.4725   &    \multicolumn{2}{c} 
{ }
\\ \hline
\multirow{4}{*}{\tab[2mm] 328 \tab[2mm]} & 2-methyl-pentane\_out\_G09 &
\multirow{2}{*}{Different} & \multirow{2}{*}{Different}
\\
& E = -237.084759065 \tab Freq =80.9103   &    &  \\ \cline{2-2}
& 3-methyl-pentane\_out\_G09\_invertion   & \multicolumn{2}{c}
 {RMS = 0.808631}
\\
& E = -237.083213438 \tab Freq =81.4725   &    \multicolumn{2}{c} 
{ }
\\ \hline
\multirow{4}{*}{\tab[2mm] 329 \tab[2mm]} & 2-methyl-pentane\_out\_G09 &
\multirow{2}{*}{Different} & \multirow{2}{*}{Different}
\\
& E = -237.084759065 \tab Freq =80.9103   &    &  \\ \cline{2-2}
& hexane\_out\_G09   & \multicolumn{2}{c}
 {RMS = 0.0426031}
\\
& E = -237.085500432 \tab Freq =72.0752   &    \multicolumn{2}{c} 
{ }
\\ \hline
\end{tabular}
\newpage

\vtab[-2cm]
\tab[-2cm]
\begin{tabular}{c|m{8cm}|c|c}
\# & Moléculas & Restultado esperado & Resultado programa \\ \hline\hline
\multirow{4}{*}{\tab[2mm] 330 \tab[2mm]} & 2-methyl-pentane\_out\_G09 &
\multirow{2}{*}{Different} & \multirow{2}{*}{Different}
\\
& E = -237.084759065 \tab Freq =80.9103   &    &  \\ \cline{2-2}
& hexane\_out\_G09\_invertion   & \multicolumn{2}{c}
 {RMS = 0.0426031}
\\
& E = -237.085500432 \tab Freq =72.0752   &    \multicolumn{2}{c} 
{ }
\\ \hline
\multirow{4}{*}{\tab[2mm] 331 \tab[2mm]} & 3-methyl-pentane\_out\_G09\_invertion &
\multirow{2}{*}{Different} & \multirow{2}{*}{Different}
\\
& E = -237.083213438 \tab Freq =81.4725   &    &  \\ \cline{2-2}
& hexane\_out\_G09   & \multicolumn{2}{c}
 {RMS = 0.766028}
\\
& E = -237.085500432 \tab Freq =72.0752   &    \multicolumn{2}{c} 
{ }
\\ \hline
\multirow{4}{*}{\tab[2mm] 332 \tab[2mm]} & 3-methyl-pentane\_out\_G09\_invertion &
\multirow{2}{*}{Different} & \multirow{2}{*}{Different}
\\
& E = -237.083213438 \tab Freq =81.4725   &    &  \\ \cline{2-2}
& hexane\_out\_G09\_invertion   & \multicolumn{2}{c}
 {RMS = 0.766028}
\\
& E = -237.085500432 \tab Freq =72.0752   &    \multicolumn{2}{c} 
{ }
\\ \hline
\multirow{4}{*}{\tab[2mm] 333 \tab[2mm]} & 3-methyl-pentane\_out\_G09 &
\multirow{2}{*}{Equal} & \multirow{2}{*}{Equal}
\\
& E = -237.083213438 \tab Freq =81.4725   &    &  \\ \cline{2-2}
& 3-methyl-pentane\_out\_G09\_invertion   & \multicolumn{2}{c}
{ RMS = 2.333815E-07}
\\
& E = -237.083213438 \tab Freq =81.4725   &    \multicolumn{2}{c} 
{ }
\\ \hline
\multirow{4}{*}{\tab[2mm] 334 \tab[2mm]} & 3-methyl-pentane\_out\_G09 &
\multirow{2}{*}{Different} & \multirow{2}{*}{Different}
\\
& E = -237.083213438 \tab Freq =81.4725   &    &  \\ \cline{2-2}
& hexane\_out\_G09   & \multicolumn{2}{c}
 {RMS = 0.766027}
\\
& E = -237.085500432 \tab Freq =72.0752   &    \multicolumn{2}{c} 
{ }
\\ \hline
\multirow{4}{*}{\tab[2mm] 335 \tab[2mm]} & 3-methyl-pentane\_out\_G09 &
\multirow{2}{*}{Different} & \multirow{2}{*}{Different}
\\
& E = -237.083213438 \tab Freq =81.4725   &    &  \\ \cline{2-2}
& hexane\_out\_G09\_invertion   & \multicolumn{2}{c}
 {RMS = 0.766027}
\\
& E = -237.085500432 \tab Freq =72.0752   &    \multicolumn{2}{c} 
{ }
\\ \hline
\multirow{4}{*}{\tab[2mm] 336 \tab[2mm]} & hexane\_out\_G09 &
\multirow{2}{*}{Equal} & \multirow{2}{*}{Equal}
\\
& E = -237.085500432 \tab Freq =72.0752   &    &  \\ \cline{2-2}
& hexane\_out\_G09\_invertion   & \multicolumn{2}{c}
{ RMS = 0}
\\
& E = -237.085500432 \tab Freq =72.0752   &    \multicolumn{2}{c} 
{ }
\\ \hline
\end{tabular}
\newpage

\vtab[-2cm]
\tab[-2cm]
\begin{tabular}{c|m{8cm}|c|c}
\# & Moléculas & Restultado esperado & Resultado programa \\ \hline\hline
\multirow{4}{*}{\tab[2mm] 337 \tab[2mm]} & neopentane\_out\_G09\_invertion &
\multirow{2}{*}{Equal} & \multirow{2}{*}{Equal}
\\
& E = -197.772988492 \tab Freq =208.9793   &    &  \\ \cline{2-2}
& neopentane\_out\_G09\_rot\_x15-y15-z15   & \multicolumn{2}{c}
{\textcolor{Red}{ RMS = 0.0989238}}
\\
& E = -197.772988492 \tab Freq =208.9793   &    \multicolumn{2}{c} 
{\textcolor{Red}{ Warning: Degeneracy}}
\\ \hline
\multirow{4}{*}{\tab[2mm] 338 \tab[2mm]} & neopentane\_out\_G09\_invertion &
\multirow{2}{*}{Different} & \multirow{2}{*}{Different}
\\
& E = -197.772988492 \tab Freq =208.9793   &    &  \\ \cline{2-2}
& tert-butylamine\_out\_G09   & \multicolumn{2}{c}
 {RMS = 0.0319157}
\\
& E = -213.803110609 \tab Freq =215.0777   &    \multicolumn{2}{c} 
{ }
\\ \hline
\multirow{4}{*}{\tab[2mm] 339 \tab[2mm]} & neopentane\_out\_G09\_invertion &
\multirow{2}{*}{Different} & \multirow{2}{*}{Different}
\\
& E = -197.772988492 \tab Freq =208.9793   &    &  \\ \cline{2-2}
& tert-butylamine\_out\_G09\_invertion   & \multicolumn{2}{c}
 {RMS = 0.0319157}
\\
& E = -213.803110609 \tab Freq =215.0777   &    \multicolumn{2}{c} 
{ }
\\ \hline
\multirow{4}{*}{\tab[2mm] 340 \tab[2mm]} & neopentane\_out\_G09\_invertion &
\multirow{2}{*}{Different} & \multirow{2}{*}{Different}
\\
& E = -197.772988492 \tab Freq =208.9793   &    &  \\ \cline{2-2}
& tert-butylamine\_out\_G09\_rot\_x15\_y15\_z15   & \multicolumn{2}{c}
 {RMS = 0.0319157}
\\
& E = -213.803110609 \tab Freq =215.0777   &    \multicolumn{2}{c} 
{ }
\\ \hline
\multirow{4}{*}{\tab[2mm] 341 \tab[2mm]} & neopentane\_out\_G09\_invertion &
\multirow{2}{*}{Different} & \multirow{2}{*}{Different}
\\
& E = -197.772988492 \tab Freq =208.9793   &    &  \\ \cline{2-2}
& tetramethylsilane\_out\_G09   & \multicolumn{2}{c}
 {RMS = 0.390641}
\\
& E = -449.191777923 \tab Freq =156.0268   &    \multicolumn{2}{c} 
{ }
\\ \hline
\multirow{4}{*}{\tab[2mm] 342 \tab[2mm]} & neopentane\_out\_G09\_invertion &
\multirow{2}{*}{Different} & \multirow{2}{*}{Different}
\\
& E = -197.772988492 \tab Freq =208.9793   &    &  \\ \cline{2-2}
& tetramethylsilane\_out\_G09\_invertion   & \multicolumn{2}{c}
 {RMS = 0.390642}
\\
& E = -449.191777923 \tab Freq =156.0268   &    \multicolumn{2}{c} 
{ }
\\ \hline
\multirow{4}{*}{\tab[2mm] 343 \tab[2mm]} & neopentane\_out\_G09\_rot\_x15-y15-z15 &
\multirow{2}{*}{Different} & \multirow{2}{*}{Different}
\\
& E = -197.772988492 \tab Freq =208.9793   &    &  \\ \cline{2-2}
& tert-butylamine\_out\_G09   & \multicolumn{2}{c}
 {RMS = 0.0670111}
\\
& E = -213.803110609 \tab Freq =215.0777   &    \multicolumn{2}{c} 
{ }
\\ \hline
\end{tabular}
\newpage

\vtab[-2cm]
\tab[-2cm]
\begin{tabular}{c|m{8cm}|c|c}
\# & Moléculas & Restultado esperado & Resultado programa \\ \hline\hline
\multirow{4}{*}{\tab[2mm] 344 \tab[2mm]} & neopentane\_out\_G09\_rot\_x15-y15-z15 &
\multirow{2}{*}{Different} & \multirow{2}{*}{Different}
\\
& E = -197.772988492 \tab Freq =208.9793   &    &  \\ \cline{2-2}
& tert-butylamine\_out\_G09\_invertion   & \multicolumn{2}{c}
 {RMS = 0.0670111}
\\
& E = -213.803110609 \tab Freq =215.0777   &    \multicolumn{2}{c} 
{ }
\\ \hline
\multirow{4}{*}{\tab[2mm] 345 \tab[2mm]} & neopentane\_out\_G09\_rot\_x15-y15-z15 &
\multirow{2}{*}{Different} & \multirow{2}{*}{Different}
\\
& E = -197.772988492 \tab Freq =208.9793   &    &  \\ \cline{2-2}
& tert-butylamine\_out\_G09\_rot\_x15\_y15\_z15   & \multicolumn{2}{c}
 {RMS = 0.0670111}
\\
& E = -213.803110609 \tab Freq =215.0777   &    \multicolumn{2}{c} 
{ }
\\ \hline
\multirow{4}{*}{\tab[2mm] 346 \tab[2mm]} & neopentane\_out\_G09\_rot\_x15-y15-z15 &
\multirow{2}{*}{Different} & \multirow{2}{*}{Different}
\\
& E = -197.772988492 \tab Freq =208.9793   &    &  \\ \cline{2-2}
& tetramethylsilane\_out\_G09   & \multicolumn{2}{c}
 {RMS = 0.489568}
\\
& E = -449.191777923 \tab Freq =156.0268   &    \multicolumn{2}{c} 
{ }
\\ \hline
\multirow{4}{*}{\tab[2mm] 347 \tab[2mm]} & neopentane\_out\_G09\_rot\_x15-y15-z15 &
\multirow{2}{*}{Different} & \multirow{2}{*}{Different}
\\
& E = -197.772988492 \tab Freq =208.9793   &    &  \\ \cline{2-2}
& tetramethylsilane\_out\_G09\_invertion   & \multicolumn{2}{c}
 {RMS = 0.489569}
\\
& E = -449.191777923 \tab Freq =156.0268   &    \multicolumn{2}{c} 
{ }
\\ \hline
\multirow{4}{*}{\tab[2mm] 348 \tab[2mm]} & neopentane\_out\_G09 &
\multirow{2}{*}{Equal} & \multirow{2}{*}{Equal}
\\
& E = -197.772988492 \tab Freq =208.9793   &    &  \\ \cline{2-2}
& neopentane\_out\_G09\_invertion   & \multicolumn{2}{c}
{\textcolor{Red}{ RMS = 0.0922148}}
\\
& E = -197.772988492 \tab Freq =208.9793   &    \multicolumn{2}{c} 
{\textcolor{Red}{ Warning: Degeneracy}}
\\ \hline
\multirow{4}{*}{\tab[2mm] 349 \tab[2mm]} & neopentane\_out\_G09 &
\multirow{2}{*}{Equal} & \multirow{2}{*}{Equal}
\\
& E = -197.772988492 \tab Freq =208.9793   &    &  \\ \cline{2-2}
& neopentane\_out\_G09\_rot\_x15-y15-z15   & \multicolumn{2}{c}
{\textcolor{Red}{ RMS = 0.00670901}}
\\
& E = -197.772988492 \tab Freq =208.9793   &    \multicolumn{2}{c} 
{\textcolor{Red}{ Warning: Degeneracy}}
\\ \hline
\multirow{4}{*}{\tab[2mm] 350 \tab[2mm]} & neopentane\_out\_G09 &
\multirow{2}{*}{Different} & \multirow{2}{*}{Different}
\\
& E = -197.772988492 \tab Freq =208.9793   &    &  \\ \cline{2-2}
& tert-butylamine\_out\_G09   & \multicolumn{2}{c}
 {RMS = 0.0603026}
\\
& E = -213.803110609 \tab Freq =215.0777   &    \multicolumn{2}{c} 
{ }
\\ \hline
\end{tabular}
\newpage

\vtab[-2cm]
\tab[-2cm]
\begin{tabular}{c|m{8cm}|c|c}
\# & Moléculas & Restultado esperado & Resultado programa \\ \hline\hline
\multirow{4}{*}{\tab[2mm] 351 \tab[2mm]} & neopentane\_out\_G09 &
\multirow{2}{*}{Different} & \multirow{2}{*}{Different}
\\
& E = -197.772988492 \tab Freq =208.9793   &    &  \\ \cline{2-2}
& tert-butylamine\_out\_G09\_invertion   & \multicolumn{2}{c}
 {RMS = 0.0603026}
\\
& E = -213.803110609 \tab Freq =215.0777   &    \multicolumn{2}{c} 
{ }
\\ \hline
\multirow{4}{*}{\tab[2mm] 352 \tab[2mm]} & neopentane\_out\_G09 &
\multirow{2}{*}{Different} & \multirow{2}{*}{Different}
\\
& E = -197.772988492 \tab Freq =208.9793   &    &  \\ \cline{2-2}
& tert-butylamine\_out\_G09\_rot\_x15\_y15\_z15   & \multicolumn{2}{c}
 {RMS = 0.0603026}
\\
& E = -213.803110609 \tab Freq =215.0777   &    \multicolumn{2}{c} 
{ }
\\ \hline
\multirow{4}{*}{\tab[2mm] 353 \tab[2mm]} & neopentane\_out\_G09 &
\multirow{2}{*}{Different} & \multirow{2}{*}{Different}
\\
& E = -197.772988492 \tab Freq =208.9793   &    &  \\ \cline{2-2}
& tetramethylsilane\_out\_G09   & \multicolumn{2}{c}
 {RMS = 0.48286}
\\
& E = -449.191777923 \tab Freq =156.0268   &    \multicolumn{2}{c} 
{ }
\\ \hline
\multirow{4}{*}{\tab[2mm] 354 \tab[2mm]} & neopentane\_out\_G09 &
\multirow{2}{*}{Different} & \multirow{2}{*}{Different}
\\
& E = -197.772988492 \tab Freq =208.9793   &    &  \\ \cline{2-2}
& tetramethylsilane\_out\_G09\_invertion   & \multicolumn{2}{c}
 {RMS = 0.48286}
\\
& E = -449.191777923 \tab Freq =156.0268   &    \multicolumn{2}{c} 
{ }
\\ \hline
\multirow{4}{*}{\tab[2mm] 355 \tab[2mm]} & neopentane\_Symmetry\_out\_G09\_invertion &
\multirow{2}{*}{Equal} & \multirow{2}{*}{Equal}
\\
& E = -197.772988495 \tab Freq =208.9816   &    &  \\ \cline{2-2}
& neopentane\_out\_G09   & \multicolumn{2}{c}
{\textcolor{Red}{ RMS = 0.092204}}
\\
& E = -197.772988492 \tab Freq =208.9793   &    \multicolumn{2}{c} 
{\textcolor{Red}{ Warning: Degeneracy}}
\\ \hline
\multirow{4}{*}{\tab[2mm] 356 \tab[2mm]} & neopentane\_Symmetry\_out\_G09\_invertion &
\multirow{2}{*}{Equal} & \multirow{2}{*}{Equal}
\\
& E = -197.772988495 \tab Freq =208.9816   &    &  \\ \cline{2-2}
& neopentane\_out\_G09\_invertion   & \multicolumn{2}{c}
{ RMS = 1.07659E-05}
\\
& E = -197.772988492 \tab Freq =208.9793   &    \multicolumn{2}{c} 
{\textcolor{Red}{ Warning: Degeneracy}}
\\ \hline
\multirow{4}{*}{\tab[2mm] 357 \tab[2mm]} & neopentane\_Symmetry\_out\_G09\_invertion &
\multirow{2}{*}{Equal} & \multirow{2}{*}{Equal}
\\
& E = -197.772988495 \tab Freq =208.9816   &    &  \\ \cline{2-2}
& neopentane\_out\_G09\_rot\_x15-y15-z15   & \multicolumn{2}{c}
{\textcolor{Red}{ RMS = 0.0989131}}
\\
& E = -197.772988492 \tab Freq =208.9793   &    \multicolumn{2}{c} 
{\textcolor{Red}{ Warning: Degeneracy}}
\\ \hline
\end{tabular}
\newpage

\vtab[-2cm]
\tab[-2cm]
\begin{tabular}{c|m{8cm}|c|c}
\# & Moléculas & Restultado esperado & Resultado programa \\ \hline\hline
\multirow{4}{*}{\tab[2mm] 358 \tab[2mm]} & neopentane\_Symmetry\_out\_G09\_invertion &
\multirow{2}{*}{Equal} & \multirow{2}{*}{Equal}
\\
& E = -197.772988495 \tab Freq =208.9816   &    &  \\ \cline{2-2}
& neopentane\_Symmetry\_out\_G09\_rot\_x45\_y45\_z60   & \multicolumn{2}{c}
{\textcolor{Red}{ RMS = 0.0214957}}
\\
& E = -197.772988495 \tab Freq =208.9816   &    \multicolumn{2}{c} 
{\textcolor{Red}{ Warning: Degeneracy}}
\\ \hline
\multirow{4}{*}{\tab[2mm] 359 \tab[2mm]} & neopentane\_Symmetry\_out\_G09\_invertion &
\multirow{2}{*}{Different} & \multirow{2}{*}{Different}
\\
& E = -197.772988495 \tab Freq =208.9816   &    &  \\ \cline{2-2}
& tert-butylamine\_out\_G09   & \multicolumn{2}{c}
 {RMS = 0.0319051}
\\
& E = -213.803110609 \tab Freq =215.0777   &    \multicolumn{2}{c} 
{ }
\\ \hline
\multirow{4}{*}{\tab[2mm] 360 \tab[2mm]} & neopentane\_Symmetry\_out\_G09\_invertion &
\multirow{2}{*}{Different} & \multirow{2}{*}{Different}
\\
& E = -197.772988495 \tab Freq =208.9816   &    &  \\ \cline{2-2}
& tert-butylamine\_out\_G09\_invertion   & \multicolumn{2}{c}
 {RMS = 0.0319051}
\\
& E = -213.803110609 \tab Freq =215.0777   &    \multicolumn{2}{c} 
{ }
\\ \hline
\multirow{4}{*}{\tab[2mm] 361 \tab[2mm]} & neopentane\_Symmetry\_out\_G09\_invertion &
\multirow{2}{*}{Different} & \multirow{2}{*}{Different}
\\
& E = -197.772988495 \tab Freq =208.9816   &    &  \\ \cline{2-2}
& tert-butylamine\_out\_G09\_rot\_x15\_y15\_z15   & \multicolumn{2}{c}
 {RMS = 0.0319051}
\\
& E = -213.803110609 \tab Freq =215.0777   &    \multicolumn{2}{c} 
{ }
\\ \hline
\multirow{4}{*}{\tab[2mm] 362 \tab[2mm]} & neopentane\_Symmetry\_out\_G09\_invertion &
\multirow{2}{*}{Different} & \multirow{2}{*}{Different}
\\
& E = -197.772988495 \tab Freq =208.9816   &    &  \\ \cline{2-2}
& tetramethylsilane\_out\_G09   & \multicolumn{2}{c}
 {RMS = 0.390652}
\\
& E = -449.191777923 \tab Freq =156.0268   &    \multicolumn{2}{c} 
{ }
\\ \hline
\multirow{4}{*}{\tab[2mm] 363 \tab[2mm]} & neopentane\_Symmetry\_out\_G09\_invertion &
\multirow{2}{*}{Different} & \multirow{2}{*}{Different}
\\
& E = -197.772988495 \tab Freq =208.9816   &    &  \\ \cline{2-2}
& tetramethylsilane\_out\_G09\_invertion   & \multicolumn{2}{c}
 {RMS = 0.390653}
\\
& E = -449.191777923 \tab Freq =156.0268   &    \multicolumn{2}{c} 
{ }
\\ \hline
\multirow{4}{*}{\tab[2mm] 364 \tab[2mm]} & neopentane\_Symmetry\_out\_G09\_rot\_x45\_y45\_z60 &
\multirow{2}{*}{Equal} & \multirow{2}{*}{Equal}
\\
& E = -197.772988495 \tab Freq =208.9816   &    &  \\ \cline{2-2}
& neopentane\_out\_G09   & \multicolumn{2}{c}
{\textcolor{Red}{ RMS = 0.1137}}
\\
& E = -197.772988492 \tab Freq =208.9793   &    \multicolumn{2}{c} 
{\textcolor{Red}{ Warning: Degeneracy}}
\\ \hline
\end{tabular}
\newpage

\vtab[-2cm]
\tab[-2cm]
\begin{tabular}{c|m{8cm}|c|c}
\# & Moléculas & Restultado esperado & Resultado programa \\ \hline\hline
\multirow{4}{*}{\tab[2mm] 365 \tab[2mm]} & neopentane\_Symmetry\_out\_G09\_rot\_x45\_y45\_z60 &
\multirow{2}{*}{Equal} & \multirow{2}{*}{Equal}
\\
& E = -197.772988495 \tab Freq =208.9816   &    &  \\ \cline{2-2}
& neopentane\_out\_G09\_invertion   & \multicolumn{2}{c}
{\textcolor{Red}{ RMS = 0.0214849}}
\\
& E = -197.772988492 \tab Freq =208.9793   &    \multicolumn{2}{c} 
{\textcolor{Red}{ Warning: Degeneracy}}
\\ \hline
\multirow{4}{*}{\tab[2mm] 366 \tab[2mm]} & neopentane\_Symmetry\_out\_G09\_rot\_x45\_y45\_z60 &
\multirow{2}{*}{Equal} & \multirow{2}{*}{Equal}
\\
& E = -197.772988495 \tab Freq =208.9816   &    &  \\ \cline{2-2}
& neopentane\_out\_G09\_rot\_x15-y15-z15   & \multicolumn{2}{c}
{\textcolor{Red}{ RMS = 0.120409}}
\\
& E = -197.772988492 \tab Freq =208.9793   &    \multicolumn{2}{c} 
{\textcolor{Red}{ Warning: Degeneracy}}
\\ \hline
\multirow{4}{*}{\tab[2mm] 367 \tab[2mm]} & neopentane\_Symmetry\_out\_G09\_rot\_x45\_y45\_z60 &
\multirow{2}{*}{Different} & \multirow{2}{*}{Different}
\\
& E = -197.772988495 \tab Freq =208.9816   &    &  \\ \cline{2-2}
& tert-butylamine\_out\_G09   & \multicolumn{2}{c}
 {RMS = 0.0533989}
\\
& E = -213.803110609 \tab Freq =215.0777   &    \multicolumn{2}{c} 
{ }
\\ \hline
\multirow{4}{*}{\tab[2mm] 368 \tab[2mm]} & neopentane\_Symmetry\_out\_G09\_rot\_x45\_y45\_z60 &
\multirow{2}{*}{Different} & \multirow{2}{*}{Different}
\\
& E = -197.772988495 \tab Freq =208.9816   &    &  \\ \cline{2-2}
& tert-butylamine\_out\_G09\_invertion   & \multicolumn{2}{c}
 {RMS = 0.0533989}
\\
& E = -213.803110609 \tab Freq =215.0777   &    \multicolumn{2}{c} 
{ }
\\ \hline
\multirow{4}{*}{\tab[2mm] 369 \tab[2mm]} & neopentane\_Symmetry\_out\_G09\_rot\_x45\_y45\_z60 &
\multirow{2}{*}{Different} & \multirow{2}{*}{Different}
\\
& E = -197.772988495 \tab Freq =208.9816   &    &  \\ \cline{2-2}
& tert-butylamine\_out\_G09\_rot\_x15\_y15\_z15   & \multicolumn{2}{c}
 {RMS = 0.0533989}
\\
& E = -213.803110609 \tab Freq =215.0777   &    \multicolumn{2}{c} 
{ }
\\ \hline
\multirow{4}{*}{\tab[2mm] 370 \tab[2mm]} & neopentane\_Symmetry\_out\_G09\_rot\_x45\_y45\_z60 &
\multirow{2}{*}{Different} & \multirow{2}{*}{Different}
\\
& E = -197.772988495 \tab Freq =208.9816   &    &  \\ \cline{2-2}
& tetramethylsilane\_out\_G09   & \multicolumn{2}{c}
 {RMS = 0.369158}
\\
& E = -449.191777923 \tab Freq =156.0268   &    \multicolumn{2}{c} 
{ }
\\ \hline
\multirow{4}{*}{\tab[2mm] 371 \tab[2mm]} & neopentane\_Symmetry\_out\_G09\_rot\_x45\_y45\_z60 &
\multirow{2}{*}{Different} & \multirow{2}{*}{Different}
\\
& E = -197.772988495 \tab Freq =208.9816   &    &  \\ \cline{2-2}
& tetramethylsilane\_out\_G09\_invertion   & \multicolumn{2}{c}
 {RMS = 0.369159}
\\
& E = -449.191777923 \tab Freq =156.0268   &    \multicolumn{2}{c} 
{ }
\\ \hline
\end{tabular}
\newpage

\vtab[-2cm]
\tab[-2cm]
\begin{tabular}{c|m{8cm}|c|c}
\# & Moléculas & Restultado esperado & Resultado programa \\ \hline\hline
\multirow{4}{*}{\tab[2mm] 372 \tab[2mm]} & neopentane\_Symmetry\_out\_G09 &
\multirow{2}{*}{Equal} & \multirow{2}{*}{Equal}
\\
& E = -197.772988495 \tab Freq =208.9816   &    &  \\ \cline{2-2}
& neopentane\_out\_G09   & \multicolumn{2}{c}
{\textcolor{Red}{ RMS = 6.59223E-06}}
\\
& E = -197.772988492 \tab Freq =208.9793   &    \multicolumn{2}{c} 
{\textcolor{Red}{ Warning: Degeneracy}}
\\ \hline
\multirow{4}{*}{\tab[2mm] 373 \tab[2mm]} & neopentane\_Symmetry\_out\_G09 &
\multirow{2}{*}{Equal} & \multirow{2}{*}{Equal}
\\
& E = -197.772988495 \tab Freq =208.9816   &    &  \\ \cline{2-2}
& neopentane\_out\_G09\_invertion   & \multicolumn{2}{c}
{\textcolor{Red}{ RMS = 0.0922214}}
\\
& E = -197.772988492 \tab Freq =208.9793   &    \multicolumn{2}{c} 
{\textcolor{Red}{ Warning: Degeneracy}}
\\ \hline
\multirow{4}{*}{\tab[2mm] 374 \tab[2mm]} & neopentane\_Symmetry\_out\_G09 &
\multirow{2}{*}{Equal} & \multirow{2}{*}{Equal}
\\
& E = -197.772988495 \tab Freq =208.9816   &    &  \\ \cline{2-2}
& neopentane\_out\_G09\_rot\_x15-y15-z15   & \multicolumn{2}{c}
{\textcolor{Red}{ RMS = 0.00670242}}
\\
& E = -197.772988492 \tab Freq =208.9793   &    \multicolumn{2}{c} 
{\textcolor{Red}{ Warning: Degeneracy}}
\\ \hline
\multirow{4}{*}{\tab[2mm] 375 \tab[2mm]} & neopentane\_Symmetry\_out\_G09 &
\multirow{2}{*}{Equal} & \multirow{2}{*}{Equal}
\\
& E = -197.772988495 \tab Freq =208.9816   &    &  \\ \cline{2-2}
& neopentane\_Symmetry\_out\_G09\_invertion   & \multicolumn{2}{c}
{\textcolor{Red}{ RMS = 0.0922106}}
\\
& E = -197.772988495 \tab Freq =208.9816   &    \multicolumn{2}{c} 
{\textcolor{Red}{ Warning: Degeneracy}}
\\ \hline
\multirow{4}{*}{\tab[2mm] 376 \tab[2mm]} & neopentane\_Symmetry\_out\_G09 &
\multirow{2}{*}{Equal} & \multirow{2}{*}{Equal}
\\
& E = -197.772988495 \tab Freq =208.9816   &    &  \\ \cline{2-2}
& neopentane\_Symmetry\_out\_G09\_rot\_x45\_y45\_z60   & \multicolumn{2}{c}
{\textcolor{Red}{ RMS = 0.113706}}
\\
& E = -197.772988495 \tab Freq =208.9816   &    \multicolumn{2}{c} 
{\textcolor{Red}{ Warning: Degeneracy}}
\\ \hline
\multirow{4}{*}{\tab[2mm] 377 \tab[2mm]} & neopentane\_Symmetry\_out\_G09 &
\multirow{2}{*}{Different} & \multirow{2}{*}{Different}
\\
& E = -197.772988495 \tab Freq =208.9816   &    &  \\ \cline{2-2}
& tert-butylamine\_out\_G09   & \multicolumn{2}{c}
 {RMS = 0.0603092}
\\
& E = -213.803110609 \tab Freq =215.0777   &    \multicolumn{2}{c} 
{ }
\\ \hline
\multirow{4}{*}{\tab[2mm] 378 \tab[2mm]} & neopentane\_Symmetry\_out\_G09 &
\multirow{2}{*}{Different} & \multirow{2}{*}{Different}
\\
& E = -197.772988495 \tab Freq =208.9816   &    &  \\ \cline{2-2}
& tert-butylamine\_out\_G09\_invertion   & \multicolumn{2}{c}
 {RMS = 0.0603092}
\\
& E = -213.803110609 \tab Freq =215.0777   &    \multicolumn{2}{c} 
{ }
\\ \hline
\end{tabular}
\newpage

\vtab[-2cm]
\tab[-2cm]
\begin{tabular}{c|m{8cm}|c|c}
\# & Moléculas & Restultado esperado & Resultado programa \\ \hline\hline
\multirow{4}{*}{\tab[2mm] 379 \tab[2mm]} & neopentane\_Symmetry\_out\_G09 &
\multirow{2}{*}{Different} & \multirow{2}{*}{Different}
\\
& E = -197.772988495 \tab Freq =208.9816   &    &  \\ \cline{2-2}
& tert-butylamine\_out\_G09\_rot\_x15\_y15\_z15   & \multicolumn{2}{c}
 {RMS = 0.0603092}
\\
& E = -213.803110609 \tab Freq =215.0777   &    \multicolumn{2}{c} 
{ }
\\ \hline
\multirow{4}{*}{\tab[2mm] 380 \tab[2mm]} & neopentane\_Symmetry\_out\_G09 &
\multirow{2}{*}{Different} & \multirow{2}{*}{Different}
\\
& E = -197.772988495 \tab Freq =208.9816   &    &  \\ \cline{2-2}
& tetramethylsilane\_out\_G09   & \multicolumn{2}{c}
 {RMS = 0.482866}
\\
& E = -449.191777923 \tab Freq =156.0268   &    \multicolumn{2}{c} 
{ }
\\ \hline
\multirow{4}{*}{\tab[2mm] 381 \tab[2mm]} & neopentane\_Symmetry\_out\_G09 &
\multirow{2}{*}{Different} & \multirow{2}{*}{Different}
\\
& E = -197.772988495 \tab Freq =208.9816   &    &  \\ \cline{2-2}
& tetramethylsilane\_out\_G09\_invertion   & \multicolumn{2}{c}
 {RMS = 0.482867}
\\
& E = -449.191777923 \tab Freq =156.0268   &    \multicolumn{2}{c} 
{ }
\\ \hline
\multirow{4}{*}{\tab[2mm] 382 \tab[2mm]} & tert-butylamine\_out\_G09\_invertion &
\multirow{2}{*}{\textcolor{Red}{\bf Equal}} & \multirow{2}{*}{\textcolor{Red}{\bf Enantiomers}}
\\
& E = -213.803110609 \tab Freq =215.0777   &    &  \\ \cline{2-2}
& tert-butylamine\_out\_G09\_rot\_x15\_y15\_z15   & \multicolumn{2}{c}
 {RMS = 0.09408316}
\\
& E = -213.803110609 \tab Freq =215.0777   &    \multicolumn{2}{c} 
{ }
\\ \hline
\multirow{4}{*}{\tab[2mm] 383 \tab[2mm]} & tert-butylamine\_out\_G09\_invertion &
\multirow{2}{*}{Different} & \multirow{2}{*}{Different}
\\
& E = -213.803110609 \tab Freq =215.0777   &    &  \\ \cline{2-2}
& tetramethylsilane\_out\_G09   & \multicolumn{2}{c}
 {RMS = 0.102172}
\\
& E = -449.191777923 \tab Freq =156.0268   &    \multicolumn{2}{c} 
{ }
\\ \hline
\multirow{4}{*}{\tab[2mm] 384 \tab[2mm]} & tert-butylamine\_out\_G09\_invertion &
\multirow{2}{*}{Different} & \multirow{2}{*}{Different}
\\
& E = -213.803110609 \tab Freq =215.0777   &    &  \\ \cline{2-2}
& tetramethylsilane\_out\_G09\_invertion   & \multicolumn{2}{c}
 {RMS = 0.102172}
\\
& E = -449.191777923 \tab Freq =156.0268   &    \multicolumn{2}{c} 
{ }
\\ \hline
\multirow{4}{*}{\tab[2mm] 385 \tab[2mm]} & tert-butylamine\_out\_G09\_rot\_x15\_y15\_z15 &
\multirow{2}{*}{Different} & \multirow{2}{*}{Different}
\\
& E = -213.803110609 \tab Freq =215.0777   &    &  \\ \cline{2-2}
& tetramethylsilane\_out\_G09   & \multicolumn{2}{c}
 {RMS = 0.266281}
\\
& E = -449.191777923 \tab Freq =156.0268   &    \multicolumn{2}{c} 
{ }
\\ \hline
\end{tabular}
\newpage

\vtab[-2cm]
\tab[-2cm]
\begin{tabular}{c|m{8cm}|c|c}
\# & Moléculas & Restultado esperado & Resultado programa \\ \hline\hline
\multirow{4}{*}{\tab[2mm] 386 \tab[2mm]} & tert-butylamine\_out\_G09\_rot\_x15\_y15\_z15 &
\multirow{2}{*}{Different} & \multirow{2}{*}{Different}
\\
& E = -213.803110609 \tab Freq =215.0777   &    &  \\ \cline{2-2}
& tetramethylsilane\_out\_G09\_invertion   & \multicolumn{2}{c}
 {RMS = 0.266281}
\\
& E = -449.191777923 \tab Freq =156.0268   &    \multicolumn{2}{c} 
{ }
\\ \hline
\multirow{4}{*}{\tab[2mm] 387 \tab[2mm]} & tert-butylamine\_out\_G09 &
\multirow{2}{*}{\textcolor{Red}{\bf Equal}} & \multirow{2}{*}{\textcolor{Red}{\bf Enantiomers}}
\\
& E = -213.803110609 \tab Freq =215.0777   &    &  \\ \cline{2-2}
& tert-butylamine\_out\_G09\_invertion   & \multicolumn{2}{c}
 {RMS = 0.09408316}
\\
& E = -213.803110609 \tab Freq =215.0777   &    \multicolumn{2}{c} 
{ }
\\ \hline
\multirow{4}{*}{\tab[2mm] 388 \tab[2mm]} & tert-butylamine\_out\_G09 &
\multirow{2}{*}{Equal} & \multirow{2}{*}{Equal}
\\
& E = -213.803110609 \tab Freq =215.0777   &    &  \\ \cline{2-2}
& tert-butylamine\_out\_G09\_rot\_x15\_y15\_z15   & \multicolumn{2}{c}
{ RMS = 1.188653E-12}
\\
& E = -213.803110609 \tab Freq =215.0777   &    \multicolumn{2}{c} 
{ }
\\ \hline
\multirow{4}{*}{\tab[2mm] 389 \tab[2mm]} & tert-butylamine\_out\_G09 &
\multirow{2}{*}{Different} & \multirow{2}{*}{Different}
\\
& E = -213.803110609 \tab Freq =215.0777   &    &  \\ \cline{2-2}
& tetramethylsilane\_out\_G09   & \multicolumn{2}{c}
 {RMS = 0.143206}
\\
& E = -449.191777923 \tab Freq =156.0268   &    \multicolumn{2}{c} 
{ }
\\ \hline
\multirow{4}{*}{\tab[2mm] 390 \tab[2mm]} & tert-butylamine\_out\_G09 &
\multirow{2}{*}{Different} & \multirow{2}{*}{Different}
\\
& E = -213.803110609 \tab Freq =215.0777   &    &  \\ \cline{2-2}
& tetramethylsilane\_out\_G09\_invertion   & \multicolumn{2}{c}
 {RMS = 0.143206}
\\
& E = -449.191777923 \tab Freq =156.0268   &    \multicolumn{2}{c} 
{ }
\\ \hline
\multirow{4}{*}{\tab[2mm] 391 \tab[2mm]} & tetramethylsilane\_out\_G09 &
\multirow{2}{*}{Equal} & \multirow{2}{*}{Equal}
\\
& E = -449.191777923 \tab Freq =156.0268   &    &  \\ \cline{2-2}
& tetramethylsilane\_out\_G09\_invertion   & \multicolumn{2}{c}
{\textcolor{Red}{ RMS = 7.27607E-07}}
\\
& E = -449.191777923 \tab Freq =156.0268   &    \multicolumn{2}{c} 
{\textcolor{Red}{ Warning: Degeneracy}}
\\ \hline
\multirow{4}{*}{\tab[2mm] 392 \tab[2mm]} & Porphin\_out\_G09\_invertion &
\multirow{2}{*}{Different} & \multirow{2}{*}{Different}
\\
& E = -988.182666824 \tab Freq =-3564.4549   &    &  \\ \cline{2-2}
& Tetrabenzoporphyrin\_out\_G09   & \multicolumn{2}{c}
 {RMS = 0.117363}
\\
& E = -1604.15793928 \tab Freq =20.2596   &    \multicolumn{2}{c} 
{ }
\\ \hline
\end{tabular}
\newpage

\vtab[-2cm]
\tab[-2cm]
\begin{tabular}{c|m{8cm}|c|c}
\# & Moléculas & Restultado esperado & Resultado programa \\ \hline\hline
\multirow{4}{*}{\tab[2mm] 393 \tab[2mm]} & Porphin\_out\_G09\_invertion &
\multirow{2}{*}{Different} & \multirow{2}{*}{Different}
\\
& E = -988.182666824 \tab Freq =-3564.4549   &    &  \\ \cline{2-2}
& Tetrabenzoporphyrin\_out\_G09\_invertion   & \multicolumn{2}{c}
 {RMS = 0.117363}
\\
& E = -1604.15793928 \tab Freq =20.2596   &    \multicolumn{2}{c} 
{ }
\\ \hline
\multirow{4}{*}{\tab[2mm] 394 \tab[2mm]} & Porphin\_out\_G09\_invertion &
\multirow{2}{*}{Different} & \multirow{2}{*}{Different}
\\
& E = -988.182666824 \tab Freq =-3564.4549   &    &  \\ \cline{2-2}
& Tetrabenzoporphyrin\_rotated02\_out\_G09   & \multicolumn{2}{c}
 {RMS = 0.167859}
\\
& E = -1604.15788879 \tab Freq =19.3954   &    \multicolumn{2}{c} 
{ }
\\ \hline
\multirow{4}{*}{\tab[2mm] 395 \tab[2mm]} & Porphin\_out\_G09\_invertion &
\multirow{2}{*}{Different} & \multirow{2}{*}{Different}
\\
& E = -988.182666824 \tab Freq =-3564.4549   &    &  \\ \cline{2-2}
& Tetrabenzoporphyrin\_rotated02\_out\_G09\_invertion   & \multicolumn{2}{c}
 {RMS = 0.167858}
\\
& E = -1604.15788879 \tab Freq =19.3954   &    \multicolumn{2}{c} 
{ }
\\ \hline
\multirow{4}{*}{\tab[2mm] 396 \tab[2mm]} & Porphin\_out\_G09\_invertion &
\multirow{2}{*}{Different} & \multirow{2}{*}{Different}
\\
& E = -988.182666824 \tab Freq =-3564.4549   &    &  \\ \cline{2-2}
& Tetrabenzoporphyrin\_rotated\_out\_G09   & \multicolumn{2}{c}
 {RMS = 0.117363}
\\
& E = -1604.15793928 \tab Freq =20.2596   &    \multicolumn{2}{c} 
{ }
\\ \hline
\multirow{4}{*}{\tab[2mm] 397 \tab[2mm]} & Porphin\_out\_G09\_invertion &
\multirow{2}{*}{Different} & \multirow{2}{*}{Different}
\\
& E = -988.182666824 \tab Freq =-3564.4549   &    &  \\ \cline{2-2}
& Tetrabenzoporphyrin\_rotated\_out\_G09\_invertion   & \multicolumn{2}{c}
 {RMS = 0.117363}
\\
& E = -1604.15793928 \tab Freq =20.2596   &    \multicolumn{2}{c} 
{ }
\\ \hline
\multirow{4}{*}{\tab[2mm] 398 \tab[2mm]} & Porphin\_out\_G09 &
\multirow{2}{*}{Equal} & \multirow{2}{*}{Equal}
\\
& E = -988.182666824 \tab Freq =-3564.4549   &    &  \\ \cline{2-2}
& Porphin\_out\_G09\_invertion   & \multicolumn{2}{c}
{ RMS = 0}
\\
& E = -988.182666824 \tab Freq =-3564.4549   &    \multicolumn{2}{c} 
{ }
\\ \hline
\multirow{4}{*}{\tab[2mm] 399 \tab[2mm]} & Porphin\_out\_G09 &
\multirow{2}{*}{Different} & \multirow{2}{*}{Different}
\\
& E = -988.182666824 \tab Freq =-3564.4549   &    &  \\ \cline{2-2}
& Tetrabenzoporphyrin\_out\_G09   & \multicolumn{2}{c}
 {RMS = 0.117363}
\\
& E = -1604.15793928 \tab Freq =20.2596   &    \multicolumn{2}{c} 
{ }
\\ \hline
\end{tabular}
\newpage

\vtab[-2cm]
\tab[-2cm]
\begin{tabular}{c|m{8cm}|c|c}
\# & Moléculas & Restultado esperado & Resultado programa \\ \hline\hline
\multirow{4}{*}{\tab[2mm] 400 \tab[2mm]} & Porphin\_out\_G09 &
\multirow{2}{*}{Different} & \multirow{2}{*}{Different}
\\
& E = -988.182666824 \tab Freq =-3564.4549   &    &  \\ \cline{2-2}
& Tetrabenzoporphyrin\_out\_G09\_invertion   & \multicolumn{2}{c}
 {RMS = 0.117363}
\\
& E = -1604.15793928 \tab Freq =20.2596   &    \multicolumn{2}{c} 
{ }
\\ \hline
\multirow{4}{*}{\tab[2mm] 401 \tab[2mm]} & Porphin\_out\_G09 &
\multirow{2}{*}{Different} & \multirow{2}{*}{Different}
\\
& E = -988.182666824 \tab Freq =-3564.4549   &    &  \\ \cline{2-2}
& Tetrabenzoporphyrin\_rotated02\_out\_G09   & \multicolumn{2}{c}
 {RMS = 0.167859}
\\
& E = -1604.15788879 \tab Freq =19.3954   &    \multicolumn{2}{c} 
{ }
\\ \hline
\multirow{4}{*}{\tab[2mm] 402 \tab[2mm]} & Porphin\_out\_G09 &
\multirow{2}{*}{Different} & \multirow{2}{*}{Different}
\\
& E = -988.182666824 \tab Freq =-3564.4549   &    &  \\ \cline{2-2}
& Tetrabenzoporphyrin\_rotated02\_out\_G09\_invertion   & \multicolumn{2}{c}
 {RMS = 0.167858}
\\
& E = -1604.15788879 \tab Freq =19.3954   &    \multicolumn{2}{c} 
{ }
\\ \hline
\multirow{4}{*}{\tab[2mm] 403 \tab[2mm]} & Porphin\_out\_G09 &
\multirow{2}{*}{Different} & \multirow{2}{*}{Different}
\\
& E = -988.182666824 \tab Freq =-3564.4549   &    &  \\ \cline{2-2}
& Tetrabenzoporphyrin\_rotated\_out\_G09   & \multicolumn{2}{c}
 {RMS = 0.117363}
\\
& E = -1604.15793928 \tab Freq =20.2596   &    \multicolumn{2}{c} 
{ }
\\ \hline
\multirow{4}{*}{\tab[2mm] 404 \tab[2mm]} & Porphin\_out\_G09 &
\multirow{2}{*}{Different} & \multirow{2}{*}{Different}
\\
& E = -988.182666824 \tab Freq =-3564.4549   &    &  \\ \cline{2-2}
& Tetrabenzoporphyrin\_rotated\_out\_G09\_invertion   & \multicolumn{2}{c}
 {RMS = 0.117363}
\\
& E = -1604.15793928 \tab Freq =20.2596   &    \multicolumn{2}{c} 
{ }
\\ \hline
\multirow{4}{*}{\tab[2mm] 405 \tab[2mm]} & Tetrabenzoporphyrin\_out\_G09\_invertion &
\multirow{2}{*}{\textcolor{Red}{\bf Equal}} & \multirow{2}{*}{\textcolor{Red}{\bf Enantiomers}}
\\
& E = -1604.15793928 \tab Freq =20.2596   &    &  \\ \cline{2-2}
& Tetrabenzoporphyrin\_rotated02\_out\_G09   & \multicolumn{2}{c}
{\textcolor{Red}{ RMS = 0.0001074955}}
\\
& E = -1604.15788879 \tab Freq =19.3954   &    \multicolumn{2}{c} 
{ }
\\ \hline
\multirow{4}{*}{\tab[2mm] 406 \tab[2mm]} & Tetrabenzoporphyrin\_out\_G09\_invertion &
\multirow{2}{*}{\textcolor{Red}{\bf Equal}} & \multirow{2}{*}{\textcolor{Red}{\bf Enantiomers}}
\\
& E = -1604.15793928 \tab Freq =20.2596   &    &  \\ \cline{2-2}
& Tetrabenzoporphyrin\_rotated02\_out\_G09\_invertion   & \multicolumn{2}{c}
{\textcolor{Red}{ RMS = 0.0001075617}}
\\
& E = -1604.15788879 \tab Freq =19.3954   &    \multicolumn{2}{c} 
{ }
\\ \hline
\end{tabular}
\newpage

\vtab[-2cm]
\tab[-2cm]
\begin{tabular}{c|m{8cm}|c|c}
\# & Moléculas & Restultado esperado & Resultado programa \\ \hline\hline
\multirow{4}{*}{\tab[2mm] 407 \tab[2mm]} & Tetrabenzoporphyrin\_out\_G09\_invertion &
\multirow{2}{*}{Equal} & \multirow{2}{*}{Equal}
\\
& E = -1604.15793928 \tab Freq =20.2596   &    &  \\ \cline{2-2}
& Tetrabenzoporphyrin\_rotated\_out\_G09   & \multicolumn{2}{c}
{ RMS = 0}
\\
& E = -1604.15793928 \tab Freq =20.2596   &    \multicolumn{2}{c} 
{ }
\\ \hline
\multirow{4}{*}{\tab[2mm] 408 \tab[2mm]} & Tetrabenzoporphyrin\_out\_G09\_invertion &
\multirow{2}{*}{Equal} & \multirow{2}{*}{Equal}
\\
& E = -1604.15793928 \tab Freq =20.2596   &    &  \\ \cline{2-2}
& Tetrabenzoporphyrin\_rotated\_out\_G09\_invertion   & \multicolumn{2}{c}
{ RMS = 0}
\\
& E = -1604.15793928 \tab Freq =20.2596   &    \multicolumn{2}{c} 
{ }
\\ \hline
\multirow{4}{*}{\tab[2mm] 409 \tab[2mm]} & Tetrabenzoporphyrin\_out\_G09 &
\multirow{2}{*}{Equal} & \multirow{2}{*}{Equal}
\\
& E = -1604.15793928 \tab Freq =20.2596   &    &  \\ \cline{2-2}
& Tetrabenzoporphyrin\_out\_G09\_invertion   & \multicolumn{2}{c}
{ RMS = 0}
\\
& E = -1604.15793928 \tab Freq =20.2596   &    \multicolumn{2}{c} 
{ }
\\ \hline
\multirow{4}{*}{\tab[2mm] 410 \tab[2mm]} & Tetrabenzoporphyrin\_out\_G09 &
\multirow{2}{*}{\textcolor{Red}{\bf Equal}} & \multirow{2}{*}{\textcolor{Red}{\bf Enantiomers}}
\\
& E = -1604.15793928 \tab Freq =20.2596   &    &  \\ \cline{2-2}
& Tetrabenzoporphyrin\_rotated02\_out\_G09   & \multicolumn{2}{c}
{\textcolor{Red}{ RMS = 0.0001074955}}
\\
& E = -1604.15788879 \tab Freq =19.3954   &    \multicolumn{2}{c} 
{ }
\\ \hline
\multirow{4}{*}{\tab[2mm] 411 \tab[2mm]} & Tetrabenzoporphyrin\_out\_G09 &
\multirow{2}{*}{\textcolor{Red}{\bf Equal}} & \multirow{2}{*}{\textcolor{Red}{\bf Enantiomers}}
\\
& E = -1604.15793928 \tab Freq =20.2596   &    &  \\ \cline{2-2}
& Tetrabenzoporphyrin\_rotated02\_out\_G09\_invertion   & \multicolumn{2}{c}
{\textcolor{Red}{ RMS = 0.0001075617}}
\\
& E = -1604.15788879 \tab Freq =19.3954   &    \multicolumn{2}{c} 
{ }
\\ \hline
\multirow{4}{*}{\tab[2mm] 412 \tab[2mm]} & Tetrabenzoporphyrin\_out\_G09 &
\multirow{2}{*}{Equal} & \multirow{2}{*}{Equal}
\\
& E = -1604.15793928 \tab Freq =20.2596   &    &  \\ \cline{2-2}
& Tetrabenzoporphyrin\_rotated\_out\_G09   & \multicolumn{2}{c}
{ RMS = 0}
\\
& E = -1604.15793928 \tab Freq =20.2596   &    \multicolumn{2}{c} 
{ }
\\ \hline
\multirow{4}{*}{\tab[2mm] 413 \tab[2mm]} & Tetrabenzoporphyrin\_out\_G09 &
\multirow{2}{*}{Equal} & \multirow{2}{*}{Equal}
\\
& E = -1604.15793928 \tab Freq =20.2596   &    &  \\ \cline{2-2}
& Tetrabenzoporphyrin\_rotated\_out\_G09\_invertion   & \multicolumn{2}{c}
{ RMS = 0}
\\
& E = -1604.15793928 \tab Freq =20.2596   &    \multicolumn{2}{c} 
{ }
\\ \hline
\end{tabular}
\newpage

\vtab[-2cm]
\tab[-2cm]
\begin{tabular}{c|m{8cm}|c|c}
\# & Moléculas & Restultado esperado & Resultado programa \\ \hline\hline
\multirow{4}{*}{\tab[2mm] 414 \tab[2mm]} & Tetrabenzoporphyrin\_rotated02\_out\_G09\_invertion &
\multirow{2}{*}{\textcolor{Red}{\bf Equal}} & \multirow{2}{*}{\textcolor{Red}{\bf Enantiomers}}
\\
& E = -1604.15788879 \tab Freq =19.3954   &    &  \\ \cline{2-2}
& Tetrabenzoporphyrin\_rotated\_out\_G09   & \multicolumn{2}{c}
{\textcolor{Red}{ RMS = 0.0001075617}}
\\
& E = -1604.15793928 \tab Freq =20.2596   &    \multicolumn{2}{c} 
{ }
\\ \hline
\multirow{4}{*}{\tab[2mm] 415 \tab[2mm]} & Tetrabenzoporphyrin\_rotated02\_out\_G09\_invertion &
\multirow{2}{*}{\textcolor{Red}{\bf Equal}} & \multirow{2}{*}{\textcolor{Red}{\bf Enantiomers}}
\\
& E = -1604.15788879 \tab Freq =19.3954   &    &  \\ \cline{2-2}
& Tetrabenzoporphyrin\_rotated\_out\_G09\_invertion   & \multicolumn{2}{c}
{\textcolor{Red}{ RMS = 0.0001075617}}
\\
& E = -1604.15793928 \tab Freq =20.2596   &    \multicolumn{2}{c} 
{ }
\\ \hline
\multirow{4}{*}{\tab[2mm] 416 \tab[2mm]} & Tetrabenzoporphyrin\_rotated02\_out\_G09 &
\multirow{2}{*}{Equal} & \multirow{2}{*}{Equal}
\\
& E = -1604.15788879 \tab Freq =19.3954   &    &  \\ \cline{2-2}
& Tetrabenzoporphyrin\_rotated02\_out\_G09\_invertion   & \multicolumn{2}{c}
{ RMS = 6.620476E-08}
\\
& E = -1604.15788879 \tab Freq =19.3954   &    \multicolumn{2}{c} 
{ }
\\ \hline
\multirow{4}{*}{\tab[2mm] 417 \tab[2mm]} & Tetrabenzoporphyrin\_rotated02\_out\_G09 &
\multirow{2}{*}{\textcolor{Red}{\bf Equal}} & \multirow{2}{*}{\textcolor{Red}{\bf Enantiomers}}
\\
& E = -1604.15788879 \tab Freq =19.3954   &    &  \\ \cline{2-2}
& Tetrabenzoporphyrin\_rotated\_out\_G09   & \multicolumn{2}{c}
{\textcolor{Red}{ RMS = 0.0001074955}}
\\
& E = -1604.15793928 \tab Freq =20.2596   &    \multicolumn{2}{c} 
{ }
\\ \hline
\multirow{4}{*}{\tab[2mm] 418 \tab[2mm]} & Tetrabenzoporphyrin\_rotated02\_out\_G09 &
\multirow{2}{*}{\textcolor{Red}{\bf Equal}} & \multirow{2}{*}{\textcolor{Red}{\bf Enantiomers}}
\\
& E = -1604.15788879 \tab Freq =19.3954   &    &  \\ \cline{2-2}
& Tetrabenzoporphyrin\_rotated\_out\_G09\_invertion   & \multicolumn{2}{c}
{\textcolor{Red}{ RMS = 0.0001074955}}
\\
& E = -1604.15793928 \tab Freq =20.2596   &    \multicolumn{2}{c} 
{ }
\\ \hline
\multirow{4}{*}{\tab[2mm] 419 \tab[2mm]} & Tetrabenzoporphyrin\_rotated\_out\_G09 &
\multirow{2}{*}{Equal} & \multirow{2}{*}{Equal}
\\
& E = -1604.15793928 \tab Freq =20.2596   &    &  \\ \cline{2-2}
& Tetrabenzoporphyrin\_rotated\_out\_G09\_invertion   & \multicolumn{2}{c}
{ RMS = 0}
\\
& E = -1604.15793928 \tab Freq =20.2596   &    \multicolumn{2}{c} 
{ }
\\ \hline
\multirow{4}{*}{\tab[2mm] 420 \tab[2mm]} & R-R-R-Fructuose\_out\_G09 &
\multirow{2}{*}{Different} & \multirow{2}{*}{Different}
\\
& E = -687.366006448 \tab Freq =36.7188   &    &  \\ \cline{2-2}
& R-R-S-Fructuose\_out\_G09   & \multicolumn{2}{c}
 {RMS = 0.449321}
\\
& E = -687.367479320 \tab Freq =43.3327   &    \multicolumn{2}{c} 
{ }
\\ \hline
\end{tabular}
\newpage

\vtab[-2cm]
\tab[-2cm]
\begin{tabular}{c|m{8cm}|c|c}
\# & Moléculas & Restultado esperado & Resultado programa \\ \hline\hline
\multirow{4}{*}{\tab[2mm] 421 \tab[2mm]} & R-R-R-Fructuose\_out\_G09 &
\multirow{2}{*}{Different} & \multirow{2}{*}{Different}
\\
& E = -687.366006448 \tab Freq =36.7188   &    &  \\ \cline{2-2}
& R-S-R-Fructuose\_out\_G09   & \multicolumn{2}{c}
 {RMS = 0.274434}
\\
& E = -687.362461360 \tab Freq =34.8257   &    \multicolumn{2}{c} 
{ }
\\ \hline
\multirow{4}{*}{\tab[2mm] 422 \tab[2mm]} & R-R-R-Fructuose\_out\_G09 &
\multirow{2}{*}{Different} & \multirow{2}{*}{Different}
\\
& E = -687.366006448 \tab Freq =36.7188   &    &  \\ \cline{2-2}
& R-S-S-Fructuose\_out\_G09   & \multicolumn{2}{c}
 {RMS = 0.404519}
\\
& E = -687.363689916 \tab Freq =34.8064   &    \multicolumn{2}{c} 
{ }
\\ \hline
\multirow{4}{*}{\tab[2mm] 423 \tab[2mm]} & R-R-R-Fructuose\_out\_G09 &
\multirow{2}{*}{Different} & \multirow{2}{*}{Different}
\\
& E = -687.366006448 \tab Freq =36.7188   &    &  \\ \cline{2-2}
& S-R-R-Fructuose\_out\_G09   & \multicolumn{2}{c}
 {RMS = 0.299091}
\\
& E = -687.363690093 \tab Freq =34.7777   &    \multicolumn{2}{c} 
{ }
\\ \hline
\multirow{4}{*}{\tab[2mm] 424 \tab[2mm]} & R-R-R-Fructuose\_out\_G09 &
\multirow{2}{*}{Different} & \multirow{2}{*}{Different}
\\
& E = -687.366006448 \tab Freq =36.7188   &    &  \\ \cline{2-2}
& S-R-S-Fructuose\_out\_G09   & \multicolumn{2}{c}
 {RMS = 0.269535}
\\
& E = -687.362461325 \tab Freq =   &    \multicolumn{2}{c} 
{ }
\\ \hline
\multirow{4}{*}{\tab[2mm] 425 \tab[2mm]} & R-R-R-Fructuose\_out\_G09 &
\multirow{2}{*}{Different} & \multirow{2}{*}{Different}
\\
& E = -687.366006448 \tab Freq =36.7188   &    &  \\ \cline{2-2}
& S-S-R-Fructuose\_out\_G09   & \multicolumn{2}{c}
 {RMS = 0.327879}
\\
& E = -687.367479374 \tab Freq =43.3371   &    \multicolumn{2}{c} 
{ }
\\ \hline
\multirow{4}{*}{\tab[2mm] 426 \tab[2mm]} & R-R-R-Fructuose\_out\_G09 &
\multirow{2}{*}{Different} & \multirow{2}{*}{Different}
\\
& E = -687.366006448 \tab Freq =36.7188   &    &  \\ \cline{2-2}
& S-S-S-Fructuose\_out\_G09   & \multicolumn{2}{c}
 {RMS = 0.402361}
\\
& E = -687.368445846 \tab Freq =30.5992   &    \multicolumn{2}{c} 
{ }
\\ \hline
\multirow{4}{*}{\tab[2mm] 427 \tab[2mm]} & R-R-S-Fructuose\_out\_G09 &
\multirow{2}{*}{Different} & \multirow{2}{*}{Different}
\\
& E = -687.367479320 \tab Freq =43.3327   &    &  \\ \cline{2-2}
& R-S-R-Fructuose\_out\_G09   & \multicolumn{2}{c}
 {RMS = 0.174888}
\\
& E = -687.362461360 \tab Freq =34.8257   &    \multicolumn{2}{c} 
{ }
\\ \hline
\end{tabular}
\newpage

\vtab[-2cm]
\tab[-2cm]
\begin{tabular}{c|m{8cm}|c|c}
\# & Moléculas & Restultado esperado & Resultado programa \\ \hline\hline
\multirow{4}{*}{\tab[2mm] 428 \tab[2mm]} & R-R-S-Fructuose\_out\_G09 &
\multirow{2}{*}{Different} & \multirow{2}{*}{Different}
\\
& E = -687.367479320 \tab Freq =43.3327   &    &  \\ \cline{2-2}
& R-S-S-Fructuose\_out\_G09   & \multicolumn{2}{c}
 {RMS = 0.0448022}
\\
& E = -687.363689916 \tab Freq =34.8064   &    \multicolumn{2}{c} 
{ }
\\ \hline
\multirow{4}{*}{\tab[2mm] 429 \tab[2mm]} & R-R-S-Fructuose\_out\_G09 &
\multirow{2}{*}{Different} & \multirow{2}{*}{Different}
\\
& E = -687.367479320 \tab Freq =43.3327   &    &  \\ \cline{2-2}
& S-R-R-Fructuose\_out\_G09   & \multicolumn{2}{c}
 {RMS = 0.15023}
\\
& E = -687.363690093 \tab Freq =34.7777   &    \multicolumn{2}{c} 
{ }
\\ \hline
\multirow{4}{*}{\tab[2mm] 430 \tab[2mm]} & R-R-S-Fructuose\_out\_G09 &
\multirow{2}{*}{Different} & \multirow{2}{*}{Different}
\\
& E = -687.367479320 \tab Freq =43.3327   &    &  \\ \cline{2-2}
& S-R-S-Fructuose\_out\_G09   & \multicolumn{2}{c}
 {RMS = 0.179786}
\\
& E = -687.362461325 \tab Freq =   &    \multicolumn{2}{c} 
{ }
\\ \hline
\multirow{4}{*}{\tab[2mm] 431 \tab[2mm]} & R-R-S-Fructuose\_out\_G09 &
\multirow{2}{*}{Enantiomers} & \multirow{2}{*}{Enantiomers}
\\
& E = -687.367479320 \tab Freq =43.3327   &    &  \\ \cline{2-2}
& S-S-R-Fructuose\_out\_G09   & \multicolumn{2}{c}
{\textcolor{Red}{ RMS = 0.0005551068}}
\\
& E = -687.367479374 \tab Freq =43.3371   &    \multicolumn{2}{c} 
{ }
\\ \hline
\multirow{4}{*}{\tab[2mm] 432 \tab[2mm]} & R-R-S-Fructuose\_out\_G09 &
\multirow{2}{*}{Different} & \multirow{2}{*}{Different}
\\
& E = -687.367479320 \tab Freq =43.3327   &    &  \\ \cline{2-2}
& S-S-S-Fructuose\_out\_G09   & \multicolumn{2}{c}
 {RMS = 0.0469601}
\\
& E = -687.368445846 \tab Freq =30.5992   &    \multicolumn{2}{c} 
{ }
\\ \hline
\multirow{4}{*}{\tab[2mm] 433 \tab[2mm]} & R-S-R-Fructuose\_out\_G09 &
\multirow{2}{*}{Different} & \multirow{2}{*}{Different}
\\
& E = -687.362461360 \tab Freq =34.8257   &    &  \\ \cline{2-2}
& R-S-S-Fructuose\_out\_G09   & \multicolumn{2}{c}
 {RMS = 0.130085}
\\
& E = -687.363689916 \tab Freq =34.8064   &    \multicolumn{2}{c} 
{ }
\\ \hline
\multirow{4}{*}{\tab[2mm] 434 \tab[2mm]} & R-S-R-Fructuose\_out\_G09 &
\multirow{2}{*}{Different} & \multirow{2}{*}{Different}
\\
& E = -687.362461360 \tab Freq =34.8257   &    &  \\ \cline{2-2}
& S-R-R-Fructuose\_out\_G09   & \multicolumn{2}{c}
 {RMS = 0.0246573}
\\
& E = -687.363690093 \tab Freq =34.7777   &    \multicolumn{2}{c} 
{ }
\\ \hline
\end{tabular}
\newpage

\vtab[-2cm]
\tab[-2cm]
\begin{tabular}{c|m{8cm}|c|c}
\# & Moléculas & Restultado esperado & Resultado programa \\ \hline\hline
\multirow{4}{*}{\tab[2mm] 435 \tab[2mm]} & R-S-R-Fructuose\_out\_G09 &
\multirow{2}{*}{Enantiomers} & \multirow{2}{*}{Enantiomers}
\\
& E = -687.362461360 \tab Freq =34.8257   &    &  \\ \cline{2-2}
& S-R-S-Fructuose\_out\_G09   & \multicolumn{2}{c}
 {RMS = 0.005799255}
\\
& E = -687.362461325 \tab Freq =   &    \multicolumn{2}{c} 
{ }
\\ \hline
\multirow{4}{*}{\tab[2mm] 436 \tab[2mm]} & R-S-R-Fructuose\_out\_G09 &
\multirow{2}{*}{Different} & \multirow{2}{*}{Different}
\\
& E = -687.362461360 \tab Freq =34.8257   &    &  \\ \cline{2-2}
& S-S-R-Fructuose\_out\_G09   & \multicolumn{2}{c}
 {RMS = 0.0534457}
\\
& E = -687.367479374 \tab Freq =43.3371   &    \multicolumn{2}{c} 
{ }
\\ \hline
\multirow{4}{*}{\tab[2mm] 437 \tab[2mm]} & R-S-R-Fructuose\_out\_G09 &
\multirow{2}{*}{Different} & \multirow{2}{*}{Different}
\\
& E = -687.362461360 \tab Freq =34.8257   &    &  \\ \cline{2-2}
& S-S-S-Fructuose\_out\_G09   & \multicolumn{2}{c}
 {RMS = 0.127928}
\\
& E = -687.368445846 \tab Freq =30.5992   &    \multicolumn{2}{c} 
{ }
\\ \hline
\multirow{4}{*}{\tab[2mm] 438 \tab[2mm]} & R-S-S-Fructuose\_out\_G09 &
\multirow{2}{*}{Enantiomers} & \multirow{2}{*}{Enantiomers}
\\
& E = -687.363689916 \tab Freq =34.8064   &    &  \\ \cline{2-2}
& S-R-R-Fructuose\_out\_G09   & \multicolumn{2}{c}
{\textcolor{Red}{ RMS = 0.0001020726}}
\\
& E = -687.363690093 \tab Freq =34.7777   &    \multicolumn{2}{c} 
{ }
\\ \hline
\multirow{4}{*}{\tab[2mm] 439 \tab[2mm]} & R-S-S-Fructuose\_out\_G09 &
\multirow{2}{*}{Different} & \multirow{2}{*}{Different}
\\
& E = -687.363689916 \tab Freq =34.8064   &    &  \\ \cline{2-2}
& S-R-S-Fructuose\_out\_G09   & \multicolumn{2}{c}
 {RMS = 0.134984}
\\
& E = -687.362461325 \tab Freq =   &    \multicolumn{2}{c} 
{ }
\\ \hline
\multirow{4}{*}{\tab[2mm] 440 \tab[2mm]} & R-S-S-Fructuose\_out\_G09 &
\multirow{2}{*}{Different} & \multirow{2}{*}{Different}
\\
& E = -687.363689916 \tab Freq =34.8064   &    &  \\ \cline{2-2}
& S-S-R-Fructuose\_out\_G09   & \multicolumn{2}{c}
 {RMS = 0.0766398}
\\
& E = -687.367479374 \tab Freq =43.3371   &    \multicolumn{2}{c} 
{ }
\\ \hline
\multirow{4}{*}{\tab[2mm] 441 \tab[2mm]} & R-S-S-Fructuose\_out\_G09 &
\multirow{2}{*}{Different} & \multirow{2}{*}{Different}
\\
& E = -687.363689916 \tab Freq =34.8064   &    &  \\ \cline{2-2}
& S-S-S-Fructuose\_out\_G09   & \multicolumn{2}{c}
 {RMS = 0.00215791}
\\
& E = -687.368445846 \tab Freq =30.5992   &    \multicolumn{2}{c} 
{ }
\\ \hline
\end{tabular}
\newpage

\vtab[-2cm]
\tab[-2cm]
\begin{tabular}{c|m{8cm}|c|c}
\# & Moléculas & Restultado esperado & Resultado programa \\ \hline\hline
\multirow{4}{*}{\tab[2mm] 442 \tab[2mm]} & S-R-R-Fructuose\_out\_G09 &
\multirow{2}{*}{Different} & \multirow{2}{*}{Different}
\\
& E = -687.363690093 \tab Freq =34.7777   &    &  \\ \cline{2-2}
& S-R-S-Fructuose\_out\_G09   & \multicolumn{2}{c}
 {RMS = 0.029556}
\\
& E = -687.362461325 \tab Freq =   &    \multicolumn{2}{c} 
{ }
\\ \hline
\multirow{4}{*}{\tab[2mm] 443 \tab[2mm]} & S-R-R-Fructuose\_out\_G09 &
\multirow{2}{*}{Different} & \multirow{2}{*}{Different}
\\
& E = -687.363690093 \tab Freq =34.7777   &    &  \\ \cline{2-2}
& S-S-R-Fructuose\_out\_G09   & \multicolumn{2}{c}
 {RMS = 0.0287884}
\\
& E = -687.367479374 \tab Freq =43.3371   &    \multicolumn{2}{c} 
{ }
\\ \hline
\multirow{4}{*}{\tab[2mm] 444 \tab[2mm]} & S-R-R-Fructuose\_out\_G09 &
\multirow{2}{*}{Different} & \multirow{2}{*}{Different}
\\
& E = -687.363690093 \tab Freq =34.7777   &    &  \\ \cline{2-2}
& S-S-S-Fructuose\_out\_G09   & \multicolumn{2}{c}
 {RMS = 0.10327}
\\
& E = -687.368445846 \tab Freq =30.5992   &    \multicolumn{2}{c} 
{ }
\\ \hline
\multirow{4}{*}{\tab[2mm] 445 \tab[2mm]} & S-R-S-Fructuose\_out\_G09 &
\multirow{2}{*}{Different} & \multirow{2}{*}{Different}
\\
& E = -687.362461325 \tab Freq =   &    &  \\ \cline{2-2}
& S-S-R-Fructuose\_out\_G09   & \multicolumn{2}{c}
 {RMS = 0.0583444}
\\
& E = -687.367479374 \tab Freq =43.3371   &    \multicolumn{2}{c} 
{ }
\\ \hline
\multirow{4}{*}{\tab[2mm] 446 \tab[2mm]} & S-R-S-Fructuose\_out\_G09 &
\multirow{2}{*}{Different} & \multirow{2}{*}{Different}
\\
& E = -687.362461325 \tab Freq =   &    &  \\ \cline{2-2}
& S-S-S-Fructuose\_out\_G09   & \multicolumn{2}{c}
 {RMS = 0.132826}
\\
& E = -687.368445846 \tab Freq =30.5992   &    \multicolumn{2}{c} 
{ }
\\ \hline
\multirow{4}{*}{\tab[2mm] 447 \tab[2mm]} & S-S-R-Fructuose\_out\_G09 &
\multirow{2}{*}{Different} & \multirow{2}{*}{Different}
\\
& E = -687.367479374 \tab Freq =43.3371   &    &  \\ \cline{2-2}
& S-S-S-Fructuose\_out\_G09   & \multicolumn{2}{c}
 {RMS = 0.0744819}
\\
& E = -687.368445846 \tab Freq =30.5992   &    \multicolumn{2}{c} 
{ }
\\ \hline
\multirow{4}{*}{\tab[2mm] 448 \tab[2mm]} & r-enan\_out\_G09\_invertion\_output &
\multirow{2}{*}{Enantiomers} & \multirow{2}{*}{Enantiomers}
\\
& E = -9590.12967553 \tab Freq =134.7377   &    &  \\ \cline{2-2}
& r-enan\_rotated\_out\_G09   & \multicolumn{2}{c}
 {RMS = 0.01735201}
\\
& E = -9590.12967553 \tab Freq =134.7377   &    \multicolumn{2}{c} 
{ }
\\ \hline
\end{tabular}
\newpage

\vtab[-2cm]
\tab[-2cm]
\begin{tabular}{c|m{8cm}|c|c}
\# & Moléculas & Restultado esperado & Resultado programa \\ \hline\hline
\multirow{4}{*}{\tab[2mm] 449 \tab[2mm]} & r-enan\_out\_G09\_invertion\_output &
\multirow{2}{*}{Equal} & \multirow{2}{*}{Equal}
\\
& E = -9590.12967553 \tab Freq =134.7377   &    &  \\ \cline{2-2}
& s-enan\_out\_G09   & \multicolumn{2}{c}
{ RMS = 1.185059E-06}
\\
& E = -9590.12967522 \tab Freq =134.6541   &    \multicolumn{2}{c} 
{ }
\\ \hline
\multirow{4}{*}{\tab[2mm] 450 \tab[2mm]} & r-enan\_out\_G09\_invertion\_output &
\multirow{2}{*}{Enantiomers} & \multirow{2}{*}{Enantiomers}
\\
& E = -9590.12967553 \tab Freq =134.7377   &    &  \\ \cline{2-2}
& s-enan\_out\_G09\_invertion\_output   & \multicolumn{2}{c}
 {RMS = 0.005068578}
\\
& E = -9590.12967522 \tab Freq =134.6541   &    \multicolumn{2}{c} 
{ }
\\ \hline
\multirow{4}{*}{\tab[2mm] 451 \tab[2mm]} & r-enan\_out\_G09\_invertion\_output &
\multirow{2}{*}{Equal} & \multirow{2}{*}{Equal}
\\
& E = -9590.12967553 \tab Freq =134.7377   &    &  \\ \cline{2-2}
& s-enan\_rotated\_out\_G09   & \multicolumn{2}{c}
{ RMS = 1.178757E-06}
\\
& E = -9590.12967522 \tab Freq =134.6541   &    \multicolumn{2}{c} 
{ }
\\ \hline
\multirow{4}{*}{\tab[2mm] 452 \tab[2mm]} & r-enan\_out\_G09 &
\multirow{2}{*}{Enantiomers} & \multirow{2}{*}{Enantiomers}
\\
& E = -9590.12967553 \tab Freq =134.7377   &    &  \\ \cline{2-2}
& r-enan\_out\_G09\_invertion\_output   & \multicolumn{2}{c}
 {RMS = 0.01735201}
\\
& E = -9590.12967553 \tab Freq =134.7377   &    \multicolumn{2}{c} 
{ }
\\ \hline
\multirow{4}{*}{\tab[2mm] 453 \tab[2mm]} & r-enan\_out\_G09 &
\multirow{2}{*}{Equal} & \multirow{2}{*}{Equal}
\\
& E = -9590.12967553 \tab Freq =134.7377   &    &  \\ \cline{2-2}
& r-enan\_rotated\_out\_G09   & \multicolumn{2}{c}
{ RMS = 0}
\\
& E = -9590.12967553 \tab Freq =134.7377   &    \multicolumn{2}{c} 
{ }
\\ \hline
\multirow{4}{*}{\tab[2mm] 454 \tab[2mm]} & r-enan\_out\_G09 &
\multirow{2}{*}{Enantiomers} & \multirow{2}{*}{Enantiomers}
\\
& E = -9590.12967553 \tab Freq =134.7377   &    &  \\ \cline{2-2}
& s-enan\_out\_G09   & \multicolumn{2}{c}
 {RMS = 0.00506593}
\\
& E = -9590.12967522 \tab Freq =134.6541   &    \multicolumn{2}{c} 
{ }
\\ \hline
\multirow{4}{*}{\tab[2mm] 455 \tab[2mm]} & r-enan\_out\_G09 &
\multirow{2}{*}{Equal} & \multirow{2}{*}{Equal}
\\
& E = -9590.12967553 \tab Freq =134.7377   &    &  \\ \cline{2-2}
& s-enan\_out\_G09\_invertion\_output   & \multicolumn{2}{c}
{ RMS = 3.832848E-06}
\\
& E = -9590.12967522 \tab Freq =134.6541   &    \multicolumn{2}{c} 
{ }
\\ \hline
\end{tabular}
\newpage

\vtab[-2cm]
\tab[-2cm]
\begin{tabular}{c|m{8cm}|c|c}
\# & Moléculas & Restultado esperado & Resultado programa \\ \hline\hline
\multirow{4}{*}{\tab[2mm] 456 \tab[2mm]} & r-enan\_out\_G09 &
\multirow{2}{*}{Enantiomers} & \multirow{2}{*}{Enantiomers}
\\
& E = -9590.12967553 \tab Freq =134.7377   &    &  \\ \cline{2-2}
& s-enan\_rotated\_out\_G09   & \multicolumn{2}{c}
 {RMS = 0.005065924}
\\
& E = -9590.12967522 \tab Freq =134.6541   &    \multicolumn{2}{c} 
{ }
\\ \hline
\multirow{4}{*}{\tab[2mm] 457 \tab[2mm]} & r-enan\_rotated\_out\_G09 &
\multirow{2}{*}{Enantiomers} & \multirow{2}{*}{Enantiomers}
\\
& E = -9590.12967553 \tab Freq =134.7377   &    &  \\ \cline{2-2}
& s-enan\_out\_G09   & \multicolumn{2}{c}
 {RMS = 0.00506593}
\\
& E = -9590.12967522 \tab Freq =134.6541   &    \multicolumn{2}{c} 
{ }
\\ \hline
\multirow{4}{*}{\tab[2mm] 458 \tab[2mm]} & r-enan\_rotated\_out\_G09 &
\multirow{2}{*}{Equal} & \multirow{2}{*}{Equal}
\\
& E = -9590.12967553 \tab Freq =134.7377   &    &  \\ \cline{2-2}
& s-enan\_out\_G09\_invertion\_output   & \multicolumn{2}{c}
{ RMS = 3.832848E-06}
\\
& E = -9590.12967522 \tab Freq =134.6541   &    \multicolumn{2}{c} 
{ }
\\ \hline
\multirow{4}{*}{\tab[2mm] 459 \tab[2mm]} & r-enan\_rotated\_out\_G09 &
\multirow{2}{*}{Enantiomers} & \multirow{2}{*}{Enantiomers}
\\
& E = -9590.12967553 \tab Freq =134.7377   &    &  \\ \cline{2-2}
& s-enan\_rotated\_out\_G09   & \multicolumn{2}{c}
 {RMS = 0.005065924}
\\
& E = -9590.12967522 \tab Freq =134.6541   &    \multicolumn{2}{c} 
{ }
\\ \hline
\multirow{4}{*}{\tab[2mm] 460 \tab[2mm]} & s-enan\_out\_G09\_invertion\_output &
\multirow{2}{*}{Enantiomers} & \multirow{2}{*}{Enantiomers}
\\
& E = -9590.12967522 \tab Freq =134.6541   &    &  \\ \cline{2-2}
& s-enan\_rotated\_out\_G09   & \multicolumn{2}{c}
 {RMS = 0.01734936}
\\
& E = -9590.12967522 \tab Freq =134.6541   &    \multicolumn{2}{c} 
{ }
\\ \hline
\multirow{4}{*}{\tab[2mm] 461 \tab[2mm]} & s-enan\_out\_G09 &
\multirow{2}{*}{Enantiomers} & \multirow{2}{*}{Enantiomers}
\\
& E = -9590.12967522 \tab Freq =134.6541   &    &  \\ \cline{2-2}
& s-enan\_out\_G09\_invertion\_output   & \multicolumn{2}{c}
 {RMS = 0.01734936}
\\
& E = -9590.12967522 \tab Freq =134.6541   &    \multicolumn{2}{c} 
{ }
\\ \hline
\multirow{4}{*}{\tab[2mm] 462 \tab[2mm]} & s-enan\_out\_G09 &
\multirow{2}{*}{Equal} & \multirow{2}{*}{Equal}
\\
& E = -9590.12967522 \tab Freq =134.6541   &    &  \\ \cline{2-2}
& s-enan\_rotated\_out\_G09   & \multicolumn{2}{c}
{ RMS = 6.301958E-09}
\\
& E = -9590.12967522 \tab Freq =134.6541   &    \multicolumn{2}{c} 
{ }
\\ \hline
\end{tabular}
\newpage

\vtab[-2cm]
\tab[-2cm]
\begin{tabular}{c|m{8cm}|c|c}
\# & Moléculas & Restultado esperado & Resultado programa \\ \hline\hline
\multirow{4}{*}{\tab[2mm] 463 \tab[2mm]} & test-R\_out\_G09\_invertion &
\multirow{2}{*}{Enantiomers} & \multirow{2}{*}{Enantiomers}
\\
& E = -386.195553794 \tab Freq =38.2428   &    &  \\ \cline{2-2}
& test-R\_out\_G09\_rot-45-45-45   & \multicolumn{2}{c}
 {RMS = 0.1261689}
\\
& E = -386.195553794 \tab Freq =38.2428   &    \multicolumn{2}{c} 
{ }
\\ \hline
\multirow{4}{*}{\tab[2mm] 464 \tab[2mm]} & test-R\_out\_G09\_invertion &
\multirow{2}{*}{Equal} & \multirow{2}{*}{Equal}
\\
& E = -386.195553794 \tab Freq =38.2428   &    &  \\ \cline{2-2}
& test-S\_out\_G09   & \multicolumn{2}{c}
{ RMS = 1.131338E-07}
\\
& E = -386.195553794 \tab Freq =38.2428   &    \multicolumn{2}{c} 
{ }
\\ \hline
\multirow{4}{*}{\tab[2mm] 465 \tab[2mm]} & test-R\_out\_G09\_invertion &
\multirow{2}{*}{Enantiomers} & \multirow{2}{*}{Enantiomers}
\\
& E = -386.195553794 \tab Freq =38.2428   &    &  \\ \cline{2-2}
& test-S\_out\_G09\_invertion   & \multicolumn{2}{c}
{\textcolor{Red}{ RMS = 7.641079E-17}}
\\
& E = -386.195553794 \tab Freq =38.2428   &    \multicolumn{2}{c} 
{ }
\\ \hline
\multirow{4}{*}{\tab[2mm] 466 \tab[2mm]} & test-R\_out\_G09\_invertion &
\multirow{2}{*}{Equal} & \multirow{2}{*}{Equal}
\\
& E = -386.195553794 \tab Freq =38.2428   &    &  \\ \cline{2-2}
& test-S\_out\_G09\_rot-45-45-45   & \multicolumn{2}{c}
{ RMS = 1.575741E-07}
\\
& E = -386.195553794 \tab Freq =38.2428   &    \multicolumn{2}{c} 
{ }
\\ \hline
\multirow{4}{*}{\tab[2mm] 467 \tab[2mm]} & test-R\_out\_G09\_rot-45-45-45 &
\multirow{2}{*}{Enantiomers} & \multirow{2}{*}{Enantiomers}
\\
& E = -386.195553794 \tab Freq =38.2428   &    &  \\ \cline{2-2}
& test-S\_out\_G09   & \multicolumn{2}{c}
 {RMS = 0.1189346}
\\
& E = -386.195553794 \tab Freq =38.2428   &    \multicolumn{2}{c} 
{ }
\\ \hline
\multirow{4}{*}{\tab[2mm] 468 \tab[2mm]} & test-R\_out\_G09\_rot-45-45-45 &
\multirow{2}{*}{Equal} & \multirow{2}{*}{Equal}
\\
& E = -386.195553794 \tab Freq =38.2428   &    &  \\ \cline{2-2}
& test-S\_out\_G09\_invertion   & \multicolumn{2}{c}
{ RMS = 1.229453E-06}
\\
& E = -386.195553794 \tab Freq =38.2428   &    \multicolumn{2}{c} 
{ }
\\ \hline
\multirow{4}{*}{\tab[2mm] 469 \tab[2mm]} & test-R\_out\_G09\_rot-45-45-45 &
\multirow{2}{*}{Enantiomers} & \multirow{2}{*}{Enantiomers}
\\
& E = -386.195553794 \tab Freq =38.2428   &    &  \\ \cline{2-2}
& test-S\_out\_G09\_rot-45-45-45   & \multicolumn{2}{c}
 {RMS = 0.1261691}
\\
& E = -386.195553794 \tab Freq =38.2428   &    \multicolumn{2}{c} 
{ }
\\ \hline
\end{tabular}
\newpage

\vtab[-2cm]
\tab[-2cm]
\begin{tabular}{c|m{8cm}|c|c}
\# & Moléculas & Restultado esperado & Resultado programa \\ \hline\hline
\multirow{4}{*}{\tab[2mm] 470 \tab[2mm]} & test-R\_out\_G09 &
\multirow{2}{*}{Enantiomers} & \multirow{2}{*}{Enantiomers}
\\
& E = -386.195553794 \tab Freq =38.2428   &    &  \\ \cline{2-2}
& test-R\_out\_G09\_invertion   & \multicolumn{2}{c}
 {RMS = 0.2451035}
\\
& E = -386.195553794 \tab Freq =38.2428   &    \multicolumn{2}{c} 
{ }
\\ \hline
\multirow{4}{*}{\tab[2mm] 471 \tab[2mm]} & test-R\_out\_G09 &
\multirow{2}{*}{Equal} & \multirow{2}{*}{Equal}
\\
& E = -386.195553794 \tab Freq =38.2428   &    &  \\ \cline{2-2}
& test-R\_out\_G09\_rot-45-45-45   & \multicolumn{2}{c}
{ RMS = 1.324403E-06}
\\
& E = -386.195553794 \tab Freq =38.2428   &    \multicolumn{2}{c} 
{ }
\\ \hline
\multirow{4}{*}{\tab[2mm] 472 \tab[2mm]} & test-R\_out\_G09 &
\multirow{2}{*}{Enantiomers} & \multirow{2}{*}{Enantiomers}
\\
& E = -386.195553794 \tab Freq =38.2428   &    &  \\ \cline{2-2}
& test-S\_out\_G09   & \multicolumn{2}{c}
{\textcolor{Red}{ RMS = 1.273513E-17}}
\\
& E = -386.195553794 \tab Freq =38.2428   &    \multicolumn{2}{c} 
{ }
\\ \hline
\multirow{4}{*}{\tab[2mm] 473 \tab[2mm]} & test-R\_out\_G09 &
\multirow{2}{*}{Equal} & \multirow{2}{*}{Equal}
\\
& E = -386.195553794 \tab Freq =38.2428   &    &  \\ \cline{2-2}
& test-S\_out\_G09\_invertion   & \multicolumn{2}{c}
{\textcolor{Red}{ RMS = 9.494985E-08}}
\\
& E = -386.195553794 \tab Freq =38.2428   &    \multicolumn{2}{c} 
{ }
\\ \hline
\multirow{4}{*}{\tab[2mm] 474 \tab[2mm]} & test-R\_out\_G09 &
\multirow{2}{*}{Enantiomers} & \multirow{2}{*}{Enantiomers}
\\
& E = -386.195553794 \tab Freq =38.2428   &    &  \\ \cline{2-2}
& test-S\_out\_G09\_rot-45-45-45   & \multicolumn{2}{c}
 {RMS = 0.2451037}
\\
& E = -386.195553794 \tab Freq =38.2428   &    \multicolumn{2}{c} 
{ }
\\ \hline
\multirow{4}{*}{\tab[2mm] 475 \tab[2mm]} & test-S\_out\_G09\_invertion &
\multirow{2}{*}{Enantiomers} & \multirow{2}{*}{Enantiomers}
\\
& E = -386.195553794 \tab Freq =38.2428   &    &  \\ \cline{2-2}
& test-S\_out\_G09\_rot-45-45-45   & \multicolumn{2}{c}
{\textcolor{Red}{ RMS = 1.575741E-07}}
\\
& E = -386.195553794 \tab Freq =38.2428   &    \multicolumn{2}{c} 
{ }
\\ \hline
\multirow{4}{*}{\tab[2mm] 476 \tab[2mm]} & test-S\_out\_G09 &
\multirow{2}{*}{Enantiomers} & \multirow{2}{*}{Enantiomers}
\\
& E = -386.195553794 \tab Freq =38.2428   &    &  \\ \cline{2-2}
& test-S\_out\_G09\_invertion   & \multicolumn{2}{c}
 {RMS = 0.2451035}
\\
& E = -386.195553794 \tab Freq =38.2428   &    \multicolumn{2}{c} 
{ }
\\ \hline
\end{tabular}
\newpage

\vtab[-2cm]
\tab[-2cm]
\begin{tabular}{c|m{8cm}|c|c}
\# & Moléculas & Restultado esperado & Resultado programa \\ \hline\hline
\multirow{4}{*}{\tab[2mm] 477 \tab[2mm]} & test-S\_out\_G09 &
\multirow{2}{*}{Equal} & \multirow{2}{*}{Equal}
\\
& E = -386.195553794 \tab Freq =38.2428   &    &  \\ \cline{2-2}
& test-S\_out\_G09\_rot-45-45-45   & \multicolumn{2}{c}
{ RMS = 2.749179E-07}
\\
& E = -386.195553794 \tab Freq =38.2428   &    \multicolumn{2}{c} 
{ }
\\ \hline
\end{tabular}



%-----------------------------------------------------------------------------% 
% Extras
%-----------------------------------------------------------------------------% 
\section{Extras}

Durante la realización de este reporte me surgió la necesidad de crear un par de scripts que puedan ayudar para el trabajo y 
el empleo de esté programa. La tarea del script {\it scriptMake\_All\_Comparations.bsh} es realizar todas las comparaciones 
posibles en una carpeta que contenga varias geometrías en el formato de archivo XYZ, guardando el resultado en otro archivo.\\
El otro script es {\it translate\_comparatation2Latex.bsh} que utiliza la salida del anterior script para crear un resumen de las comparaciones 
y traducirlas a un formato personalizado de latex y, de esta manera, poder visualizar con mayor facilidad las comparaciones.\\
Actualmente el script sigue en desarrollo y se encuentra en el repositiorio del programa.


\end{document}
